\RequirePackage{fix-cm}
\documentclass[b5paper,twoside]{article}

\usepackage[T1]{fontenc}
\usepackage[utf8]{inputenc}
\usepackage{amsmath,amssymb,amsthm,mathtools}
\usepackage{newpxtext}
\usepackage{newpxmath}
\usepackage[fontsize=13pt]{scrextend}
\usepackage{microtype}
\usepackage{enumitem}
\usepackage{xparse}
\allowdisplaybreaks[4]

\usepackage[
  paper=b5paper,
  left=0.7cm,
  right=0.7cm,
  top=0.91cm,
  bottom=0.91cm,
  includehead,
  headheight=16pt,
  headsep=7pt
]{geometry}

\usepackage{fancyhdr}
\pagestyle{fancy}
\fancyhf{}
\fancyhead[LE]{\thepage}
\fancyhead[CE]{\footnotesize\scshape F. J. Murray and J. von Neumann}
\fancyhead[CO]{\footnotesize\scshape On Rings of Operators. IV}
\fancyhead[RO]{\thepage}
\renewcommand{\headrulewidth}{0pt}
\renewcommand{\footrulewidth}{0pt}

\usepackage{xcolor}
\definecolor{MvNLink}{HTML}{1F5A87}
\usepackage{hyperref}
\hypersetup{
  colorlinks=true,
  linkcolor=MvNLink,
  citecolor=MvNLink,
  urlcolor=MvNLink,
  linktoc=all,
  hypertexnames=false,
  bookmarksnumbered=true,
  pdfstartpage=1,
  pdfstartview=FitH,
  pdftitle={On Rings of Operators. IV},
  pdfauthor={F. J. Murray and J. von Neumann}
}

\linespread{1.1}\selectfont
\setlength{\parindent}{1.25em}
\setlength{\parskip}{0pt}
\setlist{nosep,leftmargin=2.4em}
\setlength{\emergencystretch}{2em}

\newcounter{paperchapter}
\newcounter{introsection}
\renewcommand{\thesection}{\arabic{paperchapter}.\arabic{section}}
\renewcommand{\theHsection}{\arabic{paperchapter}.\arabic{section}}

\NewDocumentCommand{\PaperChapter}{m m m O{#3}}{%
  \par\medskip
  \setcounter{paperchapter}{#1}%
  \setcounter{section}{0}%
  \phantomsection\label{#2}%
  \addcontentsline{toc}{section}{\texorpdfstring{Chapter \Roman{paperchapter}. #3}{Chapter \Roman{paperchapter}}}%
  \begin{center}\bfseries\MakeUppercase{Chapter \Roman{paperchapter}. #3}\end{center}%
  \nobreak\smallskip
}

\NewDocumentCommand{\PaperSection}{m m O{#2}}{%
  \par\medskip
  \refstepcounter{section}\label{#1}%
  \addcontentsline{toc}{subsection}{\texorpdfstring{\protect\S\thesection\quad #2}{Section \thesection}}%
  \noindent\textbf{\S\thesection}\quad\ignorespaces
}

\NewDocumentCommand{\IntroductionHeading}{}{%
  \phantomsection\label{chapter:introduction}%
  \pdfbookmark[1]{Introduction}{bookmark.chapter.introduction}%
  \begin{center}\bfseries\MakeUppercase{Introduction}\end{center}%
}

\NewDocumentCommand{\IntroSection}{m}{%
  \par\medskip
  \refstepcounter{introsection}\phantomsection\label{#1}%
  \noindent\textbf{\S\arabic{introsection}.}\quad\ignorespaces
}

\newcommand{\secref}[2]{\hyperref[#1]{\S#2}}
\newcommand{\chapterref}[2]{\hyperref[#1]{Chapter~#2}}
\newcommand{\thmref}[2]{\hyperref[#1]{Theorem~#2}}
\newcommand{\lemref}[2]{\hyperref[#1]{Lemma~#2}}
\newcommand{\defref}[2]{\hyperref[#1]{Definition~#2}}

\newtheoremstyle{mvnplain}%
  {0.55\baselineskip}{0.35\baselineskip}%
  {\itshape}{}% body font and indent
  {\bfseries}{.}{0.5em}{}
\theoremstyle{mvnplain}
\newtheorem{theorem}{Theorem}
\renewcommand{\thetheorem}{\Roman{theorem}}
\newtheorem{lemma}{Lemma}[section]
\newtheorem{definition}{Definition}[section]
\newtheorem{proposition}{Proposition}[section]
\newtheorem*{corollary}{Corollary}

\NewDocumentEnvironment{namedstatement}{m m}{%
  \par\addvspace{0.55\baselineskip}%
  \noindent\textbf{#1 #2.}\ \itshape\ignorespaces
}{%
  \par\addvspace{0.35\baselineskip}%
}

\renewenvironment{proof}[1][Proof]{%
  \par\addvspace{0.35\baselineskip}%
  \noindent\textit{#1.}\ \normalfont\ignorespaces
}{%
  \par\addvspace{0.35\baselineskip}%
}

\newcommand{\M}{\mathbf M}
\newcommand{\N}{\mathbf N}
\newcommand{\R}{\mathbf R}
\newcommand{\A}{\mathbf A}
\newcommand{\Hilbert}{\mathfrak H}
\newcommand{\Metric}[1]{\left[\!\left[#1\right]\!\right]}
\newcommand{\OpNorm}[1]{\left\lvert\!\left\lvert\!\left\lvert#1\right\rvert\!\right\rvert\!\right\rvert}
\newcommand{\set}[1]{\left\{#1\right\}}
\newcommand{\given}{\,\middle|\,}

\pdfstringdefDisableCommands{%
  \def\M{M}%
  \def\N{N}%
  \def\R{R}%
  \def\A{A}%
  \def\Hilbert{H}%
  \def\Metric#1{[[#1]]}%
  \def\Gamma{Gamma}%
  \def\operatorname#1{#1}%
  \def\widetilde#1{#1}%
  \def\mathfrak#1{#1}%
}

\begin{document}
\pagenumbering{arabic}
\setcounter{page}{1}

\thispagestyle{empty}

\noindent{\small\scshape Annals of Mathematics}\\[-0.2em]
{\small Vol. 44, No. 4, October, 1943}

\vspace{1.8em}
\begin{center}
  {\Large\bfseries\MakeUppercase{On Rings of Operators. IV}}\\[1.0em]
  {\large By F. J. Murray and J. von Neumann}\\[0.8em]
  {Received February 24, 1943}
\end{center}

\vspace{0.8em}
\IntroductionHeading

\IntroSection{intro:1}
The object of this paper is the investigation of the isomorphism properties of
those operator rings which are factors and which are the substratum of several
earlier publications of the authors. (Cf. \cite{ref5,ref6,ref8}. For the
detailed definitions and also for the cases of factors---of which there will be
more to say below---cf. \cite{ref5}, p.~138, Definition 3.1.2, p.~172,
Theorem VIII.) The discrete cases---i.e. the cases (I)---have been exhaustively
dealt with before. (Cf. \cite{ref5}, p.~173, Lemma 8.6.1, p.~139,
Definitions 3.2.1 and 3.2.2 and Lemmas 3.2.1--3.2.3.) The purely infinite
case---i.e. the case (III)---is the most refractory of all and we have, at least
for the time being, scarcely any tools to investigate it. (\cite{ref8} deals
mainly with these factors.) Thus we are left with the continuous cases---i.e.
the cases (II)---and they are our main objective in this paper.

An added justification of this program may be found in the fact that among all
factors those of the finite continuous case $(\mathrm{II}_1)$ have the
strongest immediate interest. (Cf. \cite{ref5}, part V, \cite{ref6}, Chapter
IV and Appendix.)
% Editorial note E1: the source reads "finite continuous case case"; the
% duplicated word has been removed.

It will be seen however that the discrete cases can be included in our
discussion with scarcely any extra effort. So we shall deal with them, i.e.
with all cases but the purely infinite one.

For the discrete and continuous cases, the finite ones---i.e. the cases
$(\mathrm I_n)$ $(n=1,2,\ldots)$ and $(\mathrm{II}_1)$ respectively---are
basic because the infinite ones---i.e. the cases $(\mathrm I_\infty)$ and
$(\mathrm{II}_\infty)$ respectively---can be subsequently described with
their help. (Cf. \thmref{thm:IX}{IX} and \lemref{lem:3.1.6}{3.1.6}.)
Therefore we shall direct our main effort on the finite cases. And since the
discrete ones---(the cases $(\mathrm I_n)$, $(n=1,2,\ldots)$)---are just the
finite order matrix rings, this means essentially the above-mentioned
continuous finite case $(\mathrm{II}_1)$.

\IntroSection{intro:2}
Let us now state the main problems of isomorphism more precisely.

Consider an operator ring $\M$ in a Hilbert space $\Hilbert$ which contains
$1$. (We do not restrict $\M$ in any other way yet. For the significant
definitions, cf. \cite{ref2}, pp.~388--389, Definitions 1--3.) Then there exist
two kinds of notions in $\M$: First, those which can be expressed in terms of
the entity $1$ (the unit operator) and the operations $\alpha A$ ($\alpha$ any
complex number), $A^*$, $A+B$, $A\cdot B$ alone and referring only to the
operators belonging to $\M$; Second those which need other things as well,
e.g. operators outside of $\M$, elements of $\Hilbert$, etc. The former
notions are purely algebraical, the latter ones are not; we shall also call
them spatial.

These notions were already investigated by the second author in \cite{ref9}.

It seems worth while to formulate this distinction in terms of isomorphisms.
Let, for each $i=1,2$, a Hilbert space $\Hilbert_i$ and an operator ring
$\M_i$ in $\Hilbert_i$ be given. We now define

\begin{description}[style=nextline,leftmargin=1.5em]
\item[A)] Spatial isomorphism of $\M_1$ and $\M_2$. This is a one-to-one
isomorphic mapping of $\Hilbert_1$ on $\Hilbert_2$ (i.e. one which is linear
and isometric---unitary) which carries $\M_1$ into $\M_2$.
\item[B)] Algebraical ring isomorphism of $\M_1$ and $\M_2$. This is a
one-to-one mapping of $\M_1$ onto $\M_2$ which leaves the entity $1$ and the
operations $\alpha A$ ($\alpha$ any complex number), $A^*$, $A+B$, $AB$
(when passing from $\M_1$ to $\M_2$) invariant.
\end{description}

(Cf. also \cite{ref5}, p.~145, \S5.2, in particular Definition 5.2.1. We have
added the requirement that $1$ be invariant.)

Any spatial property is invariant under the spatial isomorphisms of A), while
the purely algebraical properties are characterized by their invariance under
the algebraical isomorphisms of B).

Clearly A) implies B) while the converse need not be true.

An important ring theoretical notion which is only spatial is $\M'$. (Cf.
\cite{ref2}, p.~388, Def. 3. For the spatial character, cf. the detailed
discussion of \secref{sec:3.3}{3.3}.) Nevertheless the factor property, i.e.
$\M\cap\M'=(\alpha\cdot1)$ (cf. \secref{sec:3.3}{3.3}, loc. cit., and also
the beginning of \hyperref[intro:1]{\S1}) is purely algebraical since it
states that the operators $\alpha\cdot1$ exhaust the center of $\M$ (cf.
\cite{ref5}, p.~138, Def. 3.1.2).

Now our program is subdivided as follows.

\noindent Question I: When does B) hold?\\
Question II: Under what additional conditions does B) imply A)?

Question II was already investigated in a special case in \cite{ref6}. It was
shown there that under certain conditions A) and B) are equivalent. (Cf.
\cite{ref6}, p.~244, Theorem XI.) We shall obtain a complete answer to
Question II in \thmref{thm:X}{X}.

Question I is more difficult. It coincides with Problem 6 in \cite{ref5},
p.~172, and the present paper contains what progress the authors have been
able to make in that direction. The main results are these: An extensive class
of factors of case $(\mathrm{II}_1)$ which are all isomorphic to each other in
the sense B) will be determined. These factors are called ``approximately
finite''. (Cf. \thmref{thm:XII}{XII}, \thmref{thm:XIV}{XIV}, which are based
on \defref{def:4.1.1}{4.1.1}, \defref{def:4.3.1}{4.3.1},
\defref{def:4.5.2}{4.5.2} and \defref{def:4.6.1}{4.6.1} below.) The
isomorphisms announced in \cite{ref5}, p.~229, (v) are contained in this
class. On the other hand, certain factors of case $(\mathrm{II}_1)$ which are
not isomorphic to the approximately finite ones will be constructed. (Cf.
\thmref{thm:XVI}{XVI}, \thmref{thm:XVI-prime}{XVI$'$}.)

\IntroSection{intro:3}
There are indications to the effect that the approximately finite factors are
the simplest among those of case $(\mathrm{II}_1)$ but the evidence is not
quite conclusive. It is true that every factor in case $(\mathrm{II}_1)$ has an
approximately finite sub-ring. (Cf. \thmref{thm:XIII}{XIII}.) However, this
``imbedding theorem'' does not settle the matter, since the analogue of the
Cantor--Bernstein ``equivalence theorem'' is not true: Two factors in the case
$(\mathrm{II}_1)$ might be such that each is isomorphic to a sub-ring of the
other, but the factors themselves may not be isomorphic. (An example of this
is given in the \hyperref[appendix]{appendix}.) Hence the possibility exists
that any factor in the case $(\mathrm{II}_1)$ is isomorphic to a sub-ring of
any other such factor.

\IntroSection{intro:4}
The best we can do at present concerning the isomorphism problem in the case
$(\mathrm{II}_1)$ is this: The limits of the approximately finite case are
extended rather far in \secref{sec:5.2}{5.2} and \secref{sec:5.6}{5.6}. The
existence of non-approximately-finite factors in this case is established in
\thmref{thm:XVI}{XVI}, \thmref{thm:XVI-prime}{XVI$'$}. Certain algebraical
invariants of factors in the case $(\mathrm{II}_1)$ are formed. ((1), (2), in
\secref{sec:4.8}{4.8}, and the property $\Gamma$ in
\defref{def:6.1.1}{6.1.1}) of which the first two are probably of greater
general significance, but the last one has so far been put to greater
practical use. (Cf. the remark at the beginning of \secref{sec:6.1}{6.1}.)

The isomorphism questions of the case $(\mathrm{II}_\infty)$ are reduced to
those of the case $(\mathrm{II}_1)$. (This follows from
\thmref{thm:IX}{IX}.) Those of the discrete cases---the cases (I)---have been
settled before (cf. above at the beginning of \hyperref[intro:1]{\S1}; also
$\alpha$) in \thmref{thm:IX}{IX}).

This enumeration exhausts our present program. It answers Problems 6, 7 in
\cite{ref5}, p.~172, and \cite{ref5} eod., p.~229, as far as possible at this
moment.

\IntroSection{intro:5}
We add that \secref{sec:5.3}{5.3} contains a new technique of constructing
factors in the case $(\mathrm{II}_1)$. This seems to be considerably simpler
than our previous procedures, but it is closely related to them. (Cf.
\secref{sec:5.4}{5.4}, \secref{sec:5.5}{5.5}.) It also throws some light on
the meaning of this work from the point of view of the unitary representation
theory of groups (cf. the remarks after \lemref{lem:5.3.4}{5.3.4}). Further
generalizations are probably possible and important. (Cf. the remark at the
end of \secref{sec:5.6}{5.6}.)

The result that not all factors in the case $(\mathrm{II}_1)$ are isomorphic
to each other, expressed in \thmref{thm:XVI-prime}{XVI$'$}, deserves some
further comment. From the point of view of the systematic build-up of the
paper, it seemed best to put it at the end. It is, however, intelligible and of
interest in itself, and it can be derived independently of most of the paper.
The reader who is primarily interested in this particular result need only
read \secref{sec:5.3}{5.3}, \defref{def:6.1.1}{6.1.1},
\lemref{lem:6.1.1}{6.1.1}, \lemref{lem:6.2.1}{6.2.1},
\lemref{lem:6.2.2}{6.2.2}. The entire remainder of this paper is unnecessary
for this purpose, in particular the extensive theories of algebraical and
spatial types, of genera, and of approximate finiteness and its various
equivalent forms. But we think, nevertheless, that knowledge of the whole
paper will help to see this result too in the proper perspective.

\IntroSection{intro:6}
The notations are the same as in the papers referred to in the bibliography,
particularly \cite{ref5} and \cite{ref6}.

We use the Kronecker--Weierstrass symbol in the general sense:
\[
  \delta_{a,b}=
  \begin{cases}
    1&\text{for }a=b,\\
    0&\text{otherwise},
  \end{cases}
\]
for any objects $a,b$.

\renewcommand{\refname}{\centering\MakeUppercase{Bibliography}}
\begin{thebibliography}{9}
\bibitem{ref1}
John von Neumann,
\href{https://doi.org/10.1007/BF01782338}{\emph{Allgemeine Eigenwerttheorie
Hermitescher Funktionaloperatoren}}, Math. Annalen, vol.~102 (1929),
pp.~49--131.

\bibitem{ref2}
John von Neumann,
\href{https://doi.org/10.1007/BF01782352}{\emph{Zur Algebra der
Funktionaloperatoren}}, Math. Annalen, vol.~102 (1929), pp.~370--427.

\bibitem{ref3}
John von Neumann,
\href{https://doi.org/10.2307/1968185}{\emph{\"Uber Funktionen von
Funktionaloperatoren}}, Annals of Math., vol.~32 (1931), pp.~191--226.

\bibitem{ref4}
John von Neumann,
\href{https://doi.org/10.2307/1968692}{\emph{On a certain topology for rings
of operators}}, Annals of Math., vol.~37 (1936), pp.~111--115.

\bibitem{ref5}
F. J. Murray and John von Neumann,
\href{https://doi.org/10.2307/1968693}{\emph{On rings of operators}}, Annals
of Math., vol.~37 (1936), pp.~116--229.

\bibitem{ref6}
F. J. Murray and John von Neumann,
\href{https://doi.org/10.1090/S0002-9947-1937-1501899-4}{\emph{On rings of
operators II}}, Trans. Amer. Math. Soc., vol.~41 (1937), pp.~208--248.

\bibitem{ref7}
John von Neumann,
\href{https://www.numdam.org/item/CM_1939__6__1_0/}{\emph{On infinite direct
products}}, Compositio Mathematica, vol.~6 (1938), pp.~1--77.

\bibitem{ref8}
John von Neumann,
\href{https://doi.org/10.2307/1968823}{\emph{On rings of operators III}},
Annals of Math., vol.~41 (1940), pp.~94--161.

\bibitem{ref9}
John von Neumann,
\href{https://doi.org/10.2307/1969106}{\emph{On some algebraical properties
of operator rings}}, Annals of Math., vol.~44, preceding immediately.
\end{thebibliography}

\phantomsection\label{contents}
\pdfbookmark[1]{Contents}{bookmark.contents}
\begin{center}\bfseries\MakeUppercase{Contents}\end{center}

\newcommand{\ContentsChapter}[2]{%
  \par\noindent\hyperref[#1]{\textbf{#2}}\par}
\newcommand{\ContentsSection}[3]{%
  \par\hangindent=2.8em\hangafter=1\noindent\hspace*{1.5em}%
  \hyperref[#1]{\S#2\quad #3}\par}
\newcommand{\ContentsStatement}[3]{%
  \par\hangindent=4.7em\hangafter=1\noindent\hspace*{3.3em}%
  \hyperref[#1]{#2.\quad #3}\par}

\ContentsChapter{chapter:introduction}{Introduction}
\ContentsSection{intro:1}{1}{Purpose of the paper.}
\ContentsSection{intro:2}{2}{Algebraical and spatial isomorphisms. Questions I, II.}
\ContentsSection{intro:3}{3}{Approximate finiteness. Imbedding.}
\ContentsSection{intro:4}{4}{Outline of results.}
\ContentsSection{intro:5}{5}{The role of new examples.}
\ContentsSection{intro:6}{6}{Notations.}

\ContentsChapter{chapter:I}{Chapter I. The metric of $\M$. Purely algebraical character of certain notions}
\ContentsSection{sec:1.1}{1.1}{Purely algebraical concepts; Difficulties; The operator topologies; Subrings.}
\ContentsSection{sec:1.2}{1.2}{Character of $\operatorname{Tr}_{\M}(A)$, $\Metric{A}$. Metric and restricted metric closure. Linear sets. Algebras.}
\ContentsSection{sec:1.3}{1.3}{Restricted metric convergence and strong convergence.}
\ContentsSection{sec:1.4}{1.4}{Metric closures and strong closure.}
\ContentsSection{sec:1.5}{1.5}{Continuation of the discussion of closure.}
\ContentsSection{sec:1.6}{1.6}{Conclusion of the discussion of closure.}
\ContentsStatement{thm:I}{Theorem I}{Equivalence of various closures.}
\ContentsStatement{thm:II}{Theorem II}{Subrings are purely algebraical concepts.}

\ContentsChapter{chapter:II}{Chapter II. The algebraical type and its operations. The fundamental group}
\ContentsSection{sec:2.1}{2.1}{Algebraical type $\overline{\M}$. Discussion.}
\ContentsSection{sec:2.2}{2.2}{Conjugate dual isomorphisms. $\overline{\M}^{c}$ and $\overline{\M'}$.}
\ContentsSection{sec:2.3}{2.3}{Properties of the above. $\overline{\M}^{cc}=\overline{\M}=\overline{\M}''$.}
\ContentsStatement{thm:III}{Theorem III}{Properties of $\overline{\M}^{c}$.}
\ContentsSection{sec:2.4}{2.4}{$p$-matrix isomorphisms. Construction of $\overline{\M}^{p}$.}
\ContentsStatement{thm:IV}{Theorem IV}{Properties of $\overline{\M}^{p}$.}
\ContentsSection{sec:2.5}{2.5}{One-valuedness of $\overline{\M}^{\infty}$.}
\ContentsStatement{thm:IV-prime}{Theorem IV$'$}{One-valuedness of $\overline{\M}^{\infty}$.}
\ContentsSection{sec:2.6}{2.6}{Inversion of $\overline{\M}=\overline{\N}^{p}$. Existence of $\overline{\M}^{1/p}$.}
\ContentsStatement{thm:V}{Theorem V}{Inversion of $\overline{\M}=\overline{\N}^{p}$.}
\ContentsSection{sec:2.7}{2.7}{The operation $\overline{\N}^{p}$ and the factor character, case, finiteness.}
\ContentsSection{sec:2.8}{2.8}{The operation $\overline{\M}^{\alpha}$, $(0<\alpha\leq1)$ and its properties.}
\ContentsStatement{thm:VI}{Theorem VI}{Properties of $\overline{\M}^{\alpha}$, $(0<\alpha\leq1)$.}
\ContentsSection{sec:2.9}{2.9}{The operation $\overline{\M}^{\theta}$, ($\theta$ general) and its properties.}
\ContentsStatement{thm:VII}{Theorem VII}{Properties of $\overline{\M}^{\theta}$ ($\theta$ general).}
\ContentsSection{sec:2.10}{2.10}{The fundamental group $\mathfrak G(\overline{\M})$. Special cases.}
\ContentsStatement{thm:VIII}{Theorem VIII}{Properties of $\mathfrak G(\overline{\M})$.}

\ContentsChapter{chapter:III}{Chapter III. Characterization of all spatial types in terms of algebraical types}
\ContentsSection{sec:3.1}{3.1}{Commensurability. The genus $\widetilde{\overline{\M}}$.}
\ContentsSection{sec:3.2}{3.2}{Elementary properties of the genus. Genus and fundamental group.}
\ContentsSection{sec:3.3}{3.3}{The spatial type $\widetilde{\M}$. The number $\theta=(1/c)(D_{\M'}(\Hilbert)/D_{\M}(\Hilbert))$. Characterization of $\widetilde{\M}$ in terms of $\overline{\M}$, $\overline{\M'}$ and $\theta$. Precise relationship between $\overline{\M}$, $\overline{\M'}$ and $\theta$.}
\ContentsStatement{thm:X}{Theorem X}{$\widetilde{\M}$ and $\overline{\M}$, $\overline{\M'}$ and $\theta$.}

\ContentsChapter{chapter:IV}{Chapter IV. Approximate finiteness}
\ContentsSection{sec:4.1}{4.1}{Approximate finiteness of type $[p_1,p_2,\ldots]$.}
\ContentsStatement{thm:XI}{Theorem XI}{First isomorphism theorem.}
\ContentsSection{sec:4.2}{4.2}{Further approximation results.}
\ContentsSection{sec:4.3}{4.3}{Approximate finiteness (A).}
\ContentsSection{sec:4.4}{4.4}{Comparison of (A) and type $[p_1,p_2,\ldots]$.}
\ContentsSection{sec:4.5}{4.5}{Approximate finiteness (B). Structure of finite rings.}
\ContentsSection{sec:4.6}{4.6}{Approximate finiteness (C). Equivalence results.}
\ContentsStatement{thm:XII}{Theorem XII}{Equivalence theorem.}
\ContentsSection{sec:4.7}{4.7}{Existence results.}
\ContentsStatement{thm:XIII}{Theorem XIII}{Imbedding theorem.}
\ContentsStatement{thm:XIV}{Theorem XIV}{Final isomorphism theorem.}
\ContentsSection{sec:4.8}{4.8}{The operations $\overline{\M}^{c}$ and $\overline{\M}^{\theta}$ in conjunction with approximate finiteness. The corresponding invariants.}
\ContentsStatement{thm:XV}{Theorem XV}{$\overline{\A}^{c}=\overline{\A}$, $\mathfrak G(\overline{\A})=P$.}

\ContentsChapter{chapter:V}{Chapter V. Discussion of various examples}
\ContentsSection{sec:5.1}{5.1}{General remarks on the examples.}
\ContentsSection{sec:5.2}{5.2}{Approximate finiteness of certain existing examples.}
\ContentsSection{sec:5.3}{5.3}{Construction of the new examples.}
\ContentsSection{sec:5.4}{5.4}{First connection between the two kinds of examples.}
\ContentsSection{sec:5.5}{5.5}{Second connection between the two kinds of examples.}
\ContentsSection{sec:5.6}{5.6}{Approximate finiteness among the new examples.}

\ContentsChapter{chapter:VI}{Chapter VI. Not approximately finite factors}
\ContentsSection{sec:6.1}{6.1}{The property $\Gamma$ and approximate finiteness.}
\ContentsSection{sec:6.2}{6.2}{Existence of not approximately finite factors.}
\ContentsStatement{thm:XVI}{Theorem XVI}{Existence theorem.}
\ContentsStatement{thm:XVI-prime}{Theorem XVI$'$}{Existence of several algebraical types of case $(\mathrm{II}_1)$.}
\ContentsSection{sec:6.3}{6.3}{Discussion of further examples.}

\ContentsChapter{appendix}{Appendix: Example of two non-isomorphic factors in Case $(\mathrm{II}_1)$, each of which is isomorphic to a subring of the other.}
\ContentsStatement{thm:A}{Theorem A}{Incomparability result.}

\PaperChapter{1}{chapter:I}{The Metric of $\M$. Purely Algebraical Character of Various Notions}[The Metric of M. Purely Algebraical Character of Various Notions]
\PaperSection{sec:1.1}{Purely algebraical concepts; Difficulties; The operator topologies; Subrings}
Consider a Hilbert space $\Hilbert$ and an operator ring $\M$ in $\Hilbert$.
We know from \hyperref[intro:2]{\S2 in the introduction} as well as from
\cite{ref9} that the purely algebraical character of certain notions in $\M$
is neither obvious nor certain.

It was demonstrated in \cite{ref9}, however, that the following notions in
$\M$ are purely algebraical.

Definiteness of an $A\in\M$; the numerical value $\OpNorm{A}$ of the bound of
an $A\in\M$; convergence in the sense of the strong and of the weak operator
topologies. (Convergence in the sense of the strongest topology is the same as
in the strong one. Cf. \cite{ref4}, p.~112, end of \S2.)

We did not establish the same thing for the operator topologies themselves,

i.e. for the strongest, the strong and the weak.  (For the definitions, cf.
\cite[p.~111, \S2]{ref4}, and \cite[pp.~381--382]{ref2}.)  The uniform topology
(cf. \cite[pp.~384--386]{ref2}) causes no difficulty, since
$\OpNorm{A}$ is purely algebraical, but it will play no role in our discussions.
This is very unsatisfactory since the notion of a subring of $\M$ must be based
on one of these topologies.  (Cf. \cite[p.~112, beginning of \S3]{ref4};
\cite[p.~396, end of \S3; p.~388, Definition~1]{ref2}.  Best consulted in the
reverse order.)  They cannot be replaced by convergence notions, nor by the
uniform topology.  (Cf. in particular \cite[pp.~382--383 and p.~384]{ref2}.)
And yet the notion of a subring ought to be purely algebraical.

For our present purposes it will be adequate to settle these questions for a
restricted class of rings $\M$ and this we now do.

\PaperSection{sec:1.2}{Character of $\operatorname{Tr}_{\M}(A)$ and
$\Metric{A}$. Metric and restricted metric closure. Linear sets.}
Observe first that the following notions are purely algebraical.

The factor character of $\M$ and the case to which $\M$ belongs, the numerical
value of the relative dimension function $D_{\M}(E)$ ($E$ any projection
$\in\M$) apart from its normalization and even the standard normalization in
the cases where one is defined; the numerical value of the relative trace
$\operatorname{Tr}_{\M}(A)$, ($A\in\M$) in the finite cases.

These assertions were established in \cite[p.~145, Lemma~5.2.1;
p.~173, Theorem~IX]{ref5}; \cite[p.~219, Property~IV]{ref6} (since that
characterization is purely algebraical).

We assume now that $\M$ is a factor in a finite case.\footnote{A finite case is
either a $(\mathrm I_n)$, $n=1,2,\ldots$, or a $(\mathrm{II}_1)$.  The inclusion
of the former is unnecessary, since our real interest is in case
$(\mathrm{II}_1)$.  But these considerations have the totality of all finite
cases as their natural scope.}  We shall use the relative dimension function
$D_{\M}(E)$ and it will be advantageous not to normalize it.  We shall also use
the relative trace $\operatorname{Tr}_{\M}(A)$ with its usual properties.  We
define for every $A\in\M$
\begin{equation}
  \Metric{A}
  =\sqrt{\operatorname{Tr}_{\M}(A^*A)}
  =\sqrt{\operatorname{Tr}_{\M}(AA^*)}.
  \tag{1.2.\ensuremath{\alpha}}\label{eq:1.2.alpha}
\end{equation}
(Cf. \cite[p.~241, Lemma~4.3.2]{ref6}.  The considerations of
\cite[p.~102, Definition~1.3.1 and Theorem~II]{ref8} are somewhat more general
but of the same type.)  Owing to our above remarks, the numerical value of
$\Metric{A}$ is purely algebraical too.

By \cite[p.~235, Theorem~III]{ref6}, there exists a fixed finite system
$g_1,\ldots,g_m\in\Hilbert$ (depending on $\M$ only), such that for all
$A\in\M$
\begin{equation}
  \operatorname{Tr}_{\M}(A)=\sum_{i=1}^{m}(Ag_i,g_i).
  \tag{1.2.\ensuremath{\beta}}\label{eq:1.2.beta}
\end{equation}
Combining \eqref{eq:1.2.alpha} with \eqref{eq:1.2.beta}, and remembering that
$(A^*Ag,g)=\lVert Ag\rVert^2$, $(AA^*g,g)=\lVert A^*g\rVert^2$, we obtain
\begin{equation}
  \Metric{A}
  =\sqrt{\sum_{i=1}^{m}\lVert Ag_i\rVert^2}
  =\sqrt{\sum_{i=1}^{m}\lVert A^*g_i\rVert^2}.
  \tag{1.2.\ensuremath{\gamma}}\label{eq:1.2.gamma}
\end{equation}
This formula is useful for many purposes, although it obscures the purely
algebraical character of $\Metric{A}$.

$\Metric{A}$ has all the properties of a \emph{norm} in a linear space, and
accordingly $\Metric{A-B}$ has all those of a \emph{distance}.  (Cf. for these
and for some further properties \cite[p.~242, Property~II\textsuperscript{o}]{ref6}.)
We need these for $\M$ only, and not, as loc. cit., for its extension $Q(\M)$.
All these properties can be read off the representation
\eqref{eq:1.2.gamma}---the decisive fact being, of course, that we have them
both.)  Thus $\M$ has become a topological space in a new way, with a purely
algebraical metric.\footnote{The uniform topology, already mentioned in
\secref{sec:1.1}{1.1}, is actually of the same character: $\OpNorm{A}$,
$\OpNorm{A-B}$ are a norm and a metric, just as $\Metric{A}$,
$\Metric{A-B}$.  But only the latter is useful in the theory of operator rings.
This is due to its connection with the strong topology, which will appear in
\secref{sec:1.4}{1.4} and \secref{sec:1.5}{1.5}.}

We now define
\begin{definition}\phantomsection\label{def:1.2.1}
A sequence $A_1,A_2,\ldots\in\M$ is metrically convergent to $A\in\M$ if
$\lim_{\nu\to\infty}\Metric{A_\nu-A}=0$.  It is restrictedly metrically
convergent to $A\in\M$ if $\lim_{\nu\to\infty}\Metric{A_\nu-A}=0$ and the
(numerical) sequence $\OpNorm{A_1},\OpNorm{A_2},\ldots$ is bounded.
\end{definition}

The notions of closure for subsets $\mathbf S$ of $\M$ which are induced by
these two notions of convergence, are the ``metrical closure'' and the
``restricted metrical closure'' of $\mathbf S$.\footnote{The metrical closure is
the conventional topological notion, based on the specified metric.  The
restricted metrical closure is not of that nature.  Thus it is not evident that
the set of all limits of all restrictedly metrically convergent sequences from a
given set $\mathbf S$ is necessarily closed in that sense.  Actually this is not
true for all sets $\mathbf S$ although when $\mathbf S$ is an algebra it does
follow from the results of \secref{sec:1.5}{1.5} (cf. there).}

We shall also need:
\begin{definition}\phantomsection\label{def:1.2.2}
A subset $\mathbf S$ of $\M$ is linear if $A,B\in\mathbf S$ imply
$\alpha A,A+B\in\mathbf S$.  It is an algebra if $A,B\in\mathbf S$ imply
$\alpha A,A^*,A+B,AB\in\mathbf S$.
\end{definition}
Clearly both these definitions define purely algebraical notions.

Now a subring $\mathbf S$ of $\M$ is a closed algebra, where closure may be
understood in any one of the three operator topologies: the strongest, the
strong, or the weak (cf. the end of \secref{sec:1.1}{1.1}).  We shall establish
the purely algebraical character of the notion of a subring by proving the
equivalence of the notion of closure in all these topologies to the purely
algebraical closures of \defref{def:1.2.1}{1.2.1} for a set $\mathbf S$ which is
an algebra.\footnote{Observe that restrictions on both the ring $\M$ and its
subset $\mathbf S$ are necessary in order to secure this result.  $\M$ must be a
factor in a finite case, $\mathbf S$, an algebra.}

\PaperSection{sec:1.3}{Restricted metric convergence and strong convergence.}
Observe first:
\begin{lemma}\phantomsection\label{lem:1.3.1}
For the $g_1,\ldots,g_m$ of \eqref{eq:1.2.beta} and \eqref{eq:1.2.gamma} the
$A'g_i$, ($A'\in\M'$, $i=1,\ldots,m$) span the closed linear set $\Hilbert$.
\end{lemma}
\begin{proof}
For a fixed $i$ ($=1,\ldots,m$) the $A'g_i$, $A'\in\M'$, span the closed
linear set $\mathfrak M_{g_i}^{\M'}$ (cf. \cite[p.~143,
Definition~5.1.1]{ref5}).  Put
$\mathfrak N=[\mathfrak M_{g_i}^{\M'},\ i=1,\ldots,m]$.  We must prove
$\mathfrak N=\Hilbert$.  Now $\mathfrak N\eta\M'$ (cf. loc. cit.,
Lemma~5.1.1).  Hence $\mathfrak N\eta\M$ and $E=P_{\mathfrak N}\in\M$.
Also $g_i\in\mathfrak M_{g_i}^{\M'}\subset\mathfrak N$, so
\begin{equation}
  Eg_i=g_i.
  \tag{1.3.\ensuremath{\alpha}}\label{eq:1.3.alpha}
\end{equation}
Now \eqref{eq:1.2.gamma} shows that if $A\in\M$ and $Ag_i=0$ for
$i=1,\ldots,m$, then $A=0$.  Apply this to $A=1-E$.  Then
\eqref{eq:1.3.alpha} yields $1-E=0$, $P_{\mathfrak N}=E=1$, as desired.
\end{proof}

\begin{lemma}\phantomsection\label{lem:1.3.2}
The sequence $A_1,A_2,\ldots\in\M$ is strongly convergent to $A\in\M$ if and
only if it is restrictedly metrically convergent to this $A$.
\end{lemma}
\begin{proof}
Necessity: Assume that $A_1,A_2,\ldots$ is strongly convergent to $A$.  Then
the (numerical) sequence $\OpNorm{A_1},\OpNorm{A_2},\ldots$ is bounded by
\cite[p.~382, footnote~35]{ref2}, and
$\lim_{\nu\to\infty}\Metric{A_\nu-A}=0$ by \eqref{eq:1.2.gamma}.  Hence we
have restricted metric convergence.

Sufficiency: Assume that $A_1,A_2,\ldots$ are restrictedly metrically
convergent to $A$.  We must show that they converge strongly too, i.e. that the
set $\mathfrak M$ of all $f$ with
$\lim_{\nu\to\infty}\lVert(A_\nu-A)f\rVert=0$ is $\Hilbert$.
$\mathfrak M$ is obviously a linear set.  Since the (numerical) sequence
$\OpNorm{A_1},\OpNorm{A_2},\ldots$ is bounded, $\mathfrak M$ is also closed.
Hence it suffices to show that $\mathfrak M$ spans the closed linear set
$\Hilbert$.

Now by \eqref{eq:1.2.gamma},
$\lVert(A_\nu-A)g_i\rVert\leq\Metric{A_\nu-A}$.  Hence
$\lim_{\nu\to\infty}\lVert(A_\nu-A)g_i\rVert=0$.  Therefore for every
$A'\in\M'$, $\lim_{\nu\to\infty}\lVert A'(A_\nu-A)g_i\rVert=0$.  But
$A'(A_\nu-A)=(A_\nu-A)A'$ since $A'\in\M'$ and $A_\nu,A\in\M$.
Consequently
$\lim_{\nu\to\infty}\lVert(A_\nu-A)A'g_i\rVert=0$ and all
$A'g_i\in\mathfrak M$.  Hence $\mathfrak M$ spans the closed linear set
$\Hilbert$ by \lemref{lem:1.3.1}{1.3.1}, and the proof is completed.
\end{proof}

Thus closure with respect to strong convergence is equivalent to restricted
metric closure for every subset $\mathbf S$ of $\M$.  This however is not enough
to establish the purely algebraical character of the notion of a subring, since
that involves topological closure which cannot be replaced by convergence
closure (cf. the end of \secref{sec:1.1}{1.1}).  In fact, from that point of view
we have made no progress at all, since we knew all along from \cite{ref9} that
strong convergence is a purely algebraical notion for all rings $\M$.  Besides
our present restrictions on $\M$, we shall also have to restrict $\mathbf S$ to
algebras before we can get any further (cf. the end of \secref{sec:1.2}{1.2}).

\PaperSection{sec:1.4}{Metric closures and strong closure.}
Consider first the following two simple results, valid for all subsets
$\mathbf S$ of $\M$.
\begin{lemma}\phantomsection\label{lem:1.4.1}
Strong closure of $\mathbf S$ implies the restricted metric closure.
\end{lemma}
\begin{proof}
Strong closure of $\mathbf S$ implies closure with respect to strong convergence,
and by \lemref{lem:1.3.2}{1.3.2} the latter is equivalent to restricted metric
closure.
\end{proof}

\begin{lemma}\phantomsection\label{lem:1.4.2}
Metric closure of $\mathbf S$ implies its strong closure.
\end{lemma}
\begin{proof}
Let $\mathbf S$ be metrically closed and consider a strong condensation point
$A$ of $\mathbf S$.  Since $\mathbf S\subset\M$ and $\M$ is strongly closed,
this implies $A\in\M$.

For every $\nu=1,2,\ldots$ choose an $A_\nu\in\mathbf S$ with
$\lVert(A_\nu-A)g_i\rVert\leq1/\nu$ for $i=1,\ldots,m$.  This can be done
since, by the definition of a limit point, we must have at least one
$A_\nu\in\M$ in the strong neighborhood
$U_3(A,g_1,\ldots,g_m;1/\nu)$ of $A$ (cf. \cite[\S2, p.~381]{ref2}).  Then
\eqref{eq:1.2.gamma} gives
$\Metric{A_\nu-A}\leq\sqrt{m}/\nu$ and hence
$\lim_{\nu\to\infty}\Metric{A_\nu-A}=0$.  Thus $A_1,A_2,\ldots$ converges
metrically to $A$ and since $\mathbf S$ is metrically closed, $A\in\mathbf S$.
This completes the proof.
\end{proof}

Now if we can show that the restricted metric closure implies metric closure,
the above lemmas will yield the equivalence of the strong, the metric and the
restricted metric notions of closure.  As pointed out at the end of
\secref{sec:1.3}{1.3}, this is what we need.  The assumption that $\mathbf S$ is
an algebra will be needed to secure the above implication.

\PaperSection{sec:1.5}{Continuation of the discussion of closure.}
In what follows we shall make repeated use of the notion and properties of the
functions of bounded Hermitian and of unitary operators,\footnote{Cf., e.g.,
\cite[pp.~205 and 220]{ref3}.  The real function-theory methods of \cite{ref3}
are actually much more than what is needed for the limited objectives of
\secref{sec:1.5}{1.5}, as only continuous functions (sometimes even polynomials)
of operators occur.  ((i) in \lemref{lem:1.5.1}{1.5.1} can even be proven by
elementary evaluations.)  But it would take up too much space to go into such
methodological questions systematically.} and it will be convenient to omit
specific references to that theory.

\begin{lemma}\phantomsection\label{lem:1.5.1}
\begin{enumerate}[label=(\roman*)]
\item If $A$ is bounded Hermitian, then $(A-i1)^{-1}$ can be formed and
\[
  \frac{A+i1}{A-i1}=(A-i1)^{-1}(A+i1)=(A+i1)(A-i1)^{-1}
\]
is unitary.
\item If $A,B\in\M$ and Hermitian, then
$\dfrac{A+i1}{A-i1},\dfrac{B+i1}{B-i1}\in\M$ too, and
\[
  \Metric{\frac{A+i1}{A-i1}-\frac{B+i1}{B-i1}}
  \leq2\Metric{A-B}.
\]
\end{enumerate}
\end{lemma}
\begin{proof}
Ad (i).  Consider the two functions
\[
  f(\lambda)=\frac{1}{\lambda-i},\qquad
  g(\lambda)=\frac{\lambda+i}{\lambda-i}.
\]
Since they are continuous for all real $\lambda$ we can form $f(A)$ and $g(A)$.
We also have
\[
  (\lambda-i)f(\lambda)=f(\lambda)(\lambda-i)=1,
  \qquad
  (\lambda+i)f(\lambda)=f(\lambda)(\lambda+i)=g(\lambda).
\]
Therefore $f(A)=(A-i1)^{-1}$ and
$g(A)=(A+i1)f(A)=f(A)(A+i1)$.  Thus
$g(A)=(A+i1)(A-i1)^{-1}=(A-i1)^{-1}(A+i1)$ or
$g(A)=(A+i1)/(A-i1)$ in the sense indicated above.

Since $\overline{g(\lambda)}g(\lambda)=1=g(\lambda)\overline{g(\lambda)}$ for
all real $\lambda$, $g(A)^*g(A)=1=g(A)g(A)^*$ and hence $g(A)$ is unitary.
Thus the proof of (i) is completed.  We note also that inasmuch as
$|f(\lambda)|\leq1$ for all real $\lambda$, $\OpNorm{f(A)}\leq1$ or
\begin{equation}
  \OpNorm{(A-i1)^{-1}}\leq1.
  \tag{1.5.\ensuremath{\alpha}}\label{eq:1.5.alpha}
\end{equation}

Ad (ii).  $A,B\in\M$ imply $g(A),g(B)\in\M$ or
$(A+i1)/(A-i1),(B+i1)/(B-i1)\in\M$.

We have
\begin{align*}
  \frac{A+i1}{A-i1}-\frac{B+i1}{B-i1}
  &=(A-i1)^{-1}(A+i1)-(B+i1)(B-i1)^{-1}\\
  &=(A-i1)^{-1}\{(A+i1)(B-i1)-(A-i1)(B+i1)\}(B-i1)^{-1}\\
  &=-2i(A-i1)^{-1}(A-B)(B-i1)^{-1}.
\end{align*}
Consequently
\[
  \Metric{\frac{A+i1}{A-i1}-\frac{B+i1}{B-i1}}
  \leq2\OpNorm{(A-i1)^{-1}}\Metric{A-B}\OpNorm{(B-i1)^{-1}},
\]
and by \eqref{eq:1.5.alpha}
\[
  \Metric{\frac{A+i1}{A-i1}-\frac{B+i1}{B-i1}}
  \leq2\Metric{A-B}.
\]
\end{proof}

\begin{lemma}\phantomsection\label{lem:1.5.2}
Let a function $\varphi(\xi)$ be given, defined and continuous for all $\xi$ with
$|\xi|=1$.  Then there exists a function $\omega_1(\epsilon)$ (given
$\varphi$, $\omega_1$ is determined) defined and $>0$ for all $\epsilon>0$,
with the following property: If $U,V$ are unitary and $\in\M$ then
$\varphi(U),\varphi(V)\in\M$ and
$\Metric{U-V}<\omega_1(\epsilon)$ implies
$\Metric{\varphi(U)-\varphi(V)}<\epsilon$.
\end{lemma}
\begin{proof}
That $U,V$ unitary and $\in\M$ imply $\varphi(U),\varphi(V)\in\M$ is obvious.
We shall now show that it suffices to prove our assertion for all trigonometrical
polynomials
\[
  \varphi(\xi)=\sum_{\nu=-n}^{n}a_\nu\xi^\nu.\footnote{Since $|\xi|=1$ we
  may put $\xi=e^{i\alpha}$, $\alpha$ real, modulo $2\pi$.  Now
  $\varphi(\xi)=\sum_{\nu=-n}^{n}a_\nu\xi^\nu
  =\sum_{\nu=-n}^{n}a_\nu e^{i\nu\alpha}$.  Thus these $\varphi(\xi)$ are the
  trigonometrical polynomials with their familiar approximation properties.
  Cf. Zygmund, Antoni, ``Trigonometrical Series,'' Warsaw (1935), \S3.23(iii),
  p.~47.}
\]
Assume accordingly that it is true for these and consider a general
$\varphi(\xi)$.

Consider an $\epsilon'>0$.  Choose a trigonometrical polynomial
\[
  \varphi^{\epsilon'}(\xi)
  =\sum_{\nu=-n^{\epsilon'}}^{n^{\epsilon'}}
    a_\nu^{\epsilon'}\xi^\nu
\]
such that $|\varphi(\xi)-\varphi^{\epsilon'}(\xi)|\leq\frac13\epsilon'$ for all
$\xi$ with $|\xi|=1$.  Then the unitarity of $U,V$ implies
$\OpNorm{\varphi(U)-\varphi^{\epsilon'}(U)}\leq\frac13\epsilon'$ and
$\OpNorm{\varphi(V)-\varphi^{\epsilon'}(V)}\leq\frac13\epsilon'$.  Hence,
\emph{a fortiori},
\[
  \Metric{\varphi(U)-\varphi^{\epsilon'}(U)}\leq\tfrac13\epsilon',
  \qquad
  \Metric{\varphi(V)-\varphi^{\epsilon'}(V)}\leq\tfrac13\epsilon'.
\]
Denote the $\omega_1(\epsilon)$ of $\varphi^{\epsilon'}(\xi)$ by
$\omega_1^{\epsilon'}(\epsilon)$.  Then
$\Metric{U-V}<\omega_1^{\epsilon'}(\frac13\epsilon')$ implies
\[
  \Metric{\varphi^{\epsilon'}(U)-\varphi^{\epsilon'}(V)}<\tfrac13\epsilon'
\]
and so, owing to the above,
\[
  \Metric{\varphi(U)-\varphi(V)}<\epsilon'.
\]
Thus
$\omega_1(\epsilon')=\omega_1^{\epsilon'}(\frac13\epsilon')$ meets our
requirements for the original $\varphi(\xi)$.

It only remains therefore to prove our assertion for the trigonometrical
polynomials
\[
  \varphi(\xi)=\sum_{\nu=-n}^{n}a_\nu\xi^\nu.
\]
Obviously we may even restrict ourselves to the monomials
\[
  \varphi(\xi)=\xi^\nu,\qquad(\nu=0,\pm1,\pm2,\ldots).
\]
For $\nu=0$ this is trivial, and for $\nu<0$ the unitarity of $U,V$ implies
\[
  \Metric{U^\nu-V^\nu}
  =\Metric{(U^\nu-V^\nu)^*}
  =\Metric{U^{*\nu}-V^{*\nu}}
  =\Metric{U^{-\nu}-V^{-\nu}}.
\]
Therefore we may even assume $\nu>0$.  In this case, however,
\begin{align*}
  U^\nu-V^\nu
  &=\sum_{\mu=1}^{\nu}
    \bigl(U^\mu V^{\nu-\mu}-U^{\mu-1}V^{\nu-\mu+1}\bigr)\\
  &=\sum_{\mu=1}^{\nu}U^{\mu-1}(U-V)V^{\nu-\mu},
\end{align*}
and
\begin{align*}
  \Metric{U^\nu-V^\nu}
  &=\Metric{\sum_{\mu=1}^{\nu}U^{\mu-1}(U-V)V^{\nu-\mu}}\\
  &\leq\sum_{\mu=1}^{\nu}
  \OpNorm{U^{\mu-1}}\Metric{U-V}\OpNorm{V^{\nu-\mu}}\\
  &=\sum_{\mu=1}^{\nu}\Metric{U-V}
   =\nu\Metric{U-V}.
\end{align*}
Thus $\omega_1(\epsilon)=(1/\nu)\epsilon$ meets our requirements for these
$\varphi(\xi)$.
\end{proof}

\begin{lemma}\phantomsection\label{lem:1.5.3}
Let a function $\psi(\lambda)$ be given, defined and continuous for all real
$\lambda$ and for which $\lim_{\lambda\to\pm\infty}\psi(\lambda)$
exists.\footnote{We require
$\lim_{\lambda\to\infty}\psi(\lambda)=
\lim_{\lambda\to-\infty}\psi(\lambda)$.  This condition could be removed, but
it simplifies the proof and suffices for our present purposes.}  Then there
exists a function $\omega_2(\epsilon)$ (given $\psi$, $\omega_2$ is determined)
defined and $>0$ for all $\epsilon>0$ with the following property: If $A,B$ are
Hermitian and $\in\M$ then $\psi(A),\psi(B)\in\M$ and
$\Metric{A-B}<\omega_2(\epsilon)$ implies
$\Metric{\psi(A)-\psi(B)}<\epsilon$.
\end{lemma}
\begin{proof}
That $A,B$ Hermitian and $\in\M$ imply $\psi(A),\psi(B)\in\M$ is obvious.
The correspondence
\[
  \xi=\frac{\lambda+i}{\lambda-i},
  \qquad\text{or equivalently}\qquad
  \lambda=i\frac{\xi+1}{\xi-1}
\]
is a one-to-one and bicontinuous mapping of $\xi$ with $|\xi|=1$ on all real
$\lambda$ and $\lambda=\pm\infty$ where $\xi=1$ corresponds to
$\lambda=\pm\infty$.  Hence
\[
  \varphi(\xi)=
  \begin{cases}
    \displaystyle\psi\!\left(i\frac{\xi+1}{\xi-1}\right),
      & |\xi|=1,\ \xi\neq1,\\[6pt]
    \displaystyle\lim_{\lambda\to\pm\infty}\psi(\lambda),
      & \xi=1,
  \end{cases}
\]
is defined and continuous for all $\xi$ with $|\xi|=1$.  Clearly
\[
  \psi(\lambda)=\varphi\!\left(\frac{\lambda+i}{\lambda-i}\right).
\]
We recall the $g(\lambda)=(\lambda+i)/(\lambda-i)$ used in the proof of
\lemref{lem:1.5.1}{1.5.1} and let
\[
  U=g(A)=\frac{A+i1}{A-i1},\qquad
  V=g(B)=\frac{B+i1}{B-i1}.
\]
Then the above relations imply
\[
  \psi(A)=\varphi(U),\qquad \psi(B)=\varphi(V).
\]
By \lemref{lem:1.5.1}{1.5.1}, $U$ and $V$ are unitary and $\in\M$ and
\[
  \Metric{U-V}\leq2\Metric{A-B}.
\]
Now \lemref{lem:1.5.2}{1.5.2} applies to $\varphi(\xi)$.  Consider the
resulting $\omega_1(\epsilon)$.  Then
$\Metric{A-B}<\frac12\omega_1(\epsilon)$ implies
$\Metric{U-V}<\omega_1(\epsilon)$ and hence
$\Metric{\varphi(U)-\varphi(V)}<\epsilon$ or
$\Metric{\psi(A)-\psi(B)}<\epsilon$.  Therefore
$\omega_2(\epsilon)=\frac12\omega_1(\epsilon)$ meets our requirements.
\end{proof}

\begin{lemma}\phantomsection\label{lem:1.5.4}
Let an algebra $\mathbf S$ in $\M$ be given and a sequence
$A_1,A_2,\ldots\in\mathbf S$ which converges metrically to an $A\in\M$.
Then there exists another sequence $A'_1,A'_2,\ldots\in\mathbf S$ which
converges restrictedly metrically to $A$.
\end{lemma}
\begin{proof}
$A_1,A_2,\ldots$ converges metrically to $A$; hence
$A_1^*,A_2^*,\ldots$ converges metrically to $A^*$.\footnote{Since
$\Metric{A}=\Metric{A^*}$, $A\mapsto A^*$ is an isometric mapping.}
Consequently
$\frac12(A_1+A_1^*),\frac12(A_2+A_2^*),\ldots$ converges metrically to
$\frac12(A+A^*)$, and
$\frac1{2i}(A_1-A_1^*),\frac1{2i}(A_2-A_2^*),\ldots$ converges metrically
to $\frac1{2i}(A-A^*)$.  The operators occurring in the two last assertions
are all Hermitian and $\in\mathbf S$.  Also
$A=\frac12(A+A^*)+i\frac1{2i}(A-A^*)$.

Now suppose our result has been established in the special case in which all
the operators $A_1,A_2,\ldots$ are Hermitian.  Then to treat the general case,
we should apply the special result to the above-mentioned Hermitian ``real and
imaginary parts'' and by taking the linear combination we should obtain the
desired result since the notion of ``restricted metrical convergence'' holds for
linear combinations of series if it holds for the series themselves.  Thus we
may assume in what follows that $A,A_1,A_2,\ldots$ are Hermitian operators.

Now form the function
\[
  \psi(\lambda)=
  \begin{cases}
    \lambda, & |\lambda|\leq\OpNorm{A},\\[2pt]
    \displaystyle\frac{\OpNorm{A}^{,2}}{\lambda},
      & |\lambda|\geq\OpNorm{A}.
  \end{cases}
\]
\lemref{lem:1.5.3}{1.5.3} applies to $\psi(\lambda)$ yielding
$\omega_2(\epsilon)$.

Consider a $\nu=1,2,\ldots$.  Choose a $\nu'=1,2,\ldots$ (depending on
$\nu$ as well as on $A,A_1,A_2,\ldots$) with
\[
  \Metric{A_{\nu'}-A}<\omega_2\!\left(\frac1\nu\right).
\]
Then \lemref{lem:1.5.3}{1.5.3} gives
$\Metric{\psi(A_{\nu'})-\psi(A)}<1/\nu$.  Since
$\psi(\lambda)=\lambda$ for all $\lambda$ with
$|\lambda|\leq\OpNorm{A}$, $\psi(A)=A$.\footnote{The set
$|\lambda|\leq\OpNorm{A}$ contains the entire spectrum of $A$.}  Thus we have
\begin{equation}
  \Metric{\psi(A_{\nu'})-A}<\frac1\nu.
  \tag{1.5.\ensuremath{\beta}}\label{eq:1.5.beta}
\end{equation}

Form next the function
\[
  \vartheta(\lambda)=
  \begin{cases}
    1, & |\lambda|\leq\OpNorm{A},\\[2pt]
    \displaystyle\frac{\OpNorm{A}^{,2}}{\lambda^2},
      & |\lambda|\geq\OpNorm{A}.
  \end{cases}
\]
It is clear that $\psi(\lambda)=\lambda\vartheta(\lambda)$.  Choose a
polynomial $p_\nu(\lambda)$ with
\[
  \vartheta(\lambda)\geq p_\nu(\lambda)
  \geq\operatorname{Max}\!\left(
    \vartheta(\lambda)-\frac1{\nu|\lambda|},0\right)
  \quad\text{for }|\lambda|\leq\OpNorm{A_{\nu'}}.\footnote{This is an easy
  variant of the Weierstrass (polynomial) approximation theorem.}
\]
Then the polynomial $q_\nu(\lambda)=\lambda p_\nu(\lambda)$ has no constant
term and we have
\[
  |q_\nu(\lambda)-\psi(\lambda)|\leq\frac1\nu
  \quad\text{for }|\lambda|\leq\OpNorm{A_{\nu'}}.
\]
Furthermore $0\leq q_\nu(\lambda)\leq\psi(\lambda)\leq\OpNorm{A}$ and thus
\[
  |q_\nu(\lambda)|\leq\OpNorm{A}
  \quad\text{for }|\lambda|\leq\OpNorm{A_{\nu'}}.
\]
These two inequalities imply
\[
  \OpNorm{q_\nu(A_{\nu'})-\psi(A_{\nu'})}\leq\frac1\nu,\footnote{The set
  $|\lambda|\leq\OpNorm{A_{\nu'}}$ contains the entire spectrum of
  $A_{\nu'}$.}
\]
and hence, \emph{a fortiori},
\begin{equation}
  \Metric{q_\nu(A_{\nu'})-\psi(A_{\nu'})}\leq\frac1\nu,
  \tag{1.5.\ensuremath{\gamma}}\label{eq:1.5.gamma}
\end{equation}
and
\begin{equation}
  \OpNorm{q_\nu(A_{\nu'})}\leq\OpNorm{A}.
  \tag{1.5.\ensuremath{\delta}}\label{eq:1.5.delta}
\end{equation}
Combining \eqref{eq:1.5.beta} and \eqref{eq:1.5.gamma} gives
\begin{equation}
  \Metric{q_\nu(A_{\nu'})-A}<\frac2\nu.
  \tag{1.5.\ensuremath{\epsilon}}\label{eq:1.5.epsilon}
\end{equation}
Now put $A'_\nu=q_\nu(A_{\nu'})$.  Then $A_{\nu'}\in\mathbf S$ implies
$q_\nu(A_{\nu'})\in\mathbf S$ since $q_\nu(\lambda)$ has no constant term.
Hence the sequence $A'_1,A'_2,\ldots$ consists of elements of $\mathbf S$, and
it is restrictedly metrically convergent to $A$ by \eqref{eq:1.5.delta} and
\eqref{eq:1.5.epsilon}.

Observe that in the above proof we even had
$\OpNorm{A'_\nu}\leq\OpNorm{A}$ when $A$ is Hermitian.  This could be
extended to all $A$ but we shall not pursue this matter further.
\end{proof}

\PaperSection{sec:1.6}{Conclusion of the discussion of closure.}
The desired results are now immediate.
\begin{lemma}\phantomsection\label{lem:1.6.1}
Restricted metric closure of an algebra $\mathbf S$ implies its metric closure.
\end{lemma}
\begin{proof}
Immediate by \lemref{lem:1.5.4}{1.5.4}.
\end{proof}

\begin{theorem}\phantomsection\label{thm:I}
For an algebra $\mathbf S$ within a finite $\M$ the following eight notions of
closure are equivalent to each other:
\begin{description}
  \item[$\alpha)$] Restricted metric closure
  \item[$\beta)$] Metric closure
  \item[$\gamma)$] Weak (topological) closure
  \item[$\delta)$] Weak convergence closure
  \item[$\epsilon)$] Strong (topological) closure
  \item[$\vartheta)$] Strong convergence closure
  \item[$\zeta)$] Strongest (topological) closure
  \item[$\eta)$] Strongest convergence closure.
\end{description}
\end{theorem}
\begin{proof}
The implications $\epsilon)\Rightarrow\alpha)$,
$\beta)\Rightarrow\epsilon)$, $\alpha)\Rightarrow\beta)$ were established in
\lemref{lem:1.4.1}{1.4.1}, \lemref{lem:1.4.2}{1.4.2},
\lemref{lem:1.6.1}{1.6.1}, respectively.  Consequently
$\alpha)\Leftrightarrow\beta)\Leftrightarrow\epsilon)$.  Also
\lemref{lem:1.3.2}{1.3.2} yields $\alpha)\Leftrightarrow\vartheta)$.  From
\cite[p.~396, end of \S3]{ref2}, we have
$\gamma)\Leftrightarrow\epsilon)$.  \cite[p.~112, beginning of \S3]{ref4}
yields $\gamma)\Leftrightarrow\zeta)$.  Also
\cite[p.~112, end of \S2]{ref4} implies
$\vartheta)\Leftrightarrow\eta)$.  These equivalences give together
\[
  \alpha)\Leftrightarrow\beta)\Leftrightarrow\gamma)
  \Leftrightarrow\epsilon)\Leftrightarrow\vartheta)
  \Leftrightarrow\zeta)\Leftrightarrow\eta).
\]
And obviously
\[
  \gamma)\Rightarrow\delta)\Rightarrow\vartheta).
\]
Thus $\delta)$ also is equivalent to the others.
\end{proof}

Thus all those topologizations of the space of all operators, which may be
reasonably used in this connection, turn out to be equivalent to each other for
the algebras $\mathbf S$ in a factor $\M$ which is in a finite case.

We also state explicitly
\begin{theorem}\phantomsection\label{thm:II}
The notion of a subring $\mathbf S$ of $\M$, for a finite $\M$, is purely
algebraical.
\end{theorem}
\begin{proof}
The equivalences of \thmref{thm:I}{I} now permit us to apply the argument given
at the end of \secref{sec:1.2}{1.2}.
\end{proof}
Thus the desideratum expressed at the end of \secref{sec:1.1}{1.1} is fully
realized.

\PaperChapter{2}{chapter:II}{The Algebraical Type and its Operations. The Fundamental Group}

\PaperSection{sec:2.1}{Algebraical type $\overline{\M}$. Discussion.}
We shall call the abstraction of a ring $\M$ with respect to algebraical
isomorphism, its algebraical type, i.e. two rings $\M_1$ and $\M_2$ will be said
to have the same algebraical type if and only if they are algebraically
isomorphic.  We denote the algebraic type of the ring $\M$ by
$\overline{\M}$.  Notions in $\M$ which are purely algebraical, i.e. invariant
under algebraical isomorphisms, will be said to be notions concerning the
algebraical type $\overline{\M}$ (and not $\M$ itself).\footnote{This definition
should be compared and contrasted with that of spatial type in
\secref{sec:3.3}{3.3}.}

Thus the properties enumerated in \secref{sec:1.1}{1.1} (or equally in
\cite{ref9}) are properties of $\overline{\M}$ for every ring $\M$.  The notion
of a subring of $\M$ concerns $\overline{\M}$ when $\M$ is a factor in a finite
class (cf. \thmref{thm:II}{II}).  This is also true of $D_{\M}(E)$,
$\operatorname{Tr}_{\M}(A)$, $\Metric{A}$ (cf. the beginning of
\secref{sec:1.2}{1.2}), and the various notions of closure, as enumerated in
\thmref{thm:I}{I}, when they concern an algebra $\mathbf S$.  Thus when $\M$ is
in a finite class, the properties concerning $\overline{\M}$ are particularly
extensive.

\PaperSection{sec:2.2}{Conjugate dual isomorphisms. $\overline{\M}^{c}$ and
$\overline{\M}'$.}
As in the Introduction, \secref{intro:2}{2}, let for each $i=1,2$ a Hilbert
space $\Hilbert_i$ and an operator ring $\M_i$ in $\Hilbert_i$ be given.  We
generalize B), loc. cit., as follows:

\noindent\textbf{B*) General (algebraic ring) isomorphism of $\M_1$ and
$\M_2$.} This is a one-to-one mapping of $\M_1$ on $\M_2$ which leaves the
entity $1$ and the operations $A^*$, $A+B$ invariant; while it carries the
entity $\alpha A$ ($\alpha$ any complex number) into either $\alpha A$ always
(case $a'$) or into $\overline\alpha A$ always (case $b'$) and $AB$ into $AB$
always (case $a''$) or into $BA$ always (case $b''$).  More specifically, we
talk according to which combination of the above cases occurs, of an
$[a',a'']$ (or ``proper''), $[b',a'']$ (or ``conjugate''), $[a',b'']$ (or
``dual'') or $[b',b'']$ (or ``conjugate dual'') isomorphism.

The existence of a general isomorphism, of any one of the above four kinds,
between $\M_1$ and $\M_2$ is clearly a property of the types
$\overline{\M}_1$ and $\overline{\M}_2$.  If a (proper) isomorphism holds, we
have $\overline{\M}_1=\overline{\M}_2$.  If a conjugate or a dual isomorphism
holds, then $\overline{\M}_1$ determines $\overline{\M}_2$ uniquely, assuming
its existence.  Hence we may express these relationships by
$\overline{\M}_1^{c}=\overline{\M}_2$ or
$\overline{\M}_1'=\overline{\M}_2$, respectively.

We must now deduce the existential and other properties of the operations
$\overline{\M}^{c}$ and $\overline{\M}'$.

\PaperSection{sec:2.3}{Properties of the above. $\overline{\M}'=\overline{\M}^{c}$;
$\overline{\M}^{cc}=\overline{\M}$.}
Let a Hilbert space $\Hilbert$ be given.  Modify the fundamental operations
$\alpha f$, $f+g$, $(f,g)$ in $\Hilbert$ by replacing them by
$\overline\alpha f$, $f+g$, $\overline{(f,g)}$.  Denote the set $\Hilbert$ with
the new definition of its fundamental operations by $\Hilbert_c$.  Clearly
$\Hilbert_c$ is also a Hilbert space; every (linear bounded) operator $A$ in
$\Hilbert$ is also one in $\Hilbert_c$, and every ring of operators $\M$ in
$\Hilbert$ is also one in $\Hilbert_c$.  But we shall denote the $A,\M$ of
$\Hilbert$, when considered in $\Hilbert_c$, by $A_c,\M_c$.

Now consider an operator ring $\M$ in $\Hilbert$.  The identical mapping
$\mathfrak I^{\circ}$ then maps $\M$ in $\Hilbert$ on $\M_c$ in $\Hilbert_c$
and it is in this aspect a conjugate isomorphism of $\M$ and $\M_c$.

Consider next the mapping $\mathfrak I^*$, $A\mapsto A^*$.  This maps $\M$ in
$\Hilbert$ on $\M$ in $\Hilbert$ and it is, in this aspect, a conjugate dual
isomorphism of $\M$ and $\M$.  Consequently
$\mathfrak I^*\mathfrak I^{\circ}$ (i.e. $\mathfrak I^*$ in a new aspect) is a
dual isomorphism of $\M$ and $\M_c$.

Thus we have
\begin{theorem}\phantomsection\label{thm:III}
$\overline{\M}^{c}$ and $\overline{\M}'$ always exist and are single-valued.
They are both equal to the above $\overline{\M}_c$.  We shall denote them by
$\overline{\M}^{c}$.
\end{theorem}
Clearly
\begin{equation}
  \overline{\M}^{cc}=\overline{\M}
  \tag{2.3.\ensuremath{\alpha}}\label{eq:2.3.alpha}
\end{equation}
since $\Hilbert_{c,c}=\Hilbert$, $A_{c,c}=A$, $\M_{c,c}=\M$.

Observe that the definitions of a factor and of the case to which it belongs
are unaffected by conjugate isomorphisms.  For these definitions make use of
$\alpha A$ only in defining the set of all $\alpha1$ but never for a specific
numerical value of $\alpha$.  To see this, reconsider
\cite[p.~138, Definition~3.1.2; p.~173, Theorem~IX]{ref5}.  Hence the operation
$\overline{\M}^{c}$ (viewed as $\overline{\M}'$) does not affect the factor
character of $\M$ nor the case to which it belongs.


We now are able to prove:


\PaperSection{sec:2.4}{$p$-matrix isomorphisms. Construction of
$\overline{\M}^{p}$.}
In what follows we shall consider matrices of operators.  Given a
$p=1,2,\ldots,\infty$, a $p$-order matrix $\langle A_{t,s}\rangle$ is a scheme
or array of operators $A_{t,s}$ in which the indices $s$ and $t$ range from
$1$ to $p$ for a finite $p$ and over $1,2,\ldots$ if $p=\infty$.  (This last
will be abbreviated $s,t\leq p$ in both cases.)  An
$\langle A_{t,s}\rangle$ is finite if there exists a finite
$q=1,2,\ldots$ such that $A_{t,s}=0$ if either $t>q$ or $s>q$.  For a finite
$p$ this is automatically fulfilled, but it is a real restriction for
$p=\infty$.

Now as in the Introduction, \secref{intro:2}{2}, let, for each $i=1,2$, a
Hilbert space $\Hilbert_i$ and an operator ring $\M_i$ in $\Hilbert_i$ be given.
We generalize B), loc. cit., as follows:

\noindent\textbf{$B_p)$ $p$-matrix (algebraic ring) isomorphism of $\M_1$ over
$\M_2$.} This is a one-to-one mapping of $\M_1$ on a set $\mathfrak J$ of
$p$-order matrices $\langle A_{t,s}\rangle$ over $\M_2$ (i.e.
$A_{t,s}\in\M_2$) which is algebraically ring-isomorphic in the matrix sense,
i.e. it carries $1$, $\alpha A^{\circ}$, $A^{\circ *}$,
$A^{\circ}+B^{\circ}$, $A^{\circ}B^{\circ}$ in $\M_1$ into
\[
  \langle\delta_{t,s}1\rangle,\quad
  \langle\alpha A_{t,s}\rangle,\quad
  \langle A_{s,t}^*\rangle,\quad
  \langle A_{t,s}+B_{t,s}\rangle,\quad
  \left\langle\sum_{r=1}^{p}A_{r,s}B_{t,r}\right\rangle
\]
in $\mathfrak J$.  The $\sum_{r=1}^{\infty}$ in the above expression, when
$p=\infty$, must converge in the sense of the strong operator topology,
irrespective of the order of summation.

The set $\mathfrak J$ must, at any rate, contain all finite matrices
$\langle A_{t,s}\rangle$ over $\M_2$.  (Cf. \cite[pp.~136--137,
Lemma~2.4.3]{ref5}, where the same notions are used, the same multiplication
convention obtains, as well as the same notion of convergence for
$\sum_{r=1}^{\infty}$.)

For a given $p,\M_1,\M_2$ the existence of such a $p$-matrix isomorphism of
$\M_1$ over $\M_2$ is a property of types $\overline{\M}_1,\overline{\M}_2$.
This is obvious for $p=1,2,\ldots$ while for $p=\infty$ it is only necessary to
remember the purely algebraical character of the notion of strong convergence
(cf. the beginning of \secref{sec:1.1}{1.1}, or \cite[Theorem~II]{ref9}).
Whether $\overline{\M}_1$ (assuming its existence) is uniquely determined by
$\overline{\M}_2$ or not, for a given $p=1,2,\ldots,\infty$, depends obviously
on whether the set $\mathfrak J$ in $B_p)$ is uniquely determined or not.  For
$p=1,2,\ldots$ ($p$ finite), the last part of $B_p)$ implies that the set
$\mathfrak J$ must be the set of all $\langle A_{t,s}\rangle$ with
$A_{t,s}\in\M_2$ since they are all finite.  Hence in this case
$\overline{\M}_1$ is uniquely determined.  For $p=\infty$ the question of
uniqueness is not so immediately settled.  We shall prove that $\mathfrak J$
and with it $\overline{\M}_1$ are again unique, but we delay this discussion to
the next section.  Nevertheless, for $p=1,2,\ldots,\infty$ we express the
existence of a $p$-order matrix isomorphism of $\M_1$ over $\M_2$ as a
relationship between $\overline{\M}_1$ and $\overline{\M}_2$ by
$\overline{\M}_1=\overline{\M}_2^{p}$.  But then, quite apart from the question
of existence, we may assume that $\overline{\M}^{p}$ is single-valued only for
a finite $p$.  At present we must consider the possibility that for
$p=\infty$, $\overline{\M}^{p}$ is many valued.

The general existence of $\overline{\M}^{p}$ belongs in abstract algebra, but
we shall need $\overline{\M}^{p}$ as an operator ring.  For this reason, as well
as for the sake of general orientation, we shall construct it explicitly within
the framework of \cite[pp.~135--137, \S2.4]{ref5}.

Let a Hilbert space $\Hilbert$ and a ring $\M$ in it be given, and a
$p=1,2,\ldots,\infty$.  We perform the constructions of \cite[loc. cit.]{ref5},
the $\Hilbert_2$ there being our $\Hilbert$ and the $\Hilbert_1$ there, any
$p$-dimensional unitary space.  (This is not the $\Hilbert_1$ of $B_p)$ above.
We conform to the notations of \cite[loc. cit.]{ref5}.)  Form
$\Hilbert_1\otimes\Hilbert_2$ and the ring $\mathbf B^{\otimes}$ of all its
bounded operators.  Then all elements of $\mathbf B^{\otimes}$ can be
represented as $p$-order matrices $\langle A_{t,s}\rangle$
($t,s\leq p$, $A_{t,s}$ in $\Hilbert_2=\Hilbert$).  (Cf.
\cite[p.~136, Definition~2.4.2 and Lemma~2.4.2]{ref5}.)  For the given ring
$\M$, form the set $\M_1$ consisting of all operators in
$\Hilbert_1\otimes\Hilbert_2$ which have a matrix $\langle A_{t,s}\rangle$
with all $A_{t,s}\in\M$.  $\mathfrak J$ is the set of all those
$\langle A_{t,s}\rangle$ with $A_{t,s}\in\M$ which correspond to some
operator in $\Hilbert_1\otimes\Hilbert_2$.  From this it follows immediately
that $\mathfrak J$ contains all finite matrices $\langle A_{t,s}\rangle$ with
$A_{t,s}\in\M$.\footnote{For a finite matrix $\langle A_{t,s}\rangle$ an
operator $A^{\circ}$ in $\Hilbert_1\otimes\Hilbert_2$ with
$A^{\circ}\sim\langle A_{t,s}\rangle$ can be constructed in the sense of
\cite[p.~136, Definition~2.4.2]{ref5}, without any convergence difficulties.}
Now $\M_1$ in $\Hilbert_1\otimes\Hilbert_2$ and $\M$ in
$\Hilbert=\Hilbert_2$ together with the interpretation of the matrix scheme
$\langle A_{t,s}\rangle$ which we are now using, clearly fulfill all the
requirements of $B_p)$.  (We have our
$\Hilbert_1\otimes\Hilbert_2$, $\Hilbert=\Hilbert_2$, $\M_1$, $\M$ in place
of the $\Hilbert_1$, $\Hilbert_2$, $\M_1$, $\M_2$ of $B_p)$.)

Summing up the results obtained so far, we have
\begin{theorem}\phantomsection\label{thm:IV}
$\overline{\M}^{p}$ exists for all $p=1,2,\ldots,\infty$ and it is certainly
one-valued for $p=1,2,\ldots$ ($p$ finite).
\end{theorem}
Our definitions in $B_p)$ and that of \secref{sec:2.3}{2.3} give together
\begin{equation}
  (\overline{\M}^{c})^{p}=(\overline{\M}^{p})^{c}
  \qquad(p=1,2,\ldots),
  \tag{2.4.\ensuremath{\alpha}}\label{eq:2.4.alpha}
\end{equation}
and $B_p)$ gives
\begin{equation}
  (\overline{\M}^{p})^{q}=\overline{\M}^{pq}
  \qquad(p,q=1,2,\ldots).
  \tag{2.4.\ensuremath{\beta}}\label{eq:2.4.beta}
\end{equation}
(The exponent $\infty$ will be considered in the next section.)

\PaperSection{sec:2.5}{One-valuedness of $\overline{\M}^{\infty}$.}
In this section we shall make use of the operations $A_{(\mathfrak M)}$,
$\M_{(\mathfrak M)}$ defined in \cite[p.~78]{ref1}, and
\cite[p.~186, Definition~11.3.1]{ref5}.  If $A$ is an operator in $\Hilbert$
which is reduced by the closed linear set $\mathfrak M$ then
$A_{(\mathfrak M)}$ is its part in $\mathfrak M$.  And $\M_{(\mathfrak M)}$ is
the set of all $A_{(\mathfrak M)}$ for all $A\in\M$ which are reduced by
$\mathfrak M$.  $A_{(\mathfrak M)},\M_{(\mathfrak M)}$ are in $\mathfrak M$
while $A,\M$ were in $\Hilbert$.

\begin{lemma}\phantomsection\label{lem:2.5.1}
Assume $E=P_{\mathfrak M}\in\M$.  Let $\M^{E}$ be the set of all $A\in\M$ for
which $EA=AE=A$.  Then we have
\begin{enumerate}[label=(\roman*)]
\item $A_{(\mathfrak M)}$ can be formed for all $A\in\M^{E}$.
\item $\M_{(\mathfrak M)}$ coincides with $(\M^{E})_{(\mathfrak M)}$.
\item $A\mapsto A_{(\mathfrak M)}$ is a one-to-one mapping of $\M^{E}$ on
$\M_{(\mathfrak M)}$ which carries $E$ (in $\M^{E}$, i.e. in $\Hilbert$) into
$1$ (in $\M_{(\mathfrak M)}$, i.e. in $\mathfrak M$) and leaves the operations
$\alpha A,A^*,A+B,AB$ invariant.
\end{enumerate}
\end{lemma}
\begin{proof}
Ad (i).  Every $A\in\M^{E}$ commutes with $E=P_{\mathfrak M}$, i.e. it is
reduced by $\mathfrak M$.

Ad (ii).  $A\in\M$ implies $EAE\in\M$ and thence clearly $EAE\in\M^{E}$.
Now if $\mathfrak M$ reduces $A$, then both $A$ and $EAE$ have the same part
in $\mathfrak M$, i.e. $A_{(\mathfrak M)}=(EAE)_{(\mathfrak M)}$.  Thus
$(\M^{E})_{(\mathfrak M)}\supseteq\M_{(\mathfrak M)}$.  But
$\M^{E}\subseteq\M$ hence also
$(\M^{E})_{(\mathfrak M)}\subseteq\M_{(\mathfrak M)}$.  So
$(\M^{E})_{(\mathfrak M)}=\M_{(\mathfrak M)}$.

Ad (iii).  Every $A\in\M^{E}$ is reduced by $\mathfrak M$ (cf. (i) above) and
its part in $\mathfrak M^{\perp}$ ($=\Hilbert\ominus\mathfrak M$) is clearly
$0$.  So we may view it as an operator in $\mathfrak M$ as well as in
$\Hilbert$, and in its former aspect it is $A_{(\mathfrak M)}$.  All our
assertions are now immediate.
\end{proof}

Consider now $\M_1,\M_2$ and $\Hilbert_1,\Hilbert_2$ and an $\infty$-matrix
isomorphism of $\M_1$ over $\M_2$ as in $B_\infty)$ in the last section.  For
any $q=1,2,\ldots$ form the set $\mathfrak J^q$ of all matrices
$\langle A_{t,s}\rangle$ for which all $A_{t,s}\in\M_2$ and where
$A_{t,s}=0$ when $t>q$ or $s>q$.  Form also the matrix
$\langle e^q_{t,s}\cdot1\rangle$ where $e^q_{t,s}=1$ if $t=s\leq q$ and
$e^q_{t,s}=0$ otherwise.  Then it is easy to establish that
\begin{lemma}\phantomsection\label{lem:2.5.2}
\begin{enumerate}[label=(\roman*)]
\item $\mathfrak J^q\subseteq\mathfrak J$.
\item $\langle e^q_{t,s}\cdot1\rangle\in\mathfrak J^q$ corresponds to a
projection $E^q\in\M_1$.  Put $E^q=P_{\mathfrak M^q}$.
\item $E^1\leq E^2\leq\cdots$ and $\lim_{q\to\infty}E^q=1$ in the strong
sense.
\item Every matrix $\langle A_{t,s}\rangle\in\mathfrak J^q$ can be viewed both
as a $q$th order matrix and as an $\infty$-order matrix.  It corresponds in its
former aspect to an element of a ring $\M_2^q$ with
$\overline{\M_2^q}=\overline{\M}_2^q$ and in its latter aspect to an element of
$(\M_1)^{E^q}\subseteq\M_1$.  The latter correspondence can be continued by
\lemref{lem:2.5.1}{2.5.1} to one with an element of
$(\M_1)_{(\mathfrak M^q)}$.  These elements exhaust the sets
$\M_2^q$, $(\M_1)^{E^q}$ and $(\M_1)_{(\mathfrak M^q)}$, respectively.
\item The correspondence of (iv) is an algebraic ring isomorphism of $\M_2^q$
and $(\M_1)_{(\mathfrak M^q)}$.  Hence
$\overline{(\M_1)_{(\mathfrak M^q)}}=\overline{\M_2^q}=\overline{\M}_2^q$.
\end{enumerate}
\end{lemma}
\begin{proof}
Ad (i).  Immediate by the final clause of $B_\infty)$.

Ad (ii).  Clearly $\langle e^q_{t,s}\cdot1\rangle\in\mathfrak J^q$ hence it is
$\in\mathfrak J$ and thus corresponds to an element of $\M_1$.  The rules of
matrix computation show that this element is a projection.

Ad (iii).  The rules of matrix computation show that
$E^1\leq E^2\leq\cdots$.  Now if
$f=\langle f_1,f_2,\ldots\rangle\in\Hilbert_1\otimes\Hilbert_2$ then
$E^qf=\langle f_1,\ldots,f_q,0,\ldots\rangle\to f$ as $q\to\infty$.
Hence $\lim_{q\to\infty}E^q$ is $1$.

Ad (iv).  The $\langle A_{t,s}\rangle\in\mathfrak J^q$ are clearly
characterized by
\[
  \langle e^q_{t,s}\cdot1\rangle\langle A_{t,s}\rangle
  =\langle A_{t,s}\rangle\langle e^q_{t,s}\cdot1\rangle
  =\langle A_{t,s}\rangle
\]
and thus we may add by (i), $\langle A_{t,s}\rangle\in\mathfrak J$.  Hence the
second correspondence mentioned maps the
$\langle A_{t,s}\rangle\in\mathfrak J^q$ precisely on the
$A^{\circ}\in\M_1$ with $E^qA^{\circ}=A^{\circ}E^q=A^{\circ}$, i.e. on the
$A^{\circ}\in(\M_1)^{E^q}$.  All other assertions are obvious, remembering in
particular \lemref{lem:2.5.1}{2.5.1}.

Ad (v).  The corresponding assertions for the relationship between $\M_2^q$
and $(\M_1)^{E^q}$ follow immediately from the rules of matrix computation.
They are transferred to $\M_2^q$ and $(\M_1)_{(\mathfrak M^q)}$ by using
\lemref{lem:2.5.1}{2.5.1}.
\end{proof}

We now go further:
\begin{lemma}\phantomsection\label{lem:2.5.3}
A matrix $\langle A_{t,s}\rangle\in\mathfrak J^q$ corresponds to three
operators in the sense of (iv) in \lemref{lem:2.5.2}{2.5.2}, i.e. to an operator
of $\M_2^q$, to one in $(\M_1)^{E^q}$ and to one in
$(\M_1)_{(\mathfrak M^q)}$.  All three operators have the same bound, and we
denote this common bound by $\OpNorm{\langle A_{t,s}\rangle}$.
\end{lemma}
\begin{proof}
The operator in $(\M_1)^{E^q}$ and that one in
$(\M_1)_{(\mathfrak M^q)}$ have the same bound, since we are considering the
same operator, once in $\Hilbert_1$ and once in $\mathfrak M^q$.\footnote{The
part of the first-mentioned operator in $(\mathfrak M^q)^{\perp}$ is $0$.}
The operator in $\M_2^q$ and that one in $(\M_1)_{(\mathfrak M^q)}$ have the
same bound since they obtain from one another by an algebraical ring isomorphism
(cf. (v) in \lemref{lem:2.5.2}{2.5.2}) and since the numerical value of the
bound is a purely algebraical notion (cf. \secref{sec:1.1}{1.1}).  This
completes the proof.
\end{proof}
\begin{lemma}\phantomsection\label{lem:2.5.4}
Consider a matrix $\langle A_{t,s}\rangle$, all $A_{t,s}\in\M_2$.  Define
the matrices $\langle A^q_{t,s}\rangle\in\mathfrak J^q$ by the equations
$A^q_{t,s}=A_{t,s}$ if $t\leq q$ and $s\leq q$, and $A^q_{t,s}=0$
otherwise.  Then $\langle A_{t,s}\rangle\in\mathfrak J$ if and only if the
numerical sequence
$\OpNorm{\langle A^1_{t,s}\rangle},
 \OpNorm{\langle A^2_{t,s}\rangle},\ldots$
(cf. \lemref{lem:2.5.3}{2.5.3} above) is bounded.
\end{lemma}
\begin{proof}
Necessity: Assume $\langle A_{t,s}\rangle\in\mathfrak J$ and let it
correspond to $A^{\circ}\in\M_1$.  Then, clearly,
\[
  \langle A^q_{t,s}\rangle
  =\langle e^q_{t,s}\cdot1\rangle\langle A_{t,s}\rangle
   \langle e^q_{t,s}\cdot1\rangle,
\]
i.e. $\langle A^q_{t,s}\rangle$ corresponds to the operator
$E^qA^{\circ}E^q\in\M_1$.  Clearly
$E^qA^{\circ}E^q\in(\M_1)^{E^q}$.  Thus the second definition of
\lemref{lem:2.5.3}{2.5.3} gives
$\OpNorm{\langle A^q_{t,s}\rangle}=\OpNorm{E^qA^{\circ}E^q}$.
Now we have in $\M_1$ that
\[
  \OpNorm{E^qA^{\circ}E^q}
  \leq\OpNorm{E^q}^{2}\OpNorm{A^{\circ}}
  =\OpNorm{A^{\circ}}.
\]
So $\OpNorm{\langle A^q_{t,s}\rangle}\leq\OpNorm{A^{\circ}}$.  This holds
for $q=1,2,\ldots$ and hence the sequence
$\OpNorm{\langle A^1_{t,s}\rangle},
 \OpNorm{\langle A^2_{t,s}\rangle},\ldots$ is bounded.

Sufficiency: Assume that the numerical sequence
$\OpNorm{\langle A^1_{t,s}\rangle},
 \OpNorm{\langle A^2_{t,s}\rangle},\ldots$ is bounded.
$\langle A^q_{t,s}\rangle\in\mathfrak J^q$, so it corresponds to an
$A^q\in(\M_1)^{E^q}\subseteq\M_1$.  Owing to the definition of
\lemref{lem:2.5.3}{2.5.3},
$\OpNorm{\langle A^q_{t,s}\rangle}=\OpNorm{A^q}$, hence
\begin{equation}
  \text{The (numerical) sequence }
  \OpNorm{A^1},\OpNorm{A^2},\ldots\text{ is bounded.}
  \tag{2.5.\ensuremath{\alpha}}\label{eq:2.5.alpha}
\end{equation}

For $q\geq r$ clearly
\[
  \langle e^r_{t,s}\cdot1\rangle
  \langle A^q_{t,s}\rangle
  \langle e^r_{t,s}\cdot1\rangle
  =\langle A^r_{t,s}\rangle,
\]
hence $E^rA^qE^r=A^r$.  Now consider an $f\in\mathfrak M^r$.

For $q\geq r$, let $f_q=A^qf$.  Since
$\lVert f_q\rVert^2\leq\OpNorm{A^q}^{2}\lVert f\rVert^2$, these $f_q$'s
are bounded.  Also $f\in\mathfrak M^r\subseteq\mathfrak M^q$, so
$E^qf=f$.  Hence for $q'\geq q\geq r$,
$E^qf_{q'}=E^qA^{q'}f=E^qA^{q'}E^qf=A^qf=f_q$.  So $f_{q'}-f_q$ and
$f_q$ are orthogonal,
$\lVert f_{q'}-f_q\rVert^2+\lVert f_q\rVert^2=\lVert f_{q'}\rVert^2$,
i.e.
\begin{equation}
  \lVert f_{q'}-f_q\rVert^2
  =\lVert f_{q'}\rVert^2-\lVert f_q\rVert^2
  \qquad(q'\geq q\geq r).
  \tag{2.5.\ensuremath{\beta}}\label{eq:2.5.beta}
\end{equation}
Thus the $\lVert f_r\rVert^2,\lVert f_{r+1}\rVert^2,\ldots$ form a
monotone bounded sequence, and hence the
$\lim_{q\to\infty}\lVert f_q\rVert^2$ exists.  Therefore the right-hand
side of \eqref{eq:2.5.beta} converges to zero for $q,q'\to\infty$.  Hence
\eqref{eq:2.5.beta} now implies that the strong limit $A^qf$ exists for
$q\to\infty$.  The restriction $q\geq r$ may now be dropped.

The set of all $f$ for which $A^qf$ is strongly convergent is closed, since
the $\OpNorm{A^q}$ are uniformly bounded.  This set is obviously linear and
by the above it contains
$\mathfrak M^1,\mathfrak M^2,\ldots$.  Thus we see
\begin{equation}
  \left\{
  \begin{minipage}{0.78\linewidth}
  The sequence $A^1,A^2,\ldots$ converges for all $f$ in the closed linear
  set determined by $\mathfrak M^1,\mathfrak M^2,\ldots$.
  \end{minipage}
  \right.
  \tag{2.5.\ensuremath{\gamma}}\label{eq:2.5.gamma}
\end{equation}
Since the $\mathfrak M^1,\mathfrak M^2,\ldots$ are together dense, the set
of \eqref{eq:2.5.gamma} is necessarily $\Hilbert_1\otimes\Hilbert_2$ itself.
Denote the limit of $A^1f,A^2f,\ldots$ by $Af$.  Then $A$ belongs to $\M_1$
along with $A^1,A^2,\ldots$.  Thus we have shown
\begin{equation}
  \left\{
  \begin{minipage}{0.78\linewidth}
  The sequence $A^1,A^2,\ldots$ converges for all
  $f\in\Hilbert_1\otimes\Hilbert_2$ to an $A\in\M_1$ in the strong sense.
  \end{minipage}
  \right.
  \tag{2.5.\ensuremath{\delta}}\label{eq:2.5.delta}
\end{equation}

Let $A$ correspond to the matrix $\langle B_{t,s}\rangle$ in $B_\infty)$.
Of course $\langle B_{t,s}\rangle\in\mathfrak J$.  Since for $q\geq r$,
$E^rA^qE^r=A^r$, \eqref{eq:2.5.delta} implies that $E^rAE^r=A^r$.
Hence
\[
  \langle e^r_{t,s}\cdot1\rangle
  \langle B_{t,s}\rangle
  \langle e^r_{t,s}\cdot1\rangle
  =\langle A^r_{t,s}\rangle.
\]
Thus the rules of matrix computation yield $B_{t,s}=A_{t,s}$ if
$t,s\leq r$.  Now let $t,s$ be given and choose
$r=\operatorname{Max}(t,s)$.  Then the above formula gives
$B_{t,s}=A_{t,s}$.  Since $t,s$ were arbitrary, this means
$\langle B_{t,s}\rangle=\langle A_{t,s}\rangle$ and consequently
$\langle A_{t,s}\rangle\in\mathfrak J$ too.

Thus the proof is completed.
\end{proof}

\begin{lemma}\phantomsection\label{lem:2.5.5}
The set $\mathfrak J$ can be uniquely and purely algebraically characterized
in terms of $\M_2$ alone.
\end{lemma}
\begin{proof}
A unique characterization of $\mathfrak J$ was given in
\lemref{lem:2.5.4}{2.5.4}.  By using the first definition of
\lemref{lem:2.5.3}{2.5.3} for
$\OpNorm{\langle A^q_{t,s}\rangle}$, that one in terms of $\M_2^q$, it
becomes clear that this characterization is purely algebraical in terms of
$\M_2$ and of the $\M_2^q$, $q=1,2,\ldots$.  But we have already reduced the
$\M_2^q$, i.e. the $\overline{\M}_2^q$ for $q$ finite, to $\M_2$ by
\thmref{thm:IV}{IV}.

The proof is thus completed.
\end{proof}

The remarks made in \secref{sec:2.4}{2.4}, preceding \thmref{thm:IV}{IV},
in conjunction with the above \lemref{lem:2.5.5}{2.5.5}, permit us now to
assert
\begin{namedstatement}{Theorem}{IV$'$}\phantomsection\label{thm:IV-prime}
$\overline{\M}^{\infty}$ too is one-valued.
\end{namedstatement}
Our definitions in $B_\infty)$ and that of \secref{sec:2.3}{2.3} give
together
\begin{equation}
  (\overline{\M}^{c})^{\infty}=(\overline{\M}^{\infty})^{c}.
  \tag{2.5.\ensuremath{\epsilon}}\label{eq:2.5.epsilon}
\end{equation}
(This corresponds to \eqref{eq:2.4.alpha}.  We could extend
\eqref{eq:2.4.beta} too, but we shall not need it here.  Furthermore this
extension is an easy consequence of \lemref{lem:3.1.6}{3.1.6}.)

\PaperSection{sec:2.6}{Inverse operation to $\overline{\N}^{p}$; matrix units.}
We undertake now to find the inverse operation to $\overline{\N}^{p}$.  At
first we restrict neither the ring $\N$ nor the exponent
$p=1,2,\ldots,\infty$.

\begin{definition}\phantomsection\label{def:2.6.1}
A system of $p^2$ operators $W_{\mu,\nu}$ ($\mu,\nu\leq p$) is a system of
$p$-matrix units if it possesses the following properties:
\begin{enumerate}[label=(\roman*)]
\item $W_{\mu,\nu}^{*}=W_{\nu,\mu}$,
\item
\[
  W_{\mu,\nu}W_{\sigma,\omega}
  =\begin{cases}
     W_{\mu,\omega},&\nu=\sigma,\\
     0,&\nu\ne\sigma.
   \end{cases}
\]
\end{enumerate}
\end{definition}

Some immediate properties of these systems are given in
\begin{lemma}\phantomsection\label{lem:2.6.1}
Every system of $p$-matrix units $W_{\mu,\nu}$ ($\mu,\nu\leq p$) possesses
the following properties:
\begin{enumerate}[label=(\roman*)]
\item The $W_{\mu,\mu}$ ($\mu\leq p$) are pairwise orthogonal projections.
\item Every $W_{\mu,\nu}$ is partially isometric with the initial projection
$W_{\nu,\nu}$ and the final projection $W_{\mu,\mu}$ (cf.
\cite[\S4.3]{ref5}).
\item Form $E_0=\sum_{\mu=1}^{p}W_{\mu,\mu}$.  This $\sum_{\mu=1}^{p}$ is
either a finite sum or when $p=\infty$ converges in the sense of the strong
operator topology, irrespective of the order in which the
$\mu=1,2,\ldots$ are gone through.  This $E_0$ is a projection and it is the
unit of the system $W_{\mu,\nu}$, i.e. always
\[
  E_0W_{\mu,\nu}=W_{\mu,\nu}E_0=W_{\mu,\nu}.
\]
\end{enumerate}
\end{lemma}
\begin{proof}
Ad (i), (iii): The convergence in (iii) is a consequence of (i), as is seen
from familiar considerations concerning pairwise orthogonal projections
(cf. e.g. \cite[pp.~76--78]{ref5}, or the proof of
\eqref{eq:2.5.delta} in the proof of \lemref{lem:2.5.4}{2.5.4} above).  All
other assertions are immediately verified by using (i), (ii) in
\defref{def:2.6.1}{2.6.1}.

Ad (ii).  \defref{def:2.6.1}{2.6.1} gives directly
\[
  W_{\mu,\nu}^{*}W_{\mu,\nu}=W_{\nu,\nu},
  \qquad
  W_{\mu,\nu}W_{\mu,\nu}^{*}=W_{\mu,\mu}.
\]
Since $W_{\nu,\nu}$, $W_{\mu,\mu}$ are projections by (i) above, these
formulae imply our assertions by \cite[p.~142, Lemma~4.3.2 (the last two
criteria), and p.~141, Lemma~4.3.1]{ref5}.
\end{proof}


We prove now
\begin{theorem}\phantomsection\label{thm:V}
$\overline{\M}=\overline{\N}^{p}$ is equivalent to this: There exists a
system of $p$-matrix units $W_{\mu,\nu}\in\M$ ($\mu,\nu\leq p$) with the
unit $1$ (cf. (iii) in \lemref{lem:2.6.1}{2.6.1}) and with the following
further property:

$W_{1,1}$ is a projection.  Let $\mathfrak M$ be the closed linear set with
$P_{\mathfrak M}=W_{1,1}$.  Then $\overline{\M_{(\mathfrak M)}}$ must be
algebraically isomorphic to $\overline{\N}$.
% Editorial note E2: corrected the original typo ``circumstqnces'' to ``circumstances'' in footnote 15.
\footnote{Observe that in this presentation $\overline{\N}$ is determined
by $\overline{\M}$ and $W_{1,1}$ together.  Whether $\overline{\M}$ alone
suffices to determine $\overline{\N}$ is a different question.  If, however,
$p$ is finite and $\M$ is a factor, we do have that $\overline{\M}$ and $p$
determine $\overline{\N}$ uniquely if there is such an $\overline{\N}$.
For if $\mathfrak M^{(1)}$ is the range of $W_{1,1}$ for one choice of the
$p$-matrix units and $\mathfrak M^{(2)}$ that for another, then
$D_{\M}(\mathfrak M^{(1)})=(1/p)D_{\M}(\Hilbert)
=D_{\M}(\mathfrak M^{(2)})$.  Hence \lemref{lem:2.8.1}{2.8.1} (of our
present text) implies
$\overline{\M_{(\mathfrak M^{(1)})}}
=\overline{\M_{(\mathfrak M^{(2)})}}$.  Hence
$\overline{\M_{(\mathfrak M)}}$ depends only on $\overline{\M}$ and $p$.

If however $p$ is $\infty$ then $\M$ itself is in an infinite case by (iii)
of \lemref{lem:2.7.1}{2.7.1}.  Under these circumstances,
$\overline{\M}$ fails to determine $\overline{\N}$ as may be seen from
$\delta)$ in \lemref{lem:3.1.5}{3.1.5} (write $\overline{\M}$ for its
$\overline{\M}^{\infty}$).}
\end{theorem}
(Concerning the omission of this extra condition, cf. the corollary below.)
\begin{proof}
Sufficiency: Assume the existence of $W_{\mu,\nu}$ ($\mu,\nu\leq p$) as
described.  $W_{1,1}\in\M$ is a projection.  Let $\mathfrak M$ be such that
$P_{\mathfrak M}=E=W_{1,1}$.

For every $A^{0}\in\M$ define
\begin{equation}
  A'_{t,s}=W_{1,s}A^{0}W_{t,1}\qquad(t,s\leq p).
  \tag{2.6.\ensuremath{\alpha}}\label{eq:2.6.alpha}
\end{equation}
Clearly $A'_{t,s}\in\M$ and one verifies immediately that
$EA'_{t,s}=A'_{t,s}E=A'_{t,s}$ ($E=W_{1,1}$).  Hence
$A'_{t,s}\in\M^{E}$.  Now apply \lemref{lem:2.5.1}{2.5.1} and form
\begin{equation}
  A_{t,s}=(A'_{t,s})_{(\mathfrak M)}.
  \tag{2.6.\ensuremath{\beta}}\label{eq:2.6.beta}
\end{equation}
Accordingly $A_{t,s}\in\M_{(\mathfrak M)}$.

Thus to every $A^{0}\in\M$ there corresponds a $p$-order matrix
$\langle A_{t,s}\rangle$ ($t,s\leq p$) with all
$A_{t,s}\in\M_{(\mathfrak M)}$.  Let $\mathfrak J$ be the set of these
$\langle A_{t,s}\rangle$ which correspond to an $A^{0}\in\M$.  We shall
now show that this correspondence establishes a $p$-matrix automorphism of
$\M$ over $\M_{(\mathfrak M)}$ in the sense of $B_p)$ in
\secref{sec:2.4}{2.4}.  That means
$\overline{\M}=\overline{\M_{(\mathfrak M)}}^{p}$ and as
$\overline{\M_{(\mathfrak M)}}=\overline{\N}$, so
$\overline{\M}=\overline{\N}^{p}$.  This will complete the proof of the
sufficiency.  We prove the above assertions in three successive steps.

The correspondence is one-to-one---i.e. if $A^{0},B^{0}\in\M$ have
$\langle A_{t,s}\rangle=\langle B_{t,s}\rangle$ then $A^{0}=B^{0}$.
Now $\langle A_{t,s}\rangle=\langle B_{t,s}\rangle$ means
$A_{t,s}=B_{t,s}$ for $t,s\leq p$ and hence
$(A'_{t,s})_{(\mathfrak M)}=(B'_{t,s})_{(\mathfrak M)}$ and
$A'_{t,s}=B'_{t,s}$ since both are $\in\M^{E}$.  Thus
$W_{1,s}A^{0}W_{t,1}=W_{1,s}B^{0}W_{t,1}$.  Multiplication by $W_{s,1}$ on
the left and by $W_{1,t}$ on the right gives
$W_{s,s}A^{0}W_{t,t}=W_{s,s}B^{0}W_{t,t}$.  If we sum over $s$ and $t$ we
obtain $A^{0}=B^{0}$ as desired, since by hypothesis
$E_0=1=\sum_{s=1}^{p}W_{s,s}$.

The correspondence is a matrix isomorphism.  This can be verified by the use
of \defref{def:2.6.1}{2.6.1},
$E_0=\sum_{\mu=1}^{p}W_{\mu,\mu}=1$ and the equations
\eqref{eq:2.6.alpha} and \eqref{eq:2.6.beta}.  For these show that by the use
of Equation \eqref{eq:2.6.alpha} we correspond to
$1,\alpha A^{0},A^{0*},A^{0}+B^{0},A^{0}B^{0}$ (in $\M$) matrices
\[
  \langle\delta_{t,s}E\rangle,
  \quad\langle\alpha A'_{t,s}\rangle,
  \quad\langle A_{s,t}^{\prime *}\rangle,
  \quad\langle A'_{t,s}+B'_{t,s}\rangle,
  \quad\left\langle\sum_{r=1}^{p}A'_{r,s}B'_{t,r}\right\rangle
\]
of elements of $\M^{E}$ and then by the use of Equation
\eqref{eq:2.6.beta} we correspond these to
\[
  \langle\delta_{t,s}1\rangle,
  \quad\langle\alpha A_{t,s}\rangle,
  \quad\langle A_{s,t}^{*}\rangle,
  \quad\langle A_{t,s}+B_{t,s}\rangle,
  \quad\left\langle\sum_{r=1}^{p}A_{r,s}B_{t,r}\right\rangle,
\]
whose elements are in $\M_{(\mathfrak M)}$.

$\mathfrak J$ contains all finite matrices.  Consider a finite matrix
$\langle A_{t,s}\rangle$ where $A_{t,s}\in\M_{(\mathfrak M)}$.  There is
a $q<\infty$ such that $A_{t,s}=0$ when $t>q$ or when $s>q$.  Let
$A'_{t,s}\in\M^{E}$ satisfy \eqref{eq:2.6.beta} and put
\[
  A^{0}=\sum_{t=1}^{q}\sum_{s=1}^{q}W_{s,1}A'_{t,s}W_{1,t}.
\]
Then $A^{0}\in\M$ and it follows from the rules of
\defref{def:2.6.1}{2.6.1} that \eqref{eq:2.6.alpha} holds.  Thus
$\langle A_{t,s}\rangle$ corresponds to this $A^{0}$ and therefore it
belongs to $\mathfrak J$.

Necessity: Assume $\overline{\M}=\overline{\N}^{p}$ in the sense of $B_p)$
in \secref{sec:2.4}{2.4}.  For every pair $\mu,\nu\leq p$ form the matrix
$\langle\delta_{t,\mu}\delta_{s,\nu}1\rangle$.  This matrix is clearly
finite and hence it belongs to $\mathfrak J$.  Thus it corresponds to an
element of $\M$ which we denote by $W_{\mu,\nu}$.  Then the matrix
computation rules of $B_p)$ give immediately the rules of
\defref{def:2.6.1}{2.6.1}.  Thus the results of
\lemref{lem:2.6.1}{2.6.1} are now available.  Now in the notation of (ii) in
\lemref{lem:2.5.2}{2.5.2},
$E^q=\sum_{\mu=1}^{q}W_{\mu,\mu}$ ($q=1,2,\ldots$).  Hence
\lemref{lem:2.5.2}{2.5.2} (iii) and \lemref{lem:2.6.1}{2.6.1} (iii) imply
$E_0=\sum_{\mu=1}^{p}W_{\mu,\mu}=1$.

Put $P_{\mathfrak M}=E=W_{1,1}$.  $A^{0}\in\M^{E}$ means
$EA^{0}=A^{0}E=A^{0}$ or, by the above-mentioned rules, that the matrix of
$A^{0}$ has the form
$\langle\delta_{t,1}\delta_{s,1}A\rangle$, $A\in\N$.  This correspondence
between the $A^{0}\in\M^{E}$ and the $A\in\N$ is clearly one-to-one, it
carries $E\in\M^{E}$ (which has the matrix
$\langle\delta_{t,1}\delta_{s,1}1\rangle$) into $1\in\N$ and it leaves the
operations $\alpha A,A^{*},A+B,AB$ invariant.  These results, together with
\lemref{lem:2.5.1}{2.5.1}, show therefore that
$\M_{(\mathfrak M)}$ is algebraically isomorphic to $\N$.

Thus all our requirements are fulfilled.
\end{proof}

\begin{corollary}\phantomsection\label{cor:2.6.1}
Given $\overline{\M}$ and $p$, the equation
$\overline{\M}=\overline{\N}^{p}$ possesses a solution $\overline{\N}$ if
and only if there exists a system of $p$-matrix units,
$W_{\mu,\nu}\in\M$ ($\mu,\nu\leq p$), with the unit $1$.
\end{corollary}
\begin{proof}
This becomes clear when we omit the last condition in the above theorem,
since the only effect of that condition is to determine $\overline{\N}$ in
terms of the $W_{\mu,\nu}$.
\end{proof}

Under certain conditions we may even go further.
\begin{lemma}\phantomsection\label{lem:2.6.2}
If $\M$ is a factor not in case $(\mathrm I_n)$, and $p=2,3,\ldots$ is
finite, then there is a system of $p$-matrix units
$W_{\mu,\nu}\in\M$ ($\mu,\nu\leq p$) with unit $E_0=1$.
\end{lemma}
\begin{proof}
We first obtain a set $E_1,\ldots,E_p$ of projections in $\M$ such that
$D_{\M}(E_i)=D_{\M}(E_1)$ and $\sum_{n=1}^{p}E_n=1$.  Suppose first that
$\M$ is in a case $(\mathrm{II}_1)$.  We may assume that $D_{\M}(1)=1$.
The range of $D_{\M}$ contains $1/p$ (cf. \cite[p.~172, Theorem~VIII]{ref5}),
and hence there is an $E_1\in\M$ with $D_{\M}(E_1)=1/p$.  Since
$D_{\M}(E_1)=1/p\leq1-1/p=D_{\M}(1-E_1)$, there exists a projection
$E_2\leq1-E_1$ with $D_{\M}(E_2)=1/p=D_{\M}(E_1)$ (cf.
\cite[p.~167, Lemma~8.3.1]{ref5}).  Furthermore
$D_{\M}(1-E_1-E_2)=1-2/p$.  If $p>2$ we can find an
$E_3\leq1-E_1-E_2$ such that $D_{\M}(E_3)=1/p$.  Thus we may continue until
we obtain a set of $p$ pairwise orthogonal projections
$E_1,\ldots,E_p\in\M$ with
$\sum_{n=1}^{p}E_n=1$, $D_{\M}(E_n)=1/p$.

If $\M$ is in an infinite case, we can easily modify the proof of
\cite[p.~157, Lemma~7.2.3]{ref5} to yield that there are $p$ pairwise
orthogonal, dimensionally equivalent, closed, linear sets
$\mathfrak M_1,\ldots,\mathfrak M_p$ which together span $\Hilbert$.  In
that proof it is only necessary to separate the $\mathfrak R_i$'s into $p$
sets rather than $2$.  Then if $E_i$ is the projection on $\mathfrak M_i$,
$\sum_{n=1}^{p}E_n=1$ and $D_{\M}(E_n)=\infty=D_{\M}(E_1)$.

Since $D_{\M}(E_t)=D_{\M}(E_1)$ there is a partially isometric
$W_{t,1}\in\M$ with initial set $\mathfrak M_1$ and final set
$\mathfrak M_t$ by the definition of the dimensionality function.  If we
now define $W_{t,s}$ by the equation $W_{t,s}=W_{t,1}W_{s,1}^{*}$, the
properties of the partially isometric $W_{t,1}$ (cf. \cite[\S4.3]{ref5})
readily yield the equations of \defref{def:2.6.1}{2.6.1} above.  Hence the
$W_{t,s}$ constitute a set of $p$-matrix units with $E_0=1$,
$W_{t,s}\in\M$.
\end{proof}

\begin{lemma}\phantomsection\label{lem:2.6.3}
If $\M$ is in case $(\mathrm I_n)$ then $\M$ possesses a set of $p$-matrix
units with $E_0=1$ if and only if $p$ divides $n$.
\end{lemma}
\begin{proof}
Necessity: If $\M$ has a set of $p$-matrix units, with $E_0=1$, then by
\lemref{lem:2.6.1}{2.6.1}, the $W_{1,1},\ldots,W_{p,p}$ are $p$ pairwise
orthogonal projections each with same relative dimension, and whose sum is
$1$.  It follows that
$D_{\M}(W_{t,t})=(1/p)D_{\M}(\Hilbert)$.  If we take
$D_{\M}(\Hilbert)=n$, we know that
$D_{\M}(W_{t,t})=(1/p)D_{\M}(\Hilbert)$ is a whole number and hence $p$
divides $n$ (cf. \cite[p.~172, Theorem~VIII]{ref5}).
% Editorial note E3: corrected the original typo ``p divide n'' to ``p divides n''.

Sufficiency: If $p$ divides $n$, we can find $p$ dimensionally equivalent
pairwise orthogonal projections $E_1,\ldots,E_p\in\M$ and with
$\sum_{n=1}^{p}E_n=1$.  For if $D_{\M}(\Hilbert)=n$, we can find a
projection $E_1\in\M$ with $D_{\M}(E_1)=n/p$ and we proceed as in the proof
of \lemref{lem:2.6.2}{2.6.2} to find $E_2,\ldots,E_p$.  The rest of that
proof then applies to the present situation.
\end{proof}

\thmref{thm:V}{V} yields
\begin{lemma}\phantomsection\label{lem:2.6.4}
\begin{enumerate}[label=(\roman*)]
\item If $\M$ is a factor, not in case $(\mathrm I_n)$ and $p$ is a finite
integer, then there is an $\overline{\N}$ such that
$\overline{\N}^{p}=\overline{\M}$.
\item If $\M$ is in case $(\mathrm I_n)$ there is an $\overline{\N}$ such
that $\overline{\N}^{p}=\overline{\M}$ if and only if $p$ divides $n$.
\item If there is an $\overline{\N}$ such that
$\overline{\N}^{p}=\overline{\M}$ for
$p<\infty$, $p=2,3,\ldots$, then there is a manifold $\mathfrak M$ such that
$\overline{\M_{(\mathfrak M)}}=\overline{\N}$ and
$D_{\M}(\mathfrak M)=(1/p)D_{\M}(\Hilbert)$.
\end{enumerate}
\end{lemma}

As we remarked in footnote 15, we may add
\begin{lemma}\phantomsection\label{lem:2.6.5}
If $p$ is finite, $\M$ is a factor and if $\overline{\N}$ is such that
$\overline{\N}^{p}=\overline{\M}$ then $\overline{\N}$ is uniquely
determined.
\end{lemma}

\PaperSection{sec:2.7}{Nature of $\overline{\N}$; the exponent $p=\infty$.}
It remains to consider the nature of $\overline{\N}$ and also what happens if
$p=\infty$.
\begin{lemma}\phantomsection\label{lem:2.7.1}
We have for any fixed $p=1,2,\ldots,\infty$,
\begin{enumerate}[label=(\roman*)]
\item $\overline{\N}$ is a factor if and only if
$\overline{\N}^{p}$ is one.
\item If $\overline{\N},\overline{\N}^{p}$ are factors, then they belong to
the same one of the Cases (I), (II), (III).
\item If $\overline{\N},\overline{\N}^{p}$ are factors, then
$\overline{\N}^{p}$ is in a finite case if and only if $\overline{\N}$ is
in a finite case and $p$ is finite.
\end{enumerate}
\end{lemma}
\begin{proof}
Choose an $\M$ with $\overline{\M}=\overline{\N}^{p}$.

Ad (i).  An $A^{0}\in\M$ belongs to the center of $\M$ if and only if its
matrix $\langle A_{t,s}\rangle$ (all $A_{t,s}\in\N$) commutes with all
matrices $\langle B_{t,s}\rangle\in\mathfrak J$.  For this it is necessary
that it commute with all finite matrices $\langle B_{t,s}\rangle$ (all
$B_{t,s}\in\N$, cf. $B_p)$).  This is immediately seen to imply that
$A_{t,s}=\delta_{t,s}A$ ($t,s\leq p$) for an $A$ belonging to the center of
$\N$.  This condition is also clearly sufficient.  Hence the $A^{0}$ of the
center of $\M$ are precisely the $\alpha1$ corresponding to the matrices
$\langle\alpha\delta_{t,s}1\rangle$ if and only if the $A$ of the center of
$\N$ are precisely the $\alpha1$, i.e. $\M$ is a factor if and only if
$\N$ is one.

This completes the proof.

Ad (ii).  By \thmref{thm:V}{V}, $\N$ is isomorphic to an
$\M_{(\mathfrak M)}$.  Hence our assertion follows from
\cite[p.~189, Lemma~11.4.3]{ref5}.

Ad (iii).  The $W_{\mu,\mu}$ ($\mu\leq p$) are pairwise orthogonal
projections $\in\M$ all with the same relative dimension in $\M$ owing to
(i), (ii) in \lemref{lem:2.6.1}{2.6.1}.  Hence
\refstepcounter{footnote}
\begin{equation}
  D_{\M}(1)=\sum_{\mu=1}^{p}D_{\M}(W_{\mu,\mu})
  =pD_{\M}(W_{1,1}).\text{\textsuperscript{\hyperref[fn:16]{\thefootnote}}}
  \tag{2.7.\ensuremath{\alpha}}\label{eq:2.7.alpha}
\end{equation}
\footnotetext[\value{footnote}]{\phantomsection\label{fn:16}Here as in some subsequent instances
(proof of \lemref{lem:5.3.4}{5.3.4}; III in
\secref{sec:5.5}{5.5}) it is convenient to use the convention
$\infty\cdot0=0$.}
Now $W_{1,1}$ is not zero, since if it were we should have
$W_{\mu,\mu}=0$ and $\sum_{\mu=1}^{p}W_{\mu,\mu}=0$ for
$p=1,2,\ldots,\infty$.  Hence $D_{\M}(W_{1,1})\ne0$.  Now $\M$ is in a
finite case if and only if $D_{\M}(1)$ is finite.  By the above, this is
equivalent to the statement that both $p$ and $D_{\M}(W_{1,1})$ are finite.
And the finiteness of
$D_{\M}(W_{1,1})=D_{\M}(\mathfrak M)$,
($P_{\mathfrak M}=W_{1,1}$) is, according to
\cite[p.~189, Lemma~11.4.3]{ref5}, equivalent to
$\M_{(\mathfrak M)}$ being in a finite case, i.e. by
\thmref{thm:V}{V}, to $\N$ being in a finite case.

This completes the proof.
\end{proof}
\PaperSection{sec:2.8}{Finite factors; the operation $\overline{\M}^{\alpha}$,
$0<\alpha\leq1$.}
Let $\M$ be a factor in a Hilbert space $\Hilbert$.  As usual we denote by
$D_{\M}(E)$ and by $D_{\M}(\mathfrak M)$,
($E=P_{\mathfrak M}\in\M$) a fixed relative dimension function of $\M$.
We consider again the rings $\M^{E}$ and $\M_{(\mathfrak M)}$ as described
in the beginning of \secref{sec:2.5}{2.5}.

\begin{lemma}\phantomsection\label{lem:2.8.1}
If $\mathfrak M_1,\mathfrak M_2$ are two closed linear sets $\eta\M$ with
$D_{\M}(\mathfrak M_1)=D_{\M}(\mathfrak M_2)$ then
$\M_{(\mathfrak M_1)}$ and $\M_{(\mathfrak M_2)}$ (in $\mathfrak M_1$ and
$\mathfrak M_2$, respectively) are spatially isomorphic.
\end{lemma}
\begin{proof}
Since $D_{\M}(\mathfrak M_1)=D_{\M}(\mathfrak M_2)$ there exists a
partially isometric $U\in\M$ with the initial set $\mathfrak M_1$ and final
set $\mathfrak M_2$.  This is a one-to-one isomorphic mapping of
$\mathfrak M_1$ on $\mathfrak M_2$ and it obviously carries
$\M_{(\mathfrak M_1)}$ into $\M_{(\mathfrak M_2)}$.  Hence it establishes
the desired spatial isomorphism.
\end{proof}

For the remainder of this chapter, we assume that the factor $\M$ is in a
finite case.

Consider a real number $\alpha$ and a projection $E\in\M$ or equivalently a
closed linear set $\mathfrak M\,\eta\,\M$ with $E=P_{\mathfrak M}$.  We assume
$E\ne0$, i.e. $\mathfrak M\ne\{0\}$ and
\begin{equation}
  \alpha=\frac{D_{\M}(E)}{D_{\M}(1)}.
  \tag{2.8.\ensuremath{\alpha}}\label{eq:2.8.alpha}
\end{equation}
Then the above \lemref{lem:2.8.1}{2.8.1} shows that
$\overline{\M_{(\mathfrak M)}}$ depends only on $\overline{\M}$ and
$\alpha$ but not on the $E,\mathfrak M$ themselves.  (A spatial isomorphism
implies an algebraic isomorphism.)

Consider now two algebraically isomorphic $\M_1$ and $\M_2$.  Choose
$E_1\in\M_1$ with
$\alpha=D_{\M_1}(E_1)/D_{\M_1}(1)$.  Then the isomorphism carries $E_1$
into an $E_2\in\M_2$ with
$\alpha=D_{\M_2}(E_2)/D_{\M_2}(1)$.  Let
$\mathfrak M_1,\mathfrak M_2$ be $\eta\M_1,\eta\M_2$ respectively, and such
that $P_{\mathfrak M_1}=E_1$, $P_{\mathfrak M_2}=E_2$.  The isomorphism
carries $\M_1^{E_1}$ into $\M_2^{E_2}$, in particular $E_1$ into $E_2$, and
it leaves the operations $\alpha A,A^{*},A+B,AB$ invariant.  Hence it
generates, by \lemref{lem:2.5.1}{2.5.1}, an algebraic isomorphism of
$(\M_1)_{(\mathfrak M_1)}$ and $(\M_2)_{(\mathfrak M_2)}$.

Thus $\overline{\M_{(\mathfrak M)}}$ (assuming its existence) is uniquely
determined by $\overline{\M}$ and $\alpha$.  We denote this
$\overline{\M_{(\mathfrak M)}}$ by $\overline{\M}^{\alpha}$.

This notation could conflict with that of \thmref{thm:IV}{IV} when
$p=\alpha=1$.  But in that case clearly $\overline{\M}^{1}$ exists and is
equal to $\overline{\M}$ under both definitions.

The $\alpha$, for which $\overline{\M}^{\alpha}$ can be formed, comprise
precisely the range of $D_{\M}(E)/D_{\M}(1)$ for all $E\in\M$, $E\ne0$.
This is the set $(1/n,2/n,\ldots,1)$ if $\M$ is in a case
$(\mathrm I_n)$ and the set $0<\alpha\leq1$ if $\M$ is in a case
$(\mathrm{II}_1)$.

Summing up,
\begin{theorem}\phantomsection\label{thm:VI}
$\overline{\M}^{\alpha}$ exists if and only if $\alpha$ belongs to the
following set: $(1/n,2/n,\ldots,1)$ if $\M$ is in a case $(\mathrm I_n)$,
($n=1,2,\ldots$), $(0<\alpha\leq1)$ if $\M$ is in the case
$(\mathrm{II}_1)$.
\end{theorem}
There is no conflict between this notation and that of \thmref{thm:IV}{IV}.

Since the relative dimension $D_{\M}(E)$ can be characterized in a way which
is invariant under conjugate isomorphisms (cf. the remark at the end of
\secref{sec:2.3}{2.3}), our present definition and that of
\secref{sec:2.3}{2.3} yield together
\begin{equation}
  (\overline{\M}^{c})^{\alpha}=(\overline{\M}^{\alpha})^{c}
  \qquad(\alpha\text{ in the range of \thmref{thm:VI}{VI}}).
  \tag{2.8.\ensuremath{\beta}}\label{eq:2.8.beta}
\end{equation}

Consider next two projections $E,F\in\M$, $E,F\ne0$ and $E\leq F$, or
equivalently the closed linear sets $\mathfrak M,\mathfrak N\,\eta\,\M$ with
$P_{\mathfrak M}=E$, $P_{\mathfrak N}=F$.  Hence
$\mathfrak M\ne(0)$, $\mathfrak N\ne(0)$ and
$\mathfrak M\subseteq\mathfrak N$.  To avoid misunderstanding, let $1$ be
the unit operator in $\Hilbert$ and $1'=F_{(\mathfrak N)}$ be the unit in
$\mathfrak N$.  Put
\[
  \alpha=\frac{D_{\M}(F)}{D_{\M}(1)},
  \qquad
  \beta=\frac{D_{\M_{(\mathfrak N)}}(E_{(\mathfrak N)})}
                   {D_{\M_{(\mathfrak N)}}(1')}
       =\frac{D_{\M}(E)}{D_{\M}(F)}.
\]
Clearly $\alpha\beta=D_{\M}(E)/D_{\M}(1)$ and
$\M_{(\mathfrak M)}=(\M_{(\mathfrak N)})_{(\mathfrak M)}$.  Our
definitions now yield immediately
\begin{equation}
  (\overline{\M}^{\alpha})^{\beta}=\overline{\M}^{\alpha\beta}
  \qquad(\alpha,\beta\text{ in the range of \thmref{thm:VI}{VI}}).
  \tag{2.8.\ensuremath{\gamma}}\label{eq:2.8.gamma}
\end{equation}


\PaperSection{sec:2.9}{Extension to all positive exponents
$\overline{\M}^{\theta}$.}
We continue the discussion of the preceding section.
\begin{lemma}\phantomsection\label{lem:2.9.1}
If $p$ is finite and $\M$ and $\N$ are factors in a finite case,
$\overline{\N}=\overline{\M}^{1/p}$ is equivalent to the statement
$\overline{\M}=\overline{\N}^{p}$.  (The statement
``$\overline{\M}^{1/p}=\overline{\N}$'' includes
``$\overline{\M}^{1/p}$ exists.'')
\end{lemma}
\begin{proof}
Assume $\overline{\M}=\overline{\N}^{p}$.  It follows from
\thmref{thm:V}{V} that there is an $E$ with
$D_{\M}(E)=(1/p)D_{\M}(1)$ such that
$\overline{\N}=\overline{\M_{(\mathfrak M)}}$ where $\mathfrak M$ is such
that $E=P_{\mathfrak M}$.  Hence, by the above definition,
$\overline{\M}^{1/p}=\overline{\N}$.

Conversely, suppose $\overline{\M}^{1/p}=\overline{\N}$.  It follows that
$1/p$ is in the set of \thmref{thm:VI}{VI} and hence, if $\M$ is in a case
$(\mathrm I_n)$, $p$ divides $n$.  Thus \lemref{lem:2.6.2}{2.6.2} and
\lemref{lem:2.6.3}{2.6.3} imply that there is an $\overline{\N}_0$ such
that $\overline{\N}_0^{p}=\overline{\M}$.  Hence the above result (with
$\overline{\N}_0$ in place of $\overline{\N}$) gives
$\overline{\N}_0=\overline{\M}^{1/p}$, i.e.
$\overline{\N}_0=\overline{\N}$.  Consequently
$\overline{\N}^{p}=\overline{\N}_0^{p}=\overline{\M}$.

This completes the proof.
\end{proof}

\begin{lemma}\phantomsection\label{lem:2.9.2}
If $\overline{\M}^{\alpha}$ exists (in the sense of
\thmref{thm:VI}{VI}) and if $p\alpha=q\beta$,
($p,q=1,2,\ldots$) then
$(\overline{\M}^{\alpha})^{p}=(\overline{\M}^{q})^{\beta}$.  (This is to
include the statement that $(\overline{\M}^{q})^{\beta}$ exists.)
\end{lemma}
\begin{proof}
$p\alpha=q\beta$.  Hence $\alpha/q=\beta/p$,
$\overline{\M}^{\alpha}$ exists.  From this we shall draw inferences so that
the existence of each expression that appears will be established when we
write it down.  In this way \eqref{eq:2.8.gamma} and
\lemref{lem:2.9.1}{2.9.1} give
\[
  \overline{\M}^{\alpha}
  =\bigl((\overline{\M}^{q})^{1/q}\bigr)^{\alpha}
  =(\overline{\M}^{q})^{\alpha/q}
  =(\overline{\M}^{q})^{\beta/p}.
\]
Hence by \lemref{lem:2.9.1}{2.9.1},
\[
  (\overline{\M}^{\alpha})^{p}
  =\bigl((\overline{\M}^{q})^{\beta/p}\bigr)^{p}
  =\left(\bigl((\overline{\M}^{q})^{\beta}\bigr)^{1/p}\right)^{p}
  =(\overline{\M}^{q})^{\beta}.
\]
\end{proof}

\begin{lemma}\phantomsection\label{lem:2.9.3}
For a given $\overline{\M}$ consider all $\alpha$ of the set in
\thmref{thm:VI}{VI}, and all $p=1,2,\ldots$.  Then we have
\begin{enumerate}[label=(\roman*)]
\item The range of $\theta=p\alpha$ is the set
$(1/n,2/n,\ldots)$ if $\M$ is in a case $(\mathrm I_n)$
($n=1,2,\ldots$), and the set $0<\theta<\infty$ if $\M$ is in the case
$(\mathrm{II}_1)$.
\item $(\overline{\M}^{\alpha})^{p}$ depends only on $\theta=p\alpha$ and
not on the individual values of $\alpha$ and $p$.
\end{enumerate}
\end{lemma}
\begin{proof}
Ad (i).  Immediate by \thmref{thm:V}{V}.

Ad (ii).  \lemref{lem:2.9.2}{2.9.2} yields, when $p\alpha=q\beta$,
$(\overline{\M}^{\alpha})^{p}=(\overline{\M}^{q})^{\beta}$ and
$(\overline{\M}^{\beta})^{q}=(\overline{\M}^{p})^{\alpha}$.  Hence
$(\overline{\M}^{\alpha})^{p}=(\overline{\M}^{\beta})^{q}$.
\end{proof}

We denote the above $(\overline{\M}^{\alpha})^{p}$, $\theta=p\alpha$, by
$\overline{\M}^{\theta}$.

This notation does not conflict with those of \thmref{thm:IV}{IV} and
\thmref{thm:VI}{VI}.  This can be seen by putting $\alpha=1$ or $p=1$,
respectively.

Summing up,---
\begin{theorem}\phantomsection\label{thm:VII}
$\overline{\M}^{\theta}$ exists if and only if $\theta$ belongs to the
following set: $(1/n,2/n,\ldots)$ if $\M$ is in a case $(\mathrm I_n)$,
($n=1,2,\ldots$), $(0<\theta<\infty)$ if $\M$ is in the case
$(\mathrm{II}_1)$.
\end{theorem}
There is no conflict between this notation and that of \thmref{thm:IV}{IV}
and \thmref{thm:VI}{VI}.

Combining \eqref{eq:2.4.alpha} and \eqref{eq:2.8.beta} gives
\begin{equation}
  (\overline{\M}^{c})^{\theta}=(\overline{\M}^{\theta})^{c}
  \qquad(\theta\text{ in the set of \thmref{thm:VII}{VII}}).
  \tag{2.9.\ensuremath{\alpha}}\label{eq:2.9.alpha}
\end{equation}

Also, by the use of \eqref{eq:2.4.beta}, \eqref{eq:2.8.gamma} and
\lemref{lem:2.9.2}{2.9.2}, we obtain
\[
  \left(((\overline{\M}^{\alpha})^{p})^{\beta}\right)^{q}
  =\left(((\overline{\M}^{\alpha})^{\beta})^{p}\right)^{q}
  =(\overline{\M}^{\alpha\beta})^{pq}
  =\overline{\M}^{pq\alpha\beta}.
\]
Hence if we put $\theta=p\alpha$, $\xi=q\beta$, then
\begin{equation}
  (\overline{\M}^{\theta})^{\xi}
  =(\overline{\M}^{\xi})^{\theta}
  =\overline{\M}^{\xi\theta}
  \qquad(\theta\text{ and }\xi\text{ in the set of
  \thmref{thm:VII}{VII}}).
  \tag{2.9.\ensuremath{\beta}}\label{eq:2.9.beta}
\end{equation}
(Concerning the exponent $\infty$, consider the remark at the end of
\secref{sec:2.5}{2.5}.)

\PaperSection{sec:2.10}{The fundamental group.}
Denote the set of all $\theta$ for which $\overline{\M}^{\theta}$ exists and
equals $\overline{\M}$ by
\begin{equation}
  \mathfrak G=\mathfrak G(\overline{\M}).
  \tag{2.10.\ensuremath{\alpha}}\label{eq:2.10.alpha}
\end{equation}
Obviously $1\in\mathfrak G$ and $\theta,\xi\in\mathfrak G$ imply
$\theta\xi\in\mathfrak G$.  Since
$(\overline{\M}^{\theta})^{1/\theta}$ exists and equals $\overline{\M}$,
$\overline{\M}^{\theta}=\overline{\M}$ implies
$\overline{\M}^{1/\theta}=\overline{\M}$, i.e.
$\theta\in\mathfrak G$ implies $1/\theta\in\mathfrak G$.  (All this is due
to \eqref{eq:2.9.beta} in \secref{sec:2.9}{2.9}.)

\begin{theorem}\phantomsection\label{thm:VIII}
The set $\mathfrak G$ of \eqref{eq:2.10.alpha} above is a subgroup of $P$,
the multiplication group of the real numbers $\theta$,
$0<\theta<\infty$.  We call $\mathfrak G$ the fundamental group of
$\overline{\M}$.
\end{theorem}

Some properties of $\mathfrak G$.
\begin{lemma}\phantomsection\label{lem:2.10.1}
$\overline{\M}^{\theta}=\overline{\M}^{\xi}$ (both sides are assumed to
exist) is equivalent to $\theta/\xi\in\mathfrak G$.
\end{lemma}
\begin{proof}
Sufficiency: $\theta/\xi\in\mathfrak G$ implies
$\overline{\M}^{\theta/\xi}=\overline{\M}$.  Hence
$\overline{\M}^{\theta}=(\overline{\M}^{\theta/\xi})^{\xi}
=\overline{\M}^{\xi}$.

Necessity: Assume $\overline{\M}^{\theta}=\overline{\M}^{\xi}$.
$(\overline{\M}^{\xi})^{1/\xi}$ exists and equals $\overline{\M}$.  Hence
$(\overline{\M}^{\theta})^{1/\xi}=\overline{\M}^{\theta/\xi}$ exists and
equals $\overline{\M}$.  Thus $\theta/\xi\in\mathfrak G$.
\end{proof}

\begin{lemma}\phantomsection\label{lem:2.10.2}
$\overline{\M}$, $\overline{\M}^{c}$ and all
$\overline{\M}^{\theta}$ have the same fundamental group.
\end{lemma}
\begin{proof}
Ad $\overline{\M}^{c}$.  Immediate by \eqref{eq:2.9.alpha}.

Ad $\overline{\M}^{\theta}$.  Immediate by
\lemref{lem:2.10.1}{2.10.1}.
\end{proof}

Before we conclude this chapter, we determine the behavior of the discrete
finite cases, i.e. the $(\mathrm I_n)$, $n=1,2,\ldots$.

Consider the ring $\boldsymbol\Theta$ of all operators $\alpha\cdot1$.  This
is obviously the prototype of all factors of the case $(\mathrm I_1)$, its
type $\overline{\boldsymbol\Theta}$ is that of the ring of all complex
numbers.

Now we have
\begin{lemma}\phantomsection\label{lem:2.10.3}
If $\M$ is in case $(\mathrm I_n)$, ($n=1,2,\ldots$), then
$\overline{\M}^{\alpha}$ ($\alpha$ finite) exists if and only if
$n\alpha=1,2,\ldots$, and then
$\overline{\M}^{\alpha}=\overline{\boldsymbol\Theta}^{,n\alpha}$.
\end{lemma}
\begin{proof}
The range asserted for $\alpha$ is the same as that of
\thmref{thm:VII}{VII}.  Put $n\alpha=k$.  Then
\[
  \overline{\M}^{\alpha}
  =\overline{\M}^{k/n}
  =\bigl(\overline{\M}^{1/n}\bigr)^{k}.
\]
Now consider $\overline{\M}^{1/n}$.  Apply the definition of
\secref{sec:2.8}{2.8}, preceding \thmref{thm:VII}{VII}.  For the $E$ in
question, $D_{\M}(E)/D_{\M}(1)$ assumes its minimal value, $1/n$.  Hence,
this $E$ is minimal (cf. \cite[pp.~143--144, Definition~5.1.2]{ref5}).
Consequently $\M^{E}$ is the set of all $\alpha E$ (cf.
\cite[p.~144, the last part of the proof of Lemma~5.1.3]{ref5}---or it may
be proved directly).  Hence $\M_{(\mathfrak M)}$ is the set of all
$\alpha\cdot1$ in $\mathfrak M$, i.e.
\[
  \overline{\M}^{1/n}=\overline{\M_{(\mathfrak M)}}
  =\overline{\boldsymbol\Theta}.
\]
Consequently
\[
  \overline{\M}^{\alpha}=\overline{\boldsymbol\Theta}^{,k},
  \qquad k=n\alpha,
\]
as desired.
\end{proof}

\begin{corollary}\phantomsection\label{cor:2.10.3}
If $\M$ is in a case $(\mathrm I_n)$ ($n=1,2,\ldots$) then
$\overline{\M}=\overline{\boldsymbol\Theta}^{,n}$.
\end{corollary}
\begin{proof}
Put $\alpha=1$ in the above lemma.
\end{proof}

\begin{lemma}\phantomsection\label{lem:2.10.4}
If $\M$ is in a case $(\mathrm I_n)$ ($n=1,2,\ldots$) then
$\overline{\M}^{c}=\overline{\M}$ and
$\mathfrak G=\mathfrak G(\overline{\M})=(1)$.
\end{lemma}
\begin{proof}\footnote{Both assertions could be proved directly by matrix
considerations.}
Ad $\overline{\M}^{c}=\overline{\M}$.  Clearly
$\overline{\boldsymbol\Theta}^{c}=\overline{\boldsymbol\Theta}$.  Hence the
corollary to \lemref{lem:2.10.3}{2.10.3}, and \eqref{eq:2.4.alpha} or
\eqref{eq:2.9.alpha} imply $\overline{\M}^{c}=\overline{\M}$.

Ad $\mathfrak G=(1)$.  By \lemref{lem:2.10.3}{2.10.3} and its corollary,
$\overline{\M}^{\alpha}$ is in case $(\mathrm I_k)$, $k=n\alpha$.  Hence
$\overline{\M}^{\alpha}=\overline{\M}$ implies $n\alpha=n$ or
$\alpha=1$.  So $\mathfrak G=(1)$.
\end{proof}

The validity or invalidity of the equation
$\overline{\M}^{c}=\overline{\M}$ and the fundamental group
$\mathfrak G=\mathfrak G(\overline{\M})$ are algebraical isomorphism
invariants of $\M$.  \lemref{lem:2.10.4}{2.10.4} showed us how they behave
when $\M$ is in a discrete finite case, i.e. in a case $(\mathrm I_n)$,
$n=1,2,\ldots$.  The really interesting factors are the remaining finite
ones, those in the continuous finite case $(\mathrm{II}_1)$.  We shall see
that they are not all isomorphic to each other, but we shall achieve this
differentiation with the help of other invariants.  (Cf. the end of
\secref{sec:5.6}{5.6} and the beginning of \secref{sec:6.1}{6.1}.)

For all $\M$ in case $(\mathrm{II}_1)$ for which we have succeeded in
settling this question, $\overline{\M}^{c}=\overline{\M}$ and
$\overline{\M}^{\theta}=\overline{\M}$ ($\theta$ in the set of
\thmref{thm:VII}{VII}), i.e. $\mathfrak G=P$, was found.  (Cf.
\thmref{thm:XV}{XV} and the end of \secref{sec:5.6}{5.6}.  Observe
\lemref{lem:2.10.4}{2.10.4} by comparison.)  There seems to be, however, no
reason to believe the general validity of these relations.  The general
behavior of the above invariants remains therefore an open question.

% Source PDF p. 27 (journal p. 742).
\PaperChapter{3}{chapter:III}{Characterization of All Spatial Types in Terms of Algebraical Types}

\PaperSection{sec:3.1}{}
We have classified all operator rings \(\M\) according to their types
\(\overline{\M}\), i.e. with respect to algebraical isomorphism. Now we shall
introduce a broader classification, to be called the \emph{genus}, each genus
consisting of one or more types. It will be necessary, however, to restrict this
new classification to factors \(\M\).

We need some auxiliary lemmas which lead up to the desired definition.
Throughout this section \(\M\) will be a factor in a Hilbert space \(\Hilbert\).
As before, we denote by \(D_{\M}(E)\) and \(D_{\M}(\mathfrak M)\) a fixed
relative dimension function of \(\M\) applicable to the \(E\in\M\) or the
\(\mathfrak M\eta\M\).

% Source PDF p. 28 (journal p. 743).
\begin{lemma}\phantomsection\label{lem:3.1.1}
If \(\mathfrak M\) is a closed linear set, \(\mathfrak M\eta\M\), and such
that \(\mathfrak M\ne(0)\) and
\(D_{\M}(\Hilbert\ominus\mathfrak M)=\infty\), then
\[
  \overline{\M}=\bigl(\overline{\M_{(\mathfrak M)}}\bigr)^{\infty}.
\]
\end{lemma}

\begin{proof}
Under our hypothesis, \(\mathfrak M\) ``divides''
\(\Hilbert\ominus\mathfrak M\) an infinite number of times, i.e. there are an
infinite number of pairwise orthogonal closed linear sets
\(\mathfrak M_2,\mathfrak M_3,\ldots\) such that
\(D_{\M}(\mathfrak M_i)=D_{\M}(\mathfrak M)\) and
\(\mathfrak M_i\subset\Hilbert\ominus\mathfrak M\) for
\(i=2,3,\ldots\). If \(\mathfrak M\) is relatively finite-dimensional, this is
obvious. If \(\mathfrak M\) is infinite-dimensional, we apply a variant of
\cite[p.~157, Lemma~7.2.3]{ref5} to \(\Hilbert\ominus\mathfrak M\), which
shows that \(\Hilbert\ominus\mathfrak M\) is determined by a denumerably
infinite number of pairwise orthogonal closed linear sets
\(\mathfrak M'_2,\mathfrak M'_3,\ldots\), each of which is relatively
infinite-dimensional. (To prove the variant, it is only necessary, in the proof
of \cite[Lemma~7.2.3]{ref5}, to divide the
\(\mathfrak P_1,\mathfrak P_2,\ldots\) into an infinite number of sets, each
having an infinite number of \(\mathfrak P_i\)'s, rather than into two sets.)

If we now apply \cite[p.~155, Lemma~7.1.2]{ref5} to \(\mathfrak M\) and
\(\Hilbert\ominus\mathfrak M\), we see that
\(\Hilbert\ominus\mathfrak M\) is determined by a denumerably infinite
number of sets \(\mathfrak M_2,\mathfrak M_3,\ldots\), which are pairwise
orthogonal and for which
\(D_{\M}(\mathfrak M_i)=D_{\M}(\mathfrak M)\), for
\(i=2,3,\ldots\). Let \(\mathfrak M=\mathfrak M_1\).

The second paragraph of the proof of \lemref{lem:2.6.2}{2.6.2} can now be
used to show the existence in \(\M\) of a set of \(\infty\)-matrix units
\(W_{\mu,\nu}\), with
\(W_{\mu,\mu}=P_{\mathfrak M_\mu}\),
\(W_{1,1}=P_{\mathfrak M_1}=P_{\mathfrak M}\). For the hypothesis
\(p<\infty\) is not used here. Furthermore,
\[
\begin{aligned}
E_0&=\sum_{\mu=1}^{\infty}W_{\mu,\mu}
 =W_{1,1}+\sum_{\mu=2}^{\infty}W_{\mu,\mu}\\
&=P_{\mathfrak M_1}+\sum_{\mu=2}^{\infty}P_{\mathfrak M_\mu}
 =P_{\mathfrak M}+P_{\Hilbert\ominus\mathfrak M}=1.
\end{aligned}
\]
Thus \thmref{thm:V}{V} in \secref{sec:2.6}{2.6} above yields the desired
result.
\end{proof}

\begin{lemma}\phantomsection\label{lem:3.1.2}
If \(\mathfrak M\) is a closed linear set, \(\mathfrak M\eta\M\) and
\(\mathfrak M\ne(0)\), and if \(\M\) is in an infinite case, then
\[
  \overline{\M}=\bigl(\overline{\M_{(\mathfrak M)}}\bigr)^{\infty}.
\]
\end{lemma}

\begin{proof}
If \(D_{\M}(\Hilbert\ominus\mathfrak M)=\infty\), the result is implied by
\lemref{lem:3.1.1}{3.1.1}. Hence we may assume that
\(D_{\M}(\Hilbert\ominus\mathfrak M)\) is finite. Since \(\M\) is in an
infinite case, \(D_{\M}(\Hilbert)\) is infinite and hence
\[
  D_{\M}(\mathfrak M)=D_{\M}(\Hilbert)
    -D_{\M}(\Hilbert\ominus\mathfrak M)
\]
is infinite also.

Since \(D_{\M}(\Hilbert)\) is infinite, there exists by
\cite[p.~157, Lemma~7.2.3]{ref5} a closed linear set \(\mathfrak N\) such
that \(D_{\M}(\mathfrak N)\) and
\(D_{\M}(\Hilbert\ominus\mathfrak N)\) are also infinite. Hence
\(\overline{\M_{(\mathfrak N)}}=\overline{\M_{(\mathfrak M)}}\) by
\lemref{lem:2.8.1}{2.8.1} (we need only the algebraical, not the spatial,
isomorphism), and
\(\overline{\M}=\bigl(\overline{\M_{(\mathfrak N)}}\bigr)^\infty\) by
\lemref{lem:3.1.1}{3.1.1}. These give together
\(\overline{\M}=\bigl(\overline{\M_{(\mathfrak M)}}\bigr)^\infty\).
\end{proof}

\begin{lemma}\phantomsection\label{lem:3.1.3}
A factor \(\M\) is in an infinite case if and only if
\[
  \overline{\M}=\overline{\M}^{\infty}.
\]
\end{lemma}

\begin{proof}
Sufficiency: Immediate by (iii) in \lemref{lem:2.7.1}{2.7.1}.

Necessity: Put \(\mathfrak M=\Hilbert\) in \lemref{lem:3.1.2}{3.1.2}.
\end{proof}

% Source PDF p. 29 (journal p. 744).
\begin{definition}\phantomsection\label{def:3.1.1}
If, for two factors, \(\M,\N\),
\[
  \overline{\N}=\overline{\M_{(\mathfrak M)}}
\]
for a suitable closed set \(\mathfrak M\ne(0)\) and
\(\mathfrak M\eta\M\) in the Hilbert space \(\Hilbert\) of \(\M\), then we
say that \(\overline{\M}\) is a \emph{multiple} of \(\overline{\N}\), or that
\(\overline{\N}\) is a \emph{divisor} of \(\overline{\M}\).
\end{definition}

It is clear that this relationship concerns \(\overline{\M}\) and
\(\overline{\N}\). Cf. the argument immediately preceding
\thmref{thm:VI}{VI} in \secref{sec:2.8}{2.8}. It is also obviously transitive.

\begin{lemma}\phantomsection\label{lem:3.1.4}
\(\overline{\N}\) is a divisor of \(\overline{\M}\) if and only if
\begin{enumerate}[label={}]
\item[\(\alpha\))] \(\overline{\M}=\overline{\N}\), when \(\overline{\N}\) is in an
infinite case;
\item[\(\beta\))] \(\overline{\M}=\overline{\N}^{\theta}\), with a \(\theta\geq1\) in
the range of \thmref{thm:VII}{VII} or of \thmref{thm:IV-prime}{IV$'$}
(i.e. \(\theta=\infty\)), when \(\overline{\N}\) is in a finite case.
\end{enumerate}
\end{lemma}

\begin{proof}
Ad \(\alpha\)). Sufficiency: Obvious.

Necessity: Since
\(\overline{\N}=\overline{\M_{(\mathfrak M)}}\) is in an infinite case,
\(D_{\M}(\mathfrak M)\) is infinite by
\cite[p.~189, Lemma~11.4.3]{ref5}. Hence \(\overline{\M}\) is in an infinite
case and thus \lemref{lem:3.1.2}{3.1.2} gives
\[
  \overline{\M}
  =\bigl(\overline{\M_{(\mathfrak M)}}\bigr)^\infty
  =\overline{\N}^{\infty}.
\]
On the other hand, \(\overline{\N}=\overline{\N}^{\infty}\) by
\lemref{lem:3.1.3}{3.1.3}. Thus \(\overline{\M}=\overline{\N}\).

Ad \(\beta\)). Sufficiency: \(\theta\) finite: Then
\(\overline{\N}=\overline{\M}^{\alpha}\), with
\(\alpha=1/\theta\leq1\). Hence the assertion follows by our definition of
\(\overline{\M}^{\alpha}\) for \(\alpha\leq1\) in
\secref{sec:2.8}{2.8}. \(\theta\) infinite: Immediate by
\thmref{thm:V}{V}.

Necessity: \(\overline{\M}\) in a finite case: Clearly
\(\overline{\N}=\overline{\M_{(\mathfrak M)}}=\overline{\M}^{\alpha}\),
with \(\alpha=D_{\M}(E)/D_{\M}(1)\leq1\),
\(E=P_{\mathfrak M}\). Hence
\(\overline{\M}=\overline{\N}^{\theta}\), with
\(\theta=1/\alpha\geq1\). \(\overline{\M}\) in an infinite case:
\(\overline{\M}=\bigl(\overline{\M_{(\mathfrak M)}}\bigr)^\infty
=\overline{\N}^{\infty}\) by \lemref{lem:3.1.2}{3.1.2}.
\end{proof}

\begin{lemma}\phantomsection\label{lem:3.1.5}
The four following statements are equivalent to each other:
\begin{enumerate}[label={}]
\item[\(\alpha\))] \(\overline{\M}\) is either a multiple or divisor of
\(\overline{\N}\);
\item[\(\beta\))] \(\overline{\M}\) and \(\overline{\N}\) possess a common multiple;
\item[\(\gamma\))] \(\overline{\M}\) and \(\overline{\N}\) possess a common divisor;
\item[\(\delta\))] \(\overline{\M}^{\infty}=\overline{\N}^{\infty}\).
\end{enumerate}
\end{lemma}

\begin{proof}
The implications
\[
  \alpha)\longrightarrow\beta),\qquad
  \alpha)\longrightarrow\gamma)
\]
are obvious. Furthermore,
\[
  \delta)\longrightarrow\beta)
\]
follows from \thmref{thm:V}{V}, since this theorem shows that
\(\overline{\M}^{\infty}=\overline{\N}^{\infty}\) is a multiple of
\(\overline{\M}\) and of \(\overline{\N}\).

We prove next
\[
  \beta)\longrightarrow\alpha).
\]
Let \(\overline{\mathbf L}\) be the common multiple of both \(\overline{\M}\) and
\(\overline{\N}\). Let closed linear sets \(\mathfrak M\) and
\(\mathfrak N\) be chosen in the \(\Hilbert\) associated with \(\mathbf L\), so that
\(\overline{\M}=\overline{\mathbf L_{(\mathfrak M)}}\) and
\(\overline{\N}=\overline{\mathbf L_{(\mathfrak N)}}\). Since \(\mathbf L\) is a factor,
either \(\mathfrak M\) has the same relative dimension (in \(\mathbf L\)) as a
\(\mathfrak N'\subseteq\mathfrak N\), or \(\mathfrak N\) has the same relative
dimension as a \(\mathfrak M'\subseteq\mathfrak M\). Cf.
\cite[p.~153, Lemma~6.2.3, or p.~154, Theorem~VI]{ref5}. By symmetry we may
assume the former. By \lemref{lem:2.8.1}{2.8.1} we may now replace
\(\mathfrak M\) by \(\mathfrak N'\), i.e. we can assume that
\(\mathfrak M\subseteq\mathfrak N\).

% Source PDF p. 30 (journal p. 745).
Thus
\[
  \overline{\bigl(\mathbf L_{(\mathfrak N)}\bigr)_{(\mathfrak M)}}
  =\overline{\mathbf L_{(\mathfrak M)}},
\]
so that \(\overline{\mathbf L_{(\mathfrak N)}}\) is a multiple of
\(\overline{\mathbf L_{(\mathfrak M)}}\). Hence \(\overline{\N}\) is a multiple of
\(\overline{\M}\).

Finally,
\[
  \gamma)\longrightarrow\delta).
\]
Let \(\overline{\mathbf K}\) be a common divisor of \(\overline{\M}\) and
\(\overline{\N}\). Since \(\overline{\M}\) is a divisor of
\(\overline{\M}^{\infty}\) by \thmref{thm:V}{V}, so
\(\overline{\mathbf K}\) is a divisor of \(\overline{\M}^{\infty}\) too. Now
\(\overline{\M}^{\infty}\) is in an infinite case by (iii) in
\lemref{lem:2.7.1}{2.7.1}. Hence
\(\overline{\mathbf K}^{\infty}=\overline{\M}^{\infty}\) by
\lemref{lem:3.1.2}{3.1.2}. Similarly
\(\overline{\mathbf K}^{\infty}=\overline{\N}^{\infty}\). Consequently
\(\overline{\M}^{\infty}=\overline{\N}^{\infty}\).

All these equivalences give together
\[
  \alpha)\longrightarrow\gamma)\longrightarrow\delta)
  \longrightarrow\beta)\longrightarrow\alpha),
\]
thus completing the proof.
\end{proof}

\begin{definition}\phantomsection\label{def:3.1.2}
If the equivalent conditions of \lemref{lem:3.1.5}{3.1.5} are satisfied, then
we say that \(\overline{\M},\overline{\N}\) are \emph{commensurable}.
\end{definition}

Condition \(\delta\)) of \lemref{lem:3.1.5}{3.1.5} implies that the notion of
commensurability is reflexive, symmetric and transitive.

\begin{lemma}\phantomsection\label{lem:3.1.6}
\(\overline{\M},\overline{\N}\) are commensurable if and only if
\begin{enumerate}[label={}]
\item[\(\alpha\))] \(\overline{\M}=\overline{\N}\), when \(\overline{\M},\overline{\N}\)
are both in infinite cases;
\item[\(\beta\))] \(\overline{\M}=\overline{\N}^{\infty}\), when \(\overline{\M}\) is
in an infinite case and \(\overline{\N}\) in a finite case;
\item[\(\gamma\))] \(\overline{\M}^{\infty}=\overline{\N}\), when \(\overline{\M}\) is
in a finite case and \(\overline{\N}\) in an infinite case;
\item[\(\delta\))] \(\overline{\M}=\overline{\N}^{\theta}\), with a \(\theta\) in the
range of \thmref{thm:VII}{VII}, when \(\overline{\N}\) and
\(\overline{\M}\) are both in finite cases.
\end{enumerate}
\end{lemma}

\begin{proof}
Ad \(\alpha\))--\(\gamma\)). Condition \(\delta\)) of
\lemref{lem:3.1.5}{3.1.5} for commensurability is
\(\overline{\M}^{\infty}=\overline{\N}^{\infty}\).
\lemref{lem:3.1.3}{3.1.3} shows that when \(\overline{\M}\) is in an infinite
case we may substitute \(\overline{\M}\) for \(\overline{\M}^{\infty}\) in
this condition, and correspondingly for \(\overline{\N}\). When
\(\overline{\M}\) and \(\overline{\N}\) are both infinite, this yields
\(\alpha\)), while if only one is infinite we obtain \(\beta\)) or
\(\gamma\)).

Ad \(\delta\)). By \lemref{lem:3.1.4}{3.1.4}, \(\beta\)),
\(\overline{\M}=\overline{\N}^{\theta}\), for \(\theta\geq1\) in the range
of \thmref{thm:VII}{VII}, is equivalent to \(\overline{\N}\) a divisor of
\(\overline{\M}\). Since \(\overline{\M}\) is finite, \(\theta\) is finite
by \lemref{lem:2.7.1}{2.7.1} (iii).

On the other hand, \(\overline{\M}=\overline{\N}^{\theta}\), for a
\(\theta\leq1\) in the range of \thmref{thm:VII}{VII} (for
\(\overline{\N}\)), is equivalent to
\(\overline{\M}^{1/\theta}=\overline{\N}\), with \(1/\theta\) in the range
of \thmref{thm:VII}{VII} (for \(\overline{\M}\)). (Cf.
\eqref{eq:2.9.beta}. That \(\theta\) is in the correct sets in the discrete
cases is readily shown by reference to \lemref{lem:2.10.3}{2.10.3} and its
corollary.) \lemref{lem:3.1.4}{3.1.4}, with \(\overline{\M}\) and
\(\overline{\N}\) interchanged, shows that this is equivalent to
\(\overline{\N}\) is a multiple of \(\overline{\M}\).

Thus \(\overline{\M}=\overline{\N}^{\theta}\), for a \(\theta\) in the range
of \thmref{thm:VII}{VII}, is equivalent to \(\overline{\N}\) is a divisor of
\(\overline{\M}\) or \(\overline{\N}\) is a multiple of
\(\overline{\M}\), i.e. to the condition of \lemref{lem:3.1.5}{3.1.5} for
commensurability.
\end{proof}

We call the abstraction of type \(\overline{\M}\) (for factors) with respect
to commensurability its \emph{genus}, i.e. two types
\(\overline{\M}_1\) and \(\overline{\M}_2\) will be said to have the same
genus if and only if they are commensurable. We denote the genus of the type
\(\overline{\M}\) by \(\widetilde{\overline{\M}}\). Notions invariant under
commensurability will be said to be notions concerning the genus
\(\widetilde{\overline{\M}}\) (and not \(\overline{\M}\) or \(\M\) itself).

Thus the cases I, II, III are notions concerning the genus. This follows from
(ii) in \lemref{lem:2.7.1}{2.7.1} (with \(p=\infty\)) and the
characterization \(\delta\)) in \lemref{lem:3.1.5}{3.1.5}.

% Source PDF p. 31 (journal p. 746).
Similarly \(\overline{\M}^{,c}\) is an operation concerning the genus. This
follows from \eqref{eq:2.5.epsilon} and the above-mentioned characterization.
Therefore we can define
\[
  \widetilde{\overline{\M}}^{,c}
  =\widetilde{\overline{\M}^{,c}}.
\]

Observe finally that \(\overline{\M_{(\mathfrak M)}}\),
\(\overline{\M}^{,p}\), \(p=1,2,\ldots,\infty\), and (when it is defined,
i.e. for \(\overline{\M}\) in a finite case)
\(\overline{\M}^{\,\theta}\) (\(\theta\) in the range of
\thmref{thm:VII}{VII}) have the same genus as \(\overline{\M}\) (for
\(\overline{\M_{(\mathfrak M)}}\) this follows from \(\alpha\)) in
\lemref{lem:3.1.5}{3.1.5}; for \(\overline{\M}^{,p}\) this is implied by
\thmref{thm:V}{V} and \lemref{lem:3.1.5}{3.1.5}, \(\alpha\)); and for
\(\overline{\M}^{\,\theta}\) this follows from
\lemref{lem:3.1.6}{3.1.6}, \(\delta\)).)

It is clear then that being in a finite or an infinite case does not concern
the genus (cf. \thmref{thm:IX}{IX} for details).

\PaperSection{sec:3.2}{}
Since the cases (I), (II), (III) are notions concerning the genus, each one of
these cases is the sum of one or more genera. And of course each genus is the
sum of one or more types. These are the details.

\begin{theorem}\phantomsection\label{thm:IX}
\begin{enumerate}[label={}]
\item[\(\alpha\))] Case (I) consists of exactly one genus. This genus contains exactly one
type in an infinite case: \((\mathrm I_\infty)\); and for every positive integer
\(n\) exactly one type \((\mathrm I_n)\), these being all its types in finite
cases.
\item[\(\beta\))] Case (II) consists of one or more genera. Each such genus contains
exactly one type in an infinite case \((\mathrm {II}_\infty)\) and one or more
types in finite cases \((\mathrm {II}_1)\).
\item[\(\gamma\))] Case (III) consists of one or more genera. Each such genus contains
exactly one type which is, of course, in an infinite case
\((\mathrm {III}_\infty)\).
\end{enumerate}
\end{theorem}

\begin{proof}
Ad \(\alpha\)). Each one of the cases
\((\mathrm I_1),(\mathrm I_2),\ldots,(\mathrm I_\infty)\) contains exactly
one type by \cite[p.~173, Lemma~8.6.1]{ref5}. They can obviously all be
obtained from the last one by the operations \(\M_{(\mathfrak M)}\). Hence
they are all of the same genus as this last by \(\alpha\)) of
\lemref{lem:3.1.5}{3.1.5} and \defref{def:3.1.1}{3.1.1}.

Ad \(\beta\)). Consider a genus \(\widetilde{\overline{\M}}\) of case (II).
Then there exists an \(\overline{\M_{(\mathfrak M)}}\) in a finite case, for
any \(\mathfrak M\) with \(D_{\M}(\mathfrak M)<\infty\) will do. Thus
\(\overline{\M_{(\mathfrak M)}}\) is finite and it belongs to the given genus
of \(\overline{\M}\) by condition \(\alpha\)) of
\lemref{lem:3.1.5}{3.1.5}. \(\overline{\M}^{\infty}\) belongs also to this
genus for the same reason, since \(\overline{\M}\) is a divisor of
\(\overline{\M}^{\infty}\) by \thmref{thm:V}{V} of
\secref{sec:2.6}{2.6}. The infinite case in a particular genus is unique by
\(\alpha\)) of \lemref{lem:3.1.6}{3.1.6}.

Ad \(\gamma\)). Consider a genus \(\widetilde{\overline{\M}}\) of case
\((\mathrm {III}_\infty)\). Then \(\overline{\M}\) must be in an infinite
case, and consequently it is unique by \(\alpha\)) of
\lemref{lem:3.1.6}{3.1.6}.

Observe that none of the three cases (I), (II) or (III) is empty. Cf.
\cite[p.~94, \S1, in the introduction]{ref8}.
\end{proof}

\begin{lemma}\phantomsection\label{lem:3.2.1}
Consider a genus \(\widetilde{\overline{\M}}\) which contains types
\(\overline{\M}\) in finite cases, i.e. one in the cases (I), (II) (cf.
\thmref{thm:IX}{IX} above). Then the fundamental group
\(\mathfrak G(\overline{\M})\) is a notion concerning the genus
\(\widetilde{\overline{\M}}\), i.e. it is the same for all the types
\(\overline{\M}\) of the genus which are in a finite case.
\end{lemma}

\begin{proof}
Immediate by \(\delta\)) in \lemref{lem:3.1.6}{3.1.6} and
\lemref{lem:2.10.2}{2.10.2}.
\end{proof}

We shall therefore denote the fundamental group from now on by
\(\mathfrak G(\widetilde{\overline{\M}})\). Some simple properties of
\(\mathfrak G(\widetilde{\overline{\M}})\) follow.

\begin{lemma}\phantomsection\label{lem:3.2.2}
Consider a genus \(\widetilde{\overline{\M}}\) as above. Form the ring
\(\M\) and two closed linear sets \(\mathfrak M,\mathfrak N\eta\M\), which
are both \(\ne(0)\) and of finite relative dimension. Then we have
\begin{enumerate}[label=(\roman*)]
\item
\[
  \overline{\M_{(\mathfrak M)}}
  =\overline{\M_{(\mathfrak N)}}^{\,D_{\M}(\mathfrak M)/D_{\M}(\mathfrak N)};
\]
\item
\[
  \overline{\M_{(\mathfrak M)}}=\overline{\M_{(\mathfrak N)}}
  \quad\text{if and only if}\quad
  \frac{D_{\M}(\mathfrak M)}{D_{\M}(\mathfrak N)}
  \in\mathfrak G(\widetilde{\overline{\M}}).
\]
\end{enumerate}
\end{lemma}

% Source PDF p. 32 (journal p. 747).
\begin{proof}
Ad (i). Suppose first that
\(D_{\M}(\mathfrak M)\leq D_{\M}(\mathfrak N)\). Here
\(D_{\M}(\mathfrak M)=D_{\M}(\mathfrak N')\) for a
\(\mathfrak N'\subseteq\mathfrak N\), and since by
\lemref{lem:2.8.1}{2.8.1}
\(\overline{\M_{(\mathfrak N')}}=\overline{\M_{(\mathfrak M)}}\), we may
suppose \(\mathfrak M\subseteq\mathfrak N\). If \(\mathfrak P\) is a set
\(\eta\M\) and \(\subseteq\mathfrak N\), \(D_{\M}(\mathfrak P)\) is a
dimension function for \(\M_{(\mathfrak N)}\). Since
\(\mathfrak M\subseteq\mathfrak N\), by definition
\[
  \overline{\M_{(\mathfrak M)}}
  =\bigl(\overline{\M_{(\mathfrak N)}}\bigr)^\alpha,
  \qquad
  \alpha=\frac{D_{\M}(\mathfrak M)}{D_{\M}(\mathfrak N)}.
\]
(Cf. \secref{sec:2.8}{2.8}; \(\Hilbert,\M\) are replaced by
\(\mathfrak N,\M_{(\mathfrak N)}\).)

On the other hand, if
\(D_{\M}(\mathfrak M)>D_{\M}(\mathfrak N)\), the above argument yields
\[
  \bigl(\overline{\M_{(\mathfrak M)}}\bigr)^{1/\alpha}
  =\overline{\M_{(\mathfrak N)}}.
\]
Equation \eqref{eq:2.9.beta} can now be used to show that
\(\overline{\M_{(\mathfrak M)}}
=\bigl(\overline{\M_{(\mathfrak N)}}\bigr)^\alpha\), suitable use of
\lemref{lem:2.10.3}{2.10.3} and its corollary being made in the discrete case.

Ad (ii). By the definition of
\(\mathfrak G(\overline{\M_{(\mathfrak N)}})\), (i) implies
\(\overline{\M_{(\mathfrak M)}}=\overline{\M_{(\mathfrak N)}}\) if and
only if \(\alpha\in\mathfrak G(\overline{\M_{(\mathfrak N)}})\). But
\(\mathfrak G(\overline{\M_{(\mathfrak N)}})
=\mathfrak G(\widetilde{\overline{\M}})\) by
\lemref{lem:3.2.1}{3.2.1}.
\end{proof}

\begin{lemma}\phantomsection\label{lem:3.2.3}
Consider a genus \(\widetilde{\overline{\M}}\) as above. Then it contains
precisely one type in a finite case if and only if
\(\mathfrak G(\widetilde{\overline{\M}})=P\) (cf.
\thmref{thm:VIII}{VIII}).
\end{lemma}

\begin{proof}
Sufficiency: Immediate by \(\delta\)) in \lemref{lem:3.1.6}{3.1.6}.

Necessity: Assume \(\mathfrak G(\widetilde{\overline{\M}})\ne P\).
Consider an \(\overline{\M}\) of this genus in a finite case. If we find a
\(\theta\notin\mathfrak G(\widetilde{\overline{\M}})\) which is in the range
of \thmref{thm:VII}{VII}, then \(\overline{\M}^{\theta}\) has the same genus,
is also in a finite case and yet \(\ne\overline{\M}\), thus completing the
proof.

Now in case (I),
\(\mathfrak G(\widetilde{\overline{\M}})=(1)\) by
\lemref{lem:2.10.4}{2.10.4}. Hence any \(\theta=2,3,\ldots\) will do. And in
case (II), \(P\) is the range of \thmref{thm:VII}{VII}, so
\(\mathfrak G(\widetilde{\overline{\M}})\ne P\) guarantees the existence of
the desired \(\theta\).
\end{proof}

\PaperSection{sec:3.3}{}
The moment has now come when it is opportune to introduce the notion of a
spatial type (cf. footnote~12). We shall call the abstraction of \(\M\) with
respect to spatial isomorphism its spatial type. Thus two rings \(\M_1\) and
\(\M_2\), in two Hilbert spaces \(\Hilbert_1\) and \(\Hilbert_2\),
respectively, will be said to have the same spatial type if and only if they are
spatially isomorphic. We denote the spatial type of the ring \(\M\) by
\(\widetilde{\M}\).

We shall now answer Question II in \S2 in the introduction, for the cases (I)
and (II)\footnote{In the cases (I) of course the answer is known. Nevertheless
we include them for the reason given in footnote~1.}; i.e. determine when an
algebraic isomorphism determines a spatial one. Thus we shall establish the
relationship of \(\widetilde{\M}\) to \(\overline{\M}\). For the cases (I)
and (II) we shall determine all spatial types \(\widetilde{\M}\) in terms of
the algebraical types \(\overline{\M}\).

The notions \(\overline{\M}^{\,\theta}\),
\(\mathfrak G(\widetilde{\overline{\M}})\), as well as the genus
\(\widetilde{\overline{\M}}\), will be basic in these considerations.

A spatial isomorphism of \(\M_1\) in \(\Hilbert_1\) and \(\M_2\) in
\(\Hilbert_2\) (cf. Definition A in \S2 in the Introduction) will also carry
\(\M'_1\) into \(\M'_2\). So the spatial type \(\widetilde{\M}\) of a ring
\(\M\) determines not only the algebraical type \(\overline{\M}\), but also
the algebraical type \(\overline{\M'}\). Now \(\overline{\M}\) does not
determine \(\overline{\M'}\) uniquely, not even for a factor
\(\M\)\footnote{The direct factors described in \cite[p.~139, or p.~173,
Lemma~8.6.1]{ref5} give elementary examples of this. The exhaustive results
in this respect are contained in the last part of our \thmref{thm:X}{X}.}.
Consequently a discussion of the spatial type \(\widetilde{\M}\) of a factor
\(\M\) must begin with an investigation of the relation between
\(\overline{\M}\) and \(\overline{\M'}\). At this point the notion of genus
comes in.

% Source PDF p. 33 (journal p. 748).
\begin{lemma}\phantomsection\label{lem:3.3.1}
For every factor \(\M\) in the cases (I) and (II),
\[
  \overline{\M'}=\overline{\M}^{\,c}.
\]
\end{lemma}

\begin{proof}
Due to our assumptions concerning \(\M\), there exists a closed linear set
\(\mathfrak M\eta\M\) which is not \(=(0)\) and has a finite
\(D_{\M}(\mathfrak M)\). Choose an \(f\in\mathfrak M\), with \(f\ne0\),
and form \(\mathfrak M_f^{\M'}\), \(\mathfrak M_f^{\M}\) by
\cite[p.~143, Def.~5.1.1]{ref5}. These are closed linear sets, clearly
\(\eta\M\) and \(\eta\M'\), respectively, both \(\ne(0)\), and
\(\mathfrak M_f^{\M'}\subseteq\mathfrak M\), since
\(f\in\mathfrak M\eta\M\). Hence \(D_{\M}(\mathfrak M_f^{\M'})\) is
finite.

We now normalize \(D_{\M}\) so as to make
\(D_{\M}(\mathfrak M_f^{\M'})=1\), and then \(D_{\M'}\) so as to make
\(C=1\) for the \(C\) of \cite[p.~182, Theorem~X]{ref5}. Hence
\begin{equation}
  D_{\M}(\mathfrak M_f^{\M'})
  =D_{\M'}(\mathfrak M_f^{\M})=1.
  \tag{3.3.\ensuremath{\alpha}}\label{eq:3.3.alpha}
\end{equation}
For the sake of brevity we write
\begin{equation}
  \mathfrak M_1=\mathfrak M_f^{\M'},
  \qquad
  \mathfrak M'_1=\mathfrak M_f^{\M}.
  \tag{3.3.\ensuremath{\beta}}\label{eq:3.3.beta}
\end{equation}

Let us now use the considerations of \cite[pp.~188--190, \S11.4]{ref5}. We
form, in accord with them,
\(\M_{(\mathfrak M_1\cdot\mathfrak M'_1)}\),
\(\M'_{(\mathfrak M_1\cdot\mathfrak M'_1)}\). These are coupled factors in
\(\mathfrak M_1\cdot\mathfrak M'_1\ne(0)\), just as \(\M,\M'\) are in
\(\Hilbert\). (Cf. the remarks, loc. cit., at the beginning of \S11.4, and
immediately preceding Lemma~11.4.3.) We can use \(D_{\M}\) and
\(D_{\M'}\) to obtain relative dimension functions in
\(\M_{(\mathfrak M_1\cdot\mathfrak M'_1)}\),
\(\M'_{(\mathfrak M_1\cdot\mathfrak M'_1)}\), respectively, as described in
Lemma~11.4.2, loc. cit. Our equations \eqref{eq:3.3.alpha} and
\eqref{eq:3.3.beta} are now seen to imply:

First, the relative dimension of
\(\mathfrak M_1\cdot\mathfrak M'_1\) is 1 for
\(\M_{(\mathfrak M_1\cdot\mathfrak M'_1)}\) and for
\(\M'_{(\mathfrak M_1\cdot\mathfrak M'_1)}\). Hence both factors are in a
finite case. Besides \(C=1\) (cf. above). Hence they are either both in a case
\((\mathrm I_n)\), \(n=1,2,\ldots\), with \(1/n\) times the standard
normalization, or both in case \((\mathrm {II}_1)\), with the standard
normalization.\footnote{This discrepancy is of course due to the different
ways in which we defined the standard normalization in the cases (I) and (II)
(cf. \cite[p.~172, Theorem~VIII]{ref5}).}

Second, they fulfill the requirements of \cite[p.~235, beginning of
\S4.1]{ref6}, on which the Theorem~VI on p.~239 eod. is based.\footnote{\hypertarget{fn:21}{}The
theorem referred to is stated for the case \((\mathrm {II}_1)\) only. It holds
however for the cases \((\mathrm I_n)\), \(n=1,2,\ldots\), too, and with
precisely the same proof. Our remark in footnote~1 could have been applied to
\cite{ref6} also.} Now this theorem states that the coupled factors to which it
applies are dual isomorphic.\footnote{\hypertarget{fn:22}{}Called anti-isomorphic, loc. cit.} Again,
by \cite[p.~188, Lemma~11.4.1]{ref5},
\(\M_{(\mathfrak M_1)}\),
\(\M_{(\mathfrak M_1\cdot\mathfrak M'_1)}\) are algebraically ring
isomorphic and, similarly,
\(\M'_{(\mathfrak M'_1)}\),
\(\M'_{(\mathfrak M_1\cdot\mathfrak M'_1)}\). Thus
\(\M_{(\mathfrak M_1)}\), \(\M'_{(\mathfrak M'_1)}\) are dual isomorphic,
i.e.
\begin{equation}
  \overline{\M'_{(\mathfrak M'_1)}}
  =\overline{\M_{(\mathfrak M_1)}}^{\,c}.
  \tag{3.3.\ensuremath{\gamma}}\label{eq:3.3.gamma}
\end{equation}
Consequently,
\[
  \overline{\M'}
  =\overline{\M'_{(\mathfrak M'_1)}}
  =\overline{\M_{(\mathfrak M_1)}}^{\,c}
  =\overline{\M}^{\,c},
\]
as desired.
\end{proof}

This lemma determines the direction of further analysis of our present topic.
It is easy to obtain now more detailed information.

% Source PDF p. 34 (journal p. 749).
For this purpose we define

\begin{definition}\phantomsection\label{def:3.3.1}
For every factor \(\M\) in the cases (I) and (II) we form the number
\begin{equation}
  \theta=\frac{1}{C}\frac{D_{\M'}(\Hilbert)}{D_{\M}(\Hilbert)}
  \tag{3.3.\ensuremath{\delta}}\label{eq:3.3.delta}
\end{equation}
with any normalization of \(D_{\M},D_{\M'}\), and the corresponding \(C\)
of \cite[p.~182, Theorem~X]{ref5}.
\end{definition}

The expression for \(\theta\) is clearly independent of the normalization of
\(D_{\M},D_{\M'}\). Since
\(0<C<\infty\), \(0<D_{\M}(\Hilbert),D_{\M'}(\Hilbert)\leq\infty\),
\eqref{eq:3.3.delta} characterizes \(\theta\) as a well-defined number with
\begin{equation}
  0\leq\theta\leq\infty,
  \tag{3.3.\ensuremath{\epsilon}}\label{eq:3.3.epsilon}
\end{equation}
except when \(D_{\M}(\Hilbert)=D_{\M'}(\Hilbert)=\infty\), i.e. when
\(\M,\M'\) both belong to infinite cases. If this happens, we construe
\eqref{eq:3.3.delta} to mean
\begin{equation}
  \theta=\frac{\infty}{\infty},
  \tag{3.3.\ensuremath{\zeta}}\label{eq:3.3.zeta}
\end{equation}
the symbol \(\infty/\infty\) being considered as an entity different from all
numbers of \eqref{eq:3.3.epsilon}.

\begin{lemma}\phantomsection\label{lem:3.3.2}
Let \(\M\) and \(\theta\) be as above. Then we have
\begin{enumerate}[label={}]
\item[\(\alpha\))] If \(\M,\M'\) are both in infinite cases, then
\[
  \overline{\M'}=\overline{\M}^{\,c}
  \quad\text{and}\quad \theta=\frac{\infty}{\infty}.
\]
\item[\(\beta\))] If \(\M\) is in an infinite case, and \(\M'\) is in a finite case,
then
\[
  \overline{\M}=\bigl(\overline{\M'}^{\,c}\bigr)^\infty
  \quad\text{and}\quad \theta=0,\quad\text{so }1/\theta=\infty.
\]
\item[\(\gamma\))] If \(\M\) is in a finite case and \(\M'\) is in an infinite case,
then
\[
  \overline{\M'}=\bigl(\overline{\M}^{\,c}\bigr)^\infty
  \quad\text{and}\quad \theta=\infty.
\]
\item[\(\delta\))] If \(\M,\M'\) are both in finite cases, then
\[
  \overline{\M'}=\bigl(\overline{\M}^{\,c}\bigr)^\theta,
  \quad\text{or equivalently}\quad
  \overline{\M}=\bigl(\overline{\M'}^{\,c}\bigr)^{1/\theta},
  \qquad 0<\theta<\infty.
\]
\end{enumerate}
\end{lemma}

\begin{proof}
All assertions concerning \(\theta\) are immediate by
\defref{def:3.3.1}{3.3.1}. Let us now consider the other statements.

Ad \(\alpha\)), \(\beta\)), \(\gamma\)). \lemref{lem:3.3.1}{3.3.1}
states that \(\overline{\M'}\) and \(\overline{\M}^{\,c}\) are
commensurable. Hence our \(\alpha\)), \(\beta\)), \(\gamma\)) follow from
\(\alpha\)), \(\beta\)), \(\gamma\)) in \lemref{lem:3.1.6}{3.1.6},
respectively.

Ad \(\delta\)). The equivalence of these two equations is obvious. We
consider the first one.

We choose the normalization of \(D_{\M},D_{\M'}\) exactly as in the proof of
\lemref{lem:3.3.1}{3.3.1}. Now \eqref{eq:3.3.alpha},
\eqref{eq:3.3.beta} in that proof give
\[
  \overline{\M_{(\mathfrak M_1)}}
    =\overline{\M}^{,1/D_{\M}(\Hilbert)},
  \qquad
  \overline{\M'_{(\mathfrak M'_1)}}
    =\overline{\M'}^{,1/D_{\M'}(\Hilbert)}.
\]

% Source PDF p. 35 (journal p. 750).
Then \eqref{eq:3.3.gamma} and \eqref{eq:2.9.alpha} yield
\[
  \overline{\M'}^{,1/D_{\M'}(\Hilbert)}
  =\left(\overline{\M}^{,1/D_{\M}(\Hilbert)}\right)^c
  =\left(\overline{\M}^{,c}\right)^{1/D_{\M}(\Hilbert)},
\]
or
\[
  \overline{\M'}
  =\left(\overline{\M}^{,c}\right)^{D_{\M'}(\Hilbert)/D_{\M}(\Hilbert)}
  =\left(\overline{\M}^{,c}\right)^\theta,
\]
as desired.
\end{proof}

\par\addvspace{0.35\baselineskip}
\noindent\textbf{Remark.}\ 
The first formula of \(\delta\)) holds clearly for \(\gamma\)) as well, while
it becomes meaningless for \(\alpha\)), \(\beta\)). The second formula of
\(\delta\)) holds clearly for \(\beta\)) as well, while it becomes meaningless
for \(\alpha\)), \(\gamma\)).

It is furthermore apparent from our proof that for \(\delta\)) the exponents
\(\theta,1/\theta\) are in the range of \thmref{thm:VII}{VII}.
\par\addvspace{0.35\baselineskip}

\begin{lemma}\phantomsection\label{lem:3.3.3}
Let \(\M\) and \(\theta\) be as above. Then the spatial type
\(\widetilde{\M}\) is uniquely determined by the algebraical types
\(\overline{\M},\overline{\M'}\) and the number \(\theta\).
\end{lemma}

\begin{proof}
As in the introduction, \S2, for each \(i=1,2\) a Hilbert space
\(\Hilbert_i\) and an operator ring \(\M_i\) in \(\Hilbert_i\) must be given.
We assume that both \(\M_i\) are factors in case (I) or (II), and form for each
\(\M_i\) its \(\theta_i\) in the sense of \defref{def:3.3.1}{3.3.1}. We
assume further that
\[
  \overline{\M}_1=\overline{\M}_2,\qquad
  \overline{\M'_1}=\overline{\M'_2},qquad
  \theta_1=\theta_2=\theta.
\]
Then our task is to establish \(\widetilde{\M}_1=\widetilde{\M}_2\), i.e. to
find an isomorphic mapping of \(\Hilbert_1\) on \(\Hilbert_2\) (cf. A) in the
Introduction, \S2) which carries \(\M_1\) into \(\M_2\).

In proving this we must distinguish several alternatives corresponding to
various values of \(\theta\) (the joint value of \(\theta_i\)).

Observe first that the isomorphic mapping of \(\Hilbert_1\) on
\(\Hilbert_2\) which carries \(\M_1\) into \(\M_2\) also carries
\(\M'_1\) into \(\M'_2\). Thus we can always replace each \(\M_i\) by its
\(\M'_i\). This has the consequence that
\(\theta_1=\theta_2=\theta\) is replaced by
\(1/\theta_1=1/\theta_2=1/\theta\).

Let us now consider the alternatives for \(\theta\).

First alternative: \(0<\theta<\infty\). Since we may replace \(\theta\) by
\(1/\theta\), we can even assume that \(0<\theta\leq1\). For both
\(i=1,2\), the factors \(\M_i,\M'_i\) are in finite cases. Choose the
normalizations of both \(D_{\M_i},D_{\M'_i}\) so that
\(D_{\M'_i}(\Hilbert_i)=1\), \(C_i=1\) (we use the notation of
\defref{def:3.3.1}{3.3.1} for both \(i=1,2\)). So
\(D_{\M_i}(\Hilbert_i)=1/\theta_i=1/\theta\geq1\). Hence 1 belongs to the
range of \(D_{\M_i}\) (cf. \cite[p.~182, the discussion of the ranges of
\(\Delta,\Delta_0,\Delta',\Delta'_0\) in Theorem~X and preceding it]{ref5}).
Choose accordingly a \(\mathfrak K_i\eta\M_i\), with
\(D_{\M_i}(\mathfrak K_i)=1\). There exists furthermore an \(\epsilon\) in
the range of \(D_{\M_1}\) for which there is a \(p=1,2,\ldots\), with
\(0<\epsilon\leq1\), \(p\epsilon=1/\theta\).\footnote{The discussion of
\cite[p.~182]{ref5} shows that the ranges of \(D_{\M_1}\) and
\(D_{\M'_1}\) agree for the area \(\leq\) both their maxima, i.e.
\(1,1/\theta\). Hence we can argue as follows:

Case \((\mathrm I_n)\), \(n=1,2,\ldots\): The range of \(D_{\M_1}\) is
\((0,1/n,\ldots,1)\); now \(1/\theta\geq1/n\), and \(1/n\) is the smallest
positive element of the range of \(D_{\M'_1}\). Hence that range contains only
integer multiples of \(1/n\). So \(1/\theta=p/n\), \(p=1,2,\ldots\). Thus
\(\epsilon=1/n\) meets our requirements.

Case \((\mathrm {II}_1)\). Choose \(p=1,2,\ldots\) so that
\(1/(p\theta)\leq1\). Then \(\epsilon=1/(p\theta)\) fulfills
\(p\epsilon=1/\theta\), and besides \(\epsilon\leq1,1/\theta\). Since
\(\epsilon\) belongs to the range of \(D_{\M_1}\), it belongs to the range
of \(D_{\M'_1}\) also.} Choose accordingly a
\(\mathfrak L_1\eta\M_1\), with
\(D_{\M_1}(\mathfrak L_1)=\epsilon\),
\(\mathfrak L_1\subseteq\mathfrak K_1\).

% Source PDF p. 36 (journal p. 751).
Let us now use the considerations of \cite[pp.~188--190, \S11.4]{ref5}. We
form, in accord with them,
\(\M_{i(\mathfrak K_i)},\M'_{i(\mathfrak K_i)}\). (We choose
\(\Hilbert_i,\mathfrak K_i,\mathfrak K_i\) and \(\M_i,\M'_i\) for the
\(\Hilbert,\mathfrak M,\mathfrak M'\) and \(\M,\M'\), loc. cit.) These are
coupled factors in \(\mathfrak K_i\ne(0)\), just as \(\M_i,\M'_i\) are in
\(\Hilbert_i\). We have
\[
  \overline{\M'_{i(\mathfrak K_i)}}=\overline{\M'_i}.
\]
Hence \(\overline{\M'_1}=\overline{\M'_2}\) gives
\[
  \overline{\M'_{1(\mathfrak K_1)}}
  =\overline{\M'_{2(\mathfrak K_2)}}.
\]
Besides,
\[
\begin{aligned}
D_{\M_i(\mathfrak K_i)}(\mathfrak K_i)&=D_{\M_i}(\mathfrak K_i)=1,\\
D_{\M'_i(\mathfrak K_i)}(\mathfrak K_i)&=D_{\M'_i}(\Hilbert_i)=1,
\end{aligned}
\]
and \(C_i=1\), as before. Hence \cite[p.~244, Theorem~XI]{ref6} applies with
our \(\mathfrak K_i,\M_{i(\mathfrak K_i)},\M'_{i(\mathfrak K_i)}\) in place
of its \(\Hilbert_i,\M_i,\M'_i\). There exists an isomorphic mapping
\(\mathfrak I\) of \(\mathfrak K_1\) on \(\mathfrak K_2\) which carries
\(\M_{1(\mathfrak K_1)}\) into \(\M_{2(\mathfrak K_2)}\), and
\(\M'_{1(\mathfrak K_1)}\) into \(\M'_{2(\mathfrak K_2)}\).
\(\mathfrak L_1\eta\M_1\),
\(\mathfrak L_1\subseteq\mathfrak K_1\) imply
\(\mathfrak L_1\eta\M_{1(\mathfrak K_1)}\). So \(\mathfrak I\) carries
\(\mathfrak L_1\) into an \(\mathfrak L_2\eta\M_{2(\mathfrak K_2)}\). It
follows that \(\mathfrak L_2\eta\M_2\) and
\(\mathfrak L_2\subseteq\mathfrak K_2\). Now
\[
  D_{\M_2}(\mathfrak L_2)
  =D_{\M_2(\mathfrak K_2)}(\mathfrak L_2)
  =D_{\M_1(\mathfrak K_1)}(\mathfrak L_1)
  =D_{\M_1}(\mathfrak L_1)=\epsilon.
\]
\(\mathfrak I\) carries
\(\M_{1(\mathfrak K_1)(\mathfrak L_1)}=\M_{1(\mathfrak L_1)}\) into
\(\M_{2(\mathfrak K_2)(\mathfrak L_2)}=\M_{2(\mathfrak L_2)}\). Restrict
\(\mathfrak I\) from \(\mathfrak K_1\) to \(\mathfrak L_1\), and denote
this restricted mapping by \(\mathfrak F\). Then \(\mathfrak F\) is an
isomorphic mapping of \(\mathfrak L_1\) on \(\mathfrak L_2\), which carries
\(\M_{1(\mathfrak L_1)}\) into \(\M_{2(\mathfrak L_2)}\).

We have \(p\epsilon=1/\theta\), and therefore for both \(i=1,2\),
\[
  pD_{\M_i}(\mathfrak L_i)=D_{\M_i}(\Hilbert_i),
\]
or, if we introduce the projection \(E_i=P_{\mathfrak L_i}\), then
\[
  D_{\M_i}(E_i)=\frac1pD_{\M_i}(1).
\]
Consequently we can proceed as in the proof of \lemref{lem:2.6.2}{2.6.2} or
\lemref{lem:2.6.3}{2.6.3} to obtain a set of \(p\)-order matrix units
\(W_{i;u,v}\), \(u,v=1,\ldots,p\), with
\(W_{i;1,1}=E_i\), \(\sum_{u=1}^pW_{i;u,u}=1\). Let the projection
\(F_{i;u}=W_{i;u,u}\) have range \(\mathfrak N_{i;u}\). We note that
\(\mathfrak N_{i;1}=\mathfrak L_i\), and that \(W_{i;u,v}\) is partially
isometric with initial set \(\mathfrak N_{i;v}\) and final set
\(\mathfrak N_{i;u}\), i.e. it is an isomorphic mapping of
\(\mathfrak N_{i;v}\) on \(\mathfrak N_{i;u}\).

Thus
\(W_{2;u,1}\mathfrak F W_{1;1,u}=\mathfrak F_u\) is an isomorphic mapping
of \(\mathfrak N_{1;u}\) on \(\mathfrak N_{2;u}\), for
\(u=1,2,\ldots,p\). For \(u=1\), \(W_{i;1,1}=F_{i;1}=E_i=P_{\mathfrak L_i}\),
\(\mathfrak N_{i;1}=\mathfrak L_i\), and hence \(\mathfrak F_1\) coincides
with \(\mathfrak F\). Since the \(\mathfrak N_{i;u}\),
\(u=1,\ldots,p\), are pairwise orthogonal and span together the closed linear
set \(\Hilbert_i\), we can combine the \(\mathfrak F_u\),
\(u=1,\ldots,p\), to one isomorphic mapping \(\mathfrak F^*\) of
\(\Hilbert_1\) on \(\Hilbert_2\). On \(\mathfrak L_1\), this
\(\mathfrak F^*\) coincides with \(\mathfrak F_1=\mathfrak F\). Hence
\begin{equation}
\begin{gathered}
\mathfrak F^*\text{ carries }\mathfrak L_1\text{ into }\mathfrak L_2
\text{ and in it }\M_{1(\mathfrak L_1)}\text{ into }
\M_{2(\mathfrak L_2)}.
\end{gathered}
\tag{3.3.\ensuremath{\eta}}\label{eq:3.3.eta}
\end{equation}

% Source PDF p. 37 (journal p. 752).
By virtue of the properties of \(W_{i;u,v}\) (cf.
\defref{def:2.6.1}{2.6.1}), and owing to the definition of
\(\mathfrak F^*\), we have further
\begin{equation}
\begin{gathered}
\mathfrak F^*\text{ carries }W_{1;u,v}\text{ into }W_{2;u,v},
\qquad u,v=1,\ldots,p.
\end{gathered}
\tag{3.3.\ensuremath{\vartheta}}\label{eq:3.3.theta-map}
\end{equation}
Now \(\M_i\) can be characterized in terms of
\(\M_{i(\mathfrak L_i)},\mathfrak L_i\), and the \(W_{i;u,v}\),
\(u,v=1,\ldots,p\). Indeed: \(A_i\in\M_i\) is equivalent to
\[
  \bigl(W_{i;1,v}A_iW_{i;u,1}\bigr)_{(\mathfrak L_i)}
  \in\M_{i(\mathfrak L_i)}
  \quad\text{for all }u,v=1,\ldots,p.\footnote{\hypertarget{fn:24}{}Put
  \(A_{i;u,v}=W_{i;1,v}A_iW_{i;u,1}\). The forward implication is obvious;
  the reverse follows from the formula
  \[
    A_i=\sum_{u,v=1}^pW_{i;u,1}A_{i;u,v}W_{i;1,v},
  \]
  which is easily verified.}
\]
Consequently \eqref{eq:3.3.eta} and \eqref{eq:3.3.theta-map} imply that
\(\mathfrak F^*\) carries \(\M_1\) into \(\M_2\).

This completes the proof of the first alternative.

Second alternative: \(\theta=0\) or \(\infty\). Since we may replace
\(\theta\) by \(1/\theta\), it is clear that we need only consider the case
\(\theta=0\). For both \(i=1,2\), the factor \(\M_i\) is in an infinite
case and the factor \(\M'_i\) is in a finite case. Choose the normalization of
\(D_{\M_i},D_{\M'_i}\) so that
\[
  D_{\M'_i}(\Hilbert_i)=1,\qquad C_i=1.
\]
We may parallel our discussion of the first alternative, however with
\(\epsilon=1\), \(p=\infty\). We form the \(\mathfrak K_i\) accordingly.
Since \(\epsilon=1\), \(\mathfrak L_1=\mathfrak K_1\).
\(\mathfrak I\) obtains as loc. cit.; \(\mathfrak L_1=\mathfrak K_1\) yields
\(\mathfrak L_2=\mathfrak K_2\) and \(\mathfrak F=\mathfrak I\).

Let \(E_i=P_{\mathfrak L_i}\). We proceed to obtain an infinite number of
pairwise orthogonal manifolds
\(\mathfrak N_{i;1},\mathfrak N_{i;2},\ldots\), with
\(\mathfrak N_{i;1}=\mathfrak L_i\),
\(\mathfrak N_{i;u}\eta\M_i\),
\(D_{\M_i}(\mathfrak N_{i;u})=1\), for \(u=1,2,\ldots\), and finally
\(\Hilbert_i=[\mathfrak N_{i;1},\mathfrak N_{i;2},\ldots]\) (cf.
\cite[p.~155, Lemma~7.1.2]{ref5}). We obtain the \(W_{i;u,v}\) as in the
second part of the proof of \lemref{lem:2.6.2}{2.6.2}.

As in the discussion of the preceding alternative, we now obtain
\(\mathfrak F_u\), \(u=1,2,\ldots\), and finally \(\mathfrak F^*\); and by
a literal repetition of that argument, the isomorphic mapping
\(\mathfrak F^*\) of \(\Hilbert_1\) on \(\Hilbert_2\) is seen to carry
\(\M_1\) into \(\M_2\).\footnote{\hyperlink{fn:24}{Footnote~24} applies to this case too. The
convergence of \(\sum_{u,v=1}^{\infty}\) in footnote~24 is readily seen to
be a consequence of the convergence of the sum
\(\sum_{u=1}^{\infty}W_{i;u,u}\). And in connection with the latter
convergence, the remarks made in the proof of (iii) in
\lemref{lem:2.6.1}{2.6.1} apply again.}

Thus the second alternative too is settled.

Third alternative: \(\theta=\infty/\infty\). For both \(i=1,2\), the
factors \(\M_i,\M'_i\) are in infinite cases. Choose a finite
\(\mathfrak K_1\eta\M_1\), \(\mathfrak K_1\ne(0)\), and put
\(P_{\mathfrak K_1}=E_1\in\M_1\). Under the (algebraic ring) isomorphism of
\(\M_1,\M_2\) (due to \(\overline{\M}_1=\overline{\M}_2\)), this projection
\(E_1\in\M_1\) corresponds to a projection \(E_2\in\M_2\). Let
\(\mathfrak K_2\) be the range of \(E_2\). Then
\(\mathfrak K_2\eta\M_2\), \(\mathfrak K_2\ne(0)\). Then the isomorphism of
\(\M_1,\M_2\) also induces one of
\(\M_{1(\mathfrak K_1)},\M_{2(\mathfrak K_2)}\), owing to
\cite[p.~187, Lemma~11.3.3(i)]{ref5}. Thus
\begin{equation}
  \overline{\M_{1(\mathfrak K_1)}}
  =\overline{\M_{2(\mathfrak K_2)}}.
  \tag{3.3.\ensuremath{\iota}}\label{eq:3.3.iota}
\end{equation}
Also by (ii) eod.,
\(\overline{\M'_{i(\mathfrak K_i)}}=\overline{\M'_i}\); hence
\(\overline{\M'_1}=\overline{\M'_2}\) gives
\begin{equation}
  \overline{\M'_{1(\mathfrak K_1)}}
  =\overline{\M'_{2(\mathfrak K_2)}}.
  \tag{3.3.\ensuremath{\kappa}}\label{eq:3.3.kappa}
\end{equation}

% Source PDF p. 38 (journal p. 753).
Now choose the normalization of \(D_{\M_i},D_{\M'_i}\) so that
\(D_{\M_i}(\mathfrak K_i)=1\), \(C_i=1\). Automatically
\(D_{\M'_i}(\Hilbert_i)=\infty\). We may then write
\[
\begin{aligned}
D_{\M_{i(\mathfrak K_i)}}(\mathfrak K_i)
  &=D_{\M_i}(\mathfrak K_i)=1,\\
D_{\M'_{i(\mathfrak K_i)}}(\mathfrak K_i)
  &=D_{\M'_i}(\Hilbert_i)=\infty,
\end{aligned}
\]
and \(C_i=1\). Apply \defref{def:3.3.1}{3.3.1} to
\(\mathfrak K_i,\M_{i(\mathfrak K_i)},\M'_{i(\mathfrak K_i)}\) in place of
\(\Hilbert,\M,\M'\), and denote the resulting number by \(\theta_i^0\). The
above equations give
\begin{equation}
  \theta_1^0=\theta_2^0=\infty.
  \tag{3.3.\ensuremath{\lambda}}\label{eq:3.3.lambda}
\end{equation}
Equations \eqref{eq:3.3.iota}, \eqref{eq:3.3.kappa}, and
\eqref{eq:3.3.lambda} permit us to apply the second alternative, which we have
already disposed of, to
\(\mathfrak K_i,\M_{i(\mathfrak K_i)},\M'_{i(\mathfrak K_i)}\). There exists
an isomorphic mapping \(\mathfrak I\) of \(\mathfrak K_1\) on
\(\mathfrak K_2\), which carries \(\M_{1(\mathfrak K_1)}\) into
\(\M_{2(\mathfrak K_2)}\), and \(\M'_{1(\mathfrak K_1)}\) into
\(\M'_{2(\mathfrak K_2)}\).

We can continue from here on exactly as in the discussion of the second
alternative, i.e. paralleling the considerations of the first alternative. We
have the \(\mathfrak K_i\) already. Again we put \(\epsilon=1\),
\(p=\infty\). Since \(\epsilon=1\),
\(\mathfrak L_1=\mathfrak K_1\). We already have \(\mathfrak I\).
\(\mathfrak L_1=\mathfrak K_1\) gives
\(\mathfrak L_2=\mathfrak K_2\) and \(\mathfrak F=\mathfrak I\). We have
\(E_i=P_{\mathfrak L_i}=P_{\mathfrak K_i}\). The \(F_{i;u}\),
\(u=1,2,\ldots\), and \(W_{i;u,v}\), \(u,v=1,2,\ldots\), obtain as loc.
cit. So do the \(\mathfrak F_u\), \(u=1,2,\ldots\), and finally
\(\mathfrak F^*\); and exactly as loc. cit., the isomorphic mapping
\(\mathfrak F^*\) of \(\Hilbert_1\) on \(\Hilbert_2\) is seen to carry
\(\M_1\) into \(\M_2\).

Thus the third alternative too is settled and the entire proof is completed.
\end{proof}

We settle now the question of existence.

\begin{lemma}\phantomsection\label{lem:3.3.4}
A factor \(\M\) in the cases (I) or (II), with given algebraical types
\(\overline{\M},\overline{\M'}\) and given number \(\theta\) (cf.
\defref{def:3.3.1}{3.3.1}), exists if and only if these fulfill the alternative
conditions \(\alpha\))--\(\delta\)) in \lemref{lem:3.3.2}{3.3.2}.
\end{lemma}

\begin{proof}
Necessity: This is merely the assertion of \lemref{lem:3.3.2}{3.3.2}.

Sufficiency: We shall continue to denote the two prescribed algebraical types
by \(\overline{\M},\overline{\M'}\), although we have not yet identified them
as belonging in this way to any coupled factors \(\M,\M'\). (In fact, this is
what we propose to prove.) The relations \(\alpha\))--\(\delta\)) in
\lemref{lem:3.3.2}{3.3.2}, which we assumed, imply for the genera
\begin{equation}
  \widetilde{\overline{\M'}}
  =\widetilde{\overline{\M}}^{,c}.
  \tag{3.3.\ensuremath{\mu}}\label{eq:3.3.mu}
\end{equation}
(This cannot be deduced from \lemref{lem:3.3.1}{3.3.1}, since
\(\M,\M'\) are fictitious, cf. above.) The genus
\(\widetilde{\overline{\M}}\) contains a (unique) infinite type
\(\overline{\R}\) (cf. \thmref{thm:IX}{IX})---\(\R\) being an actual factor
in a Hilbert space \(\mathfrak R\).

Consider now the factor \(\R'\) in \(\mathfrak R\). We repeat the construction
in \secref{sec:2.4}{2.4}, immediately preceding
\thmref{thm:IV-prime}{IV$'$}, with our \(\mathfrak R,\R'\) in place of its
\(\Hilbert,\M\), and with its \(p=\infty\). Thus, using the notations of
\secref{sec:2.4}{2.4}, loc. cit., our \(\R'\) is a factor in
\(\mathfrak R=\Hilbert=\Hilbert_2\), and with the (\(p=\infty\)-dimensional)
Hilbert space \(\Hilbert_1\), we obtain the Hilbert space
\(\Hilbert_1\otimes\Hilbert_2\), and in it the ring \(\M_1\) of all
operators \(\langle A'_{t,s}\rangle\), with all \(A'_{t,s}\in\R'\) (cf. the
similar argument before \thmref{thm:IV}{IV} in \secref{sec:2.4}{2.4}). Put
\(\R_1=\M'_1\), i.e. \(\M_1=\R'_1\). It is obvious, by the usual matrix
computations, that \(\R_1=\M'_1\) is the set of all operators
\(\langle A\delta_{t,s}\rangle\), \(A\in\R\). Hence
\(\overline{\R}_1=\overline{\R}\). On the other hand, the discussions of
\secref{sec:2.4}{2.4} show that
\(\overline{\R'_1}=\overline{\M}_1=\overline{\R'}^{\infty}\). Thus
\(\overline{\R'_1}\) is in an infinite case (by
\lemref{lem:2.7.1}{2.7.1}).

% Source PDF p. 39 (journal p. 754).
Now \(\overline{\R}_1=\overline{\R}\) permits us to replace
\(\mathfrak R,\R\) by \(\Hilbert_1\otimes\Hilbert_2,\R_1\). So we have
shown: It is permissible to assume that \(\overline{\R'}\) is in an infinite
case. We do this, but we continue using the original notations
\(\mathfrak R,\R\).

By \lemref{lem:3.3.1}{3.3.1},
\(\widetilde{\overline{\R'}}=\widetilde{\overline{\R}}^{,c}\). Now
\(\widetilde{\overline{\R}}=\widetilde{\overline{\M}}\); hence
\(\widetilde{\overline{\R'}}=\widetilde{\overline{\M}}^{,c}\). By
\eqref{eq:3.3.mu}, this means
\(\widetilde{\overline{\R'}}=\widetilde{\overline{\M'}}\). Since
\(\overline{\R},\overline{\R'}\) are both infinite, we have (cf.
\thmref{thm:IX}{IX})
\begin{equation}
\begin{gathered}
\overline{\R},\overline{\R'}\text{ are the (unique) infinite types of the
genera }\widetilde{\overline{\M}},\widetilde{\overline{\M'}},
\text{ respectively.}
\end{gathered}
\tag{3.3.\ensuremath{\nu}}\label{eq:3.3.nu}
\end{equation}

Choose an \(f_0\ne0\) in \(\mathfrak R\), and form the closed linear sets
\(\mathfrak M_{f_0}^{\R'}\eta\R\),
\(\mathfrak M_{f_0}^{\R}\eta\R'\) (cf.
\cite[p.~143, Def.~5.1.1]{ref5}). We may assume that
\(\mathfrak M_{f_0}^{\R'}\) is finite dimensional. (Choose a finite
dimensional \(\mathfrak M\eta\R\), \(\mathfrak M\ne(0)\), and then
\(f_0\in\mathfrak M\). Clearly
\(\mathfrak M_{f_0}^{\R'}\subseteq\mathfrak M\).) In this case,
\(\mathfrak M_{f_0}^{\R}\) is automatically finite dimensional. (Use
\cite[p.~182, Theorem~X]{ref5}, for this and for the next step.) Hence both
\(D_{\R}(\mathfrak M_{f_0}^{\R'})\),
\(D_{\R'}(\mathfrak M_{f_0}^{\R})\) are \(>0,<\infty\), and
\[
  C=\frac{D_{\R'}(\mathfrak M_{f_0}^{\R})}
          {D_{\R}(\mathfrak M_{f_0}^{\R'})}.
\]
Now normalize \(D_{\R},D_{\R'}\) so that
\(D_{\R}(\mathfrak M_{f_0}^{\R'})
=D_{\R'}(\mathfrak M_{f_0}^{\R})\). The above formula for \(C\) gives
\begin{equation}
  C=1.
  \tag{3.3.\ensuremath{o}}\label{eq:3.3.omicron}
\end{equation}
Write \(\mathfrak M_1=\mathfrak M_{f_0}^{\R'}\),
\(\mathfrak M'_1=\mathfrak M_{f_0}^{\R}\); then
\begin{equation}
  D_{\R}(\mathfrak M_1)=D_{\R'}(\mathfrak M'_1).
  \tag{3.3.\ensuremath{\xi}}\label{eq:3.3.xi}
\end{equation}

We shall now again make use of the considerations of
\cite[pp.~188--190, \S11.4]{ref5}. In accord with them we form
\(\R_{(\mathfrak M_1\cdot\mathfrak M'_1)}\),
\(\R'_{(\mathfrak M_1\cdot\mathfrak M'_1)}\). These are coupled factors in
\(\mathfrak M_1\cdot\mathfrak M'_1\ne(0)\), just as \(\R,\R'\) are in
\(\Hilbert\), and \eqref{eq:3.3.xi} gives
\[
  D_{\R_{(\mathfrak M_1\cdot\mathfrak M'_1)}}
    (\mathfrak M_1\cdot\mathfrak M'_1)
  =D_{\R'_{(\mathfrak M_1\cdot\mathfrak M'_1)}}
    (\mathfrak M_1\cdot\mathfrak M'_1).
\]
This, together with \eqref{eq:3.3.omicron} and a normalization, shows that the
requirements of \cite[p.~239, Theorem~VI]{ref6} are fulfilled.\textsuperscript{\hyperlink{fn:21}{21}} Hence
\(\R_{(\mathfrak M_1\cdot\mathfrak M'_1)}\),
\(\R'_{(\mathfrak M_1\cdot\mathfrak M'_1)}\) are dual
isomorphic\textsuperscript{\hyperlink{fn:22}{22}}, i.e.
\[
  \overline{\R'_{(\mathfrak M_1\cdot\mathfrak M'_1)}}
  =\overline{\R_{(\mathfrak M_1\cdot\mathfrak M'_1)}}^{\,c}.
\]
Now
\(\overline{\R_{(\mathfrak M_1)}}
=\overline{\R_{(\mathfrak M_1\cdot\mathfrak M'_1)}}\),
\(\overline{\R'_{(\mathfrak M'_1)}}
=\overline{\R'_{(\mathfrak M_1\cdot\mathfrak M'_1)}}\). Consequently
\begin{equation}
  \overline{\R'_{(\mathfrak M'_1)}}
  =\overline{\R_{(\mathfrak M_1)}}^{\,c}.
  \tag{3.3.\ensuremath{\pi}}\label{eq:3.3.pi}
\end{equation}

Since \(\widetilde{\overline{\R}}=\widetilde{\overline{\M}}\),
\(\overline{\R},\overline{\M}\) are commensurable. \(\overline{\R}\) is
in an infinite case. Hence \(\alpha\)), \(\beta\)) in
\lemref{lem:3.1.6}{3.1.6} and \lemref{lem:3.1.4}{3.1.4} give together that
\(\overline{\M}\) is a divisor of \(\overline{\R}\). Hence by
\defref{def:3.1.1}{3.1.1},
\begin{equation}
  \overline{\M}=\overline{\R_{(\mathfrak L)}}
  \quad\text{for a suitable }\mathfrak L\eta\R.
  \tag{3.3.\ensuremath{\rho}}\label{eq:3.3.rho}
\end{equation}
If \(\mathfrak L\) is of finite relative dimension, we may assume by
\cite[Theorem~X]{ref5} that \(\mathfrak L=\mathfrak M_{f_0}^{\R'}\) for some
\(f_0\). Similarly,
\begin{equation}
  \overline{\M'}=\overline{\R'_{(\mathfrak L')}}
  \quad\text{for a suitable }\mathfrak L'\eta\R',\quad
  \mathfrak L'=\mathfrak M_{g_0}^{\R}\text{ when finite}.
  \tag{3.3.\ensuremath{\sigma}}\label{eq:3.3.sigma}
\end{equation}

From here on, it is again convenient to distinguish several alternatives,
corresponding to various values of \(\theta\).

% Source PDF p. 40 (journal p. 755).
First alternative: \(0<\theta<\infty\). We have the case \(\delta\)) in
\lemref{lem:3.3.2}{3.3.2}. Hence
\(\overline{\M'}=(\overline{\M}^{,c})^\theta\), or by
\eqref{eq:3.3.rho} and \eqref{eq:3.3.sigma},
\[
  \overline{\R'_{(\mathfrak L')}}
  =\left(\overline{\R_{(\mathfrak L)}}^{,c}\right)^\theta.
\]
Now \eqref{eq:3.3.pi} and (i) in \lemref{lem:3.2.2}{3.2.2} give
\[
  \overline{\R'_{(\mathfrak L')}}
  =\left(\overline{\R_{(\mathfrak L)}}^{,c}\right)^\beta,
\]
where \(\beta=D_{\R'}(\mathfrak L')/D_{\R}(\mathfrak L)\). Hence (ii) in
\lemref{lem:3.2.2}{3.2.2} gives
\[
  \alpha=\frac{\beta}{\theta}
  =\frac1\theta\frac{D_{\R'}(\mathfrak L')}{D_{\R}(\mathfrak L)}
  \in\mathfrak G(\widetilde{\overline{\R}})
  =\mathfrak G(\overline{\R_{(\mathfrak L)}}).
\]
The interchange of \(\overline{\M}\) and \(\overline{\M'}\), and with them
of \(\theta\) and \(1/\theta\), as well as \(\R,\mathfrak L\) and
\(\R',\mathfrak L'\), replaces \(\alpha\) by \(1/\alpha\). So we may
assume \(\alpha\leq1\). Hence
\(\alpha\in\mathfrak G(\overline{\R_{(\mathfrak L)}})\) secures the
constructability of \(\overline{\R_{(\mathfrak L)}}^{\alpha}\) in the original
sense of \secref{sec:2.8}{2.8}. Thus there exists an
\(\mathfrak L''\subseteq\mathfrak L\), with
\(\mathfrak L''\eta\R_{(\mathfrak L)}\), and
\[
  \alpha
  =\frac{D_{\R_{(\mathfrak L)}}(\mathfrak L'')}
          {D_{\R_{(\mathfrak L)}}(\mathfrak L)}
  =\frac{D_{\R}(\mathfrak L'')}{D_{\R}(\mathfrak L)}.
\]
(\(\mathfrak L''\eta\R_{(\mathfrak L)}\) implies
\(\mathfrak L''\eta\R\).) The membership
\(\alpha\in\mathfrak G(\overline{\R_{(\mathfrak L)}})\) gives
\[
  \overline{\R_{(\mathfrak L'')}}
  =\overline{\R_{(\mathfrak L)}}^{\alpha}
  =\overline{\R_{(\mathfrak L)}}=\overline{\M}.
\]
Thus we may replace in all our considerations \(\mathfrak L\) by
\(\mathfrak L''\). This carries \(\alpha=1\). Writing again
\(\mathfrak L\) for \(\mathfrak L''\), we see: We may assume that
\(\alpha=1\), i.e. that
\[
  \theta=\frac{D_{\R'}(\mathfrak L')}{D_{\R}(\mathfrak L)}.
\]
Now consider the coupled factors
\(\R_{(\mathfrak L\cdot\mathfrak L')}\),
\(\R'_{(\mathfrak L\cdot\mathfrak L')}\) in
\(\mathfrak L\cdot\mathfrak L'\ne(0)\). Combining our results, we see that
we have
\[
\begin{aligned}
\overline{\R_{(\mathfrak L\cdot\mathfrak L')}}
  &=\overline{\R_{(\mathfrak L)}}=\overline{\M},\\
\overline{\R'_{(\mathfrak L\cdot\mathfrak L')}}
  &=\overline{\R'_{(\mathfrak L')}}=\overline{\M'}.
\end{aligned}
\]
Thus the spatial type
\(\widetilde{\R}_{(\mathfrak L\cdot\mathfrak L')}\) meets all our
requirements.

Second alternative: \(\theta=0\) or \(\theta=\infty\). Since we may
interchange \(\overline{\M}\) and \(\overline{\M'}\), and with them
\(\theta\) and \(1/\theta\), we can even assume that \(\theta=0\). Then we
have the case \(\beta\)) in \lemref{lem:3.3.2}{3.3.2}.

Thus \(\overline{\M}\) is in an infinite case. Hence \eqref{eq:3.3.nu} gives
\(\overline{\R}=\overline{\M}\). Now consider the coupled factors
\(\R_{(\mathfrak L')}\), \(\R'_{(\mathfrak L')}\) in
\(\mathfrak L'\ne(0)\). Combining our results, we see that we have
\[
\begin{aligned}
\overline{\R_{(\mathfrak L')}}&=\overline{\R}=\overline{\M},\\
\overline{\R'_{(\mathfrak L')}}&=\overline{\M'},\\
D_{\R_{(\mathfrak L')}}(\mathfrak L')&=D_{\R}(\Hilbert)=\infty,\\
D_{\R'_{(\mathfrak L')}}(\mathfrak L')&=D_{\R'}(\mathfrak L')<\infty.
\end{aligned}
\]

% Source PDF p. 41 (journal p. 756).
Thus the spatial type \(\widetilde{\R}_{(\mathfrak L')}\) meets all our
requirements.

Third alternative: \(\theta=\infty/\infty\). We have the case \(\alpha\)) in
\lemref{lem:3.3.2}{3.3.2}. Hence \(\overline{\M},\overline{\M'}\) are both
in infinite cases, and so \eqref{eq:3.3.nu} gives
\[
  \overline{\R}=\overline{\M},
  \qquad
  \overline{\R'}=\overline{\M'}.
\]
Also
\[
  D_{\R}(\Hilbert)=D_{\R'}(\Hilbert)=\infty.
\]
Thus the spatial type \(\widetilde{\R}\) meets all our requirements.

This completes the proof.
\end{proof}

Summing up the results of \lemref{lem:3.3.3}{3.3.3} and
\lemref{lem:3.3.4}{3.3.4},

\begin{theorem}\phantomsection\label{thm:X}
Consider factors \(\M\) in the cases (I) and (II). Then the spatial type
\(\widetilde{\M}\) is uniquely determined by the algebraical types
\(\overline{\M},\overline{\M'}\), and the number \(\theta\) (cf.
\defref{def:3.3.1}{3.3.1}).

\(\widetilde{\M}\) exists if and only if \(\overline{\M},\overline{\M'}\)
and \(\theta\) fulfill the alternative conditions \(\alpha\))--\(\delta\)) in
\lemref{lem:3.3.2}{3.3.2}.
\end{theorem}

\par\addvspace{0.35\baselineskip}
\noindent\textbf{Remark.}\ 
In some cases, two of \(\overline{\M},\overline{\M'}\), and \(\theta\)
suffice to determine the third one and thus \(\widetilde{\M}\). Specifically,
\begin{enumerate}[label={}]
\item[\(\alpha\))] If \(\overline{\M}\) is in a finite case, then \(\overline{\M}\) and
\(\theta\) determine \(\overline{\M'}\) by \(\gamma\)), \(\delta\)) in
\lemref{lem:3.3.2}{3.3.2}.
\item[\(\beta\))] If \(\overline{\M'}\) is in a finite case, then \(\overline{\M'}\) and
\(\theta\) determine \(\overline{\M}\) by \(\beta\)), \(\delta\)) in
\lemref{lem:3.3.2}{3.3.2}.
\item[\(\gamma\))] If either \(\overline{\M}\) or \(\overline{\M'}\) is in an infinite
case, then \(\overline{\M},\overline{\M'}\) determine \(\theta\) by
\(\alpha\)), \(\beta\)), or \(\gamma\)) in \lemref{lem:3.3.2}{3.3.2}.
\item[\(\delta\))] If \(\overline{\M},\overline{\M'}\) are both in finite cases, then
\(\overline{\M},\overline{\M'}\) determine \(\theta\) precisely up to a
factor belonging to the group
\(\mathfrak G(\widetilde{\overline{\M}})\), owing to \(\delta\)) in
\lemref{lem:3.3.2}{3.3.2}, and to \eqref{eq:2.10.alpha} and
\lemref{lem:3.2.1}{3.2.1}. Consequently they determine \(\theta\) outright if
and only if \(\mathfrak G(\widetilde{\overline{\M}})=(1)\).\footnote{Hence
\(\theta\) is never needed in the cases (I), by
\lemref{lem:2.10.4}{2.10.4}. The situation is different in the cases (II).
Cf. \thmref{thm:XV}{XV} and the remark at the end of \secref{sec:5.6}{5.6}.}
\end{enumerate}
\par\addvspace{0.35\baselineskip}

Thus the only remaining objects for our inquiries in the cases (I) and (II)
are the algebraical types \(\overline{\M}\). \thmref{thm:IX}{IX} and
\thmref{thm:X}{X} (or \lemref{lem:3.2.1}{3.2.1} and
\lemref{lem:3.2.2}{3.2.2} in place of the latter) show that the knowledge of
these is equivalent to the knowledge of their genera
\(\widetilde{\overline{\M}}\). The case (I) is completely clear; hence the
above concepts must be studied in the case (II). And we can analyze a genus
\(\widetilde{\overline{\M}}\) in case (II) by analyzing any one of its types
\(\overline{\M}\) in a finite case. Consequently we need only be concerned
with (algebraical) types \(\overline{\M}\) in case \((\mathrm {II}_1)\).

\PaperChapter{4}{chapter:IV}{Approximate Finiteness}

\PaperSection{sec:4.1}{}
In accordance with the above we begin the systematic investigation of the
algebraical types \(\overline{\M}\) in case \((\mathrm {II}_1)\). Throughout
this chapter \(\M\) will be a factor in case \((\mathrm {II}_1)\) in Hilbert
space \(\Hilbert\). We use \(D_{\M}(E)\), (but now in the standard
normalization) \(\operatorname {Tr}_{\M}(A)\), \(\Metric{A}\) as in
\secref{sec:1.2}{1.2}. We shall also consider other factors in \(\Hilbert\),
but they will be denoted by different letters. All our considerations however
will take place in \(\Hilbert\).

% Source PDF p. 42 (journal p. 757).
We begin by stating some familiar facts about factors in the case
\((\mathrm I_n)\), \(n=1,2,\ldots\).

\begin{lemma}\phantomsection\label{lem:4.1.1}
Let \(\N\) be a factor in a case \((\mathrm I_p)\),
\(p=1,2,\ldots\). Then there exists a one-to-one correspondence between the
two following sets of entities,
\begin{enumerate}[label={}]
\item[\((\alpha)\)] The (algebraical ring) isomorphisms \(\mathfrak K\) between \(\N\) and
the system of all complex \(p\)-order matrices
\[
  a=\{a_{t,s}\},\qquad (t,s=1,\ldots,p).\footnote{The (complex numerical)
  matrices which we consider now should be viewed as endowed with the same
  operations and defined in the same way as the (operator) matrices of
  \(\mathbf B_p\) in \secref{sec:2.4}{2.4}.}
\]
\item[\((\beta)\)] The matrix bases of \(\N\), i.e. the systems \(W_{u,v}\),
\((u,v=1,\ldots,p)\), of \(p\)-order matrix units, that are contained in
\(\N\) and have
\begin{equation}\label{eq:4.1.alpha}
  E_0=\sum_{u=1}^{p}W_{u,u}=1.
  \tag{4.1.\ensuremath{\alpha}}
\end{equation}
\end{enumerate}
The correspondence is this:
\begin{enumerate}[label=\((\roman*)\)]
\item \(\mathfrak K\) is given:
\[
  W_{u,v}=\mathfrak K^{-1}\!\left(\{\delta_{t,v}\delta_{s,u}\}\right).
\]
\item The \(W_{u,v}\) are given:
\[
  \mathfrak K A=a=\{a_{t,s}\}
\]
is equivalent to
\[
  A=\sum_{u=1}^{p}\sum_{v=1}^{p}a_{v,u}W_{u,v}.
\]
\end{enumerate}
\end{lemma}

\begin{proof}\footnote{Given only for the sake of completeness.}
(i) leads from (\(\alpha\)) to (\(\beta\)). The conditions of
\defref{def:2.6.1}{2.6.1} as well as \eqref{eq:4.1.alpha} are verified by
direct matrix computation.

(ii) leads from (\(\beta\)) to (\(\alpha\)): The rules of matrix computation
are verified from \defref{def:2.6.1}{2.6.1} and \eqref{eq:4.1.alpha}.

(i) implies (ii): Immediate by matrix computation.

(ii) implies (i). Put \(a_{u,v}=\delta_{t,v}\delta_{s,u}\) in (ii).
\end{proof}

\par\addvspace{0.35\baselineskip}
\noindent\textbf{Remark 1.}\ 
It is clear by (ii) that the ring \(\N\) is generated by the \(W_{u,v}\),
\(u,v=1,\ldots,p\),
\[
  \N=\R(W_{u,v};\ u,v=1,\ldots,p).
\]

\noindent\textbf{Remark 2.}\ 
The two sets (\(\alpha\)), (\(\beta\)) are of course not empty. This follows
in numerous ways from past results; the equivalent was even stated explicitly
in the Corollary to \lemref{lem:2.10.3}{2.10.3}.

\begin{lemma}\phantomsection\label{lem:4.1.2}
Let \(\N,\mathbf O\) be two factors in the cases \((\mathrm I_p)\),
\((\mathrm I_q)\) respectively, \(p,q=1,2,\ldots\). Assume that
\[
  \N\subseteq\mathbf O.
\]
Then we have
\begin{enumerate}[label=\((\roman*)\)]
\item \(p\) is a divisor of \(q\): \(q=pr\), \(r=1,2,\ldots\).
\item To every matrix basis \(W_{u,v}\), \(u,v=1,\ldots,p\), of \(\N\)
there exists a matrix base \(X_{o,w}\), \(o,w=1,\ldots,q\), of \(\mathbf O\)
such that
\[
  W_{u,v}=\sum_{i=1}^{r}X_{(u-1)r+i,(v-1)r+i}.
\]
\item If we pass from the matrix bases \(W_{u,v}\) and \(X_{o,w}\) in (ii)
to the isomorphisms \(\mathfrak K\) and \(\mathfrak L\) which correspond to
them by \lemref{lem:4.1.1}{4.1.1}, then the correlation of (ii) is equivalent
to this:

For \(A\in\N\subseteq\mathbf O\), both
\(\mathfrak K A=a=\{a_{t,s}\}\), \(t,s=1,\ldots,p\), and
\(\mathfrak L A=\bar a=\{\bar a_{k,l}\}\), \(k,l=1,\ldots,q\), are
defined and
\[
  \bar a_{k,l}=a_{t,s}\delta_{i,j}
  \quad\text{for}\quad
  k=(t-1)r+i,\qquad l=(s-1)r+j.
\]
\end{enumerate}
\end{lemma}

% Source PDF p. 43 (journal p. 758).
\begin{proof}
Ad (i). \(D_{\mathbf O}(E)\) will do as \(D_{\N}(E)\) for the projections
\(E\in\N\subseteq\mathbf O\) in the sense of
\cite[p.~165, Definition~8.2.1]{ref5}\footnote{Consider also
\cite[p.~170, Lemma~8.3.5]{ref5}. We are using the fact that \(\mathbf O\)
belongs to a finite case.} and standard normalization. And since the dimension
is uniquely characterized by those properties, therefore
\(D_{\mathbf O}(E)=D_{\N}(E)\) for these \(E\).

Hence the total range of \(D_{\mathbf O}(E)\) contains the total range of
\(D_{\N}(E)\), i.e. the set \((0,1/q,\ldots,1)\) contains the set
\((0,1/p,\ldots,1)\). Consequently, \(p\) is a divisor of \(q\).

Ad (ii). Consider a matrix base \(W_{u,v}\) of \(\N\) and an arbitrary matrix
base \(X'_{o,w}\) of \(\mathbf O\). Put
\begin{equation}\label{eq:4.1.beta}
  W'_{u,v}=\sum_{i=1}^{r}X'_{(u-1)r+i,(v-1)r+i}.
  \tag{4.1.\ensuremath{\beta}}
\end{equation}
Then one verifies by explicit computation (using
\defref{def:2.6.1}{2.6.1}) that the \(W'_{u,v}\) also form a system of
\(p\)-order matrix units. Also
\[
  \sum_{u=1}^{p}W'_{u,u}=\sum_{o=1}^{q}X'_{o,o}=1.
\]
Since \(D_{\mathbf O}(W'_{u,u})=D_{\mathbf O}(W'_{1,1})\), we have
\(D_{\mathbf O}(W'_{1,1})=1/p\), and similarly
\(D_{\mathbf O}(W_{1,1})=1/p\). Therefore there exists a partially isometric
\(U\) in \(\mathbf O\) with the initial projection \(W'_{1,1}\) and the final
projection \(W_{1,1}\). Now put
\[
  V=\sum_{u=1}^{p}W_{u,1}UW'_{1,u}.
\]
Clearly \(V\in\mathbf O\), and one verifies by direct computation
\[
\begin{array}{lll}
(1)&V^*V=VV^*=1,&\text{i.e. \(V\) is unitary},\\
(2)&VW'_{u,v}=W_{u,v}V,&
  \text{i.e. \(W_{u,v}=VW'_{u,v}V^{-1}\)}.
\end{array}
\]
Hence replacement of \(X'_{u,v}\) by
\(VX'_{u,v}V^{-1}=X_{u,v}\) will leave all its properties intact, but replace
\(W'_{u,v}\) by \(W_{u,v}=VW'_{u,v}V^{-1}\). And now
\eqref{eq:4.1.beta} yields the desired equation of (ii).

Ad (iii). Apply (ii) in \lemref{lem:4.1.1}{4.1.1} to \(\mathfrak K\) (with
\(W_{u,v}\)) and \(\mathfrak L\) (with \(X_{o,w}\)). Then the equation of
(ii) goes over into the desired relation of (iii).
\end{proof}

% Source PDF p. 44 (journal p. 759).
We now define

\begin{definition}\phantomsection\label{def:4.1.1}
\(\M\) is approximately finite of type \([p_1,p_2,\ldots]\) if this is true.
There exists a sequence of factors \(\N_1,\N_2,\ldots\) with the following
properties:
\begin{enumerate}[label=\((\roman*)\)]
\item \(\N_n\) is in case \((\mathrm I_{p_n})\).
\item \(\N_1\subseteq\N_2\subseteq\cdots\).
\item The ring \(\M\) is generated by the \(\N_1,\N_2,\ldots\),
\[
  \M=\R(\N_n;\ n=1,2,\ldots).
\]
\end{enumerate}
\end{definition}

We have immediately:

\begin{lemma}\phantomsection\label{lem:4.1.3}
Under the conditions of \defref{def:4.1.1}{4.1.1}, we have necessarily
\begin{enumerate}[label=\((\roman*)\)]
\item \(p_n\) is a divisor of \(p_{n+1}\),
\(p_{n+1}=p_nr_n\), \((r_n=1,2,\ldots)\).
\item \(p_n\to\infty\) for \(n\to\infty\), i.e. infinitely many
\(r_n>1\).
\end{enumerate}
\end{lemma}

\begin{proof}
Ad (i). As \(\N_n\subseteq\N_{n+1}\), this follows from (i) in
\lemref{lem:4.1.2}{4.1.2}.

Ad (ii). Owing to (i), \(p_1\le p_2\le\cdots\), so failure of (ii) implies
\(p_m=p_{m+1}=p_{m+2}=\cdots\) for some \(m\). Hence equally
\(\N_m=\N_{m+1}=\N_{m+2}=\cdots\)\footnote{For \(l\ge m\), \(\N_m\) and
\(\N_l\) contain the same number of linearly independent elements:
\(p_l^2=p_m^2\). As \(\N_m\subseteq\N_l\), this implies
\(\N_m=\N_l\).} and so (iii) in \defref{def:4.1.1}{4.1.1} gives
\[
  \M=\R(\N_n;\ n=1,2,\ldots)=\N_m.
\]
But this is impossible since \(\M\) is in the case \((\mathrm {II}_1)\) and
\(\N_m\) in the case \((\mathrm I_{p_m})\).
\end{proof}

For every sequence \([p_1,p_2,\ldots]\) which meets these requirements there
exists an \(\M\) which is approximately finite of that type. But the concept
introduced by \defref{def:4.1.1}{4.1.1} is a preliminary one and we prefer to
prove this existence only in conjunction with the final reformulation of the
concept of approximate finiteness in \thmref{thm:XII}{XII} and
\thmref{thm:XIV}{XIV}. We now proceed in a different direction.

\begin{lemma}\phantomsection\label{lem:4.1.4}
For any sequence of rings \(\N_1,\N_2,\ldots\) which fulfills conditions (ii),
(iii) in \defref{def:4.1.1}{4.1.1}, denote by \(\mathbf S\) the set theoretical
sum of that sequence. Then we have
\begin{enumerate}[label=\((\roman*)\)]
\item \(\mathbf S\) is an algebra.
\item \(\M\) is the set of elements of \(\mathbf B\) which are limits of a
convergent sequence from \(\mathbf S\).
\end{enumerate}
\end{lemma}

\begin{proof}
Ad (i): Every \(\N_n\) is a ring; hence an algebra. Also
\(\N_1\subseteq\N_2\subseteq\cdots\). Hence the sum \(\mathbf S\) is an
algebra also.

Ad (ii): Denote the set of limits of all metrically convergent sequences from
\(\mathbf S\) by \(\mathbf S_1\). We must prove that \(\mathbf S_1=\M\).

Now
\[
  \M=\R(\N_n;\ n=1,2,\ldots)=\R(\mathbf S)
\]
and \(\mathbf S\subseteq\mathbf S_1\subseteq\M\). Hence
\(\mathbf S_1=\M\) is established if we can show that \(\mathbf S_1\) is a
ring.

\(\mathbf S\) is an algebra and \(\alpha A\), \(A^*\) (as one-variable
functions), \(A+B\) (as a two-variable function), and \(AB\) (as a
one-variable function with respect to either variable) are continuous for
metric convergence. (This follows from the evaluations of
\cite[p.~242, property~II\(^{\mathrm o}\)]{ref6}.) Hence \(\mathbf S_1\) is
also an algebra. Therefore its closure with respect to metric convergence
implies its closure in the weak topology by \thmref{thm:I}{I}. Consequently,
\(\mathbf S_1\) is a ring.

This completes the proof.
\end{proof}

% Source PDF p. 45 (journal p. 760).
We can now prove

\begin{theorem}\phantomsection\label{thm:XI}
The (algebraical) type \(\overline{\M}\) of an \(\M\) which is approximately
finite of type \([p_1,p_2,\ldots]\) depends on the sequence
\([p_1,p_2,\ldots]\) alone and not on the particular choice of \(\M\).

(Provided such an \(\M\) exists at all. Cf. above.)
\end{theorem}

\begin{proof}
Consider two such \(\M^{(h)}\), \(h=1,2\), and form their
\(\N_1^{(h)},\N_2^{(h)},\ldots\) by \defref{def:4.1.1}{4.1.1} and their
\(\mathbf S^{(h)}\) by \lemref{lem:4.1.4}{4.1.4}. We wish to prove that
\(\M^{(1)}\) and \(\M^{(2)}\) are algebraically isomorphic.

We shall choose by induction for every \(n=1,2,\ldots\) an (algebraical)
isomorphism \(\mathfrak K_n^{(h)}\) between \(\N_n^{(h)}\) and the system of
all (complex) \(p_n\)-order matrices in the sense of (\(\alpha\)) in
\lemref{lem:4.1.1}{4.1.1}. For \(n=1\) we choose a \(\mathfrak K_1^{(h)}\)
which meets this requirement (cf. Remark 2 after
\lemref{lem:4.1.1}{4.1.1}). If for any \(n=2,3,\ldots\),
\(\mathfrak K_{n-1}^{(h)}\) has already been chosen, then we choose
\(\mathfrak K_n^{(h)}\) in the following way: Apply
\lemref{lem:4.1.2}{4.1.2} to \(\N_{n-1}^{(h)},\N_n^{(h)},
\mathfrak K_{n-1}^{(h)}\), in place of its \(\N,\mathbf O,\mathfrak K\)
(with \(\M^{(h)}\) in place of \(\M\)) and put
\(\mathfrak K_n^{(h)}=\mathfrak L\). Thus \(\mathfrak K_{n-1}^{(h)}\) and
\(\mathfrak K_n^{(h)}\) are related to each other as \(\mathfrak K\) and
\(\mathfrak L\) in (iii) in \lemref{lem:4.1.2}{4.1.2}.

Now form the (algebraical) isomorphism
\(\mathfrak J_n=(\mathfrak K_n^{(2)})^{-1}\mathfrak K_n^{(1)}\), of
\(\N_n^{(1)}\) and \(\N_n^{(2)}\). In view of the relationship given in the
last sentence of the preceding paragraph, \(\mathfrak J_n\) is a part of
\(\mathfrak J_{n+1}\). Thus the isomorphisms \(\mathfrak J_n\) of
\(\N_n^{(1)}\) and \(\N_n^{(2)}\), for \(n=1,2,\ldots\), are compatible with
each other and they merge to an (algebraical) isomorphism \(\mathfrak J\) of
\(\mathbf S^{(1)}\) and \(\mathbf S^{(2)}\). In \(\N_n^{(1)}\), our
\(\mathfrak J\) agrees with \(\mathfrak J_n\). Hence \(\mathfrak J\) maps
\(\N_n^{(1)}\) on \(\N_n^{(2)}\).

In \(\N_n^{(h)}\), \(\operatorname {Tr}_{\N_n^{(h)}}(A)\) is purely
algebraical (cf. \secref{sec:1.2}{1.2}). Clearly
\(\operatorname {Tr}_{\M^{(h)}}(A)\) will do as
\(\operatorname {Tr}_{\N_n^{(h)}}(A)\) for the \(A\in\N_n^{(h)}\) in the
sense of \cite[p.~219, property~IV]{ref6}. And as the trace is uniquely
determined by that property, so
\(\operatorname {Tr}_{\M^{(h)}}(A)=\operatorname {Tr}_{\N_n^{(h)}}(A)\) for
\(A\in\N_n^{(h)}\). Hence
\[
  \Metric{A}
  =\bigl(\operatorname {Tr}_{\M^{(h)}}(A^*A)\bigr)^{1/2}
  =\bigl(\operatorname {Tr}_{\N_n^{(h)}}(A^*A)\bigr)^{1/2}
  \quad\text{for }A\in\N_n^{(h)},
\]
and therefore \(\Metric{A}\), and with it \(\Metric{A-B}\), is purely
algebraical in \(\N_n^{(h)}\).

Thus \(\mathfrak J\) leaves \(\Metric{A-B}\) invariant in \(\N_n^{(h)}\)
for all \(n=1,2,\ldots\). Hence the same is true in all \(\mathbf S^{(h)}\).

Summing up: \(\mathfrak J\) is an isomorphism of \(\mathbf S^{(1)}\) and
\(\mathbf S^{(2)}\) with respect to
\begin{equation}\label{eq:4.1.gamma}
  1,\quad \alpha A,\quad A^*,\quad A+B,\quad AB,
  \tag{4.1.\ensuremath{\gamma}}
\end{equation}
and
\begin{equation}\label{eq:4.1.delta}
  \Metric{A-B}.
  \tag{4.1.\ensuremath{\delta}}
\end{equation}
By \eqref{eq:4.1.delta}, \(\mathfrak J\) is isometric. Now extend
\(\mathfrak J\) as far as possible in \(\mathbf B\) by metric convergence,
thus obtaining a new isometric (and hence one-to-one) correspondence
\(\mathfrak F\). By (ii) in \lemref{lem:4.1.4}{4.1.4}, \(\mathfrak F\) maps
all of \(\M^{(1)}\) on \(\M^{(2)}\).

The continuity of the operations \eqref{eq:4.1.gamma} (cf. the proof of (ii)
in \lemref{lem:4.1.4}{4.1.4}) guarantees that \(\mathfrak F\) too leaves
\eqref{eq:4.1.gamma} invariant. Hence \(\mathfrak F\) is an (algebraical)
isomorphism of \(\M^{(1)}\) and \(\M^{(2)}\). This completes the proof.
\end{proof}
 
\PaperSection{sec:4.2}{}
As a preliminary before introducing other definitions of approximate
finiteness, we develop further the ideas which appear in
\secref{sec:1.5}{1.5}.

% Source PDF p. 46 (journal p. 761).
\begin{lemma}\phantomsection\label{lem:4.2.1}
For any \(\epsilon>0\) there exists an
\(w_3=w_3(\epsilon)>0\) with the following property:

Consider a ring \(\N\subseteq\M\) and a projection \(E\in\M\). If there
exists an \(A\in\N\) with \(\Metric{A-E}<w_3\), then there exists also a
projection \(F\in\N\) with \(\Metric{F-E}<\epsilon\).
\end{lemma}

\begin{proof}
Consider first the hypothetical \(A\). We have
\(\Metric{A^*-E}=\Metric{(A-E)^*}=\Metric{A-E}<w_3\). Since
\[
  \tfrac12(A+A^*)-E=\tfrac12\bigl((A-E)+(A^*-E)\bigr),
\]
these inequalities imply \(\Metric{\tfrac12(A+A^*)-E}<w_3\) by the triangle
inequality. Thus we may replace \(A\) by \(\tfrac12(A+A^*)\), i.e. we may
assume that \(A\) is Hermitian.

Now form the following functions
\[
\psi_1(\lambda)=
\begin{cases}
1,&|\lambda|\ge1,\\
2|\lambda|-1,&\tfrac12\le|\lambda|<1,\\
0,&|\lambda|\le\tfrac12,
\end{cases}
\qquad
\psi_2(\lambda)=
\begin{cases}
1,&|\lambda|\ge\tfrac12,\\
2|\lambda|,&|\lambda|<\tfrac12.
\end{cases}
\]
Clearly \lemref{lem:1.5.3}{1.5.3} applies to both of them. We form the
\(w_2(\epsilon)\) of that lemma for both functions and denote their minimum by
\(w'_2(\epsilon)\). We put
\begin{equation}\label{eq:4.2.alpha}
  w_3(\epsilon)=w'_2\!\left(\frac{\epsilon}{2}\right).
  \tag{4.2.\ensuremath{\alpha}}
\end{equation}
Form the further function
\[
\varphi(\lambda)=
\begin{cases}
1,&|\lambda|\ge\tfrac12,\\
0,&|\lambda|<\tfrac12.
\end{cases}\footnote{Observe that \(\psi_1(\lambda)\), \(\psi_2(\lambda)\)
are continuous, while \(\varphi(\lambda)\) is not.}
\]
The following facts are clear:
\(\overline{\varphi(\lambda)}=\varphi(\lambda)=\varphi(\lambda)^2\). Hence
% Editorial note E20: The source prints "M" here; "N" is required by the
% conclusion of the lemma and follows from functional calculus inside N.
\(F=\varphi(A)\) is a projection which belongs to \(\N\) along with \(A\).
\(\psi_1(\lambda)=\psi_2(\lambda)=\varphi(\lambda)=\lambda\) for
\(\lambda=0,1\), the entire spectrum of \(E\). Hence
\begin{equation}\label{eq:4.2.beta}
  \psi_1(E)=\psi_2(E)=\varphi(E)=E.
  \tag{4.2.\ensuremath{\beta}}
\end{equation}
Finally
\((1-\varphi(\lambda))\psi_1(\lambda)
 +\varphi(\lambda)\psi_2(\lambda)=\varphi(\lambda)\). Hence
\begin{equation}\label{eq:4.2.gamma}
  (1-F)\psi_1(A)+F\psi_2(A)=F.
  \tag{4.2.\ensuremath{\gamma}}
\end{equation}
Now \(\Metric{A-E}<w_3(\epsilon)\) implies by
\lemref{lem:1.5.3}{1.5.3} and \eqref{eq:4.2.alpha} that
\[
  \Metric{\psi_1(A)-\psi_1(E)}<\frac{\epsilon}{2},
  \qquad
  \Metric{\psi_2(A)-\psi_2(E)}<\frac{\epsilon}{2}.
\]
Hence \eqref{eq:4.2.beta} now yields that
\(\Metric{\psi_1(A)-E}<\epsilon/2\) and
\(\Metric{\psi_2(A)-E}<\epsilon/2\). Since \(\OpNorm{F}\) and
\(\OpNorm{1-F}\) are \(\le1\)\footnote{Actually we have an equality in the
case of each of these operators except when it is zero.} we can therefore
conclude that
\[
  \Metric{(1-F)(\psi_1(A)-E)+F(\psi_2(A)-E)}<\epsilon.
\]
This may be rewritten
\[
  \Metric{(1-F)\psi_1(A)+F\psi_2(A)-E}<\epsilon.
\]
With \eqref{eq:4.2.gamma} this gives \(\Metric{F-E}<\epsilon\), as desired.
\end{proof}

% Source PDF p. 47 (journal p. 762).
\begin{lemma}\phantomsection\label{lem:4.2.2}
For any \(\epsilon>0\), there exists an
\(w_4=w_4(\epsilon)>0\) with the following property.

Consider two projections \(E,F\in\M\) with the closed linear sets
\(\mathfrak M,\mathfrak N\). If \(\Metric{E-F}<w_4\), then there exists a
partially isometric \(W'\in\M\) with an initial set
\(\mathfrak M'\subseteq\mathfrak M\), a final set
\(\mathfrak N'\subseteq\mathfrak N\), and with
\(\Metric{W'-E}<\epsilon\).
\end{lemma}

\begin{proof}
Form the canonical decomposition of \(A=FE\) in the sense of
\cite[p.~143, Definition~4.4.1]{ref5}:
\begin{equation}\label{eq:4.2.delta}
  FE=A=W'B,qquad EF=A^*=BW'^*.
  \tag{4.2.\ensuremath{\delta}}
\end{equation}
Here \(B\) is Hermitian and definite, \(W'\) is partially isometric with
initial set
\[
  \mathfrak M'=[\operatorname {Range}A^*]
  =[\operatorname {Range}EF]
  \subseteq[\operatorname {Range}E]=\mathfrak M
\]
and final set
\[
  \mathfrak N'=[\operatorname {Range}A]
  =[\operatorname {Range}FE]
  \subseteq[\operatorname {Range}F]=\mathfrak N.
\]
The range of \(W'^*\) is
\(\mathfrak M'\subseteq\mathfrak M=[\operatorname {Range}E]\). Thus
\(EW'^*=W'^*\), and taking adjoints, we obtain
\begin{equation}\label{eq:4.2.epsilon}
  W'E=W'.
  \tag{4.2.\ensuremath{\epsilon}}
\end{equation}
For any canonical decomposition \(B^2=A^*A\). In our case
\(A^*A=EF\cdot FE=EFE\), so
\begin{equation}\label{eq:4.2.zeta}
  B^2=EFE.
  \tag{4.2.\ensuremath{\zeta}}
\end{equation}
Observe also that due to the character of these operators
\begin{equation}\label{eq:4.2.eta}
  \OpNorm{E}\le1,\qquad \OpNorm{F}\le1,\qquad \OpNorm{W'}\le1.
  \tag{4.2.\ensuremath{\eta}}
\end{equation}
Thus \eqref{eq:4.2.zeta} implies \(\OpNorm{B^2}\le1\), and so
\begin{equation}\label{eq:4.2.theta}
  \OpNorm{B}\le1.
  \tag{4.2.\ensuremath{\theta}}
\end{equation}
Now form the function
\[
\psi(\lambda)=
\begin{cases}
1,&|\lambda|\ge1,\\
\sqrt{|\lambda|},&|\lambda|\le1.
\end{cases}
\]
Clearly \lemref{lem:1.5.3}{1.5.3} applies to it. We form the
\(w_2(\epsilon)\) of that lemma for this function and put
\begin{equation}\label{eq:4.2.iota}
  w_4(\epsilon)=\operatorname {Min.}
  \left(w_2\!\left(\frac{\epsilon}{2}\right),\frac{\epsilon}{2}\right).
  \tag{4.2.\ensuremath{\iota}}
\end{equation}
It is clear that \(\psi(\lambda^2)=\lambda\) for \(0\le\lambda\le1\),
which includes the entire spectrum of \(B\), since \(B\) is definite and
\eqref{eq:4.2.theta} holds. Thus \(\psi(B^2)=B\), or using
\eqref{eq:4.2.zeta} we obtain
\begin{equation}\label{eq:4.2.kappa}
  \psi(EFE)=B.
  \tag{4.2.\ensuremath{\kappa}}
\end{equation}

% Source PDF p. 48 (journal p. 763).
On the other hand, \(\psi(\lambda)=\lambda\) for \(\lambda=0,1\), the
spectrum of \(E\), and hence
\begin{equation}\label{eq:4.2.lambda}
  \psi(E)=E.
  \tag{4.2.\ensuremath{\lambda}}
\end{equation}
Now our assumption, \(\Metric{E-F}<w_4\), implies by
\eqref{eq:4.2.eta} that \(\Metric{(E-F)E}<w_4\) and
\(\Metric{E(E-F)E}<w_4\), or
\begin{equation}\label{eq:4.2.mu}
  \Metric{E-FE}<w_4,qquad \Metric{E-EFE}<w_4.
  \tag{4.2.\ensuremath{\mu}}
\end{equation}
\eqref{eq:4.2.mu}, \lemref{lem:1.5.3}{1.5.3}, and
\eqref{eq:4.2.iota} together imply
\[
  \Metric{\psi(E)-\psi(EFE)}<\frac{\epsilon}{2}.
\]
This and \eqref{eq:4.2.kappa} and \eqref{eq:4.2.lambda} yield
\(\Metric{E-B}<\epsilon/2\). This and \eqref{eq:4.2.eta} imply
\(\Metric{W'(E-B)}<\epsilon/2\). Hence \eqref{eq:4.2.epsilon} and
\eqref{eq:4.2.delta} give
\begin{equation}\label{eq:4.2.nu}
  \Metric{W'-FE}<\frac{\epsilon}{2}.
  \tag{4.2.\ensuremath{\nu}}
\end{equation}
\eqref{eq:4.2.mu} and \eqref{eq:4.2.iota} imply
\(\Metric{E-FE}<\epsilon/2\). Hence \eqref{eq:4.2.nu} yields
\(\Metric{W'-E}<\epsilon\), as desired.
\end{proof}

\begin{lemma}\phantomsection\label{lem:4.2.3}
If \(U\in\M\) is unitary and \(W\in\M\) is partially isometric, and if
\(U\) agrees with \(W\) on its initial set, then
\[
  \Metric{U-W}\le\bigl(2\Metric{W-1}\bigr)^{1/2}.
\]
\end{lemma}

\begin{proof}
The initial set of \(W\) is the range or final set of \(W^*\). Therefore
\(UW^*=WW^*\). This implies for the adjoints \(WU^*=WW^*\). Hence
\[
\begin{aligned}
(U-W)(U-W)^*
 &=UU^*-UW^*-WU^*+WW^*\\
 &=1-WW^*-WW^*+WW^*=1-WW^*.
\end{aligned}
\]
Therefore
\(\Metric{U-W}^2
=\operatorname {Tr}_{\M}((U-W)(U-W)^*)
=\operatorname {Tr}_{\M}(1-WW^*)\), or
\begin{equation}\label{eq:4.2.omicron}
  \Metric{U-W}=\bigl(\operatorname {Tr}_{\M}(1-WW^*)\bigr)^{1/2}.
  \tag{4.2.\ensuremath{o}}
\end{equation}
We now make use of Schwartz's inequality in the sense of
\cite[p.~241, Theorem~VIII]{ref6}, remembering that \(\OpNorm{W}\le1\), and
so \(\Metric{W}\le1\). Thus
\[
\begin{aligned}
\operatorname {Tr}_{\M}(1-WW^*)
 &=\operatorname {Tr}_{\M}\bigl(-W(W-1)^*-(W-1)1^*\bigr)\\
 &\le2\Metric{W-1}.
\end{aligned}
\]
This and \eqref{eq:4.2.omicron} yield the desired inequality of the lemma.
\end{proof}

\begin{lemma}\phantomsection\label{lem:4.2.4}
For any \(\epsilon>0\), there exists an
\(w_5=w_5(\epsilon)>0\) with the following property:

Consider two projections \(E,F\in\M\) with
\(D_{\M}(E)=D_{\M}(F)\). If \(\Metric{E-F}<w_5\), then there exists a
unitary \(U\in\M\) with \(F=UEU^{-1}\) and with
\(\Metric{U-1}<\epsilon\).
\end{lemma}

\begin{proof}
Put \(w_5(\epsilon)=w_4(\epsilon')\), where
\(\epsilon'=\epsilon'(\epsilon)>0\) will be chosen later.

Apply \lemref{lem:4.2.2}{4.2.2} to \(E,F\) and their closed linear sets
\(\mathfrak M,\mathfrak N\). It gives a partially isometric \(W'\in\M\)
with an initial set \(\mathfrak M'\subseteq\mathfrak M\), a final set
\(\mathfrak N'\subseteq\mathfrak N\), and with
\begin{equation}\label{eq:4.2.xi}
  \Metric{W'-E}<\epsilon'.
  \tag{4.2.\ensuremath{\xi}}
\end{equation}

% Source PDF p. 49 (journal p. 764).
In view of the relation of \(W'\) to \(\mathfrak M'\) and \(\mathfrak N'\),
\(D_{\M}(\mathfrak M')=D_{\M}(\mathfrak N')\). By hypothesis,
\(D_{\M}(\mathfrak M)=D_{\M}(\mathfrak N)\). Hence
\[
  D_{\M}(\mathfrak M\ominus\mathfrak M')
  =D_{\M}(\mathfrak N\ominus\mathfrak N').
\]
Choose a partially isometric \(V'\in\M\) with initial set
\(\mathfrak M\ominus\mathfrak M'\) and final set
\(\mathfrak N\ominus\mathfrak N'\). Then \(U'=W'+V'\in\M\) is also
partially isometric, with initial set \(\mathfrak M\), final set
\(\mathfrak N\), and it agrees with \(W'\) on the initial set
\(\mathfrak M'\) of \(W'\).

Apply all this to \(1-E,1-F\) and
\(\mathfrak M^{\perp}(=\Hilbert\ominus\mathfrak M),\mathfrak N^{\perp}\)
in place of \(E,F\) and \(\mathfrak M,\mathfrak N\). (Clearly
\(\Metric{(1-E)-(1-F)}=\Metric{E-F}\).) Thus we obtain a partially
isometric \(W''\in\M\) with an initial set
\(\mathfrak M''\subseteq\mathfrak M^{\perp}\), a final set
\(\mathfrak N''\subseteq\mathfrak N^{\perp}\), and with
\begin{equation}\label{eq:4.2.pi}
  \Metric{W''-(1-E)}<\epsilon'.
  \tag{4.2.\ensuremath{\pi}}
\end{equation}
Furthermore, there exists a partially isometric \(U''\in\M\) with the initial
set \(\mathfrak M^{\perp}\), final set \(\mathfrak N^{\perp}\), and agreeing
with \(W''\) on the initial set \(\mathfrak M''\) of \(W''\).

Now put
\[
  W=W'+W'',\qquad U=U'+U''.
\]
Then \(W,U\in\M\) are partially isometric, with initial sets
\([\mathfrak M',\mathfrak M'']\),
\([\mathfrak M,\mathfrak M^{\perp}]=\Hilbert\), respectively, and with the
respective final sets \([\mathfrak N',\mathfrak N'']\) and
\([\mathfrak N,\mathfrak N^{\perp}]=\Hilbert\). So \(U\) is unitary.

\(U\) agrees with \(U'\) on the initial set \(\mathfrak M\) of \(U'\). Hence
it maps \(\mathfrak M\) on \(\mathfrak N\). Therefore
\begin{equation}\label{eq:4.2.rho}
  F=UEU^{-1}.
  \tag{4.2.\ensuremath{\rho}}
\end{equation}
\(U\) agrees with \(W\) on the initial set
\([\mathfrak M',\mathfrak M'']\) of \(W\). Therefore by
\lemref{lem:4.2.3}{4.2.3},
\begin{equation}\label{eq:4.2.sigma}
  \Metric{U-W}=\bigl(2\Metric{W-1}\bigr)^{1/2}.
  \tag{4.2.\ensuremath{\sigma}}
\end{equation}
\eqref{eq:4.2.xi} and \eqref{eq:4.2.pi} imply
\begin{equation}\label{eq:4.2.tau}
  \Metric{W-1}<2\epsilon'.
  \tag{4.2.\ensuremath{\tau}}
\end{equation}
\eqref{eq:4.2.sigma} and \eqref{eq:4.2.tau} give together
\(\Metric{U-W}<(4\epsilon')^{1/2}\). This and \eqref{eq:4.2.tau} yield
\begin{equation}\label{eq:4.2.upsilon}
  \Metric{U-1}<(4\epsilon')^{1/2}+2\epsilon'.
  \tag{4.2.\ensuremath{\upsilon}}
\end{equation}
Hence if we choose \(\epsilon'\) so that
\((4\epsilon')^{1/2}+2\epsilon'\le\epsilon\),
\eqref{eq:4.2.upsilon} yields
\begin{equation}\label{eq:4.2.phi}
  \Metric{U-1}<\epsilon.
  \tag{4.2.\ensuremath{\phi}}
\end{equation}
In view, therefore, of \eqref{eq:4.2.rho} and \eqref{eq:4.2.phi}, the \(w_5\)
which we have specified satisfies the condition of our lemma.
\end{proof}

\PaperSection{sec:4.3}{}
We shall recast the definition of approximate finiteness in a series of steps,
i.e. Definitions 4.3.1, 4.5.2 and 4.6.1, below:

\begin{definition}\phantomsection\label{def:4.3.1}
\(\M\) is approximately finite (A) if this is true:

Given any \(A_1,\ldots,A_m\in\M\) and any \(\epsilon>0\), there exists an
\(n=n_1(A_1,\ldots,A_m,\epsilon)\) such that for every \(q\ge n\) there
exists a factor
\(\N=\N_1(q,A_1,\ldots,A_m,\epsilon)\) with the following properties:
\begin{enumerate}[label=\((\roman*)\)]
\item \(\N\) is in case \((\mathrm I_q)\).
\item \(\N\subseteq\M\).
\item There exist \(B_1,\ldots,B_m\in\N\) with
\[
  \Metric{B_h-A_h}<\epsilon\qquad\text{for }h=1,\ldots,m.
\]
\end{enumerate}
\end{definition}

% Source PDF p. 50 (journal p. 765).
In the present section we are interested in certain consequences of this
definition, in particular what other properties of the approximating factor
\(\N\) may be assumed. It will be assumed therefore in the five lemmas of
this section that \(\M\) is approximately finite (A).

\begin{lemma}\phantomsection\label{lem:4.3.1}
Given any \(A_1,\ldots,A_m\in\M\), any projection \(E\in\M\) with
\(D_{\M}(E)=1/p\) for \(p=1,2,\ldots\), and any \(\epsilon>0\), there
exists an \(n_2=n_2(A_1,\ldots,A_m,E,\epsilon)\) such that for every
\(q\ge n_2\) which is divisible by \(p\) there exists a factor
\(\N=\N_2(q,A_1,\ldots,A_m,E,\epsilon)\) with the following properties:
\begin{enumerate}[label=\((\roman*)\)]
\item \(\N\) is in case \((\mathrm I_q)\).
\item \(\N\subseteq\M\).
\item There exist \(B_1,\ldots,B_m\in\N\) with
\[
  \Metric{B_h-A_h}<\epsilon\qquad\text{for }h=1,\ldots,m.
\]
\item There exists a projection \(F\in\N\) with
\(D_{\M}(F)=D_{\M}(E)=1/p\), and
\(\Metric{F-E}<\epsilon\).\footnote{\hypertarget{fn:33}{}Observe that this is a
strengthened form of \defref{def:4.3.1}{4.3.1}.}
\end{enumerate}
\end{lemma}

\begin{proof}
Apply \defref{def:4.3.1}{4.3.1} with \(A_1,\ldots,A_m,E\) in place of its
\(A_1,\ldots,A_m\), and with \(\epsilon'\) in place of \(\epsilon\), where
\(\epsilon'=\epsilon'(\epsilon)\) will be chosen later.

Put
\[
\begin{aligned}
n_2(A_1,\ldots,A_m,E,\epsilon)
  &=n_1(A_1,\ldots,A_m,E,\epsilon'),\\
\N_2(q,A_1,\ldots,A_m,E,\epsilon)
  &=\N_1(q,A_1,\ldots,A_m,E,\epsilon').
\end{aligned}
\]
Now consider a \(q\ge n_2\) which is divisible by \(p\), and its
\(\N=\N_2\).

In view of the above definition of \(\N_2\), (i) and (ii) of the present
lemma follow from (i) and (ii) of \defref{def:4.3.1}{4.3.1}. Furthermore,
(iii) of \defref{def:4.3.1}{4.3.1} states that
\begin{equation}\label{eq:4.3.alpha}
  \Metric{B_h-A_h}<\epsilon'\qquad\text{for }h=1,\ldots,m,
  \tag{4.3.\ensuremath{\alpha}}
\end{equation}
and
\begin{equation}\label{eq:4.3.beta}
  \Metric{B_{m+1}-E}<\epsilon'.
  \tag{4.3.\ensuremath{\beta}}
\end{equation}
Put
\begin{equation}\label{eq:4.3.gamma}
  \epsilon'=\operatorname {Min.}\bigl(w_3(\epsilon''),\epsilon\bigr),
  \tag{4.3.\ensuremath{\gamma}}
\end{equation}
where \(\epsilon''=\epsilon''(\epsilon)\) will be chosen later. Now
\eqref{eq:4.3.alpha} and \eqref{eq:4.3.gamma} imply (iii) in the lemma, so
we must only derive (iv) from \eqref{eq:4.3.beta}.

% Source PDF p. 51 (journal p. 766).
By \lemref{lem:4.2.1}{4.2.1}, \eqref{eq:4.3.beta} and
\eqref{eq:4.3.gamma} imply the existence of a projection \(F_1\in\N\) with
\begin{equation}\label{eq:4.3.delta}
  \Metric{F_1-E}<\epsilon''.
  \tag{4.3.\ensuremath{\delta}}
\end{equation}
This implies
\begin{equation}\label{eq:4.3.epsilon}
  \bigl|D_{\M}(F_1)-D_{\M}(E)\bigr|<\epsilon''.\footnote{For
  any two projections \(E,F\in\M\),
  \[
    \bigl|D_{\M}(F)-D_{\M}(E)\bigr|\le\Metric{F-E}.
  \]
  \textit{Proof.} We have
  \[
    \bigl|D_{\M}(F)-D_{\M}(E)\bigr|
    =\bigl|\operatorname {Tr}_{\M}(F)-\operatorname {Tr}_{\M}(E)\bigr|
    =\bigl|\operatorname {Tr}_{\M}(F-E)\bigr|.
  \]
  This is \(\le\Metric{1}\cdot\Metric{F-E}\) (cf.
  \cite[p.~241, Theorem~VIII]{ref6}); hence \(\le\Metric{F-E}\), as
  desired.}
  \tag{4.3.\ensuremath{\epsilon}}
\end{equation}
As \(\N\) is in case \((\mathrm I_q)\), and \(q\) is divisible by \(p\),
there exists a projection \(F\in\N\) with
\begin{equation}\label{eq:4.3.zeta}
  D_{\M}(F)=D_{\M}(E)=1/p,
  \tag{4.3.\ensuremath{\zeta}}
\end{equation}
and we can even choose \(F_1\gtreqless F\). This latter implies
\[
  \Metric{F_1-F}
  =\bigl(\lvert D_{\M}(F_1)-D_{\M}(F)\rvert\bigr)^{1/2}.\footnote{
  For any two projections \(E,F\in\M\) with
  \(E\gtreqless F\),
  \[
    \Metric{F-E}=\bigl(\lvert D_{\M}(F)-D_{\M}(E)\rvert\bigr)^{1/2}.
  \]
  \textit{Proof.} By symmetry we may assume \(E\le F\). Then \(F-E\) is a
  projection, \(D_{\M}(F)-D_{\M}(E)=D_{\M}(F-E)\). Put \(G=F-E\). We
  must prove \(\Metric{G}=(D_{\M}(G))^{1/2}\). Now
  \[
    \Metric{G}^2=\operatorname {Tr}_{\M}(G^*G)
    =\operatorname {Tr}_{\M}(G)=D_{\M}(G).
  \]}
\]
This equation, \eqref{eq:4.3.epsilon}, and \eqref{eq:4.3.zeta} together
imply \(\Metric{F_1-F}<(\epsilon'')^{1/2}\). Hence
\eqref{eq:4.3.delta} now implies
\begin{equation}\label{eq:4.3.eta}
  \Metric{F-E}<(\epsilon'')^{1/2}+\epsilon''.
  \tag{4.3.\ensuremath{\eta}}
\end{equation}
Thus if we choose \(\epsilon''=\epsilon''(\epsilon)>0\), so that
\((\epsilon'')^{1/2}+\epsilon''\le\epsilon\), we can conclude from
\eqref{eq:4.3.eta} that \(\Metric{F-E}<\epsilon\). In view of
\eqref{eq:4.3.zeta} and \(F\in\N\), our last inequality shows that \(\N\)
has property (iv), as desired.
\end{proof}

\begin{lemma}\phantomsection\label{lem:4.3.2}
Given any \(A_1,\ldots,A_m\in\M\), any projection \(E\in\M\) with
\(D_{\M}(E)=1/p\) for \(p=1,2,\ldots\), and any \(\epsilon>0\). Then there
exists an \(n_3=n_3(A_1,\ldots,A_m,E,\epsilon)\) such that for every
\(q\ge n_3\) which is divisible by \(p\) there exists a factor
\(\N=\N_3(q,A_1,\ldots,A_m,E,\epsilon)\) with the properties (i), (ii),
(iii) of \lemref{lem:4.3.1}{4.3.1}, and in addition
\begin{enumerate}[label=\((\roman*')\),start=4]
\item \(E\in\N\).\textsuperscript{\hyperlink{fn:33}{33}}
\end{enumerate}
\end{lemma}

\begin{proof}
Our hypotheses are identical with those of \lemref{lem:4.3.1}{4.3.1}, and we
apply the latter to \(A_1,\ldots,A_m,E\), but with \(\epsilon'\) in place of
\(\epsilon\), where
\begin{equation}\label{eq:4.3.theta}
  \epsilon'=\operatorname {Min.}\bigl(w_5(\epsilon''),\epsilon''\bigr),
  \tag{4.3.\ensuremath{\theta}}
\end{equation}
and \(\epsilon''=\epsilon''(\epsilon)>0\) will be chosen later.

Denote the \(n_2,\N_2\), and \(B_1,\ldots,B_m,F\) which we obtain in this
way by \(n'_2,\N_2^+\), and \(B'_1,\ldots,B'_m,F'\). We put
\(n_3=n'_2\) and \(\N^+=\N_2^+\). In view of \eqref{eq:4.3.theta},
\lemref{lem:4.3.1}{4.3.1}(iii) gives
\begin{equation}\label{eq:4.3.iota}
  \Metric{B'_h-A_h}<\epsilon''\qquad\text{for }h=1,\ldots,m,
  \tag{4.3.\ensuremath{\iota}}
\end{equation}
while (iv) of that lemma yields
\begin{equation}\label{eq:4.3.kappa}
  \Metric{F'-E}<w_5(\epsilon'')
  \tag{4.3.\ensuremath{\kappa}}
\end{equation}
and
\begin{equation}\label{eq:4.3.lambda}
  D_{\M}(F')=D_{\M}(E).
  \tag{4.3.\ensuremath{\lambda}}
\end{equation}
\eqref{eq:4.3.kappa} and \eqref{eq:4.3.lambda} allow us to apply
\lemref{lem:4.2.4}{4.2.4} with \(F',E\) in place of its \(E,F\).
Consequently there exists a unitary \(U\in\M\) with
\begin{equation}\label{eq:4.3.mu}
  \Metric{U-1}<\epsilon''
  \tag{4.3.\ensuremath{\mu}}
\end{equation}
and
\begin{equation}\label{eq:4.3.nu}
  E=UF'U^{-1}.
  \tag{4.3.\ensuremath{\nu}}
\end{equation}

% Source PDF p. 52 (journal p. 767).
\eqref{eq:4.3.mu} implies
\[
\begin{aligned}
\Metric{A_h-U^{-1}A_hU}
 &=\Metric{A_h-A_hU+A_hU-U^{-1}A_hU}\\
 &\le\OpNorm{A_h}\Metric{1-U}
   +\Metric{(1-U^{-1})A_hU}\\
 &\le2\OpNorm{A_h}\epsilon''.
\end{aligned}
\]
Hence by \eqref{eq:4.3.iota},
\(\Metric{B'_h-U^{-1}A_hU}<(2\OpNorm{A_h}+1)\epsilon''\), and so
\begin{equation}\label{eq:4.3.omicron}
  \Metric{UB'_hU^{-1}-A_h}<(2\OpNorm{A_h}+1)\epsilon''.
  \tag{4.3.\ensuremath{o}}
\end{equation}
Now we choose \(\epsilon''=\epsilon''(\epsilon)>0\) as
\begin{equation}\label{eq:4.3.xi}
  \epsilon''=
  \frac{\epsilon}{2\operatorname {Max}_{h=1,\ldots,m}\OpNorm{A_h}+1}.
  \tag{4.3.\ensuremath{\xi}}
\end{equation}
Put \(\N=\N_3=U\N^+U^{-1}\). Then \(\N\) is a factor in the case
\((\mathrm I_q)\) along with \(\N^+\). Thus (i) holds; (ii) holds since
\(\N^+\subseteq\M\), \(U\in\M\) implies \(\N\subseteq\M\).
\(B'_h\in\N^+\) gives \(B_h=UB'_hU^{-1}\in\N\), while
\eqref{eq:4.3.omicron} and \eqref{eq:4.3.xi} imply
\(\Metric{B_h-A_h}<\epsilon\), so (iii) holds. And finally
\(F'\in\N^+\) and \eqref{eq:4.3.nu} imply
\(E=UF'U^{-1}\in\N\), so that (iv') holds.
\end{proof}

\begin{lemma}\phantomsection\label{lem:4.3.3}
If we add to the assumptions of \lemref{lem:4.3.2}{4.3.2} that
\begin{equation}\label{eq:4.3.local-alpha}
  EA_h=A_hE=A_h,
  \tag{\ensuremath{\alpha}}
\end{equation}
then we can add to its assertions that
\begin{equation}\label{eq:4.3.local-beta}
  EB_h=B_hE=B_h.\footnote{\hypertarget{fn:36}{}With the same \(n_3,\N_3\) as in
  that lemma.}
  \tag{\ensuremath{\beta}}
\end{equation}
\end{lemma}

\begin{proof}
Apply \lemref{lem:4.3.2}{4.3.2}, but write \(B'_h\) for its \(B_h\). Now
put \(B_h=EB'_hE\). Then \eqref{eq:4.3.local-beta} is satisfied and
\(B'_h,E\in\N\) give \(B_h\in\N\). By \eqref{eq:4.3.local-alpha},
\(EA_hE=A_h\). By (iii) of \lemref{lem:4.3.2}{4.3.2},
\(\Metric{B'_h-A_h}<\epsilon\). Hence
\(\Metric{E(B'_h-A_h)E}<\epsilon\), or
\(\Metric{B_h-A_h}<\epsilon\). Thus (iii) holds for \(B_h\) as well as for
\(B'_h\).

The other assertions ((i), (ii), (iv')) are not affected by our replacing
\(B'_h\) by \(B_h\). Thus the proof is completed.
\end{proof}

% Source PDF p. 53 (journal p. 768).
\begin{lemma}\phantomsection\label{lem:4.3.4}
If we add to the assumption of \lemref{lem:4.3.3}{4.3.3} that
\(E\in\N^0\), where \(\N^0\) is a factor in case \((\mathrm I_p)\), then
we can add to its assertions
\begin{enumerate}[label=\((\roman*)\),start=5]
\item \(\N^0\subseteq\N\).\textsuperscript{\hyperlink{fn:33}{33},
\hyperlink{fn:36}{36}}
\end{enumerate}
\end{lemma}

\begin{proof}
Apply \lemref{lem:4.3.3}{4.3.3}, but write \(\N^+\) for its \(\N\).

The considerations at the beginning of the proof of (i) in
\lemref{lem:4.1.2}{4.1.2} show again that
\(D_{\M}(G)=D_{\N^0}(G)\) for all projections
\(G\in\N^0(\subseteq\M)\).

Hence in particular \(D_{\N^0}(E)=D_{\M}(E)=1/p\). Then the argument used
in the proof of \lemref{lem:2.6.2}{2.6.2} establishes the existence of a
system of \(p\)-order matrix units \(W^0_{u,v}\),
\(u,v=1,\ldots,p\), in \(\N^0\), with
\(\sum_{u=1}^{p}W^0_{u,u}=1\) and \(W^0_{1,1}=E\). Thus in terms of
(\(\beta\)) of \lemref{lem:4.1.1}{4.1.1},
\begin{equation}\label{eq:4.3.pi}
\begin{gathered}
\text{There exists a matrix base }W^0_{u,v},\quad u,v=1,\ldots,p,
  \text{ of }\N^0,\\
\text{with }W^0_{1,1}=E.
\end{gathered}
  \tag{4.3.\ensuremath{\pi}}
\end{equation}

On the other hand, the argument of \lemref{lem:2.6.3}{2.6.3} applied to
\(\N^+\) establishes the existence of a system of \(p\)-order matrix units
\(W^+_{u,v}\), \(u,v=1,\ldots,p\), in \(\N^+\), with
\(\sum_{u=1}^{p}W^+_{u,u}=1\) and \(W^+_{1,1}=E\). (Observe that this is not
a matrix base of \(\N^+\), since \(\N^+\) is in case \((\mathrm I_q)\) and
not \((\mathrm I_p)\)!)

Now put
\[
  U=\sum_{u=1}^{p}W^0_{u,1}W^+_{1,u}.
\]
As \(W^0_{u,1}\in\N^0\subseteq\M\), \(W^+_{1,u}\in\N^+\subseteq\M\), so
\(U\in\M\). One verifies by direct computation that \(U^*U=UU^*=1\),
i.e. \(U\) is unitary. Also
\begin{equation}\label{eq:4.3.rho}
  UW^+_{u,v}=W^0_{u,v}U
  \quad\text{or}\quad
  W^0_{u,v}=UW^+_{u,v}U^{-1},
  \tag{4.3.\ensuremath{\rho}}
\end{equation}
and, since \(W^0_{1,1}=W^+_{1,1}=E\),
\begin{equation}\label{eq:4.3.sigma}
  EU=UE=E.
  \tag{4.3.\ensuremath{\sigma}}
\end{equation}
Apply the mapping
\begin{equation}\label{eq:4.3.tau}
  X\longmapsto UXU^{-1}
  \tag{4.3.\ensuremath{\tau}}
\end{equation}
to \(\N^+\). This carries \(\N^+\) into \(\N=U\N^+U^{-1}\). Now we
see: \(\N\) is a factor in case \((\mathrm I_q)\) along with \(\N^+\); thus
(i) holds. \(\N^+\subseteq\M\), \(U\in\M\) imply \(\N\subseteq\M\); so
(ii) holds. By \eqref{eq:4.3.local-beta} in
\lemref{lem:4.3.3}{4.3.3} and \eqref{eq:4.3.sigma}, the mapping
\eqref{eq:4.3.tau} leaves all \(B_h\) fixed; hence they belong to \(\N\) as
well as to \(\N^+\). Now (iii) and (iv') in \lemref{lem:4.3.2}{4.3.2} are
unaffected and hence hold in the new case. By \eqref{eq:4.3.rho} the mapping
\eqref{eq:4.3.tau} carries \(W^+_{u,v}\) into \(W^0_{u,v}\). Hence
\(W^+_{u,v}\in\N^+\) gives \(W^0_{u,v}\in\N\). Hence by
\eqref{eq:4.3.pi} and (ii) in \lemref{lem:4.1.1}{4.1.1}, every
\(A\in\N^0\) has the form
\[
  A=\sum_{u=1}^{p}\sum_{v=1}^{p}a_{v,u}W^0_{u,v},
\]
so \(A\in\N\). Thus \(\N^0\subseteq\N\), which is the desired relation
(v).

Thus the proof is completed.
\end{proof}

% Fragment: Lemma 4.3.5 through the end of Chapter V.

\begin{lemma}\phantomsection\label{lem:4.3.5}
Assume that \(\M\) is approximately finite (A). Given any
\(A_1,\ldots,A_m\in\M\), any \(p=1,2,\ldots\), and any \(\epsilon>0\),
then there exists an
\[
 n_4=n_4(A_1,\ldots,A_m,p,\epsilon)
\]
such that for every \(q\geq n_4\) which is divisible by \(p\) and every factor
\(\N^0\subseteq\M\) in case \((\mathrm I_p)\), there exists a factor
\[
 \N=\N_4(q,A_1,\ldots,A_m,\N^0,p,\epsilon)
\]
with the following properties:
\begin{enumerate}[label=(\roman*)]
\item \(\N\) is in case \((\mathrm I_q)\).
\item \(\N\subseteq\M\).
\item There exist \(B_1,\ldots,B_m\in\N\) with
      \(\Metric{B_h-A_h}<\epsilon\) for \(h=1,\ldots,m\).
\setcounter{enumi}{4}
\item \(\N^0\subseteq\N\).
\end{enumerate}
(Note. This lemma is quite similar to \lemref{lem:4.3.4}{4.3.4}. However,
(\(\alpha\)) has been dropped from the assumptions and (\(\beta\)) from the
conclusion. Thus \(E\) is no longer mentioned.)
\end{lemma}

\begin{proof}
Choose a matrix base \(W_{u,v}\), \(u,v=1,\ldots,p\), of \(\N^0\), and let
\(E=W_{1,1}\). Put
\[
 A_{u,v;h}=W_{1,u}A_hW_{v,1}.
\]
We have \(D_{\M}(E)=D_{\M}(W_{1,1})=1/p\) (cf. the argument of
\eqref{eq:2.7.alpha}). Clearly
\[
 EA_{u,v;h}=A_{u,v;h}E=A_{u,v;h}.
\]
Now apply \lemref{lem:4.3.4}{4.3.4} with the \(mp^2\) operators
\(A_{u,v;h}\), \(h=1,\ldots,m\), \(u,v=1,\ldots,p\), in place of
\(A_1,\ldots,A_m\), and \(\epsilon/p^2\) in place of \(\epsilon\).

Denote the \(n_3,\N_3\) which obtain in this way by \(n_4,\N_4\). Write
\(\N=\N_4\). Denote the operators which correspond to the \(A_{u,v;h}\)
(as the \(B_h\) do to the \(A_h\) in \lemref{lem:4.3.4}{4.3.4}) by
\(B_{u,v;h}\). Clearly:
\begin{equation}\tag{4.3.\(\nu\)}\label{eq:4.3-nu}
 A_h=\sum_{u=1}^{p}\sum_{v=1}^{p}W_{u,1}A_{u,v;h}W_{1,v}.
\end{equation}
Let
\begin{equation}\tag{4.3.\(\phi\)}\label{eq:4.3-phi}
 B_h=\sum_{u=1}^{p}\sum_{v=1}^{p}W_{u,1}B_{u,v;h}W_{1,v}.
\end{equation}
Now we argue as follows. (i), (ii), and (v) in \lemref{lem:4.3.4}{4.3.4}
are unaffected. \(B_{u,v;h}\in\N\) and \(W_{u,v}\in\N^0\subseteq\N\)
imply \(B_h\in\N\). We have
\(\Metric{B_{u,v;h}-A_{u,v;h}}<\epsilon/p^2\). Hence
\[
 \Metric{W_{u,1}(B_{u,v;h}-A_{u,v;h})W_{1,v}}<\epsilon/p^2
\]
and
\[
 \Metric{\sum_{u=1}^{p}\sum_{v=1}^{p}
 W_{u,1}(B_{u,v;h}-A_{u,v;h})W_{1,v}}<\epsilon.
\]
Thus \eqref{eq:4.3-nu}, \eqref{eq:4.3-phi}, and this inequality yield
\(\Metric{B_h-A_h}<\epsilon\). So (iii) holds and the proof is complete.
\end{proof}

\PaperSection{sec:4.4}{Comparison of (A) and type \([p_1,p_2,\ldots]\).}
We now use the results of the preceding two sections to compare the two notions
of approximate finiteness which we have introduced.

\begin{lemma}\phantomsection\label{lem:4.4.1}
Assume that \(\M\) is approximately finite (A) and that the sequence
\([p_1,p_2,\ldots]\) fulfills (i), (ii) in \lemref{lem:4.1.3}{4.1.3}. Then
\([p_1,p_2,\ldots]\) possesses a subsequence
\([p'_1,p'_2,\ldots]\) such that \(\M\) is approximately finite of type
\([p'_1,p'_2,\ldots]\).
\end{lemma}

\begin{proof}
By \cite{ref2}, pp.~387--388, there exists a sequence
\(A_1,A_2,\ldots\in\M\) which is everywhere dense in \(\M\) in the sense of
strong convergence. Put \(p'_0=1\) and \(\N_0=(\alpha\mathbin{\cdot}1)\).
This is a factor \(\subseteq\M\) in case \((\mathrm I_1)\). We shall choose by
induction, for every \(n=1,2,\ldots\), a \(p'_n\) and a factor
\(\N_n\subseteq\M\) in the case \((\mathrm I_{p'_n})\). If, for any
\(n=1,2,\ldots\), \(p'_{n-1}\) and \(\N_{n-1}\) have already been chosen,
then we choose \(p'_n\) and \(\N_n\) in the following way. Apply
\lemref{lem:4.3.5}{4.3.5} to
\[
 A_1,\ldots,A_n,\quad p'_{n-1},\quad 1/n,\quad \N_{n-1}
\]
in place of its \(A_1,\ldots,A_m,p,\epsilon,\N^0\). Choose \(p'_n\) from
the sequence \([p_1,p_2,\ldots]\) with
\(p'_n>\operatorname{Max}(n_4,p'_{n-1})\). Thus \(p'_n\) is divisible by
\(p'_{n-1}\). Put \(q=p'_n\), and then put \(\N_n=\N\) (from the above
application of \lemref{lem:4.3.5}{4.3.5}).

Thus \([p'_1,p'_2,\ldots]\) is a subsequence of \([p_1,p_2,\ldots]\). We
have \(\N_{n-1}\subseteq\N_n\), hence
\begin{equation}\tag{4.4.\(\alpha\)}\label{eq:4.4-alpha}
 \N_1\subseteq\N_2\subseteq\cdots.
\end{equation}
\(\N_n\) is a factor \(\subseteq\M\) of class \((\mathrm I_{p'_n})\). Clearly
\begin{equation}\tag{4.4.\(\beta\)}\label{eq:4.4-beta}
 \R(\N_n;n=1,2,\ldots)\subseteq\M.
\end{equation}
Now for \(n\geq h\), there exists a \(B_h^{(n)}\in\N_n\) with
\(\Metric{B_h^{(n)}-A_h}<1/n\). So \(A_h\) is a metric limit of
\(\R(\N_n;n=1,2,\ldots)\) and hence, by \thmref{thm:I}{I},
\(A_h\in\R(\N_n;n=1,2,\ldots)\). Thus
\(A_1,A_2,\ldots\in\R(\N_n;n=1,2,\ldots)\). Since the
\(A_1,A_2,\ldots\) are everywhere dense in \(\M\) in the sense of strong
convergence, this and \eqref{eq:4.4-beta} yield
\begin{equation}\tag{4.4.\(\gamma\)}\label{eq:4.4-gamma}
 \R(\N_n;n=1,2,\ldots)=\M.
\end{equation}
Equations \eqref{eq:4.4-alpha}, \eqref{eq:4.4-gamma}, and our observations
establish that \(\M\) is approximately finite of type
\([p'_1,p'_2,\ldots]\), as stated.
\end{proof}

We need one more auxiliary lemma. In view of a later application, we shall
prove it in a stronger form than that immediately needed.

\begin{lemma}\phantomsection\label{lem:4.4.2}
Let \(\N,\mathbf O\) be two factors in the cases \((\mathrm I_p)\) and
\((\mathrm I_q)\) or \((\mathrm{II}_1)\), respectively. Let \(p\) be a divisor
of \(r\), and (if \(q\) exists) \(r\) a divisor of \(q\). Suppose
\(\N\subseteq\mathbf O\). Then there exists a factor \(\R\) in case
\((\mathrm I_r)\) such that \(\N\subseteq\R\subseteq\mathbf O\).
\end{lemma}

\begin{proof}
By \lemref{lem:2.6.2}{2.6.2} or \lemref{lem:2.6.3}{2.6.3}, since
\(\mathbf O\) is in case \((\mathrm{II}_1)\), or in the case
\((\mathrm I_q)\) with \(q\) divisible by \(r\), there exists a factor
\(\R^+\subseteq\mathbf O\) in the case \((\mathrm I_r)\) with a matrix base
\(W^+_{o,w}\), \(o,w=1,\ldots,r\). Choose also a matrix base
\(W^0_{u,v}\), \(u,v=1,\ldots,p\), of \(\N\). Thus \(W^+_{o,w}\) is a
system of \(r\)-order matrix units with \(\sum_{o=1}^{r}W^+_{o,o}=1\), and
\(W^0_{u,v}\) is a system of \(p\)-order matrix units with
\(\sum_{u=1}^{p}W^0_{u,u}=1\).

Put
\begin{equation}\tag{4.4.\(\delta\)}\label{eq:4.4-delta}
 W^{++}_{u,v}=\sum_{i=1}^{r/p}
 W^+_{(u-1)r/p+i,(v-1)r/p+i}.
\end{equation}
Then one verifies by direct calculation that the \(W^{++}_{u,v}\),
\(u,v=1,\ldots,p\), form a system of \(p\)-order units with
\(\sum_{u=1}^{p}W^{++}_{u,u}=1\). One can readily prove that
\(D_{\mathbf O}(W^{++}_{1,1})=1/p=D_{\mathbf O}(W^0_{1,1})\) in the standard
normalization. Thus there exists a partially isometric \(V\in\mathbf O\) such
that \(V^*V=W^{++}_{1,1}\), \(VV^*=W^0_{1,1}\). Now put
\[
 U=\sum_{u=1}^{p}W^0_{u,1}VW^{++}_{1,u}.
\]
Clearly \(U\in\mathbf O\). One verifies by direct computation that
\(U^*U=UU^*=1\), and hence that \(U\) is unitary. Also
\begin{equation}\tag{4.4.\(\epsilon\)}\label{eq:4.4-epsilon}
 UW^{++}_{u,v}=W^0_{u,v}U,qquad
 W^0_{u,v}=UW^{++}_{u,v}U^{-1}.
\end{equation}
Hence the replacement of \(\R^+\) by \(U\R^+U^{-1}\), and of the system
\(W^+_{o,w}\) by the system \(UW^+_{o,w}U^{-1}\), leaves all their
properties unaffected, but it carries \(W^{++}_{u,v}\) into
\(W^0_{u,v}=UW^{++}_{u,v}U^{-1}\). Consequently, we may assume that we
had \(W^{++}_{u,v}=W^0_{u,v}\), i.e.
\begin{equation}\tag{4.4.\(\zeta\)}\label{eq:4.4-zeta}
 W^0_{u,v}=\sum_{i=1}^{r/p}
 W^+_{(u-1)r/p+i,(v-1)r/p+i},
\end{equation}
to begin with.

Thus all \(W^0_{u,v}\in\R\), hence \(\N\subseteq\R\). We know already that
\(\R\) is a factor in the case \((\mathrm I_r)\) and that
\(\R\subseteq\mathbf O\). Thus the proof is completed.
\end{proof}

We can now conclude our deductions from \defref{def:4.3.1}{4.3.1}.

\begin{lemma}\phantomsection\label{lem:4.4.3}
Under the same assumptions as \lemref{lem:4.4.1}{4.4.1}, \(\M\) is
approximately finite of type \([p_1,p_2,\ldots]\).
\end{lemma}

\begin{proof}
Apply \lemref{lem:4.4.1}{4.4.1}. Denote its \(\N_1,\N_2,\ldots\) by
\(\N_{1'},\N_{2'},\ldots\). Write \(0'=0\),
\(\N_{0'}=\N_0=(\alpha\mathbin{\cdot}1)\).

Thus \(\N_n\) is now defined for \(n=0',1',2',\ldots\) only. Consider a
fixed \(m'=1',2',\ldots\). We have
\(\N_{(m-1)'}\subseteq\N_{m'}\), and \(\N_{(m-1)'},\N_{m'}\) are factors in
the cases \((\mathrm I_{p_{(m-1)'}})\), \((\mathrm I_{p_{m'}})\), respectively.
Each of
\[
 p_{(m-1)'},p_{(m-1)'+1},\ldots,p_{m'-1},p_{m'}
\]
is a divisor of its successor. Hence we obtain, by repeated applications of
\lemref{lem:4.4.2}{4.4.2}, a succession of factors
\(\N_{(m-1)'+1},\N_{(m-1)'+2},\ldots,\N_{m'-1}\), with these properties:
\[
 \N_{(m-1)'}\subseteq\N_{(m-1)'+1}\subseteq\cdots
 \subseteq\N_{m'-1}\subseteq\N_{m'},
\]
and \(\N_n\), \(n=(m-1)'+1,\ldots,m'-1\), is in the case
\((\mathrm I_{p_n})\).

So we have defined \(\N_n\) for all \(n=1,2,\ldots\). We have
\begin{equation}\tag{4.4.\(\eta\)}\label{eq:4.4-eta}
 \N_1\subseteq\N_2\subseteq\cdots,
\end{equation}
and \(\N_n\) is a factor in the case \((\mathrm I_{p_n})\). Finally, by
\eqref{eq:4.4-eta} and \lemref{lem:4.4.1}{4.4.1}, we have
\begin{equation}\tag{4.4.\(\theta\)}\label{eq:4.4-theta}
 \R(\N_n;n=1,2,\ldots)=\R(\N_{m'};m=1,2,\ldots)=\M.
\end{equation}
Equations \eqref{eq:4.4-eta}, \eqref{eq:4.4-theta}, and our other observations
establish that \(\M\) is approximately finite of type
\([p_1,p_2,\ldots]\), as stated.
\end{proof}

\PaperSection{sec:4.5}{Approximate finiteness (B). Structure of finite rings.}
We proceed now to the second operation announced at the beginning of
\secref{sec:4.3}{4.3}.

\begin{definition}\phantomsection\label{def:4.5.1}
A ring \(\N\) is of finite order if it possesses a finite linear basis, i.e. a
(fixed finite) system \(A^0_1,\ldots,A^0_q\) such that \(\N\) is the set of all
\[
 A=\sum_{u=1}^{q}x_uA^0_u
 \qquad(x_1,\ldots,x_q\text{ any complex numbers}).
\]
\end{definition}

\begin{definition}\phantomsection\label{def:4.5.2}
\(\M\) is approximately finite (B) if this is true: Given any
\(A_1,\ldots,A_m\in\M\) and any \(\epsilon>0\), there exists a ring \(\N\)
with the following properties:
\begin{enumerate}[label=(\roman*)]
\item \(\N\) is of finite order;
\item \(\N\subseteq\M\);
\item There exist \(B_1,\ldots,B_m\in\N\) with
\[
 \Metric{B_h-A_h}<\epsilon\qquad\text{for }h=1,\ldots,m.
\]
\end{enumerate}
\end{definition}

\begin{lemma}\phantomsection\label{lem:4.5.1}
If \(\M\) is approximately finite of type \([p_1,p_2,\ldots]\), then it is
also approximately finite (B).
\end{lemma}

\begin{proof}
Suppose \(A_1,\ldots,A_m\in\M\) and an \(\epsilon>0\) are given. Form the
\(\N_1,\N_2,\ldots\) of \defref{def:4.1.1}{4.1.1}. By (ii) in
\lemref{lem:4.1.4}{4.1.4}, there exists for each \(h=1,\ldots,m\) an
\(n_h\) such that, for some \(B_h\in\N_{n_h}\),
\begin{equation}\tag{4.5.\(\alpha\)}\label{eq:4.5-alpha}
 \Metric{B_h-A_h}<\epsilon.
\end{equation}
Put \(n=\operatorname{Max}(n_1,\ldots,n_m)\); then
\(B_h\in\N_{n_h}\subseteq\N_n\), i.e.
\begin{equation}\tag{4.5.\(\beta\)}\label{eq:4.5-beta}
 B_1,\ldots,B_m\in\N_n.
\end{equation}
Now \(\N_n\) is a factor in the case \((\mathrm I_{p_n})\). By (ii) in
\lemref{lem:4.1.1}{4.1.1} it is a ring of finite order. This, together with
\eqref{eq:4.5-alpha} and \eqref{eq:4.5-beta}, completes the proof.
\end{proof}

Comparison of \lemref{lem:4.4.3}{4.4.3} and \lemref{lem:4.5.1}{4.5.1}
shows that if we can derive approximate finiteness (A) from (B), then the
equivalence of all kinds of approximate finiteness is established. This we
accomplish in the next section. In the present section we obtain certain
preliminary lemmas on rings of finite order; indeed we shall determine the
structure of these.\footnote{The four lemmas which follow are merely an
\emph{ad hoc} derivation of Wedderburn's structure theorems for semi-simple
algebras---specialized to the present situation. They can also be interpreted as
a summary of certain results of unitary representation theory including the
theorem of Burnside--Frobenius--Schur for that case.} Thus for the present
section \(\N\) is to be a ring of finite order.

\begin{lemma}\phantomsection\label{lem:4.5.2}
\begin{enumerate}[label=(\roman*)]
\item There is a maximum \(r\) such that there is a system
\(E_1,\ldots,E_r\) of mutually orthogonal projections \(\ne0\) and in \(\N\).
\item For such a system (with \(r\) maximum),
\(E=\sum_{s=1}^{r}E_s\) is the unit of \(\N\).
\item Every \(E_s\) in such a system (with \(r\) maximum) is minimal (with
respect to \(\N\)) in the sense of \cite{ref5}, p.~143,
Definition 5.1.2.
\end{enumerate}
\end{lemma}

\begin{proof}
Ad (i): Any such system is necessarily linearly independent. Hence the number
of its elements must be \(\leq\) the \(q\) of \defref{def:4.5.1}{4.5.1}.
Hence the numbers for which there is such a system have a maximum.

Ad (ii): Let \(E\) be the unit of \(\N\). \(\sum_{s=1}^{r}E_s\) is a
projection in \(\N\); hence \(\sum_{s=1}^{r}E_s\leq E\). So
\(E_{r+1}=E-\sum_{s=1}^{r}E_s\) is a projection in \(\N\) which is
orthogonal to the original \(E_s\). Hence the assumed maximum property of
\(r\) necessitates \(E_{r+1}=0\), i.e. \(E=\sum_{s=1}^{r}E_s\).

Ad (iii): Consider a projection \(F\leq E_s\) in \(\N\). Then
\(E'_s=F\) and \(E''_s=E_s-F\) are mutually orthogonal projections in
\(\N\) which are also orthogonal to all \(E_{s'}\), \(s'\ne s\). Hence the
assumed maximum property of \(r\) necessitates \(E'_s=0\) or
\(E''_s=0\), i.e. \(F=0\) or \(E_s\). This, however, is the defining
property of minimality.
\end{proof}

\noindent\textsc{Remark.} The center \(\N\mathbin{\cdot}\N'\) of \(\N\)
is of finite order along with \(\N\). Hence we can apply the above lemma to it,
thus obtaining a system \(\bar E_1,\ldots,\bar E_{\bar r}\).

By this lemma, \(\bar E=\sum_{s=1}^{\bar r}\bar E_s\) is the unit of
\(\N\mathbin{\cdot}\N'\); hence by \cite{ref2}, p.~393,
Theorems 3 and 4, it is also the unit of \(\N\).

\begin{lemma}\phantomsection\label{lem:4.5.3}
\begin{enumerate}[label=(\roman*)]
\item \(\bar E_s\) (cf. above) is such that its range \(\bar{\mathfrak M}_s\)
reduces \(\N,\N',\N\mathbin{\cdot}\N'\).
\item
\[
 (\N\mathbin{\cdot}\N')_{(\bar{\mathfrak M}_s)}
   =(\bar E_s)_{(\bar{\mathfrak M}_s)}.
\]
\end{enumerate}
\end{lemma}

\begin{proof}
Ad (i): This is implied by \(\bar E_s\in\N\mathbin{\cdot}\N'\).

Ad (ii): This is equivalent to the statement
\(\bar E_sA=A\bar E_s=\alpha\bar E_s\) for all
\(A\in\N\mathbin{\cdot}\N'\). Put
\(\bar E_sA=A\bar E_s=A'\). Clearly
\(A'\in\N\mathbin{\cdot}\N'\),
\(\bar E_sA'=A'\bar E_s=A'\). Since \(\bar E_s\) is a minimal projection
in \(\N\mathbin{\cdot}\N'\), this implies \(A'=\alpha\bar E_s\) (cf. the
last part of the proof of Lemma 5.1.3 in \cite{ref5}, p.~144).
\end{proof}

\begin{lemma}\phantomsection\label{lem:4.5.4}
\(\N_{(\bar{\mathfrak M}_s)}\) is a factor (in \(\bar{\mathfrak M}_s\)) in a
case \((\mathrm I_{q_s})\), \(q_s=1,2,\ldots\).
\end{lemma}

\begin{proof}
If \(A'\) is in the center of \(\N_{(\bar{\mathfrak M}_s)}\),
\(A'\bar E_s\) defines an operator \(A\) in \(\Hilbert\) which is readily seen
to be in the center of \(\N\), and for which
\(A\bar E_s=\bar E_sA=A\). Thus the center of
\(\N_{(\bar{\mathfrak M}_s)}\) is included in
\((\N\mathbin{\cdot}\N')_{(\bar{\mathfrak M}_s)}\). The converse inclusion
is obvious.

It follows that the center of \(\N_{(\bar{\mathfrak M}_s)}\) is
\((\N\mathbin{\cdot}\N')_{(\bar{\mathfrak M}_s)}\), and hence
\(\N_{(\bar{\mathfrak M}_s)}\) is a factor by (ii) in
\lemref{lem:4.5.3}{4.5.3}.

\(\N_{(\bar{\mathfrak M}_s)}\) is of finite order along with \(\N\).
Therefore \lemref{lem:4.5.2}{4.5.2} implies that
\(\N_{(\bar{\mathfrak M}_s)}\) contains a minimal projection. Thus
\(\N_{(\bar{\mathfrak M}_s)}\) is a direct factor (cf.
\cite{ref5}, pp.~150, 173, Theorem IV and Lemma 8.6.1).
Obviously, since it is of finite order, \(\N_{(\bar{\mathfrak M}_s)}\) is not
in a case \((\mathrm I_\infty)\), and hence it is in a case
\((\mathrm I_{q_s})\), \(q_s=1,2,\ldots\). This completes the proof.
\end{proof}

\begin{lemma}\phantomsection\label{lem:4.5.5}
\(\N\) possesses a finite linear base which consists of the operators
\(W^s_{u,v}\), \(s=1,\ldots,\bar r\), \(u,v=1,\ldots,q_s\), where for any
fixed \(s\) the \(W^s_{u,v}\), \(u,v=1,\ldots,q_s\), are a system of
\(q_s\)-order matrix units with
\(\bar E_s=\sum_{u=1}^{q_s}W^s_{u,u}\). Thus for the \(q\) of
\defref{def:4.5.1}{4.5.1},
\[
 q=\sum_{s=1}^{\bar r}(q_s)^2.
\]
\end{lemma}

\begin{proof}
In virtue of \lemref{lem:4.5.4}{4.5.4}, we may apply
\lemref{lem:4.1.1}{4.1.1} and Remark 2 after it to
\(\N_{(\bar{\mathfrak M}_s)}\). So \(\N_{(\bar{\mathfrak M}_s)}\)
possesses a finite linear base of operators \(\widetilde W^s_{u,v}\) (in
\(\bar{\mathfrak M}_s\)), \(u,v=1,\ldots,q_s\), which form a system of
\(q_s\)-order matrix units in \(\bar{\mathfrak M}_s\), with
\((\bar E_s)_{(\bar{\mathfrak M}_s)}=
\sum_{u=1}^{q_s}\widetilde W^s_{u,u}\).

Since \(\widetilde W^s_{u,v}\in\N_{(\bar{\mathfrak M}_s)}\), there is a
\(W^{s0}_{u,v}\) in \(\N\) for which the part in
\(\bar{\mathfrak M}_s\) is \(\widetilde W^s_{u,v}\). Let
\(W^s_{u,v}=W^{s0}_{u,v}\bar E_s\). Since
\(\bar E_s\in\N\mathbin{\cdot}\N'\), we have \(W^s_{u,v}\in\N\), and also
\(W^s_{u,v}=W^{s0}_{u,v}\bar E_s=\bar E_sW^{s0}_{u,v}\). Clearly
\((W^s_{u,v})_{(\bar{\mathfrak M}_s)}=\widetilde W^s_{u,v}\), and one can
readily show that the \(W^s_{u,v}\) form a system of \(q_s\)-order matrix
units in \(\Hilbert\) with \(\bar E_s=\sum_{u=1}^{q_s}W^s_{u,u}\).

Consider an \(A_s\in\N\) with \(\bar E_sA_s=A_s\bar E_s=A_s\). Then
\(A_s\) has the part \((A_s)_{(\bar{\mathfrak M}_s)}\) in
\(\bar{\mathfrak M}_s\) and 0 in \(\bar{\mathfrak M}_s^\perp\). In the
sense, then, that \(\bar E_s\) is a transformation from \(\Hilbert\) to
\(\bar{\mathfrak M}_s\),
\(A_s=(A_s)_{(\bar{\mathfrak M}_s)}\bar E_s\), and also
\(W^s_{u,v}=\widetilde W^s_{u,v}\bar E_s\). As
\((A_s)_{(\bar{\mathfrak M}_s)}\in\N_{(\bar{\mathfrak M}_s)}\), we have,
from the first paragraph of the proof,
\[
 (A_s)_{(\bar{\mathfrak M}_s)}
   =\sum_{u=1}^{q_s}\sum_{v=1}^{q_s}a^s_{u,v}\widetilde W^s_{u,v},
\]
and hence
\[
 A_s=(A_s)_{(\bar{\mathfrak M}_s)}\bar E_s
   =\sum_{u=1}^{q_s}\sum_{v=1}^{q_s}a^s_{u,v}
      \widetilde W^s_{u,v}\bar E_s
   =\sum_{u=1}^{q_s}\sum_{v=1}^{q_s}a^s_{u,v}W^s_{u,v}.
\]

Consider now any \(A\in\N\). Put \(A_s=\bar E_sA\), which also equals
\(A\bar E_s\), since \(\bar E_s\in\N\mathbin{\cdot}\N'\). By the above,
\(A_s=\sum_{u=1}^{q_s}\sum_{v=1}^{q_s}a^s_{u,v}W^s_{u,v}\). By the
remark after \lemref{lem:4.5.2}{4.5.2}, we have
\[
 A=\bar EA=\sum_{s=1}^{\bar r}\bar E_sA
  =\sum_{s=1}^{\bar r}\sum_{u=1}^{q_s}\sum_{v=1}^{q_s}
      a^s_{u,v}W^s_{u,v}.
\]
Thus the \(W^s_{u,v}\), \(s=1,\ldots,\bar r\),
\(u,v=1,\ldots,q_s\), form indeed a finite linear basis of \(\N\).
\end{proof}

We now prove a lemma which possesses a certain interest of its own. It applies
to any ring of finite order \(\N\subseteq\M\). We assume that a representation
of \(\N\) in the sense of \lemref{lem:4.5.5}{4.5.5} is given, with the
\(W^s_{u,v}\) as described there.

\begin{lemma}\phantomsection\label{lem:4.5.6}
Given a (fixed) \(q=1,2,\ldots\), there exists a factor \(\mathbf O\) in case
\((\mathrm I_q)\), with \(\N\subseteq\mathbf O\subseteq\M\), if and only if
this is true for \(\N\):
\begin{equation}\tag{4.5.\(\gamma\)}\label{eq:4.5-gamma}
 \text{All }qD_{\M}(W^s_{1,1})\text{ are integers.}
\end{equation}
\end{lemma}

\begin{proof}
Assume that an \(\mathbf O\) with the specified properties exists. The
considerations at the beginning of the proof of (i) in
\lemref{lem:4.1.2}{4.1.2} show again that
\(D_{\M}(G)=D_{\mathbf O}(G)\) for all projections
\(G\in\mathbf O\;(\subseteq\M)\). Hence, in particular,
\(D_{\M}(W^s_{1,1})=D_{\mathbf O}(W^s_{1,1})\), since \(W^s_{1,1}\) is a
projection in \(\N\subseteq\mathbf O\). But as \(\mathbf O\) is in case
\((\mathrm I_q)\), \(D_{\M}(W^s_{1,1})=D_{\mathbf O}(W^s_{1,1})\) must have
a value \(0,1/q,\ldots,1\); i.e. \(qD_{\M}(W^s_{1,1})\) must be an integer.
Thus \eqref{eq:4.5-gamma} holds.

Sufficiency: Assume that \(\N\) fulfills \eqref{eq:4.5-gamma}. Considering
\lemref{lem:4.5.5}{4.5.5}, we must prove this:
\begin{equation}\tag{4.5.\(\delta\)}\label{eq:4.5-delta}
\begin{minipage}{0.80\linewidth}
Let a system of mutually orthogonal projections \(\ne0\) from \(\M\) be given:
\(\bar E_1,\ldots,\bar E_{\bar r}\). For each \(s\), let a system of
\(q_s\)-order matrix units \(W^s_{u,v}\) from \(\M\) be given, with
\(\bar E_s=\sum_{u=1}^{q_s}W^s_{u,u}\) and \(q_s=1,2,\ldots\). Let, for a
given \(q\), every \(qD_{\M}(W^s_{1,1})\) be an integer. Then there exists a
factor \(\mathbf O\subseteq\M\) in a case \((\mathrm I_q)\) such that all
\(W^s_{u,v}\in\mathbf O\).
\end{minipage}
\end{equation}
We therefore disregard \(\N\) completely and deal with
\eqref{eq:4.5-delta} alone.

\textit{Proof of \eqref{eq:4.5-delta}.} We have
\(D_{\M}(W^s_{1,1})=p_s/q\), \(p_s=1,2,\ldots\). The projections
\(W^s_{1,1}\) and \(W^s_{u,u}\) are equivalent in virtue of the partially
isometric \(W^s_{u,1}\) (cf. \lemref{lem:2.6.1}{2.6.1}). Hence
\begin{equation}\tag{4.5.\(\epsilon\)}\label{eq:4.5-epsilon}
 D_{\M}(W^s_{u,u})=p_s/q,qquad
 D_{\M}(\bar E_s)=p_sq_s/q.
\end{equation}
We may assume that
\begin{equation}\tag{4.5.\(\zeta\)}\label{eq:4.5-zeta}
 \sum_{s=1}^{\bar r}\bar E_s=1.
\end{equation}
Otherwise we could put \(\bar E_{\bar r+1}=1-\sum_{s=1}^{\bar r}\bar E_s\).
Then \(\bar E_{\bar r+1}\) is a projection \(\ne0\). Also
\(qD_{\M}(\bar E_{\bar r+1})=q-
\sum_{s=1}^{\bar r}qD_{\M}(\bar E_s)\) is an integer. Thus we could add
\(\bar E_{\bar r+1}\) to the system \(\bar E_1,\ldots,\bar E_{\bar r}\),
with \(q_{\bar r+1}=1\) and
\(W^{\bar r+1}_{1,1}=\bar E_{\bar r+1}\).

By \eqref{eq:4.5-epsilon} and \eqref{eq:4.5-zeta},
\(\sum_{s=1}^{\bar r}p_sq_s=q\). We put
\(l_s=\sum_{t=1}^{s}p_tq_t\), for \(s=0,1,\ldots,\bar r\). Then
\begin{equation}\tag{4.5.\(\eta\)}\label{eq:4.5-eta}
 \begin{cases}
 l_s=l_{s-1}+p_sq_s,&s=1,\ldots,\bar r,\\
 l_0=0,&l_{\bar r}=q.
 \end{cases}
\end{equation}
As \(\M\) is in case \((\mathrm{II}_1)\), by \lemref{lem:2.6.2}{2.6.2}
there exists a factor \(\mathbf O^+\subseteq\M\) in the case
\((\mathrm I_q)\), with its matrix base \(W^+_{o,w}\),
\(o,w=1,\ldots,q\). Thus \(W^+_{o,w}\) is a system of \(q\)-order matrix
units, with \(\sum_{u=1}^{q}W^+_{u,u}=1\).

Put
\begin{equation}\tag{4.5.\(\theta\)}\label{eq:4.5-theta}
 E^{+s}=\sum_{u=l_{s-1}+1}^{l_s}W^+_{u,u}
 \qquad(s=1,\ldots,\bar r),
\end{equation}
\begin{equation}\tag{4.5.\(\iota\)}\label{eq:4.5-iota}
 W^{+s}_{u,v}=\sum_{i=1}^{p_s}
 W^+_{l_{s-1}+(u-1)p_s+i,\,l_{s-1}+(v-1)p_s+i}
 \quad\begin{matrix}s=1,\ldots,\bar r,\\u,v=1,\ldots,q_s.\end{matrix}
\end{equation}
Then one verifies by a direct computation that the \(E^{+s}\) are mutually
orthogonal projections \(\ne0\), with \(\sum_{s=1}^{\bar r}E^{+s}=1\), and
(using \defref{def:2.6.1}{2.6.1}) that the \(W^{+s}_{u,v}\) form, for any
fixed \(s\), a system of \(q_s\)-order matrix units with
\(\sum_{u=1}^{q_s}W^{+s}_{u,u}=E^{+s}\). All these operators are in \(\M\).

The \(W^+_{u,u}\) are equivalent (due to the \(W^+_{u,v}\)), and since
\(\sum_{u=1}^{q}W^+_{u,u}=1\), \(D_{\M}(W^+_{u,u})=1/q\). Hence by
\eqref{eq:4.5-iota}, \(D_{\M}(W^{+s}_{1,1})=p_s/q\), and by
\eqref{eq:4.5-epsilon},
\(D_{\M}(W^{+s}_{1,1})=D_{\M}(W^s_{1,1})\). Consequently there exists a
partially isometric operator \(V_s\in\M\), with initial projection
\(W^{+s}_{1,1}\) and final projection \(W^s_{1,1}\).

Now put
\[
 U=\sum_{s=1}^{\bar r}\sum_{u=1}^{q_s}W^s_{u,1}V_sW^{+s}_{1,u}.
\]
Clearly \(U\in\M\). One verifies, by direct computation,
\(U^*U=UU^*=1\), which implies that \(U\) is unitary, and similarly that
\begin{equation}\tag{4.5.\(\kappa\)}\label{eq:4.5-kappa}
 UW^{+s}_{u,v}=W^s_{u,v}U,qquad
 W^s_{u,v}=UW^{+s}_{u,v}U^{-1}.
\end{equation}
\(\mathbf O=U\mathbf O^+U^{-1}\) is a factor \(\subseteq\M\) and in case
\((\mathrm I_q)\), along with \(\mathbf O^+\). By \eqref{eq:4.5-iota},
\(W^{+s}_{u,v}\in\mathbf O^+\), and hence by \eqref{eq:4.5-kappa},
\(W^s_{u,v}=UW^{+s}_{u,v}U^{-1}\in U\mathbf O^+U^{-1}=\mathbf O\). Thus
all \(W^s_{u,v}\in\mathbf O\).
\end{proof}

\PaperSection{sec:4.6}{Approximate finiteness (C). Equivalence results.}
We are now in a position to complete our proof of the relationships between the
various kinds of approximate finiteness.

\begin{lemma}\phantomsection\label{lem:4.6.1}
Suppose \(\N,\M,\bar E_1,\ldots,\bar E_{\bar r}\), and the
\(W^s_{u,v}\) are as in \lemref{lem:4.5.5}{4.5.5}, with
\(\bar E=\sum_{s=1}^{\bar r}\bar E_s=1\). For any \(\epsilon>0\) there
exists an \(n_5=n_5(\N,\epsilon)\) such that for every \(q\geq n_5\) there
exists a projection \(E=E(q,\N,\epsilon)\) which is in
\(\M\mathbin{\cdot}\N'\) and has the following properties:
\begin{enumerate}[label=(\roman*)]
\item \(D_{\M}(1-E)<\epsilon\);
\item all \(qD_{\M}(EW^s_{1,1})\) are \(\ne0\) and integers.
\end{enumerate}
\end{lemma}

\begin{proof}
Given a \(q=1,2,\ldots\), we wish to choose for each
\(s=1,\ldots,\bar r\) a \(k_s=1,2,\ldots\) with
\begin{equation}\tag{4.6.\(\alpha\)}\label{eq:4.6-alpha}
 q\left(D_{\M}(W^s_{1,1})-\frac{\epsilon}{\sum_{s=1}^{\bar r}q_s}\right)
 <k_s\leq qD_{\M}(W^s_{1,1}).
\end{equation}
This is certainly feasible if \(qD_{\M}(W^s_{1,1})\geq1\) and
\(q\epsilon/\sum_{s=1}^{\bar r}q_s\geq1\), for
\(s=1,\ldots,\bar r\). We assume therefore that \(q\geq n_5\), with
\begin{equation}\tag{4.6.\(\beta\)}\label{eq:4.6-beta}
 n_5=n_5(\N,\epsilon)=
 \operatorname{Max}\left(\frac{\sum_{s=1}^{\bar r}q_s}{\epsilon},
 \frac{1}{D_{\M}(W^s_{1,1})},\ s=1,\ldots,\bar r\right).
\end{equation}
Now \eqref{eq:4.6-alpha} and \(k_s\ne0\) give
\[
 D_{\M}(W^s_{1,1})-\frac{\epsilon}{\sum_{s=1}^{\bar r}q_s}
 <\frac{k_s}{q}\leq D_{\M}(W^s_{1,1}),\qquad \frac{k_s}{q}\ne0.
\]
As \(\M\) is in case \((\mathrm{II}_1)\), this implies that we can find a
projection \(G^s\in\M\), with
\begin{equation}\tag{4.6.\(\gamma\)}\label{eq:4.6-gamma}
 D_{\M}(G^s)=k_s/q\ne0,
\end{equation}
and
\begin{equation}\tag{4.6.\(\delta\)}\label{eq:4.6-delta}
 G^s\leq W^s_{1,1},
\end{equation}
and that then
\begin{equation}\tag{4.6.\(\epsilon\)}\label{eq:4.6-epsilon}
 D_{\M}(W^s_{1,1}-G^s)<\frac{\epsilon}{\sum_{s=1}^{\bar r}q_s}.
\end{equation}

\(W^s_{u,1}\) is a partially isometric operator with initial projection
\(W^s_{1,1}\) and final projection \(W^s_{u,u}\). Its mapping carries
therefore the \(G^s\leq W^s_{1,1}\) into a projection \(G^s_u\in\M\), with
\begin{equation}\tag{4.6.\(\zeta\)}\label{eq:4.6-zeta}
 G^s_u\leq W^s_{u,u}.
\end{equation}
The same mapping carries \eqref{eq:4.6-epsilon} into
\begin{equation}\tag{4.6.\(\eta\)}\label{eq:4.6-eta}
 D_{\M}(W^s_{u,u}-G^s_u)<\frac{\epsilon}{\sum_{s=1}^{\bar r}q_s}.
\end{equation}
Finally, \(W^s_{u,v}W^s_{v,1}=W^s_{u,1}\) implies that the mapping of
\(W^s_{u,v}\) carries \(G^s_v\) into \(G^s_u\). This may be written
\begin{equation}\tag{4.6.\(\theta\)}\label{eq:4.6-theta}
 W^s_{u,v}G^s_v=G^s_uW^s_{u,v}.
\end{equation}
For \(s\ne t\) or \(v\ne w\), the initial projection \(W^s_{v,v}\) of
\(W^s_{u,v}\) is orthogonal to \(W^t_{w,w}\), hence, a fortiori, to
\(G^t_w\leq W^t_{w,w}\). So we have
\begin{align}
 W^s_{u,v}G^t_w&=0 &&\text{for }s\ne t\text{ or }v\ne w,
   \tag{4.6.\(\iota\)}\label{eq:4.6-iota}\\
 G^t_wW^s_{u,v}&=0 &&\text{for }s\ne t\text{ or }u\ne w.
   \tag{4.6.\(\kappa\)}\label{eq:4.6-kappa}
\end{align}
Now put \(E=\sum_{s=1}^{\bar r}\sum_{u=1}^{q_s}G^s_u\). Clearly
\(E\in\M\), and since
\(\sum_{s=1}^{\bar r}\sum_{u=1}^{q_s}W^s_{u,u}=
\sum_{s=1}^{\bar r}\bar E_s=1\), we have
\(1-E=\sum_{s=1}^{\bar r}\sum_{u=1}^{q_s}(W^s_{u,u}-G^s_u)\). Hence
\eqref{eq:4.6-eta} yields
\begin{equation}\tag{4.6.\(\lambda\)}\label{eq:4.6-lambda}
 D_{\M}(1-E)<\epsilon.
\end{equation}
Equations \eqref{eq:4.6-theta}, \eqref{eq:4.6-iota}, and
\eqref{eq:4.6-kappa} give
\[
 W^s_{u,v}E=EW^s_{u,v}
  \;(=W^s_{u,v}G^s_v=G^s_uW^s_{u,v}).
\]
So \(E\) commutes with \(W^s_{u,v}\), and hence, by
\lemref{lem:4.5.5}{4.5.5}, with all of \(\N\). Since \(E\in\M\), this yields
\begin{equation}\tag{4.6.\(\mu\)}\label{eq:4.6-mu}
 E\in\M\mathbin{\cdot}\N'.
\end{equation}
Finally, direct computation shows that
\begin{equation}\tag{4.6.\(\nu\)}\label{eq:4.6-nu}
 EW^s_{1,1}=G^s_1=G^s.
\end{equation}
Our assertions now follow immediately. Equation \eqref{eq:4.6-mu} gives the
preliminary one, \eqref{eq:4.6-lambda} implies (i), and
\eqref{eq:4.6-nu} and \eqref{eq:4.6-gamma} yield (ii). Thus the lemma is
proven.
\end{proof}

It is now possible to take the final step.

\begin{lemma}\phantomsection\label{lem:4.6.2}
If \(\M\) is approximately finite (B), then it is also approximately finite (A).
\end{lemma}

\begin{proof}
We must verify the criteria of \defref{def:4.3.1}{4.3.1} for \(\M\).

Let therefore a system \(A_1,\ldots,A_m\in\M\) and an \(\epsilon>0\) be
given. Apply \defref{def:4.5.2}{4.5.2} to these \(A_1,\ldots,A_m\) and to
\(\epsilon/2\), thus obtaining the ring \(\N\) and
\(B_1,\ldots,B_m\in\N\) described there. Thus
\begin{equation}\tag{4.6.\(\mathrm{o}\)}\label{eq:4.6-omicron}
 \Metric{B_s-A_s}<\epsilon/2.
\end{equation}

\(\N\) is a ring of finite order \(\subseteq\M\). We form
\(\bar E_1,\ldots,\bar E_{\bar r}\) and the \(W^s_{u,v}\) in the sense of
\lemref{lem:4.5.5}{4.5.5} for this \(\N\). To these and to
\begin{equation}\tag{4.6.\(\pi\)}\label{eq:4.6-pi}
 \epsilon'=\left(\frac{\epsilon}{2\operatorname{Max}
 (\OpNorm{B_s};\ s=1,\ldots,m)}\right)^2
\end{equation}
we apply \lemref{lem:4.6.1}{4.6.1}, thus obtaining
\(n_5=n_5(\N,\epsilon')\). For any \(q\geq n_5\), we form the projection
\(E=E(q,\N,\epsilon')\) of \lemref{lem:4.6.1}{4.6.1}.

Accordingly we have \(E\in\M\mathbin{\cdot}\N'\), and
\(\N_E=E\mathbin{\cdot}\N=\N E\) is a finite-order ring \(\subseteq\M\)
along with \(\N\). Thus \(E\bar E_1,\ldots,E\bar E_{\bar r}\) and the
\(EW^s_{u,v}\) perform for \(\N_E\) the functions of
\(\bar E_1,\ldots,\bar E_{\bar r}\) and the \(W^s_{u,v}\) for \(\N\)
(cf. \lemref{lem:4.5.5}{4.5.5}). By (ii) in \lemref{lem:4.6.1}{4.6.1}, all
\(qD_{\M}(EW^s_{1,1})\) are integers. Thus \lemref{lem:4.5.6}{4.5.6} is
applicable to \(\N_E\), and we obtain a factor \(\mathbf O\) in case
\((\mathrm I_q)\), with
\begin{equation}\tag{4.6.\(\rho\)}\label{eq:4.6-rho}
 \N_E\subseteq\mathbf O\subseteq\M.
\end{equation}
\(B_s\in\N\) implies
\begin{equation}\tag{4.6.\(\sigma\)}\label{eq:4.6-sigma}
 C_s=EB_s\in\N_E\subseteq\mathbf O.
\end{equation}
Now (i) in \lemref{lem:4.6.1}{4.6.1} gives
\(D_{\M}(1-E)<\epsilon'\). Hence
\(\Metric{1-E}<\epsilon'^{1/2}\), and therefore
\[
 \Metric{B_s-C_s}=\Metric{B_s(1-E)}
 <\OpNorm{B_s}\epsilon'^{1/2}.
\]
Hence by \eqref{eq:4.6-pi},
\begin{equation}\tag{4.6.\(\tau\)}\label{eq:4.6-tau}
 \Metric{B_s-C_s}<\epsilon/2.
\end{equation}
Combining \eqref{eq:4.6-omicron} and \eqref{eq:4.6-tau} gives
\begin{equation}\tag{4.6.\(\upsilon\)}\label{eq:4.6-upsilon}
 \Metric{C_s-A_s}<\epsilon.
\end{equation}
Thus we can satisfy \defref{def:4.3.1}{4.3.1} by putting
\[
 n_1=n_1(A_1,\ldots,A_m;\epsilon)=n_5(\N,\epsilon'),
\]
and then, for \(q\geq n_1=n_5\),
\(\N_1=\N_1(q,A_1,\ldots,A_m,\epsilon)=\mathbf O\).\footnote{Observe
that the (indirectly defined) expression at the extreme right has the proper
dependence---i.e. that one indicated by the expression in the middle.} Then our
\(C_1,\ldots,C_m\) replace the \(B_1,\ldots,B_m\) of
\defref{def:4.3.1}{4.3.1}, and (i)--(iii) there obtain as follows: (i) holds by
virtue of the construction of \(\mathbf O\), (ii) follows from
\eqref{eq:4.6-rho}, and (iii) coincides with \eqref{eq:4.6-upsilon}. Thus the
proof is completed.
\end{proof}

We could now make the final steps as indicated after \lemref{lem:4.5.1}{4.5.1}
above, but we introduce one more variant of the concept of approximate
finiteness.

\begin{definition}\phantomsection\label{def:4.6.1}
\(\M\) is approximately finite (C) if there exists a sequence of rings
\(\N_1,\N_2,\ldots\) with the following properties:
\begin{enumerate}[label=(\roman*)]
\item \(\N_n\) is of finite order;
\item \(\N_1\subseteq\N_2\subseteq\cdots\);
\item the ring \(\M\) is generated by the \(\N_1,\N_2,\ldots\):
\[
 \M=\R(\N_n;n=1,2,\ldots).
\]
\end{enumerate}
\end{definition}

\begin{lemma}\phantomsection\label{lem:4.6.3}
If \(\M\) is approximately finite of type \([p_1,p_2,\ldots]\), then it is
also approximately finite (C).
\end{lemma}

\begin{proof}
Immediate by comparison of Definitions \defref{def:4.1.1}{4.1.1} and
\defref{def:4.6.1}{4.6.1}.
\end{proof}

\begin{lemma}\phantomsection\label{lem:4.6.4}
If \(\M\) is approximately finite (C), then it is also approximately finite (B).
\end{lemma}

\begin{proof}
(ii) and (iii) in \defref{def:4.6.1}{4.6.1} coincide with those in
\defref{def:4.1.1}{4.1.1}. Hence \lemref{lem:4.1.4}{4.1.4} applies to our
\(\N_1,\N_2,\ldots\).

Let then the \(A_1,\ldots,A_m\in\M\) and the \(\epsilon>0\) of
\defref{def:4.5.2}{4.5.2} be given. Owing to \lemref{lem:4.1.4}{4.1.4},
\(A_k\) is the limit of a metrically convergent sequence from the
set-theoretical sum of the \(\N_1,\N_2,\ldots\). Hence there exists a
\(B_k\in\N_{n_k}\) with \(\Metric{B_k-A_k}<\epsilon\).

Now put \(n=\operatorname{Max}(n_1,\ldots,n_m)\). Then every
\(B_k\in\N_{n_k}\subseteq\N_n\). Thus \(\N=\N_n\) meets all the
requirements (i)--(iii) of \defref{def:4.5.2}{4.5.2}.
\end{proof}

We can now prove:

\begin{theorem}\phantomsection\label{thm:XII}
All kinds of approximate finiteness are equivalent to each other: (A), (B),
(C), and every type \([p_1,p_2,\ldots]\) (for every sequence
\([p_1,p_2,\ldots]\) which fulfills (i), (ii) of
\lemref{lem:4.1.3}{4.1.3}).

We shall therefore use the expression approximately finite from now on,
without any further qualification.
\end{theorem}

\begin{proof}
The implications
\[
 (A)\longrightarrow\text{type }[p_1,p_2,\ldots]
 \longrightarrow(B)\longrightarrow(A)
\]
were established in \lemref{lem:4.4.3}{4.4.3},
\lemref{lem:4.5.1}{4.5.1}, and \lemref{lem:4.6.2}{4.6.2}, respectively.
Consequently
\[
 \text{type }[p_1,p_2,\ldots]\Longleftrightarrow(A)\Longleftrightarrow(B).
\]
Further,
\[
 \text{type }[p_1,p_2,\ldots]\longrightarrow(C)\longrightarrow(B)
\]
by \lemref{lem:4.6.3}{4.6.3}, \lemref{lem:4.6.4}{4.6.4}, respectively;
hence
\[
 \text{type }[p_1,p_2,\ldots]\Longleftrightarrow(C)
\]
too.
\end{proof}

\PaperSection{sec:4.7}{Existence of approximately finite factors.}
We have not shown so far that approximately finite factors exist at all. In
\secref{sec:5.2}{\S\S5.2}, \secref{sec:5.3}{5.3} and
\secref{sec:5.6}{5.6} we shall give various examples of such factors. Actually
it is more difficult to show that there are factors in case (II\(_1\)) which
are not approximately finite.

Our immediate goal is to prove a general theorem which yields the existence
of approximately finite factors as a byproduct.

A preparatory lemma is necessary:

\begin{lemma}\phantomsection\label{lem:4.7.1}
Let a factor \(\N\subseteq\M\) in a case (I\(_p\)),
\(p=1,2,\ldots\), be given. If an \(A\in\N\) possesses these two properties
\begin{enumerate}[label=(\roman*)]
\item \(\Metric{AX-XA}\leq\epsilon\OpNorm{X}\) for every \(X\in\N\),
\item \(\operatorname{Tr}_{\M}(A)=0\),
\end{enumerate}
then \(\Metric{A}\leq\epsilon\).
\end{lemma}

\begin{proof}
For any unitary \(U\in\N\) we have
\(\OpNorm{U}=\OpNorm{U^{-1}}=1\). Hence by (i),
\[
 \Metric{UAU^{-1}-A}
 =\Metric{U(AU^{-1}-U^{-1}A)}
 \leq\Metric{AU^{-1}-U^{-1}A}\leq\epsilon.
\]
Also by (ii)
\[
 \operatorname{Tr}_{\M}(UAU^{-1})=\operatorname{Tr}_{\M}(A)=0.
\]
If \(U_1,\ldots,U_m\) are unitary and \(\in\N\), then we have by the above
\(\Metric{U_iAU_i^{-1}-A}\leq\epsilon\),
\(\operatorname{Tr}_{\M}(U_iAU_i^{-1})=0\). So for
\begin{equation}\tag{4.7.\(\alpha\)}\label{eq:4.7-alpha}
 B=\frac1m\left(\sum_{i=1}^{m}U_iAU_i^{-1}\right)
\end{equation}
the evaluations
\begin{align}
 \Metric{B-A}&\leq\epsilon,
   \tag{4.7.\(\beta\)}\label{eq:4.7-beta}\\
 \operatorname{Tr}_{\M}(B)&=0
   \tag{4.7.\(\gamma\)}\label{eq:4.7-gamma}
\end{align}
obtain.

Now choose an isomorphism \(\mathcal K\) between \(\N\) and the system of
all (complex) \(p\)-order matrices, in the sense of (\(\alpha\)) in
\lemref{lem:4.1.1}{4.1.1}. Put
\[
 \mathcal KA=a=\{a_{t,s}\},\qquad
 \mathcal KB=b=\{b_{t,s}\}.
\]
Form next the operators \(V,X_1,\ldots,X_p\in\N\) with
\[
 \mathcal KV=v=\{v_{s,t}\},\qquad
 v_{s,t}=\begin{cases}
 1,&s-t\equiv1\pmod p,\\
 0,&\text{otherwise},
 \end{cases}
\]
\[
 \mathcal KX_r=x_r=\{x^{r}_{s,t}\},\qquad
 x^{r}_{s,t}=\begin{cases}
 1,&t=s\ne r,\\
 -1,&t=s=r,\\
 0,&\text{otherwise}.
 \end{cases}
\]
The matrices \(v,x_r\) are unitary. Hence the operators \(V,X_r\) are also
unitary.

Now put \(m=p2^p\) and let \(U_1,\ldots,U_m\) be the operators
\(V^hX_1^{k_1}\cdots X_p^{k_p}\),
\(h=0,1,\ldots,p-1\); \(k_1,\ldots,k_p=0,1\). Then a combination of
these formulas gives, due to \eqref{eq:4.7-alpha} by direct
computation,\footnote{This computation may be outlined as follows. Let
\[
 \epsilon_{s,t}^{(r)}=(-1)^{\delta_{s,r}+\delta_{r,t}},
\]
i.e.
\[
 \epsilon_{s,t}^{(r)}=\begin{cases}
 1,&r\ne t\text{ and }s\ne r\text{ or if }r=s=t,\\
 -1,&\text{if either }r=s\text{ and }r\ne t\text{ or }r\ne s\text{ and }r=t.
 \end{cases}
\]
Then if \(\mathcal KX_rAX_r^{-1}=c^{(r)}=(c_{s,t}^{(r)})\),
\[
 c_{s,t}^{(r)}=\epsilon_{s,t}^{(r)}a_{s,t}.
\]
Let
\[
 B_0=\sum_{k_1,\ldots,k_p=0,1}
 (X_1^{k_1}\cdots X_p^{k_p})A
 (X_1^{k_1}\cdots X_p^{k_p})^{-1}
\]
and \(\mathcal KB_0=d=(d_{t,s})\). Then
\begin{align*}
 d_{t,s}
 &=\sum_{k_1,\ldots,k_p=0,1}
   (\epsilon_{s,t}^{(1)})^{k_1}\cdots
   (\epsilon_{s,t}^{(p)})^{k_p}a_{s,t}\\
 &=\bigl((\epsilon_{s,t}^{(1)})^0+(\epsilon_{s,t}^{(1)})^1\bigr)
   \cdots
   \bigl((\epsilon_{s,t}^{(p)})^0+(\epsilon_{s,t}^{(p)})^1\bigr)a_{s,t}\\
 &=(1+\epsilon_{s,t}^{(1)})\cdots
   (1+\epsilon_{s,t}^{(p)})a_{s,t}.
\end{align*}
Unless \(s=t\), \(\epsilon_{s,t}^{(s)}=-1\) and \(d_{s,t}=0\). If
\(s=t\), \(d_{s,s}=2^pa_{s,s}\). Thus \(B_0\) is a diagonal matrix.

The effect of \(V\) on a diagonal matrix is just to permute the diagonal
terms cyclicly. Hence if \(B\) and \(b_{s,t}\) are as in the text,
\[
 B=\frac1{p2^p}\sum_{h=0}^{p-1}V^hB_0V^{-h}
\]
and
\[
 b_{s,t}=0\quad\text{if }s\ne t,
 \qquad b_{s,s}=\frac1p\sum_{u=1}^{p}a_{u,u}.
\]}
\[
 b_{t,s}=\begin{cases}
 \dfrac1p\displaystyle\sum_{u=1}^{p}a_{u,u}=a_0,&t=s,\\
 0,&\text{otherwise}.
 \end{cases}
\]
Consequently
\begin{equation}\tag{4.7.\(\delta\)}\label{eq:4.7-delta}
 B=a_0\cdot1.
\end{equation}
\eqref{eq:4.7-delta} and \eqref{eq:4.7-gamma} give together
\(0=\operatorname{Tr}_{\M}(B)=a_0\), i.e. \(B=0\). Hence
\eqref{eq:4.7-beta} becomes \(\Metric{A}\leq\epsilon\) as desired.
\end{proof}

We are able to prove:

\begin{theorem}\phantomsection\label{thm:XIII}
If \(\M\) is a factor in case (II\(_1\)), \(\M\) contains a subring
\(\mathbf O\) which is a factor in case (II\(_1\)), and approximately finite.
\end{theorem}

\begin{proof}
Choose a sequence \([p_1,p_2,\ldots]\) which fulfills (i), (ii) in
\lemref{lem:4.1.3}{4.1.3}, e.g. \(p_n=2^n\).

Put \(p_0=1\) and \(\N_0=(\alpha1)\). This is a factor \(\subseteq\M\) and
in case (I\(_1\)). We shall choose by induction for every
\(n=1,2,\ldots\) a factor \(\N_n\subseteq\M\) in the case
(I\(_{p_n}\)). If for any \(n=1,2,\ldots\), \(\N_{n-1}\) has already
been chosen, we choose \(\N_n\) in the following way: \(\N_{n-1},\M\) are
factors in the cases (I\(_{p_{n-1}}\)), (II\(_1\)) respectively;
\(\N_{n-1}\subseteq\M\), \(p_{n-1}\) is a divisor of \(p_n\). Apply
therefore \lemref{lem:4.4.2}{4.4.2} to \(\N_{n-1},\M,p_{n-1},p_n\) in
place of its \(\N,\mathbf O,p,r\). Put \(\N_n=\R\). Then \(\N_n\) is a
factor \(\subseteq\M\) in the case (I\(_{p_n}\)).

Thus
\begin{equation}\tag{4.7.\(\epsilon\)}\label{eq:4.7-epsilon}
 \N_1\subseteq\N_2\subseteq\cdots\subseteq\M.
\end{equation}
Denote by \(S\) the set-theoretical sum of the sequence
\(\N_1,\N_2,\ldots\). Owing to \eqref{eq:4.7-epsilon}, \(S\) is an
algebra along with \(\N_1,\N_2,\ldots\).

Denote the set of the limits of all metrically convergent sequences from
\(S\) by \(S_1\). The same argument as in the proof of
\lemref{lem:4.1.4}{4.1.4} shows that \(S_1\) is an algebra along with \(S\)
and then that it is a ring. Hence \(S\subseteq S_1\) implies
\(\R(S)\subseteq S_1\). On the other hand \(\R(S)\) is closed in the sense
of metric convergence by \thmref{thm:I}{I}. Hence \(S\subseteq\R(S)\)
implies \(S_1\subseteq\R(S)\). Consequently \(\R(S)=S_1\) and so
\begin{equation}\tag{4.7.\(\zeta\)}\label{eq:4.7-zeta}
 \M\supseteq\R(\N_n;n=1,2,\ldots)=\R(S)=S_1.
\end{equation}

Next we prove that the ring \(S_1\) is a factor, i.e. that every
\(A\in S_1\cdot S_1'\) belongs to \((\alpha\cdot1)\).

Consider an \(A\in S_1\cdot S_1'\). Choose \(\epsilon>0\). As
\(A\in S_1\), there exists a \(B_\epsilon\in\N_{p_{n_\epsilon}}\) with
\begin{equation}\tag{4.7.\(\eta\)}\label{eq:4.7-eta}
 \Metric{B_\epsilon-A}<\epsilon.
\end{equation}
Put \(C_\epsilon=B_\epsilon-\operatorname{Tr}_{\M}(B_\epsilon)\cdot1\).
Then \(C_\epsilon\in\N_{p_{n_\epsilon}}\) too, and
\begin{equation}\tag{4.7.\(\theta\)}\label{eq:4.7-theta}
 \operatorname{Tr}_{\M}(C_\epsilon)=0.
\end{equation}
Consider an \(X\in\N_{p_{n_\epsilon}}\). As \(A\in S_1'\),
\(X\in S_1\), so \(AX-XA=0\). Also, clearly
\(B_\epsilon X-XB_\epsilon=C_\epsilon X-XC_\epsilon\). Hence
\(C_\epsilon X-XC_\epsilon=(B_\epsilon-A)X-X(B_\epsilon-A)\) and so
\(\Metric{C_\epsilon X-XC_\epsilon}\leq
2\OpNorm{X}\Metric{B_\epsilon-A}\). This and \eqref{eq:4.7-eta} yield
\begin{equation}\tag{4.7.\(\iota\)}\label{eq:4.7-iota}
 \Metric{C_\epsilon X-XC_\epsilon}\leq2\OpNorm{X}\epsilon.
\end{equation}
\eqref{eq:4.7-theta} and \eqref{eq:4.7-iota} permit us to apply
\lemref{lem:4.7.1}{4.7.1} to \(C_\epsilon,2\epsilon\) in place of its
\(A,\epsilon\). Thus
\[
 \Metric{C_\epsilon}\leq2\epsilon.
\]
Put \(\operatorname{Tr}_{\M}(B_\epsilon)=a_\epsilon\). Then the above
equation means
\[
 \Metric{B_\epsilon-a_\epsilon\cdot1}\leq2\epsilon.
\]
Hence by \eqref{eq:4.7-eta}
\begin{equation}\tag{4.7.\(\kappa\)}\label{eq:4.7-kappa}
 \Metric{A-a_\epsilon\cdot1}<3\epsilon.
\end{equation}
As \eqref{eq:4.7-kappa} holds for all \(\epsilon>0\), it proves that \(A\)
is the limit of a metrically convergent sequence from \((\alpha\cdot1)\).
As \((\alpha\cdot1)\) is obviously closed in this sense, we have
\(A\in(\alpha\cdot1)\) as desired.

Thus \(S_1\) is a factor. The considerations at the beginning of the proof of
(i) in \lemref{lem:4.1.2}{4.1.2} show again that
\(D_{S_1}(E)=D_{\M}(E)\) for all projections \(E\in S_1\). Hence \(S_1\)
is in a finite case, along with \(\M\). If \(S_1\) were in a case (I\(_q\)),
\(q=1,2,\ldots\), we could argue as follows: \(\N_{p_n}\subseteq S_1\) and
\(\N_{p_n}\) is a factor in case (I\(_{p_n}\)). Hence \(p_n\) is a divisor
of \(q\) by (i) in \lemref{lem:4.1.2}{4.1.2}. But this is impossible, since
\(p_n\to\infty\) as \(n\to\infty\). Therefore \(S_1\) must be in the case
(II\(_1\)).

Now \eqref{eq:4.7-epsilon}, \eqref{eq:4.7-zeta} prove that \(S_1\) is
approximately finite of type \([p_1,p_2,\ldots]\)---i.e. approximately finite
by \thmref{thm:XII}{XII}.

Thus \(\mathbf O=S_1\) meets all our requirements.
\end{proof}

And to conclude:

\begin{theorem}\phantomsection\label{thm:XIV}
Approximately finite factors \(\M\) (in the case II\(_1\)) exist and they
have all the same algebraical type \(\overline{\M}\).

We denote this algebraical type by \(\overline{\mathbf A}\).
\end{theorem}

\begin{proof}
The existence follows from \thmref{thm:XIII}{XIII}.

Consider two approximately finite \(\M_1,\M_2\). Choose a sequence
\([p_1,p_2,\ldots]\) which fulfills (i), (ii) in
\lemref{lem:4.1.3}{4.1.3}, e.g. \(p_n=2^n\). Then \(\M_1,\M_2\) are both
approximately finite of type \([p_1,p_2,\ldots]\) by
\thmref{thm:XII}{XII}. Hence they are algebraically isomorphic by
\thmref{thm:XI}{XI}, i.e. \(\overline{\M}_1=\overline{\M}_2\).

Thus \(\overline{\M}\) (\(\M\) approximately finite) is uniquely determined
and the definition of \(\overline{\mathbf A}\) is justified.
\end{proof}

\PaperSection{sec:4.8}{Algebraical isomorphism invariants of the approximately finite type.}
The algebraical isomorphism invariants described at the end of
\secref{sec:2.10}{\S2.10} will now be studied for the algebraical type
\(\overline{\mathbf A}\). These invariants are
\begin{enumerate}[label=(\arabic*)]
\item Validity of \(\overline{\mathbf A}^{c}=\overline{\mathbf A}\),
\item The fundamental group:
\(\mathfrak G=\mathfrak G(\overline{\mathbf A})\).
\end{enumerate}

\begin{lemma}\phantomsection\label{lem:4.8.1}
\(\overline{\mathbf A}^{c}=\overline{\mathbf A}\).
\end{lemma}

\begin{proof}
By \thmref{thm:XIV}{XIV} on one hand and \thmref{thm:III}{III} on the other,
it suffices to show that a definition of approximate finiteness is unaffected
by a conjugate or by a dual isomorphism. This is obviously true, even for all
the definitions given in \thmref{thm:XII}{XII} and for both types of general
isomorphisms (cf. \(B^*\)) in \secref{sec:2.2}{\S2.2}) mentioned above.
\end{proof}

\begin{lemma}\phantomsection\label{lem:4.8.2}
\(\mathfrak G=\mathfrak G(\overline{\mathbf A})\) contains all rational
numbers \(\theta\) in \(0<\theta<\infty\).
\end{lemma}

\begin{proof}
Since \(\mathfrak G\) is a group, it suffices to prove that it contains all
\(1/p\), \(p=2,3,\ldots\).

Consider such a \(1/p\). Choose a sequence \([p_1,p_2,\ldots]\) which
fulfills (i), (ii) in \lemref{lem:4.1.3}{4.1.3} and with all \(p_m\) divisible
by \(p\), e.g. \(p_n=p^n\).

Choose an approximately finite \(\M\) and the \(\N_1,\N_2,\ldots\), for this
\(\M\) and the sequence \([p_1,p_2,\ldots]\) in the sense of
\defref{def:4.1.1}{4.1.1}.

The considerations at the beginning of the proof of (i) in
\lemref{lem:4.1.2}{4.1.2} show again that
\(D_{\M}(G)=D_{\N_n}(G)\) for all projections \(G\in\N_n\;(\subseteq\M)\).

As \(\N_1\) is in the case (I\(_{p_1}\)), \(p_1\) divisible by \(p\), there
exists a projection \(E\in\N_1\) with \(D_{\N_1}(E)=1/p\). Thus
\begin{equation}\tag{4.8.\(\alpha\)}\label{eq:4.8-alpha}
 E\in\N_1\subseteq\N_2\subseteq\cdots\subseteq\M
\end{equation}
while \(D_{\N_1}(E)=1/p\) yields
\begin{equation}\tag{4.8.\(\beta\)}\label{eq:4.8-beta}
 \frac{D_{\mathbf O}(E)}{D_{\mathbf O}(1)}=D_{\mathbf O}(E)=\frac1p
 \quad\text{for }\mathbf O=\N_1,\N_2,\ldots,\M.
\end{equation}
Denote the range of \(E\) by \(\mathfrak M\).

By \eqref{eq:4.8-alpha} we can form the factors
\(\N_{1(\mathfrak M)},\N_{2(\mathfrak M)},\ldots,\M_{(\mathfrak M)}\) in
\(\mathfrak M\), and we have
\begin{equation}\tag{4.8.\(\gamma\)}\label{eq:4.8-gamma}
 \N_{1(\mathfrak M)}\subseteq\N_{2(\mathfrak M)}\subseteq\cdots
 \subseteq\M_{(\mathfrak M)}.
\end{equation}
By \eqref{eq:4.8-beta},
\(\overline{\N_{n(\mathfrak M)}}=\overline{\N}_n^{1/p}\),
\(\overline{\M_{(\mathfrak M)}}=\overline{\M}^{1/p}\) (cf. the discussion
before \thmref{thm:VI}{VI}). The former implies by
\lemref{lem:2.10.3}{2.10.3} that
\begin{equation}\tag{4.8.\(\delta\)}\label{eq:4.8-delta}
 \N_{n(\mathfrak M)}\text{ is in the case }(\mathrm I_{p_n/p}).
\end{equation}
We restate the latter:
\begin{equation}\tag{4.8.\(\epsilon\)}\label{eq:4.8-epsilon}
 \overline{\M_{(\mathfrak M)}}=\overline{\M}^{1/p}.
\end{equation}
We know that \(\M\) is the closure of the set-theoretical sum of
\(\N_1,\N_2,\ldots\) in the strong topology.\footnote{That closure is
obviously \(=\R(\N_n,n=1,2,\ldots)\). Hence it is \(\M\) by (iii) in
\defref{def:4.1.1}{4.1.1}.} Hence it follows from \eqref{eq:4.8-alpha} that
\(\M_{(\mathfrak M)}\) is the closure of the set-theoretical sum of
\(\N_{1(\mathfrak M)},\N_{2(\mathfrak M)},\ldots\) in the strong topology.
This and \eqref{eq:4.8-gamma} give together
\begin{equation}\tag{4.8.\(\zeta\)}\label{eq:4.8-zeta}
 \M_{(\mathfrak M)}=\R(\N_{n(\mathfrak M)};n=1,2,\ldots).
\end{equation}
Now \eqref{eq:4.8-delta}, \eqref{eq:4.8-gamma} and
\eqref{eq:4.8-zeta} coincide with (i), (ii), (iii) in
\defref{def:4.1.1}{4.1.1} as applied to
\(\M_{(\mathfrak M)},\allowbreak\N_{1(\mathfrak M)},\allowbreak
\N_{2(\mathfrak M)},\allowbreak\ldots,\allowbreak
p_1/p,\allowbreak p_2/p,\allowbreak\ldots\) in place of its
\(\M,\allowbreak\N_1,\allowbreak\N_2,\allowbreak\ldots,\allowbreak
p_1,\allowbreak p_2,\allowbreak\ldots\). Hence \(\M_{(\mathfrak M)}\) is
approximately finite.

As \(\M\) is approximately finite too, this shows that
\(\overline{\M_{(\mathfrak M)}}=\overline{\M}\) by
\thmref{thm:XIV}{XIV}. Owing to \eqref{eq:4.8-epsilon}, this gives
\(\overline{\M}^{1/p}=\overline{\M}\). Hence
\(1/p\in\mathfrak G\) as desired.
\end{proof}

\begin{lemma}\phantomsection\label{lem:4.8.3}
If there exists for every \(\epsilon>0\) an \(\alpha\) in
\(1-\epsilon<\alpha<1\) such that \(\overline{\M}^{\alpha}\) is
approximately finite, then \(\M\) itself is approximately finite.
\end{lemma}

\begin{proof}
Assume that \(\M\) possesses the property described above. We shall establish
its approximate finiteness in the sense of \defref{def:4.5.2}{4.5.2}.

Let accordingly \(A_1,\ldots,A_m\in\M\) and an \(\epsilon>0\) be given. We
shall construct an \(\N\) which fulfills (i)--(iii) in that definition.

Apply first our hypothesis concerning \(\M\) to
\begin{equation}\tag{4.8.\(\eta\)}\label{eq:4.8-eta}
 \epsilon'=\left(
 \frac{\epsilon}{4\operatorname{Max}(\OpNorm{A_s},\ s=1,\ldots,m)}
 \right)^2
\end{equation}
in place of its \(\epsilon\). Choose the \(\alpha\) mentioned there with
\(1-\epsilon'<\alpha<1\). Then choose a projection \(E\in\M\) with
\(D_{\M}(E)=\alpha\). Let \(\mathfrak M\) be the range of \(E\). Form
\(\M_{(\mathfrak M)}\). Then we have
\(\overline{\M_{(\mathfrak M)}}=\overline{\M}^{\alpha}\) (cf. the discussion
before \thmref{thm:VI}{VI}). It follows from our assumptions concerning
\(\alpha\) that \(\M_{(\mathfrak M)}\) is approximately finite.

By \cite{ref5}, p.~187, (i), in Lemma 11.3.3,
\[
 A\longleftrightarrow A_{(\mathfrak M)}
\]
is a one-to-one correspondence between all \(A\in\M\) with
\(EA=AE=A\) and all \(A_{(\mathfrak M)}\in\M_{(\mathfrak M)}\). Now we can
repeat, with a slight variation an argument that was used in the proof of
\thmref{thm:XI}{XI}: Clearly \(c\operatorname{Tr}_{\M}(A)\) will do as
\(\operatorname{Tr}_{\M_{(\mathfrak M)}}(A_{(\mathfrak M)})\) in the sense
of \cite{ref6}, p.~219, Property IV, for the
\(A,A_{(\mathfrak M)}\) described above, if we choose the normalizing factor
\(c>0\) so that \(c\operatorname{Tr}_{\M}(A)=1\) when
\(A_{(\mathfrak M)}\) is the unit of \(\M_{(\mathfrak M)}\). This means
\(A=E\), so we want \(c\operatorname{Tr}_{\M}(E)=1\). Since
\(\operatorname{Tr}_{\M}(E)=D_{\M}(E)=\alpha\) this means \(c=1/\alpha\).
And as the trace is uniquely characterized by the above property,
\[
 \operatorname{Tr}_{\M_{(\mathfrak M)}}(A_{(\mathfrak M)})
 =\frac1\alpha\operatorname{Tr}_{\M}(A).
\]
From this
\[
 \frac1\alpha\operatorname{Tr}_{\M}(A^*A)
 =\operatorname{Tr}_{\M_{(\mathfrak M)}}((A^*A)_{(\mathfrak M)})
 =\operatorname{Tr}_{\M_{(\mathfrak M)}}
  ((A^*)_{(\mathfrak M)}A_{(\mathfrak M)})
\]
for the same \(A,A_{(\mathfrak M)}\), hence
\begin{equation}\tag{4.8.\(\theta\)}\label{eq:4.8-theta}
 \Metric{A_{(\mathfrak M)}}_{\M_{(\mathfrak M)}}^2
 =\frac1{\alpha^2}\Metric{A}_{\M}^2.
\end{equation}
\footnote{As a rule we did not indicate the factor (in a finite case) with
respect to which \(\Metric{A}\) is formed. However, here it is desirable.}

Now put \(B_s=EA_sE\), \(s=1,\ldots,m\). Then we have
\(D_{\M}(1-E)<\epsilon'\). Hence
\(\Metric{1-E}_{\M}<\epsilon'^{1/2}\) and therefore
\[
 \Metric{A_s-B_s}_{\M}
 =\Metric{A_s(1-E)+(1-E)A_sE}_{\M}
 <\OpNorm{A_s}\epsilon'^{1/2}.
\]
Hence by \eqref{eq:4.8-eta}
\begin{equation}\tag{4.8.\(\iota\)}\label{eq:4.8-iota}
 \Metric{A_s-B_s}_{\M}<\epsilon/2.
\end{equation}

Now \(B_s\in\M\) and \(EB_s=B_sE=B_s\) imply
\(B_{s(\mathfrak M)}\in\M_{(\mathfrak M)}\). We now apply
\defref{def:4.5.2}{4.5.2} to the approximately finite
\(\M_{(\mathfrak M)}\) with
\(B_{1(\mathfrak M)},\ldots,B_{m(\mathfrak M)},\epsilon/2\alpha^2\) in
place of its \(A_1,\ldots,A_m\) and \(\epsilon\). This yields a ring
\(\N^+\subseteq\M_{(\mathfrak M)}\) of finite order with
\(C_1^+,\ldots,C_m^+\in\N^+\) such that
\begin{equation}\tag{4.8.\(\kappa\)}\label{eq:4.8-kappa}
 \Metric{C_s^+-B_{s(\mathfrak M)}}_{\M_{(\mathfrak M)}}
 <\epsilon/2\alpha^2.
\end{equation}
As \(\N^+\subseteq\M_{(\mathfrak M)}\) we extend every \(C^+\in\N^+\) to a
\(C\in\M\) with \(EC=CE=C\) and \(C^+=C_{(\mathfrak M)}\).\footnote{
\(C\) is reduced by \(\mathfrak M\), its part in \(\mathfrak M\) is \(C^+\)
and its part in \(\mathfrak M^{\perp}\) is 0.} These \(C\) form obviously a
ring \(\N\subseteq\M\) of finite order with
\(\N^+=\N_{(\mathfrak M)}\). By this process the
\(C_1^+,\ldots,C_m^+\in\N^+\) give rise to the
\(C_1,\ldots,C_m\in\N\). Now \eqref{eq:4.8-theta} gives
\[
 \Metric{C_s^+-B_{s(\mathfrak M)}}_{\M_{(\mathfrak M)}}
 =\Metric{(C_s-B_s)_{(\mathfrak M)}}_{\M_{(\mathfrak M)}}
 =\frac1{\alpha^2}\Metric{C_s-B_s}_{\M}
\]
and so \eqref{eq:4.8-kappa} becomes
\begin{equation}\tag{4.8.\(\lambda\)}\label{eq:4.8-lambda}
 \Metric{C_s-B_s}_{\M}<\epsilon/2.
\end{equation}
\eqref{eq:4.8-iota} and \eqref{eq:4.8-lambda} together give
\begin{equation}\tag{4.8.\(\mu\)}\label{eq:4.8-mu}
 \Metric{C_s-A_s}_{\M}<\epsilon.
\end{equation}
Thus the ring \(\N\subseteq\M\) of finite order meets the requirements
(i)--(iii) in \defref{def:4.5.2}{4.5.2}. Hence \(\M\) is approximately finite
as asserted.
\end{proof}

\begin{lemma}\phantomsection\label{lem:4.8.4}
\(\mathfrak G=\mathfrak G(\overline{\mathbf A})\) is the set \(P\) of all
(real) numbers \(\theta\) in \(0<\theta<\infty\).
\end{lemma}

\begin{proof}
Consider a \(\theta\) in \(0<\theta<\infty\). By
\lemref{lem:4.8.2}{4.8.2} we have for every rational \(\xi\) in
\(0<\xi<\infty\),
\(\overline{\mathbf A}^{\xi}=\overline{\mathbf A}\). Hence
\(\overline{\mathbf A}^{\xi}\) is approximately finite. Therefore
\((\overline{\mathbf A}^{\theta})^\alpha
=\overline{\mathbf A}^{\theta\alpha}\) is approximately finite if
\(\theta\alpha\) is rational.

Now it is obviously possible to find for any given \(\epsilon>0\) an
\(\alpha\) in \(1-\epsilon<\alpha<1\) for which \(\theta\alpha\) is rational.
This makes \((\overline{\mathbf A}^{\theta})^\alpha\) approximately finite.
Hence \(\overline{\mathbf A}^{\theta}\) itself is approximately finite by
\lemref{lem:4.8.3}{4.8.3} with \(\overline{\mathbf A}^{\theta}\) in place of
its \(\M\). Consequently \thmref{thm:XIV}{XIV} gives
\(\overline{\mathbf A}^{\theta}=\overline{\mathbf A}\). Hence
\(\theta\in\mathfrak G\) as desired.
\end{proof}

We restate the results of Lemmas \lemref{lem:4.8.1}{4.8.1} and
\lemref{lem:4.8.4}{4.8.4}.

\begin{theorem}\phantomsection\label{thm:XV}
The algebraical isomorphism invariants in (1), (2) at the beginning of this
section, behave for the approximate finite type \(\overline{\mathbf A}\) as
follows:
\begin{enumerate}[label=(\arabic*)]
\item \(\overline{\mathbf A}^{c}=\overline{\mathbf A}\),
\item \(\mathfrak G(\overline{\mathbf A})=P\).
\end{enumerate}
\end{theorem}

It is worth pointing out that the two invariants (1), (2) are not sufficient to
characterize an algebraical type \(\overline{\M}\) in the case (II\(_1\)). A
proof of this is only possible if we show first that not all \(\overline{\M}\)
are approximately finite: If all were, then \thmref{thm:XIV}{XIV} would give
\(\overline{\M}=\overline{\mathbf A}\), i.e. only one \(\overline{\M}\) in
case (II\(_1\)) would exist and thus no invariants for these
\(\overline{\M}\) would be needed at all.

Now we shall prove in Theorems \thmref{thm:XVI}{XVI},
\thmref{thm:XVI-prime}{XVI\(^{\prime}\)} that not approximately finite
\(\overline{\M}\) in the case (II\(_1\)) exist. By making advance use of
this result we prove:

\begin{lemma}\phantomsection\label{lem:4.8.5}
The algebraical isomorphism invariants in (1), (2) at the beginning of this
paragraph are not sufficient to characterize the algebraical type
\(\overline{\M}\) in the case (II\(_1\)).
\end{lemma}

\begin{proof}
Assume the opposite.

Consider a type \(\overline{\M}\) in the case (II\(_1\)).

Owing to the formulae \eqref{eq:2.3.alpha} and \eqref{eq:2.8.beta} the
validity of \(\overline{\M}^{c}=\overline{\M}\) is unaffected if we replace
\(\overline{\M}\) by \(\overline{\M}^{c}\) or by an
\(\overline{\M}^{\theta}\;(0<\theta<\infty)\). By
\lemref{lem:2.10.2}{2.10.2} \(\mathfrak G(\overline{\M})\) is unchanged if
we replace \(\overline{\M}\) by \(\overline{\M}^{c}\) or by an
\(\overline{\M}^{\theta}\;(0<\theta<\infty)\). In other words: The
invariants (1) and (2) are unaffected by a replacement of \(\overline{\M}\)
by \(\overline{\M}^{c}\) or by an \(\overline{\M}^{\theta}\).

Therefore our assumptions concerning these invariants permit us to conclude
\(\overline{\M}=\overline{\M}^{c}\) and
\(\overline{\M}=\overline{\M}^{\theta}\). The latter equations imply
\(\theta\in\mathfrak G(\overline{\M})\) and therefore
\(\mathfrak G(\overline{\M})=P\).

Thus the two invariants (1) and (2) are uniquely determined, independently of
the choice of \(\overline{\M}\) in the case (II\(_1\)).

Hence a second application of our assumption concerning these invariants
yields that there exists only one type \(\overline{\M}\) in the case
(II\(_1\)). \(\overline{\mathbf A}\) is such a type; hence all
\(\overline{\M}\) in the case (II\(_1\)) must equal
\(\overline{\mathbf A}\), i.e. be approximately finite.

However, as stated above, the opposite will be proven in Theorems
\thmref{thm:XVI}{XVI}, \thmref{thm:XVI-prime}{XVI\(^{\prime}\)}, and we
propose to use this now. Thus we have a contradiction and the
proof is completed.
\end{proof}

We conclude with the observation that the existence of not approximately
finite factors \(\M\) in the case (II\(_1\)) will be established with the help
of algebraic invariants other than (1), (2) above (cf. the beginning of
\secref{sec:6.1}{\S6.1}).

\PaperChapter{5}{chapter:V}{On Various Examples}

\PaperSection{sec:5.1}{General remarks on the examples.}
We established the existence of approximately finite factors (i.e. of their
type \(\overline{\mathbf A}\)) by the indirect procedure of Theorems
\thmref{thm:XIII}{XIII} and \thmref{thm:XIV}{XIV}.\footnote{The decisive
passage from Theorem \thmref{thm:XIII}{XIII} to Theorem
\thmref{thm:XIV}{XIV} makes use
of the existence of factors in the case (II\(_1\)), for instance the examples
of \cite{ref5}.} It is therefore desirable to furnish examples of
this by explicit and direct construction and also to investigate those examples
of factors in the case (II\(_1\)) which we possess already as to their
approximate finiteness. The latter investigation is also necessary in order to
find not approximately finite factors in the case (II\(_1\)).

In connection with this work we shall develop a new technique to construct
examples of factors in the case (II\(_1\)), which is a simplification of the
original one of \cite{ref5}, pp.~192--209, Part IV.

We begin by establishing the approximate finiteness of some of our old
examples: certain ones in \cite{ref5} and that one in
\cite{ref7}, pp.~72--77, \S7.5.

\PaperSection{sec:5.2}{Approximate finiteness of certain existing examples.}

\begin{lemma}\phantomsection\label{lem:5.2.1}
The ring \(\mathbf C^{*1}\) of \cite{ref7}, p.~72, Lemma 7.5.1 is
approximately finite.
\end{lemma}

\begin{proof}
Using the notations loc. cit., particularly pp.~68, 72, we have
\[
 \mathbf C^{*1}=\R(\overline{\mathbf B}_{(n,1)};n=1,2,\ldots).
\]
Hence
\begin{equation}\tag{5.2.\(\alpha\)}\label{eq:5.2-alpha}
 \mathbf C^{*1}=\R(\N_m;m=1,2,\ldots)
\end{equation}
with
\begin{equation}\tag{5.2.\(\beta\)}\label{eq:5.2-beta}
 \N_m=\R(\overline{\mathbf B}_{(n,1)};n=1,\ldots,m).
\end{equation}
Thus \(\N_1\subseteq\N_2\subseteq\cdots\) and therefore
\(\mathbf C^{*1}\) is approximately finite by \defref{def:4.6.1}{4.6.1} if
the \(\N_m\) are shown to be of finite order.\footnote{\phantomsection\label{fn:44}It is not difficult to
show that \(\N_n\) is a factor in the case (I\(_{2^n}\)). Hence
\defref{def:4.1.1}{4.1.1} would also be applicable, yielding approximate
finiteness of type \([2,4,8,\ldots]\).}

These considerations belong in
\(\prod_{n=1,2,\ldots}(\Hilbert_{(n,1)}\otimes\Hilbert_{(n,2)})\) where
\(\mathbf C^{*1}\) is a factor in the case (II\(_1\)), but in establishing
the finite order of \(\N_m\), we may as well consider the entire space
\(\prod_{n=1,2,\ldots}(\Hilbert_{(n,1)}\otimes\Hilbert_{(n,2)})\).

We proceed as in the lower part of p.~68, loc. cit.
\(B_{(n,1)}\) is in \(\Hilbert_{(n,1)}\). It is first extended to
\(\overline B_{(n,1)}\) in
\(\Hilbert_{(n,1)}\otimes\Hilbert_{(n,2)}\) and then to
\(\overline{\mathbf B}_{(n,1)}\) in
\(\prod_{n=1,2,\ldots}(\Hilbert_{(n,1)}\otimes\Hilbert_{(n,2)})\). Hence the
\(\overline{\mathbf B}_{(n,1)}\) commute with each other. Therefore the
\(\N_m\) of \eqref{eq:5.2-beta} is of finite order if the
\(\overline{\mathbf B}_{(n,1)}\) are. Now \(B_{(n,1)}\) is of finite order
because it is in the finite 2-dimensional space \(\Hilbert_{(n,1)}\) and so
\(\overline{\mathbf B}_{(n,1)}\) is of finite order too.
\end{proof}

\begin{lemma}\phantomsection\label{lem:5.2.2}
The ring \(\M\) of \cite{ref5}, pp.~203--204, Theorem XI, is
approximately finite if the group \(\mathfrak G\) possesses this property:
\begin{enumerate}[label=(\roman*)]
\item \(\mathfrak G\) is the set-theoretical sum of a sequence
\(\mathfrak G_1\subseteq\mathfrak G_2\subseteq\cdots\) of finite groups.
\end{enumerate}
(We assume, of course, that the requirements of loc. cit.
\cite{ref5}, p.~206, Lemma 13.1.2, are fulfilled, placing \(\M\) in
the case (II\(_1\)).)
\end{lemma}

\begin{proof}
Using the notations loc. cit., particularly those on pp.~198--200, we have
\begin{equation}\tag{5.2.\(\gamma\)}\label{eq:5.2-gamma}
 \M=\R(\overline U_{a_0},\overline L_{\varphi(x)};
 a_0\in\mathfrak G,\ \varphi(x)\text{ bounded and measurable}).
\end{equation}
A familiar argument shows that the \(\varphi(x)\) in \eqref{eq:5.2-gamma}
can be restricted considerably without changing the ring \(\M\). Indeed:

First. It suffices for reasons of continuity to consider the \(\varphi(x)\)
with finite ranges only.

Second. These \(\varphi(x)\) are finite linear combinations of the
characteristic functions
\[
 \chi_S(x)=\begin{cases}1,&x\in S,\\0,&\text{otherwise}.
 \end{cases}
\]

Third. Let \(T_1,T_2,\ldots\) be a countable basis for the \(T\)
(\(\subseteq S\) and measurable).\footnote{The space of all measurable sets
\(T\) (disregarding sets of measure zero) can be metricized by the distance
\[
 \overline{T'T''}=\mu(T'+T'')-\mu(T'T'').
\]
This is obviously equal to
\[
 \lvert\chi_{T'}-\chi_{T''}\rvert^2
 =\int_S\lvert\chi_{T'}(x)-\chi_{T''}(x)\rvert^2\,dx.
\]
Since \(\Hilbert_S\) is separable, there exists a sequence
\(\chi_{T_1},\chi_{T_2},\ldots\) dense in the set of all \(\chi_T\), and
correspondingly the sequence \(T_1,T_2,\ldots\) is dense in the set of all
\(T\). This is what we mean by a basis.} Then it suffices for reasons of
continuity to consider these \(\varphi(x)=\chi_{T_i}(x)\),
\(i=1,2,\ldots\), only.

Finally, these properties of the system \(T_1,T_2,\ldots\) are not lost if we
replace it by another countable system which contains it. We do this in such a
way that the increased system \(\overline T_1,\overline T_2,\ldots\)
possesses these further properties:
\begin{enumerate}[label=(\arabic*)]
\item It contains \(T'\cdot T''\) along with \(T',T''\).
\item It contains \(a_0T'\) along with \(T'\) for any
\(a_0\in\mathfrak G\).
\end{enumerate}
These being understood, we have
\begin{equation}\tag{5.2.\(\delta\)}\label{eq:5.2-delta}
 \M=\R(\overline U_{a_0},\overline L_{\chi_{\overline T_i}(x)};
 a_0\in\mathfrak G,\ i=1,2,\ldots).
\end{equation}

Consider now the finite groups \(\mathfrak G_1,\mathfrak G_2,\ldots\) of
(i). \(\overline T_1,\ldots,\overline T_n\) generates, with the help of the
operations (1) and (2) when \(a_0\) in the latter is restricted to
\(\mathfrak G_n\), a finite subsystem of
\(\overline T_1,\overline T_2,\ldots\),---say the set
\((\overline T_i;i\in I_n)\). So \(I_n\) is a finite set. Clearly
\(I_1\subseteq I_2\subseteq\cdots\). \(I_n\) contains
\(1,2,\ldots,n\). Hence \((1,2,\ldots)\) is the set-theoretical sum of
\(I_1,I_2,\ldots\).

Consequently \eqref{eq:5.2-delta} gives
\begin{equation}\tag{5.2.\(\epsilon\)}\label{eq:5.2-epsilon}
 \M=\R(\N_n;n=1,2,\ldots)
\end{equation}
with
\begin{equation}\tag{5.2.\(\zeta\)}\label{eq:5.2-zeta}
 \N_n=\R(\overline U_{a_0},\overline L_{\chi_{\overline T_i}(x)};
 a_0\in\mathfrak G_n,\ i\in I_n).
\end{equation}
Thus \(\N_1\subseteq\N_2\subseteq\cdots\) and therefore \(\M\) is
approximately finite by \defref{def:4.6.1}{4.6.1} if the \(\N_n\) are shown
to be of finite order.\footnote{This time we really must use
\defref{def:4.6.1}{4.6.1} and it would not be possible to replace it by
\defref{def:4.1.1}{4.1.1}---in contrast with footnote~\ref{fn:44}.}

Now the group character of \(\mathfrak G_n\) and the construction of \(I_n\)
show that the operators
\begin{equation}\tag{5.2.\(\eta\)}\label{eq:5.2-eta}
 \overline U_{a_0}\overline L_{\chi_{\overline T_i}(x)},
 \qquad a_0\in\mathfrak G_n,\quad i\in I_n
\end{equation}
are reproduced by multiplication. As \(\mathfrak G_n,I_n\) are finite, these
operators \eqref{eq:5.2-eta} form a finite set too. Hence they are a finite
basis of \(\N_n\), completing the proof.
\end{proof}

This last result establishes the approximate finiteness of the examples
(\(\beta\)), (\(\gamma\)) in \cite{ref5}, p.~208. Thus our
\thmref{thm:XIV}{XIV} establishes their algebraic isomorphism and
\cite{ref6}, p.~244, Theorem XI, their spatial isomorphism as
announced by (v) in \cite{ref5}, p.~209.
% Editorial note E41: The source prints ``p. 299''; corrected to ``p. 209'',
% since reference [5] ends at p. 229 and the cited announcement is on p. 209.
In order to include (\(\alpha\)), p.~208, eod., we need this:

\begin{lemma}\phantomsection\label{lem:5.2.3}
The condition (i) in \lemref{lem:5.2.2}{5.2.2} can be replaced by this
condition:
\begin{enumerate}[label=(\roman*)]
\setcounter{enumi}{1}
\item \(\mathfrak G\) is abelian.
\end{enumerate}
\end{lemma}

The proof of this lemma is somewhat complicated. It requires some rather deep
results on the decompositions of mappings of measurable sets, which will be
published elsewhere. We shall not pursue this matter further on this occasion.

\PaperSection{sec:5.3}{A simplified construction of factors from groups.}
We proceed to expound the new (simplified) method for the construction of
examples of factors in the case (II\(_1\)), announced in
\secref{sec:5.1}{\S5.1}.

Let a group \(\mathfrak G\) be given. For the moment \(\mathfrak G\) could be
finite or countably infinite, but we shall be forced to assume the latter at a
subsequent stage (cf. \lemref{lem:5.3.4}{5.3.4} and the Remark after it).

Form a unitary vector space \(\Hilbert\) by using the elements of
\(\mathfrak G\) as indices for the vector components, i.e. \(\Hilbert\) is the
set of all complex vectors \(f=[x_a;a\in\mathfrak G]\) with a finite
\(\sum_{a\in\mathfrak G}|x_a|^2\) and for
\(f=[x_a;a\in\mathfrak G]\) and \(g=[y_a;a\in\mathfrak G]\),
\((f,g)=\sum_{a\in\mathfrak G}x_a\overline{y_a}\).\footnote{This
\(\Hilbert\) coincides with the \(\Hilbert_{\mathfrak G}\) of
\cite{ref5}, p.~194, Definition 12.1.4. The only difference is that
we now write \(x_a\) in place of the \(f(a)\), loc. cit.}

Now we introduce and discuss some operators in \(\Hilbert\).

\begin{lemma}\phantomsection\label{lem:5.3.1}
Define the operators
\begin{align*}
\text{(i)}\quad U_{a_0}[x_a;a\in\mathfrak G]
 &= [x_{aa_0};a\in\mathfrak G],\\
\text{(ii)}\quad V_{a_0}[x_a;a\in\mathfrak G]
 &= [x_{a_0^{-1}a};a\in\mathfrak G],\\
\text{(iii)}\quad W[x_a;a\in\mathfrak G]
 &= [x_{a^{-1}};a\in\mathfrak G]
 \qquad\text{for any }a_0\in\mathfrak G.
\end{align*}
These operators are unitary and
\[
 W=W^{-1},\qquad WU_{a_0}W=V_{a_0}.\footnote{\phantomsection\label{fn:48}One also verifies immediately
\[
 U_{a_1}U_{a_0}=U_{a_1a_0},\qquad
 V_{a_1}V_{a_0}=V_{a_1a_0},
\]
i.e. \(a\mapsto U_a\) and \(a\mapsto V_a\) are both representations of
\(\mathfrak G\). They correspond in fact to the regular representation of
\(\mathfrak G\).}
\]
In other words, \(f\mapsto Wf\) is an involutory spatial automorphism of
\(\Hilbert\) which interchanges \(U_{a_0}\) with \(V_{a_0}\).
\end{lemma}

\begin{proof}
All assertions are immediately verified.
\end{proof}

We introduce, as usual, the complete normalized orthogonal system
\begin{equation}\tag{5.3.\(\alpha\)}\label{eq:5.3-alpha}
 \varphi_a,\qquad a\in\mathfrak G,
\end{equation}
in \(\Hilbert\), where
\(\varphi_a=[\delta_{a,b};b\in\mathfrak G]\). Thus
\(f=[x_a;a\in\mathfrak G]\) is equivalent to
\(f=\sum_{a\in\mathfrak G}x_a\varphi_a\) and
\begin{equation}\tag{5.3.\(\beta\)}\label{eq:5.3-beta}
 U_{a_0}\varphi_a=\varphi_{aa_0^{-1}},\qquad
 V_{a_0}\varphi_a=\varphi_{a_0a},\qquad
 W\varphi_a=\varphi_{a^{-1}}.
\end{equation}

\begin{lemma}\phantomsection\label{lem:5.3.2}
Let \(I\) be the set of all \(U_{a_0}\) and \(J\) the set of all
\(V_{a_0}\). Form for each bounded operator \(A\) in \(\Hilbert\), its
matrix \(\{\alpha_{a,b}\}\;(a,b\in\mathfrak G)\) in the usual
way.\footnote{\phantomsection\label{fn:49}The definition is
\[
 \tag{49,\(\alpha\)}\label{eq:49-alpha}
 \alpha_{a,b}=(A\varphi_a,\varphi_b)
              =(\varphi_a,A^*\varphi_b).
\]
Hence
\[
 \OpNorm{A\varphi_a}^{2}
 =\sum_{b\in\mathfrak G}|(A\varphi_a,\varphi_b)|^2
 =\sum_{b\in\mathfrak G}|\alpha_{a,b}|^2,
\]
\[
 \OpNorm{A^*\varphi_b}^{2}
 =\sum_{a\in\mathfrak G}|(\varphi_a,A^*\varphi_b)|^2
 =\sum_{a\in\mathfrak G}|\alpha_{a,b}|^2,
\]
i.e.
\[
 \tag{49,\(\beta\)}\label{eq:49-beta}
 \text{all }\sum_{b\in\mathfrak G}|\alpha_{a,b}|^2,
 \quad\sum_{a\in\mathfrak G}|\alpha_{a,b}|^2\text{ are finite}.
\]
Finally
\[
 (Af,\varphi_b)=(f,A^*\varphi_b)
 =\sum_{a\in\mathfrak G}(f,\varphi_a)(\varphi_a,A^*\varphi_b)
 =\sum_{a\in\mathfrak G}\alpha_{a,b}(f,\varphi_a),
\]
i.e.
\[
 \tag{49,\(\gamma\)}\label{eq:49-gamma}
 \left.
 \begin{aligned}
 f&=\sum_{a\in\mathfrak G}x_a\varphi_a,\\
 Af&=\sum_{a\in\mathfrak G}y_a\varphi_a
 \end{aligned}
 \right\}
 \quad\Longrightarrow\quad
 y_b=\sum_{a\in\mathfrak G}\alpha_{a,b}x_a.
\]
The matrix operations for these (complex numerical) matrices are defined in
the same way as for the (operator) matrices of \(B_p\) in
\secref{sec:2.4}{\S2.4}.}

Then \(A\in I'\) if and only if \(\alpha_{a,b}\) has the form
\(\alpha_{a,b}=\eta_{ab^{-1}}\) and \(A\in J'\) if and only if
\(\alpha_{a,b}\) has the form \(\alpha_{a,b}=\eta_{a^{-1}b}\), i.e.
\(\alpha_{a,b}\) depends on the value of \(ab^{-1}\) only or on the value of
\(a^{-1}b\) only, respectively. (We do not determine for which systems
\(\eta_c,c\in\mathfrak G\) a bounded \(A\) to which the given \(\eta_c\)
corresponds, actually does exist.)
\end{lemma}

\begin{proof}
If \(A\sim\{\alpha_{a,b}\}\), then \eqref{eq:5.3-beta} gives
\(U_{a_0}^{-1}AU_{a_0}\sim
\{\alpha_{aa_0^{-1},ba_0^{-1}}\}\). Now \(A\in I'\) means that \(A\)
commutes with all \(U_{a_0}\), i.e. that
\begin{equation}\tag{5.3.\(\gamma\)}\label{eq:5.3-gamma}
 \alpha_{a,b}=\alpha_{aa_0^{-1},ba_0^{-1}}
 \quad\text{for all }a_0\in\mathfrak G.
\end{equation}
Write \(\eta_c=\alpha_{c,1}\). Then \eqref{eq:5.3-gamma} with \(a_0=b\)
gives
\begin{equation}\tag{5.3.\(\delta\)}\label{eq:5.3-delta}
 \alpha_{a,b}=\eta_{ab^{-1}}.
\end{equation}
Conversely \eqref{eq:5.3-delta} implies \eqref{eq:5.3-gamma}. Thus
\eqref{eq:5.3-delta} is characteristic for \(A\in I'\). This proves our
statement concerning \(I'\).

If \(A\sim\{\alpha_{a,b}\}\), then
\(WAW\sim\{\alpha_{a^{-1},b^{-1}}\}\).\footnote{\phantomsection\label{fn:50}Use
\eqref{eq:49-alpha} and \eqref{eq:5.3-beta}. Then
\begin{align*}
 (U_{a_0}^{-1}AU_{a_0}\varphi_a,\varphi_b)
 &=(AU_{a_0}\varphi_a,U_{a_0}\varphi_b)\\
 &=(A\varphi_{aa_0^{-1}},\varphi_{ba_0^{-1}}),\\
 (WAW\varphi_a,\varphi_b)
 &=(AW\varphi_a,W\varphi_b)
  =(A\varphi_{a^{-1}},\varphi_{b^{-1}}).
\end{align*}
Hence \(A\sim\{\alpha_{a,b}\}\) implies
\(U_{a_0}^{-1}AU_{a_0}\sim
\{\alpha_{aa_0^{-1},ba_0^{-1}}\}\) and
\(WAW\sim\{\alpha_{a^{-1},b^{-1}}\}\).} Hence the automorphism
\(f\mapsto Wf\) of \(\Hilbert\) carries \eqref{eq:5.3-delta} into
\begin{equation}\tag{5.3.\(\epsilon\)}\label{eq:5.3-epsilon}
 \alpha_{a,b}=\eta_{a^{-1}b}.
\end{equation}
Since it carries \(I\) into \(J\) it takes \(I'\) into \(J'\). Therefore
\eqref{eq:5.3-epsilon} is characteristic for \(A\in J'\). This proves our
statement concerning \(J'\).
\end{proof}

\begin{lemma}\phantomsection\label{lem:5.3.3}
\(\R(I)=J'\), \(\R(J)=I'\).
\end{lemma}

\begin{proof}
Clearly \(V_{a_0}\) commutes with all \(U_{a_0}\). So \(V_{a_0}\in I'\)
and hence \(J\subseteq I'\). As \(I'\) is a ring, this implies
\(\R(J)\subseteq I'\).

If \(A\in I'\), \(B\in J'\), then
\(A\sim\{\eta_{ab^{-1}}\}\), \(B\sim\{\theta_{a^{-1}b}\}\). Direct
computation shows that \(AB=BA\) (use the matrix computation rules as
mentioned at the end of footnote~\ref{fn:49}). Thus \(A\in J''\) and hence
\(I'\subseteq J''\). Since \(1\in J\), \(J''=\R(J)\) (cf.
\cite{ref2}, p.~397). Therefore \(I'\subseteq\R(J)\).

The results of the two preceding paragraphs imply \(\R(J)=I'\). The spatial
automorphism \(f\mapsto Wf\) of \(\Hilbert\) interchanges \(I\) with \(J\)
and thus we obtain \(\R(I)=J'\) too. The proof is now complete.
\end{proof}

\begin{lemma}\phantomsection\label{lem:5.3.4}
Put \(\M=\R(I)=J'\). Then \(\M'=\R(J)\). \(I'\) is a factor if and only if
\(\mathfrak G\) fulfills this condition:
\begin{enumerate}[label=(\roman*)]
\item For every \(a\in\mathfrak G\), \(a\ne1\), the class of \(a\),
\[
 \mathfrak L_a=(c^{-1}ac;c\in\mathfrak G),
\]
is infinite.
\end{enumerate}
\end{lemma}

\begin{proof}
Clearly \(\M'=(\R(I))'=I'\).

Now consider an \(A\in\M\cdot\M'\). Put
\(A\sim\{\alpha_{a,b}\}\). \(A\in\M\cdot\M'\) implies that both
\begin{equation}\tag{5.3.\(\zeta\)}\label{eq:5.3-zeta}
 \alpha_{a,b}=\eta_{a^{-1}b}
 \qquad\text{and}\qquad
 \alpha_{a,b}=\theta_{ab^{-1}}
\end{equation}
hold. Put \(b=1\). This gives \(\eta_{a^{-1}}=\theta_a\). Thus
\eqref{eq:5.3-zeta} is equivalent to
\begin{equation}\tag{5.3.\(\eta\)}\label{eq:5.3-eta}
 \alpha_{a,b}=\eta_{a^{-1}b}=\eta_{ba^{-1}}.
\end{equation}
The second equation in \eqref{eq:5.3-eta} may be rewritten by replacing
\(a,b\) by \(b,ab\). Then it becomes \(\eta_{b^{-1}ab}=\eta_a\), i.e. it
states that \(\eta_d\) is a constant for all \(d\in\mathfrak L_a\). So
\eqref{eq:5.3-eta} is equivalent to this:
\begin{equation}\tag{5.3.\(\theta\)}\label{eq:5.3-theta}
 \alpha_{a,b}=\eta_{a^{-1}b},\qquad
 \eta_d\text{ is constant for all }d\in\mathfrak L_a.
\end{equation}

We now prove that (i) is characteristic.

Sufficiency: Assume that (i) is true. If \(a\ne1\) then
\(\mathfrak L_a\) is infinite. Hence the finiteness of
\(\sum_{d\in\mathfrak G}|\eta_d|^2\)\footnote{The finiteness of
\(\sum_{d\in\mathfrak G}|\eta_d|^2\) follows immediately from the first
formula of \eqref{eq:49-beta} with \(a=1\):
\[
 \sum_{b\in\mathfrak G}|\alpha_{1,b}|^2
 =\sum_{b\in\mathfrak G}|\eta_b|^2.
\]} together with \eqref{eq:5.3-theta} imply \(\eta_a=0\) if \(a\ne1\).
Thus
\[
 \alpha_{a,b}=\eta_{a^{-1}b}
 =\begin{cases}\eta_1,&a=b,\\0,&\text{otherwise},\end{cases}
\]
or \(A=\eta_1\cdot1\), and so \(\M\) is a factor.

Necessity: Assume that (i) is not true. Choose an \(a_0\ne1\) with
\(\mathfrak L_{a_0}\) finite. Define
\begin{equation}\tag{5.3.\(\iota\)}\label{eq:5.3-iota}
 \alpha_{a,b}=\eta_{a^{-1}b},\qquad
 \eta_d=\begin{cases}1,&d\in\mathfrak L_{a_0},\\0,&\text{otherwise}.
 \end{cases}
\end{equation}
One verifies that the \(A\sim\{\alpha_{a,b}\}\) may be expressed as
\(A=\sum_{d\in\mathfrak L_{a_0}}U_d\).\footnote{Observe that the sum can
be formed because \(\mathfrak L_{a_0}\) is finite.} Thus
\(A\in\M\cdot\M'\) by \eqref{eq:5.3-theta} and \eqref{eq:5.3-iota}.

If \(\M\) were a factor, this would necessitate \(A=\alpha\cdot1\) or
\(A\sim\{\alpha\delta_{a,b}\}\), where \(\delta_{a,b}=1\) for \(a=b\)
and \(\delta_{a,b}=0\) otherwise. Consequently
\begin{equation}\tag{5.3.\(\kappa\)}\label{eq:5.3-kappa}
 \alpha_{a,b}=\alpha\delta_{a,b}.
\end{equation}
Now \eqref{eq:5.3-iota} and \eqref{eq:5.3-kappa} give for \(a=1\) and
\(b=a_0\) (remember \(a_0\ne1\)) \(\alpha_{1,a_0}=1\) and
\(\alpha_{1,a_0}=0\) respectively. This contradiction shows that \(\M\) is
not a factor.
\end{proof}

This last result forces us to assume that \(\mathfrak G\) is infinite, since
even its classes \(\mathfrak L_a\), \(a\ne1\), are to be infinite. Groups
\(\mathfrak G\) with this property are easy to construct. We shall give some
examples at the end of \secref{sec:5.5}{\S5.5} and in
\secref{sec:5.6}{\S5.6}. Clearly no such group can be abelian.

There is, however, another circumstance worth pointing out before we go
further. Our construction is obviously nothing but an extension of
Frobenius's group numbers to infinite groups.\footnote{We treat these groups
as discrete groups and use everywhere sums not integrals. The familiar
generalization of Peter--Weyl to continuous groups and of von Neumann to
almost-periodic groups where integrals and means are used, do not show the
peculiarities of our present discussion. In those cases the infinite group
behaves almost entirely like the finite group.} The unitary space \(\Hilbert\)
and its operators are only technical devices to take care of convergence
difficulties. Both the \(U_a\) and \(V_a\) furnish representations of
\(\mathfrak G\) (cf. footnote~\ref{fn:48}), and so both \(\M=\R(I)\) and
\(\M'=\R(J)\) are the equivalents of Frobenius's group numbers. But in the
case of finite \(\mathfrak G\) this procedure does not give a factor. The group
numbers notoriously possess other central elements than the \(\alpha\cdot1\).
There are as many linearly independent ones as there are different classes
\(\mathfrak L_a\) in \(\mathfrak G\). The real meaning of our
\lemref{lem:5.3.4}{5.3.4} appears to be this: The role of all classes in
\(\mathfrak G\), when \(\mathfrak G\) is finite, is now taken over by the
finite classes.\footnote{Concerning the connection between factors and
representation theory, cf. also \cite{ref5}, pp.~118--120,
Introduction, \S3.} For finite \(\mathfrak G\) this distinction is vacuous,
but for the infinite \(\mathfrak G\) it is now seen to be decisive.

In other words: We have essentially merely extended the construction of
Frobenius group numbers to infinite \(\mathfrak G\). For finite
\(\mathfrak G\) this can never give a factor; indeed, the usefulness of the
group numbers lies in that case in the opposite direction. For infinite
\(\mathfrak G\) however, our treatment of the convergence questions makes
this construction perfectly adequate to produce factors.

We now determine the case to which these factors \(\M,\M'\) belong, and a few
other characteristics.

The condition (i) in \lemref{lem:5.3.4}{5.3.4} is assumed to be fulfilled.

\begin{lemma}\phantomsection\label{lem:5.3.5}
\begin{enumerate}[label=(\roman*)]
\item \(\M,\M'\) are coupled factors both in case (II\(_1\)).
\item In their standard normalization, the \(C\) of
\cite{ref5}, p.~182, Theorem X, is 1.
\item If \(A\in\M\), (\(A\in\M'\)) then
\(\operatorname{Tr}_{\M}(A)\), (\(\operatorname{Tr}_{\M'}(A)\)) has the
value \(\eta_1\) for the \(\eta_c\) of \lemref{lem:5.3.2}{5.3.2}.
\end{enumerate}
\end{lemma}

\begin{proof}
We proceed in a somewhat changed order.

Define for \(A\in\M\), a quantity \(T'(A)=\eta_1\) for the \(\eta_c\) of
\lemref{lem:5.3.2}{5.3.2}. Let us check for this \(T'(A)\) the properties
(i)--(vi\(^{\prime}\)) in \cite{ref6}, p.~218.

(i)--(iii), (v) loc. cit., are obvious.

If \(A\sim\{\alpha_{a,b}\}=\{\eta_{a^{-1}b}\}\),
\(B\sim\{\beta_{a,b}\}=\{\theta_{a^{-1}b}\}\), then clearly
\(AB\sim\{\gamma_{a,b}\}=\{\zeta_{a^{-1}b}\}\), where
\(\zeta_a=\sum_{c\in\mathfrak G}\eta_c\theta_{ac^{-1}}\). (Use the matrix
computation rules indicated at the end of footnote~\ref{fn:49}.) Hence
\(T'(AB)=\zeta_1=\sum_{c\in\mathfrak G}\eta_c\theta_{c^{-1}}\). This
expression is symmetric in \(\eta_a\) and \(\theta_a\). Hence
\begin{equation}\tag{5.3.\(\mu\)}\label{eq:5.3-mu}
 T'(AB)=T'(BA)=\sum_{c\in\mathfrak G}\eta_c\theta_{c^{-1}}.
\end{equation}
This proves (vi), loc. cit., and (vi\(^{\prime}\)) eod., by replacing \(A,B\) by
\(U^{-1}A,U\).

As \(A^*=\{\overline{\alpha}_{b,a}\}=\{\xi_{a^{-1}b}\}\) with
\(\xi_a=\overline{\eta}_{a^{-1}}\), \eqref{eq:5.3-mu} yields for \(B=A^*\)
\begin{equation}\tag{5.3.\(\nu\)}\label{eq:5.3-nu}
 T'(A^*A)=\sum_{c\in\mathfrak G}|\eta_c|^2.
\end{equation}
This proves (iv), (iv\(^{\prime}\)) loc. cit.

Thus all desired relations are established.

If we consider \(T'(E)\) only for projections \(E\in\M\) then the above
relations allow us to assert the validity of (ii), (iii) in
\cite{ref5}, p.~165, Definition 8.2.1. Also
\(T'(1)=1\ne0,\infty\). Hence, by \cite{ref5}, p.~170, Lemma
8.3.5, this \(T'(E)\) is a relative dimension function of \(\M\) and \(\M\)
is in a finite case.

The cases (I\(_n\)), \(n=1,2,\ldots\), are excluded, since \(\M\) is not a
finite ring. Hence \(\M\) is in the case (II\(_1\)).

Now we conclude by \cite{ref6}, p.~218, Property IV, that
\(\operatorname{Tr}_{\M}(A)=T'(A)\), i.e. that
\(\operatorname{Tr}_{\M}(A)=\eta_1\).

Thus our (i), (iii) hold for \(\M\). The spatial automorphism
\(f\mapsto Wf\) in \(\Hilbert\) shows that they hold for \(\M'\) in view of
\lemref{lem:5.3.1}{5.3.1}.

Let us now consider our (ii). Observe that by \eqref{eq:5.3-beta}
\(\mathfrak M^{\M}_{\varphi_1}\) as well as
\(\mathfrak M^{\M'}_{\varphi_1}\) (cf. \cite{ref5}, p.~143,
Definition 5.1.1) contain all \(\varphi_a,a\in\mathfrak G\). Hence both are
\(\Hilbert\). It follows that in the standard normalization,
\[
 D_{\M}(\mathfrak M^{\M'}_{\varphi_1})
 =D_{\M'}(\mathfrak M^{\M}_{\varphi_1})=1.
\]
Therefore \cite{ref5}, p.~182, Theorem X, gives \(C=1\) as
asserted.
\end{proof}

\begin{lemma}\phantomsection\label{lem:5.3.6}
Consider an \(A\in\M\) (\(A\in\M'\)) and its \(\eta_c\) by
\lemref{lem:5.3.2}{5.3.2}. Then we have
\begin{enumerate}[label=(\roman*)]
\item \(\displaystyle\Metric{A}=
\left(\sum_{c\in\mathfrak G}|\eta_c|^2\right)^{1/2}\),
\item \(\displaystyle A=\sum_{c\in\mathfrak G}\eta_cU_c\),
(\(\displaystyle A=\sum_{c\in\mathfrak G}\eta_cV_c\))
in the sense of metric convergence in \(\M\), (\(\M'\)) irrespective of the
order in which the \(c\in\mathfrak G\) are gone through.
\end{enumerate}
\end{lemma}

\begin{proof}
The spatial automorphism \(f\mapsto Wf\) of \(\Hilbert\) interchanges
\(\M\) and \(\M'\). Therefore it suffices to consider the \(A\in\M\).

Ad (i): By Equation \eqref{eq:5.3-nu} in the proof of
\lemref{lem:5.3.5}{5.3.5},
\[
 \Metric{A}^2=\operatorname{Tr}_{\M}(A^*A)=T'(A^*A)
 =\sum_{c\in\mathfrak G}|\eta_c|^2.
\]
Hence \(\Metric{A}=(\sum_{c\in\mathfrak G}|\eta_c|^2)^{1/2}\).

Ad (ii): Choose any enumeration \(a^{(1)},a^{(2)},\ldots\) of
\(\mathfrak G\). One verifies immediately that for
\(A^{(n)}=\sum_{k=1}^{n}\eta_{a^{(k)}}U_{a^{(k)}}\),
\[
 A^{(n)}\sim\{\eta^{(n)}_{a^{-1}b}\},\qquad
 \eta_c^{(n)}=\begin{cases}
 \eta_c,&c=a^{(1)},\ldots,a^{(n)},\\
 0,&\text{otherwise}.
 \end{cases}
\]
Hence
\[
 A-A^{(n)}\sim\{\eta^{(n)1}_{a^{-1}b}\},\qquad
 \eta_c^{(n)1}=\begin{cases}
 0,&c=a^{(1)},\ldots,a^{(n)},\\
 \eta_c,&\text{otherwise}.
 \end{cases}
\]
Consequently (i) above gives
\[
 \Metric{A-A^{(n)}}
 =\left(\sum_{c\in\mathfrak G}|\eta_c^{(n)1}|^2\right)^{1/2}
 =\left(\sum_{c\in\mathfrak G,\ c\ne a^{(1)},\ldots,a^{(n)}}
 |\eta_c|^2\right)^{1/2}
 =\left(\sum_{i=n+1}^{\infty}|\eta_{a^{(i)}}|^2\right)^{1/2}.
\]
This is equivalent to
\begin{equation}\tag{5.3.\(\mathrm{o}\)}\label{eq:5.3-omicron}
 \Metric{A-\sum_{k=1}^{n}\eta_{a^{(k)}}U_{a^{(k)}}}
 =\left(\sum_{k=n+1}^{\infty}|\eta_{a^{(k)}}|^2\right)^{1/2}.
\end{equation}
Now \(\sum_{k=1}^{\infty}|\eta_{a^{(k)}}|^2
=\sum_{c\in\mathfrak G}|\eta_c|^2\) is finite by (i) above. Hence the
right hand side of \eqref{eq:5.3-omicron} tends to 0 as \(n\to\infty\). This
means that \(\sum_{k=1}^{\infty}\eta_{a^{(k)}}U_{a^{(k)}}\) converges
metrically to \(A\).

Since this holds for any enumeration \(a^{(1)},a^{(2)},\ldots\) of
\(\mathfrak G\), our assertion is established.
\end{proof}

\PaperSection{sec:5.4}{Relation with the earlier construction.}
As we pointed out, the construction of \secref{sec:5.3}{\S5.3} is closely
related to that of \cite{ref5}, pp.~192--209, Part IV. The group
\(\mathfrak G\) plays the same role in both cases, and the difference consists
in the disappearance of the space \(S\).

This space played an important part in the construction referred to above:
Every \(a\in\mathfrak G\) had to induce a mapping \(x\mapsto xa\) in \(S\)
(cf. loc. cit., p.~195, Definition 12.1.5). The absence of \(S\) from our
present construction may also be interpreted by saying that \(S\) is now a
one-element set: \(S=(x_0)\) and that for every \(a\in\mathfrak G\),
\(x_0a=x_0\).

Now the difference between our two constructions can be expressed as follows:
In \cite{ref5}, \(\mathfrak G\) was entirely unrestricted, but \(S\)
had to obey (among other things) (iii) in Definition 12.1.5, loc. cit., p.~195:
for \(a\ne1\) and \(x\in S\) always \(xa\ne x\).\footnote{Except for an
\(x\) set of measure zero.} In \secref{sec:5.3}{\S5.3}, \(\mathfrak G\) had
to fulfill (i) of \lemref{lem:5.3.4}{5.3.4}, but the (iii) referred to was most
flagrantly violated.\footnote{We have \(x_0a=x_0\) and
\(\mu((x_0))=\mu(S)\ne0\).}\footnote{There is no difference between the two
constructions as regards the other conditions of Definition 12.1.5, loc. cit.,
p.~195. In particular, \(\mathfrak G\) is ergodic in our present construction,
since \(S=(x_0)\) is a one-element set.} Thus our two procedures correspond
to two different ways of securing the factor character of \(\M\) in what is
fundamentally the same construction. In one case the whole burden of the
necessary restrictions was thrown on \(S\), in the other case, on
\(\mathfrak G\).

It would be possible to use a more general arrangement of which both these
procedures are special cases. The necessary restriction would then be of a
mixed type, affecting the structure of \(S\) and \(\mathfrak G\).

We shall not consider this question in more detail at this time.

\PaperSection{sec:5.5}{A connection between the two constructions.}
A second connection between our two constructions which deserves some
attention can be obtained as follows.

Let an arbitrary countably infinite group \(\mathfrak G\) be given. We shall
try to use it for the construction of \cite{ref5} referred to above,
that is to find an \(S\) which fulfills all preliminary conditions of that
construction in conjunction with the given \(\mathfrak G\).

This construction runs as follows:\footnote{The construction (I)--(V) which
follows is in many respects analogous to the construction (I)--(VII) in
\cite{ref7}, pp.~72--77. The main differences are these: We use
\(\mathfrak G\) and not \(\mathfrak S\) as the group of mappings in \(S\),
and the latter parts of the two constructions pursue different aims.}

(I) Let \(S\) be the set of all systems \(x=[\alpha_a;a\in\mathfrak G]\)
where each \(\alpha_a=0,1\). Let \(\mathfrak S\) be the set of those
\(x=[\alpha_a;a\in\mathfrak G]\) for which \(\alpha_a=0\) except for a
finite number of \(a\)'s.

Define in \(S\): If \(x=[\alpha_a;a\in\mathfrak G]\),
\(y=[\beta_a;a\in\mathfrak G]\), then
\(x\oplus y=[\gamma_a;a\in\mathfrak G]\), where
\(\gamma_a=\alpha_a+\beta_a\pmod2\), \(\gamma_a=0\) or 1. \(S\) is
clearly a (commutative) group with the ``unit''
\(0=[0;a\in\mathfrak G]\), and \(\mathfrak S\) is an enumerably infinite
subgroup of \(S\).

Define further in \(S\): If \(x=[\alpha_a;a\in\mathfrak G]\) and
\(a_0\in\mathfrak G\), then
\(xa_0=[\alpha_{a_0a};a\in\mathfrak G]\). Then the mappings
\(x\mapsto xa_0\) represent the group \(\mathfrak G\) by permutations of
\(S\), all of which carry \(\mathfrak S\) into itself.

Now choose an arbitrary but fixed enumeration
\(a^{(1)},a^{(2)},\ldots\) of \(\mathfrak G\). For \(S\) we use the mapping
\[
 \Xi:\quad x=[\alpha_a;a\in\mathfrak G]\longmapsto
 \xi(x)=\sum_{m=1}^{\infty}\frac{\alpha_{a^{(m)}}}{2^m}
\]
of \(S\) on the numerical interval \(0\leq\xi\leq1\). Except for the
\(\Xi\) image of \(\mathfrak S\), the set of all dyadically rational numbers,
which is a set of Lebesgue measure 0, this mapping is one-to-one. So the common
(exterior) Lebesgue measure in \(0\leq\xi\leq1\) is mapped by the inverse of
\(\Xi\) on a Lebesgue measure in \(S\) with all its usual properties.

We shall now consider \(\mathfrak G\) as the ``group'' and \(S\) (with the
``mappings'' \(x\mapsto xa_0\) for \(x\in S,a_0\in\mathfrak G\)) as the
``space'' described in \cite{ref5}, pp.~192--195. In the sense of
\cite{ref5}, p.~195, Definition 12.1.5, \(\mathfrak G\) is an
\(m\)-group and ergodic in \(S\).

First we show the \(m\)-group character. Ad (i), loc. cit. This is concerned
with the preservation of outer measure of the measure determining sets
\(\{T_i\}\). In view of the definition of measure in \(S\) given in the
previous paragraph, we may take for a set of \(T_i\)'s the set of images, under
the inverse of \(\Xi\), of \(\xi\)-intervals
\(k/2^m\leq\xi<(k+1)/2^m\). An image of a \(\xi\)-interval
\(k/2^m\leq\xi<(k+1)/2^m\) is given by specifying the values of
\(\alpha_{a^{(1)}},\ldots,\alpha_{a^{(m)}}\). Hence it is mapped by
\(x\mapsto xa_0\) onto a set, in which the values of
\(\alpha_{a_0^{-1}a^{(1)}},\ldots,\alpha_{a_0^{-1}a^{(m)}}\) are specified.
Choose \(n\) so that
\(a_0^{-1}a^{(1)},\ldots,a_0^{-1}a^{(m)}\) occur among the
\(a^{(1)},\ldots,a^{(n)}\). Then the above set is the sum of \(2^{n-m}\)
sets, each of which is defined by specifying the values of
\(\alpha_{a^{(1)}},\ldots,\alpha_{a^{(n)}}\), i.e. each of which corresponds
to an interval \(l/2^n\leq\xi<(l+1)/2^n\). Thus the set corresponding to
the \(\xi\)-interval \(k/2^m\leq\xi<(k+1)/2^m\), having measure
\(1/2^m\), is mapped on a set of measure
\(2^{n-m}/2^n=2^{-m}\). Thus the measure of each such is conserved. Ad (ii),
which states that \((xa_0)b_0=x(a_0b_0)\), can be verified by a direct
computation.\footnote{Let \(x=[\alpha_a;a\in\mathfrak G]\),
\(xa_0=[\beta_a;a\in\mathfrak G]\),
\((xa_0)b_0=[\gamma_a;a\in\mathfrak G]\). Then
\(\beta_a=\alpha_{a_0a}\) and
\(\gamma_a=\beta_{b_0a}=\alpha_{a_0(b_0a)}=\alpha_{(a_0b_0)a}\). Hence
\((xa_0)b_0=[\alpha_{(a_0b_0)a};a\in\mathfrak G]=x(a_0b_0)\).} Ad (iii)
states that if \(a_0\ne1\), \(xa_0=x\) holds only for a set of measure zero.
For \(x=[\alpha_a;a\in\mathfrak G]\), \(xa_0=x\) is equivalent to
\(\alpha_{a_0a}=\alpha_a\) for all \(a\). But since
\(a_0\in\mathfrak G\), \(a_0\ne1\), we can construct an infinite sequence
\(a_1,a_2,\ldots\) such that the
\(a_1,a_2,\ldots,a_0a_1,a_0a_2,\ldots\) are all different.\footnote{\phantomsection\label{fn:60}If the
\(a_0,a_0^2,\ldots\) are all distinct, let
\(a_i=a_0^{2i-1}\). The result
readily follows in this case. Otherwise, \(a_0\) is of finite order,
\(p=2,3,\ldots\) (\(p\ne1\) since \(a_0\ne1\)). For each \(a\) we can
define \(S_a\) as the set of \(b\in\mathfrak G\) for which \(b=a_0^ka\) for
some \(k=0,1,\ldots,p-1\). It is easily seen how we can choose a sequence
\(a_1,a_2,\ldots\) such that the \(S_{a_i}\) are mutually exclusive. For this
sequence we have that the \(a_1,a_2,\ldots,a_0a_1,a_0a_2,\ldots\) are all
different. This readily follows from the facts that the \(S_{a_i}\) are
mutually exclusive and \(a_0\ne1\).} For any \(m\), let \(n=n(m)\) be such
that \(a^{(1)},\ldots,a^{(n)}\) includes
\(a_1,\ldots,a_m,a_0a_1,\ldots,a_0a_m\). Then the \(x\)-set \(S^{(m)}\) for
which \(\alpha_{a_0a_i}=\alpha_{a_i}\), \(i=1,\ldots,m\), consists of
\(2^{n-m}\) sets \(I^{(n)}\) in each of which the
\(\alpha_{a^{(1)}},\ldots,\alpha_{a^{(n)}}\) have specified values. Such a
set \(I^{(n)}\) is the image of a \(\xi\)-interval of measure \(2^{-n}\), and
hence \(S^{(m)}\) has measure \(2^{n-m}\cdot2^{-n}=2^{-m}\). Now the
\(x\)-set for which \(\alpha_{a_0a}=\alpha_a\) for all
\(a\in\mathfrak G\) is in \(S^{(m)}\) for every \(m\). Thus it has measure
zero. Since this is also the set for which \(xa_0=x\), we have shown Ad (iii).

The ergodicity will be established in (III) below.

(II) Form for these \(S,\mathfrak G\) the spaces \(\Hilbert_S\) and
\(\Hilbert_{\mathfrak G S}\) of all complex-valued functions \(f(x)\),
respectively \(F(x,a)\), \(x\in S,a\in\mathfrak G\), which are measurable
in \(x\) and with a finite
\[
 \int_S|f(x)|^2\,dx,
 \qquad\text{respectively}\qquad
 \sum_{a\in\mathfrak G}\int_S|F(x,a)|^2\,dx.
\]
Thus
\[
 (f,g)=\int_Sf(x)\overline{g(x)}\,dx,
 \qquad
 (F,G)=\sum_{a\in\mathfrak G}\int_SF(x,a)\overline{G(x,a)}\,dx
\]
(cf. \cite{ref5}, p.~194). Form the bounded operators
\[
 U_{a_0}f(x)=f(xa_0),\qquad
 \overline U_{a_0}F(x,a)=F(xa_0,aa_0)
 \qquad(a_0\in\mathfrak G),
\]
\[
 \overline L_{\varphi(x)}F(x,a)=\varphi(x)F(x,a)
 \qquad(\varphi(x)\text{ bounded and measurable})
\]
(loc. cit., pp.~196, 198--199) and the ring \(\M\) in
\(\Hilbert_{\mathfrak G S}\) which the \(\overline U_{a_0}\),
\(\overline L_{\varphi(x)}\) generate (loc. cit., p.~200):
\begin{equation}\tag{5.5.\(\alpha\)}\label{eq:5.5-alpha}
 \M=\R(\overline U_{a_0},\overline L_{\varphi(x)};
 a_0\in\mathfrak G,\ \varphi(x)\text{ bounded and measurable}).
\end{equation}

Before we go on we form the following system of functions in \(\Hilbert_S\):
\[
 \omega_{\bar x}(x)=(-1)^{\sum_{a\in\mathfrak G}\bar\alpha_a\alpha_a}
 \quad
 (\bar x=[\bar\alpha_a;a\in\mathfrak G]\in\mathfrak S,
 \ x=[\alpha_a;a\in\mathfrak G]\in S).
\]
% Editorial note E40: The printed source jumps directly from footnote 60 to
% footnote 62; the counter is advanced here to preserve the source numbering.
\addtocounter{footnote}{1}
\footnote{Owing to \(\bar x=[\bar\alpha_a;a\in\mathfrak G]\in\mathfrak S\)
we have \(\bar\alpha_a\ne0\) only for a finite set of \(a\)'s. Hence the sum
\(\sum_{a\in\mathfrak G}\bar\alpha_a\alpha_a\) is really finite.}
It is easy to verify that these functions (for all \(\bar x\in\mathfrak S\))
form a complete normalized orthogonal system in \(\Hilbert_S\).\footnote{In
this particular case the group nature of \(\mathfrak G\) is irrelevant. Hence
we can replace \(\mathfrak G\) by any other countably infinite set. We replace
the \(a\in\mathfrak G\) by \(a=1,2,\ldots\). After this substitution our
\(\omega_{\bar x}(x)\) become the well-known complete normalized orthogonal
set of Walsh--Rademacher functions. It was also considered in
\cite{ref7}, p.~74, under the name of \(\omega_a(x)\).}
Consequently the functions
\[
 F_{\bar x,\bar a}(x,a)=\omega_{\bar x}(x)\delta_{a,\bar a}
 \qquad(x\in S,\ \bar a\in\mathfrak G)
\]
form a complete normalized orthogonal system in
\(\Hilbert_{\mathfrak G S}\).

A familiar argument shows that the \(\varphi(x)\) in \eqref{eq:5.5-alpha}
can be restricted considerably without changing the ring \(\M\). Indeed:

First it suffices, for reasons of continuity, to consider those \(\varphi(x)\)
only which are finite linear aggregates of the
\(\omega_{\bar x}(x)\).\footnote{Cf. the argument of
\cite{ref7}, pp.~73--74 (II).} Second, we may now restrict the
\(\varphi(x)\) to the \(\omega_{\bar x}(x)\) themselves. So we have
\begin{equation}\tag{5.5.\(\beta\)}\label{eq:5.5-beta}
 \M=\R(\overline U_{a_0},\overline L_{\omega_{\bar x}(x)};
 a_0\in\mathfrak G,\ \bar x\in\mathfrak S).
\end{equation}
And as 1 occurs among the \(\overline U_{a_0}\) (for \(a_0=1\)) and among
the \(\overline L_{\omega_{\bar x}(x)}\) (for \(\bar x=0\)), we can equally
write:
\begin{equation}\tag{5.5.\(\gamma\)}\label{eq:5.5-gamma}
 \M=\R(\overline U_{a_0}\overline L_{\omega_{\bar x}(x)};
 a_0\in\mathfrak G,\ \bar x\in\mathfrak S).
\end{equation}

(III) Clearly
\[
 U_{a_0}\omega_{\bar x}(x)=\omega_{\bar x}(xa_0)
 =\omega_{\bar x a_0^{-1}}(x).
\]
Now consider any \(f\in\Hilbert_S\) with \(U_{a_0}f=f\) for all
\(a_0\in\mathfrak G\). Use the orthogonal expansion
\(f(x)=\sum_{\bar x\in\mathfrak S}\lambda_{\bar x}omega_{\bar x}(x)\).
Then the above formula gives
\[
 f=\sum_{\bar x\in\mathfrak S}\lambda_{\bar x}\omega_{\bar x},
 \qquad
 U_{a_0}f
 =\sum_{\bar x\in\mathfrak S}\lambda_{\bar x}
     \omega_{\bar x a_0^{-1}}
 =\sum_{\bar x\in\mathfrak S}\lambda_{\bar x a_0}\omega_{\bar x}.
\]
Hence \(\lambda_{\bar x}=\lambda_{\bar x a_0}\). As
\(\OpNorm{f}^2=\sum_{\bar x\in\mathfrak S}|\lambda_{\bar x}|^2\) is finite,
this implies that \(\lambda_{\bar x}=0\) if there are infinitely many distinct
\(\bar x a_0\), \(a_0\in\mathfrak G\). Now this is the case for all
\(\bar x\in\mathfrak S\), \(\bar x\ne0\).\footnote{\phantomsection\label{fn:65}We prove this as follows:
Put \(\bar x=[\bar\alpha_a;a\in\mathfrak G]\). Since \(\bar x\ne0\) there
is an \(a^*\) such that \(\bar\alpha_{a^*}\ne0\). Now as in
footnote~\ref{fn:60}, we
can construct an infinite sequence of the \(a\)'s \(\in\mathfrak G\) such
that \(a_1a^*,a_2a^*,\ldots\) are all different. Then
\(\bar x a_i^{-1}=[\bar\alpha_{a_i^{-1}a};a\in\mathfrak G]\) has for its
\(a_ia^*\) component \(\bar\alpha_{a^*}=1\). Thus if we consider the
\(\bar x a_1^{-1},\bar x a_2^{-1},\ldots\), we find that there is an infinite
number of \(a\)'s \(\in\mathfrak G\) such that the \(a\) component is not
zero for at least one of the \(\bar x a_i^{-1}\). Since each
\(\bar x a_i^{-1}\) can have at most a finite number of components different
from zero, there must be an infinite number of distinct \(\bar x a_i^{-1}\).}
So \(\bar x\ne0\) implies \(\lambda_{\bar x}=0\). Thus
\(f(x)=\lambda_0\omega_0(x)=\lambda_0\).

So \(U_{a_0}f=f\) (for all \(a_0\in\mathfrak G\)) implies that \(f\) is a
constant. Hence \(\mathfrak G\) is ergodic.

(IV) Clearly
\[
 \omega_{\bar x}(x)\omega_{\bar y}(x)=
 \omega_{\bar x\oplus\bar y}(x),
\]
\[
 \overline L_{\omega_{\bar x}(x)}F_{\bar y,\bar b}(x,a)
 =\omega_{\bar x}(x)F_{\bar y,\bar b}(x,a)
 =F_{\bar x\oplus\bar y,\bar b}(x,a).
\]
Next
\[
 \overline U_{\bar a}F_{\bar y,\bar b}(x,a)
 =F_{\bar y,\bar b}(x\bar a,a\bar a)
 =F_{\bar y\bar a^{-1},\bar b\bar a^{-1}}(x,a).
\]
Combining these gives
\begin{equation}\tag{5.5.\(\epsilon\)}\label{eq:5.5-epsilon}
 \overline U_{\bar a}\overline L_{\omega_{\bar x}(x)}
 F_{\bar y,\bar b}(x,a)
 =F_{(\bar x\oplus\bar y)\bar a^{-1},\bar b\bar a^{-1}}(x,a).
\end{equation}

Now treat the index \(\bar y,\bar b\) of \(F_{\bar y,\bar b}\) as one pair.
Make these pairs into a group \(\{\mathfrak S,\mathfrak G\}\) by defining
\begin{equation}\tag{5.5.\(\zeta\)}\label{eq:5.5-zeta}
 \{\bar y,\bar b\}\{\bar z,\bar c\}
 =\{\bar y\bar c\oplus\bar z,\bar b\bar c\}.
\end{equation}
As \(\bar x\oplus\bar x=0\), \eqref{eq:5.5-zeta} implies that
\[
 \{\bar y,\bar b\}\{\bar x,\bar a\}^{-1}
 =\{(\bar x\oplus\bar y)\bar a^{-1},\bar b\bar a^{-1}\}.
\]
Hence \eqref{eq:5.5-epsilon} becomes
\begin{equation}\tag{5.5.\(\eta\)}\label{eq:5.5-eta}
 \overline U_{\bar a}\overline L_{\omega_{\bar x}(x)}
 F_{\bar y,\bar b}(x,a)
 =F_{\{\bar y,\bar b\}\{\bar x,\bar a\}^{-1}}(x,a).
\end{equation}

(V) Apply the construction of \secref{sec:5.3}{\S5.3} to the group
\(\{\mathfrak S,\mathfrak G\}\) of \eqref{eq:5.5-zeta} above, instead of
its \(\mathfrak G\). This is possible since
\(\{\mathfrak S,\mathfrak G\}\) fulfills (i) in
\lemref{lem:5.3.4}{5.3.4}; i.e. if
\(\{\bar x,\bar a\}\ne\{0,1\}\) then there exist infinitely many different
\(\{\bar y,\bar b\}^{-1}\{\bar x,\bar a\}\{\bar y,\bar b\}\). Indeed, if
\(\bar x\ne0\), then form
\[
 \{0,\bar b\}^{-1}\{\bar x,\bar a\}\{0,\bar b\}
 =\{\bar x\bar b,\bar b^{-1}\bar a\bar b\}.
\]
There are infinitely many distinct \(\bar x\bar b\) for
\(\bar b\in\mathfrak G\). Cf. footnote~\ref{fn:65}. Thus our statement is true for
\(\bar x\ne0\). If however \(\bar x=0\), then necessarily \(\bar a\ne1\).
Form
\[
 \{\bar y,1\}^{-1}\{0,\bar a\}\{\bar y,1\}
 =\{\bar y\bar a\oplus\bar y,\bar a\}.
\]
There exist infinitely many different \(\bar y\bar a\oplus\bar y\).\footnote{
Put \(\bar y_b=[\delta_{a,b};a\in\mathfrak G]\in\mathfrak S\). As
\(\bar a\ne1\),
\(\bar y_b\bar a\oplus\bar y_b=\bar y_c\bar a\oplus\bar y_c\) means either
\(b=c\) or \(\bar ab=c,\bar ac=b\). In other words: If \(b\) runs over all
\(\mathfrak G\) there can never coincide more than two
\(\bar y_b\bar a\oplus\bar y_b\). Now \(\mathfrak G\) is infinite. Hence
there are infinitely many different \(\bar y_b\bar a\oplus\bar y_b\).}

Now write \(\Hilbert_I,\M_I\) for its \(\Hilbert,\M\). We have
\begin{equation}\tag{5.5.\(\theta\)}\label{eq:5.5-theta}
 \M_I=\R(U_{\{\bar x,\bar a\}};
 \bar x\in\mathfrak S,\ \bar a\in\mathfrak G)
\end{equation}
and by \eqref{eq:5.3-beta}
\begin{equation}\tag{5.5.\(\iota\)}\label{eq:5.5-iota}
 U_{\{\bar x,\bar a\}}\varphi_{\{\bar y,\bar b\}}
 =\varphi_{\{\bar y,\bar b\}\{\bar x,\bar a\}^{-1}}.
\end{equation}

We can map \(\Hilbert_{\mathfrak G S}\) isomorphically on \(\Hilbert_I\) by
carrying the complete orthonormalized system \(F_{\{\bar y,\bar b\}}\) of the
former into the set of \(\varphi_{\{\bar y,\bar b\}}\) of the latter. Then
comparison of \eqref{eq:5.5-iota} and \eqref{eq:5.5-eta} shows that this
isomorphism carries
\(\overline U_{\bar a}\overline L_{\omega_{\bar x}(x)}\) into
\(U_{\{\bar x,\bar a\}}\). Hence \eqref{eq:5.5-gamma} and
\eqref{eq:5.5-theta} show that it carries \(\M\) into \(\M_I\). Thus \(\M\)
and \(\M_I\) are spatially isomorphic.

Summing up:

\begin{lemma}\phantomsection\label{lem:5.5.1}
Given any countably infinite group \(\mathfrak G\), define
\(S,\mathfrak S\) as described in (I) above, and the group
\(\mathfrak G_I=\{\mathfrak S,\mathfrak G\}\) as described by
\eqref{eq:5.5-zeta} (IV) above. Then the construction of
\cite{ref5}, pp.~192--209, Part IV, can be applied to
\(\mathfrak G\) and \(S\), and our construction of
\secref{sec:5.3}{\S5.3} to \(\{\mathfrak S,\mathfrak G\}\) (in place of its
\(\mathfrak G\)).

These two constructions produce spatially isomorphic factors \(\M\).
\end{lemma}

The mappings
\begin{equation}\tag{5.5.\(\kappa\)}\label{eq:5.5-kappa}
 x\{\bar y,\bar b\}=x\bar b\oplus\bar y
\end{equation}
provide a convenient representation of the group
\(\{\mathfrak S,\mathfrak G\}\) of \eqref{eq:5.5-zeta} in (V). If we
specify \(x\in S\), then \eqref{eq:5.5-kappa} appears as a permutation of
\(S\). But as it maps \(\mathfrak S\) on itself, we can also prescribe
\(x\in\mathfrak S\) and view \eqref{eq:5.5-kappa} as a permutation of
\(\mathfrak S\).

\eqref{eq:5.5-kappa} makes it clear that
\(\{\mathfrak S,\mathfrak G\}\) is the simplest combination of
\(\mathfrak S\) and \(\mathfrak G\) apart from the ``direct group product''.

Our above result can also be interpreted in this way:

We start with an arbitrary (countably infinite) group \(\mathfrak G\) and
should like to apply the construction of \secref{sec:5.3}{\S5.3}. Since
\(\mathfrak G\) may not fulfill the condition (i) of
\lemref{lem:5.3.4}{5.3.4}, this cannot always be done directly. We therefore
propose to expand \(\mathfrak G\) to a group \(\mathfrak G_I\) which fulfills
that condition. \eqref{eq:5.5-zeta} provides a very simple generally valid
mechanism to do just that. This has the further interesting consequence that
within \(\M\), the set of group numbers for any \(\mathfrak G\), the metric
closure of a subalgebra is equivalent to the topological closure for any of the
ring topologies. For we extend \(\M\) to \(\M_I\) as above, and
\thmref{thm:I}{I} (of \secref{sec:1.6}{\S1.6} above) shows that our statement
is true in the extended space \(\Hilbert_I\). Since any subalgebra of \(\M\)
will be reduced by the original space \(\Hilbert\), topological closure in
\(\Hilbert_I\) will yield topological closure. Similarly the discussion of
Chapter I may be carried through on the assumption that \(\M\) is a group
algebra.

\PaperSection{sec:5.6}{Approximate finiteness of the new examples.}
The analogue of \lemref{lem:5.2.2}{5.2.2} is true for the construction of
\secref{sec:5.3}{\S5.3} too, and this time it is even easier to prove.

\begin{lemma}\phantomsection\label{lem:5.6.1}
The ring \(\M\) of \secref{sec:5.3}{\S5.3} is approximately finite if
\(\mathfrak G\) possesses the property (i) of \lemref{lem:5.2.2}{5.2.2}
(besides the property (i) of \lemref{lem:5.3.4}{5.3.4}).
\end{lemma}

\begin{proof}
Form the \(\mathfrak G_1,\mathfrak G_2,\ldots\) of (i) in
\lemref{lem:5.2.2}{5.2.2}. Now
\[
 \M=\R(U_a;a\in\mathfrak G).
\]
Hence
\[
 \M=\R(\N_n;n=1,2,\ldots)
\]
with
\[
 \N_n=\R(U_a;a\in\mathfrak G_n).
\]
Thus \(\N_1\subseteq\N_2\subseteq\cdots\) and \(\N_n\) is of finite order
because \(\mathfrak G_n\) is finite. Therefore \(\M\) is approximately finite
by \defref{def:4.6.1}{4.6.1} (cf. the author's note following
\eqref{eq:5.2-zeta}).
\end{proof}

Observe that if \(\mathfrak G\) fulfills (i) in \lemref{lem:5.2.2}{5.2.2}
(but not necessarily (i) in \lemref{lem:5.3.4}{5.3.4}), then the group
\(\mathfrak G_I=\{\mathfrak S,\mathfrak G\}\) of
\lemref{lem:5.5.1}{5.5.1} possesses that property too. Put
\(\mathfrak G=\mathfrak G_1+\mathfrak G_2+\cdots\),
\(\mathfrak G_1\subseteq\mathfrak G_2\subseteq\cdots\), all
\(\mathfrak G_n\) finite. Let \(\mathfrak S_n\) be the set of all
\(\bar x=[\bar\alpha_a;a\in\mathfrak G]\) with \(\bar\alpha_a=0\), when
\(a\) is not in \(\mathfrak G_n\). Then
\(\mathfrak S_1\subseteq\mathfrak S_2\subseteq\cdots\) and all the
\(\mathfrak S_n\) are finite. So
\[
 \{\mathfrak S,\mathfrak G\}
 =\{\mathfrak S_1,\mathfrak G_1\}
  +\{\mathfrak S_2,\mathfrak G_2\}+\cdots,
 \qquad
 \{\mathfrak S_1,\mathfrak G_1\}
 \subseteq\{\mathfrak S_2,\mathfrak G_2\}\subseteq\cdots,
\]
and all the \(\{\mathfrak S_n,\mathfrak G_n\}\) are finite.

Thus it is easy to form groups \(\mathfrak G\) which fulfill both (i) in
\lemref{lem:5.2.2}{5.2.2} and (i) in \lemref{lem:5.3.4}{5.3.4}. The following
example is more direct.

\begin{lemma}\phantomsection\label{lem:5.6.2}
The group \(\mathfrak G\) of all those permutations of
\((1,2,\ldots)\) which move only a finite number of elements, fulfills both
(i) in \lemref{lem:5.2.2}{5.2.2} and (i) in
\lemref{lem:5.3.4}{5.3.4}.
\end{lemma}

\begin{proof}
Ad (i) in \lemref{lem:5.3.4}{5.3.4}: Given an
\(a\in\mathfrak G\), \(a\ne1\), choose \(n\) so that \(a\) moves no number
other than \(1,\ldots,n\). Assume that \(a\) does move \(i\) (which is
therefore \(=1,\ldots,n\)). Let \(c_j\) be the transposition of \(i\) and
\(j\), where \(j=n+1,n+2,\ldots\). Then the only
\(k=n+1,n+2,\ldots\) moved by \(c_j^{-1}ac_j\) is \(k=j\). Thus the
\(c_j^{-1}ac_j\), \(j=n+1,n+2,\ldots\), are pairwise different and so
\(\mathfrak L_a\) is infinite.

Ad (i) in \lemref{lem:5.2.2}{5.2.2}: Let \(\mathfrak G_n\) be the set of
those permutations which move no number other than \(1,\ldots,n\). Then
\(\mathfrak G=\mathfrak G_1+\mathfrak G_2+\cdots\),
\(\mathfrak G_1\subseteq\mathfrak G_2\subseteq\cdots\), and each
\(\mathfrak G_n\) is a finite subset of \(\mathfrak G\).
\end{proof}

The analysis of the invariants (1), (2) of \secref{sec:4.8}{\S4.8}, gives no
new information. As to (1), all examples of \cite{ref5} as well as
those of \secref{sec:5.3}{\S5.3} have
\(\overline{\M}^{c}=\overline{\M}\). Indeed: Neither of these definitions
ever mentions any specific non-real (complex) number. Hence
\(\Hilbert_c,\M_c\) coincide with \(\Hilbert,\M\), so that
\(\overline{\M}^{c}=\overline{\M_c}=\overline{\M}\) (cf.
\secref{sec:2.3}{\S2.3}).

Thus we possess no examples of a factor \(\M\) in the case (II\(_1\)) with
\(\overline{\M}^{c}\ne\overline{\M}\). Possibly a generalization of the
construction of \secref{sec:5.3}{5.3}, with the introduction of complex
numerical factors in (i), (ii) of \lemref{lem:5.3.1}{5.3.1} (and
\eqref{eq:5.3-beta} after that lemma) might achieve this.\footnote{
Corresponding to the technique of ``multipliers'' in forming skew fields over
a given center in abstract algebra.} But we have not obtained so far any
conclusive results in this respect.

As to (2), \(\mathfrak G(\overline{\M})\), we know nothing beyond
\thmref{thm:XV}{XV}.

\PaperChapter{6}{chapter:VI}{Non-Approximately Finite Factors}

\PaperSection{sec:6.1}{The property $\Gamma$ and approximate finiteness}[The property Gamma and approximate finiteness]
As pointed out at the end of \secref{sec:4.8}{4.8} and again at the end of
\secref{sec:5.6}{5.6}, we have not succeeded so far to establish the not
approximately finite character of any factor with the help of the invariants
(1), (2) of \secref{sec:4.8}{4.8}. We shall now prove the existence of not
approximately finite factors, but it is necessary to introduce a new invariant
in order to achieve this aim.

Throughout what follows, $\M$ is again a factor in case $(\mathrm{II}_1)$.

\begin{definition}\phantomsection\label{def:6.1.1}
$\M$ has the property $\Gamma$ if this is true:

Given any system $A_1,\ldots,A_m\in\M$ and any $\epsilon>0$ there exists a
unitary
\[
  U=U(A_1,\ldots,A_m,\epsilon)\in\M
\]
with
\[
  \operatorname{Tr}_{\M}(U)=0
\]
and
% Editorial note E62: the source prints A_k in the defining inequality while
% quantifying h; A_h has been restored.
\[
  \Metric{U^{-1}A_hU-A_h}<\epsilon
  \qquad\text{for }h=1,\ldots,m.
\]
\footnote{The requirement $\operatorname{Tr}_{\M}(U)=0$ is necessary,
since otherwise we could choose $U=1$. Outright commutativity of $U$ with
$A_1,\ldots,A_m$ (instead of the last inequality) would not do for the
following reason: If $\M=\mathbf R(A_1,\ldots,A_m)$ then this commutativity
would mean $U\in\M'$. Hence $U\in\M\cap\M'$, $U=\alpha\cdot1$. Since
$\operatorname{Tr}_{\M}(U)=0$ this would mean $U=0$, contradicting the
unitarity of $U$. Hence no $\M=\mathbf R(A_1,\ldots,A_m)$ could then have
this property. Now it can be shown that an approximately finite
$\M=\mathbf R(A_1,A_2)$ for two suitable $A_1,A_2$. Hence an approximately
finite $\M$ would not have this property---but we need Lemma 6.1.2.}
\footnote{This property could be further subdivided, according to the values
of $m$ and of expressions like
$\epsilon/\operatorname{Max}(\OpNorm{A_1},\ldots,\OpNorm{A_m})$. We could
also require $A_1,\ldots,A_m$ to be unitary, etc. For our present purpose,
however, these refinements are not needed.}
\end{definition}

This property $\Gamma$ is clearly a purely algebraic property---i.e. one
concerning the algebraic type $\overline{\M}$ only.

Observe first this
\begin{lemma}\phantomsection\label{lem:6.1.1}
Let a group $\mathfrak G$ fulfilling (i) in \lemref{lem:5.3.4}{5.3.4} be
given, which possesses this property:
\begin{enumerate}[label=(i),leftmargin=3em]
\item Given any system $a^{(1)},\ldots,a^{(n)}\in\mathfrak G$ there exists a
$c_0\in\mathfrak G$, $c_0\ne1$ which commutes with
$a^{(1)},\ldots,a^{(n)}$.
\end{enumerate}
Then the $\M$ of \secref{sec:5.3}{5.3} possesses the property $\Gamma$.
\end{lemma}

\begin{proof}
Let $A_1,\ldots,A_m\in\M$ and $\epsilon>0$ be given as indicated in
\defref{def:6.1.1}{6.1.1}. Following (ii) in \lemref{lem:5.3.6}{5.3.6}, we
can choose $a^{(1)},\ldots,a^{(n)}\in\mathfrak G$ and
$\zeta_{h,1},\ldots,\zeta_{h,n}$ (these are the
$\eta_{a^{(1)}},\ldots,\eta_{a^{(n)}}$ mentioned there for $A=A_h$, for
$h=1,\ldots,m$) so that
\begin{equation}
  \Metric{A_h-\sum_{i=1}^{n}\zeta_{h,i}U_{a^{(i)}}}<\epsilon/2
  \qquad\text{for }h=1,\ldots,m.
  \tag{6.1.\ensuremath{\alpha}}\label{eq:6.1.alpha}
\end{equation}
Now apply our assumption (i) to these $a^{(1)},\ldots,a^{(n)}$, obtaining a
$c_0\in\mathfrak G$, $c_0\ne1$, which commutes with them.

Then $U_{c_0}$ and $U_{a^{(i)}}$ commute
(cf. \hyperref[fn:48]{footnote~48}) and so
\[
  U_{c_0}^{-1}\left(A_h-\sum_{i=1}^{n}\zeta_{h,i}U_{a^{(i)}}\right)U_{c_0}
  =U_{c_0}^{-1}A_hU_{c_0}-\sum_{i=1}^{n}\zeta_{h,i}U_{a^{(i)}}.
\]
% Editorial note E60: the source omits the subscript h from A on the right here
% and in (6.1.beta); it has been restored from the surrounding argument.
Hence \eqref{eq:6.1.alpha} yields
\begin{equation}
  \Metric{U_{c_0}^{-1}A_hU_{c_0}
  -\sum_{i=1}^{n}\zeta_{h,i}U_{a^{(i)}}}<\epsilon/2
  \tag{6.1.\ensuremath{\beta}}\label{eq:6.1.beta}
\end{equation}
and \eqref{eq:6.1.alpha} and \eqref{eq:6.1.beta} yield together
\begin{equation}
  \Metric{U_{c_0}^{-1}A_hU_{c_0}-A_h}<\epsilon.
  \tag{6.1.\ensuremath{\gamma}}\label{eq:6.1.gamma}
\end{equation}
Clearly
\[
  U_{c_0}\sim\{\eta'_{ab^{-1}}\},
  \qquad
  \eta'_c=
  \begin{cases}
    1,&c=c_0,\\
    0,&\text{otherwise}.
  \end{cases}
\]
As $c_0\ne1$ this implies $\eta'_1=0$, i.e.
\begin{equation}
  \operatorname{Tr}_{\M}(U_{c_0})=0.
  \tag{6.1.\ensuremath{\delta}}\label{eq:6.1.delta}
\end{equation}
Thus \eqref{eq:6.1.gamma}, \eqref{eq:6.1.delta} show that $U=U_{c_0}$ meets
all our requirements.
\end{proof}

We can now prove
\begin{lemma}\phantomsection\label{lem:6.1.2}
Approximate finiteness implies the Property $\Gamma$.
\end{lemma}

\begin{proof}\footnote{There exists also a relatively simple direct proof,
based on Def. 4.1.1.}
It suffices to find a group $\mathfrak G$ which fulfills the condition (i) in
\lemref{lem:5.3.4}{5.3.4} (so that the $\M$ of \secref{sec:5.3}{5.3} is a
factor, of course, in case $(\mathrm{II}_1)$), (i) in
\lemref{lem:5.2.2}{5.2.2} (so that $\M$ is approximately finite by
\lemref{lem:5.6.1}{5.6.1}), and (i) in \lemref{lem:6.1.1}{6.1.1} (so that
$\M$ possesses the property $\Gamma$). Then \thmref{thm:XIV}{XIV} and our
remark after \defref{def:6.1.1}{6.1.1} take care of everything.

Now the group $\mathfrak G=\overline{\mathfrak G}$ of
\lemref{lem:5.6.2}{5.6.2} possesses the first two properties. As to (i) in
\lemref{lem:6.1.1}{6.1.1}: Let $a^{(1)},\ldots,a^{(n)}$ be the given
elements of $\mathfrak G$. Choose $p$ so that none of the
$a^{(1)},\ldots,a^{(n)}$ move any element other than $1,\ldots,p$. Let $c_0$
be the transposition of $p+1$ with $p+2$. Then $c_0\ne1$ and it commutes with
$a^{(1)},\ldots,a^{(n)}$.
\end{proof}

\PaperSection{sec:6.2}{Existence of not approximately finite factors}
We now proceed to establish the decisive negative result.

\begin{lemma}\phantomsection\label{lem:6.2.1}
Let a group $\mathfrak G$ fulfilling (i) in \lemref{lem:5.3.4}{5.3.4} be
given, which possesses this property:
% Editorial note E63: corrected the source's ``Hausdorf's'' to ``Hausdorff's''.
\footnote{This proof copies to a certain
extent Hausdorff's famous $1/2-1/3$ division of the sphere. In this connection
the use of the free group in Lemma 6.2.2 should be noted.}

\begin{enumerate}[label=(i),leftmargin=3em]
\item There exists a set $\mathfrak F\subseteq\mathfrak G$ with these
properties:
  \begin{enumerate}[label=(i\textsubscript{\arabic*}),leftmargin=3em]
  \item There exists a $c_1\in\mathfrak G$ such that
  \[
    \mathfrak F+c_1\mathfrak F c_1^{-1}=\mathfrak G-(1).
  \]
  \item There exists a $c_2\in\mathfrak G$, such that the three sets
  $c_2^l\mathfrak F c_2^{-l}$, $l=0,\pm1$, are disjoint.
  \end{enumerate}
\end{enumerate}
Then the $\M$ of \secref{sec:5.3}{5.3} does not possess the property $\Gamma$.
\end{lemma}

\begin{proof}
Assume the opposite, i.e. that $\M$ possesses the property $\Gamma$. Apply
\defref{def:6.1.1}{6.1.1} with $n=2$. Put $A_1=U_{c_1}$,
$A_2=U_{c_2}$, while $\epsilon>0$ will be chosen subsequently. Form the
$U=U(A_1,A_2;\epsilon)$ described there.

Then we have:
\begin{align}
  \operatorname{Tr}_{\M}(U)&=0,
    \tag{6.2.\ensuremath{\alpha}}\label{eq:6.2.alpha}\\
  \Metric{U^{-1}U_{c_h}U-U_{c_h}}&<\epsilon
    \qquad\text{for }h=1,2.
    \tag{6.2.\ensuremath{\beta}}\label{eq:6.2.beta}
\end{align}
As $U_{c_h},U$ are both unitary and
\[
  U_{c_h}^{-1}U(U^{-1}U_{c_h}U-U_{c_h})
  =U-U_{c_h}^{-1}UU_{c_h},
\]
\eqref{eq:6.2.beta} is equivalent to
\begin{equation}
  \Metric{U-U_{c_h}^{-1}UU_{c_h}}<\epsilon
  \qquad\text{for }h=1,2.
  \tag{6.2.\ensuremath{\gamma}}\label{eq:6.2.gamma}
\end{equation}

Now determine the $\eta_c$ in the sense of \lemref{lem:5.3.2}{5.3.2} for
$U$, $U_{c_h}^{-1}UU_{c_h}$, $U-U_{c_h}^{-1}UU_{c_h}$ in succession. If the
first is $\theta_c$ it is easy to verify that the second is
$\theta_{c_hcc_h^{-1}}$\textsuperscript{\hyperref[fn:50]{50}} and hence the third is
$\theta_c-\theta_{c_hcc_h^{-1}}$. Therefore the application of (i) in
\lemref{lem:5.3.6}{5.3.6} to $U$ and to $U-U_{c_h}^{-1}UU_{c_h}$ gives
\begin{align*}
  \Metric{U}^{2}&=\sum_{c\in\mathfrak G}|\theta_c|^2,\\
  \Metric{U-U_{c_h}^{-1}UU_{c_h}}^{2}
    &=\sum_{c\in\mathfrak G}
      |\theta_c-\theta_{c_hcc_h^{-1}}|^2.
\end{align*}
As $U$ is unitary, $\Metric{U}^{2}=1$. Considering this and
\eqref{eq:6.2.gamma} these equations yield
\begin{align}
  \sum_{c\in\mathfrak G}|\theta_c|^2&=1,
    \tag{6.2.\ensuremath{\delta}}\label{eq:6.2.delta}\\
  \left(\sum_{c\in\mathfrak G}
  |\theta_c-\theta_{c_hcc_h^{-1}}|^2\right)^{1/2}&<\epsilon
    \qquad\text{for }h=1,2.
    \tag{6.2.\ensuremath{\epsilon}}\label{eq:6.2.epsilon}
\end{align}

After these preparations we introduce a measure in $\mathfrak G$ by defining
\[
  \nu(\mathfrak A)=\sum_{c\in\mathfrak A}|\theta_c|^2
  \qquad\text{for }\mathfrak A\subseteq\mathfrak G.
\]
Then \eqref{eq:6.2.delta} becomes
\begin{equation}
  \nu(\mathfrak G)=1,
  \tag{6.2.\ensuremath{\zeta}}\label{eq:6.2.zeta}
\end{equation}
and \eqref{eq:6.2.alpha} means $\theta_1=0$, i.e.
\begin{equation}
  \nu((1))=0.
  \tag{6.2.\ensuremath{\eta}}\label{eq:6.2.eta}
\end{equation}
The triangle inequality in infinitely many dimensions gives
\[
 \left|\left(\sum_{c\in\mathfrak A}|\theta_c|^2\right)^{1/2}
 -\left(\sum_{c\in\mathfrak A}|\theta_{c_hcc_h^{-1}}|^2\right)^{1/2}\right|
 \leq
 \left(\sum_{c\in\mathfrak G}
 |\theta_c-\theta_{c_hcc_h^{-1}}|^2\right)^{1/2}.
\]
The left-hand side is clearly
$|\nu(\mathfrak A)^{1/2}-\nu(c_h\mathfrak A c_h^{-1})^{1/2}|$. The
right-hand side is
$\leq(\sum_{c\in\mathfrak G}|\theta_c-
\theta_{c_hcc_h^{-1}}|^2)^{1/2}$, which is $<\epsilon$ by
\eqref{eq:6.2.epsilon}. So we have
\begin{equation}
  |\nu(\mathfrak A)^{1/2}
   -\nu(c_h\mathfrak A c_h^{-1})^{1/2}|<\epsilon.
  \tag{6.2.\ensuremath{\theta}}\label{eq:6.2.theta}
\end{equation}
Now by \eqref{eq:6.2.zeta}, $\nu(\mathfrak A)$ and
$\nu(c_h\mathfrak A c_h^{-1})$ are $\leq\nu(\mathfrak G)=1$. Hence
\begin{align*}
 |\nu(\mathfrak A)-\nu(c_h\mathfrak A c_h^{-1})|
 &=|\nu(\mathfrak A)^{1/2}-
   \nu(c_h\mathfrak A c_h^{-1})^{1/2}|
   (\nu(\mathfrak A)^{1/2}+
   \nu(c_h\mathfrak A c_h^{-1})^{1/2})\\
 &\leq2|\nu(\mathfrak A)^{1/2}-
   \nu(c_h\mathfrak A c_h^{-1})^{1/2}|.
\end{align*}
Therefore \eqref{eq:6.2.theta} becomes
\begin{equation}
 |\nu(\mathfrak A)-\nu(c_h\mathfrak A c_h^{-1})|<2\epsilon.
 \tag{6.2.\ensuremath{\iota}}\label{eq:6.2.iota}
\end{equation}

Let us apply \eqref{eq:6.2.iota} to
$\mathfrak F,c_1$; $\mathfrak F,c_2$;
$c_2^{-1}\mathfrak F c_2,c_2$ in place of its $\mathfrak A,c_h$. Then
\begin{align}
 |\nu(\mathfrak F)-\nu(c_1\mathfrak F c_1^{-1})|&<2\epsilon,
  \tag{6.2.\ensuremath{\kappa}}\label{eq:6.2.kappa}\\
 |\nu(\mathfrak F)-\nu(c_2\mathfrak F c_2^{-1})|&<2\epsilon,
  \tag{6.2.\ensuremath{\lambda}}\label{eq:6.2.lambda}\\
 |\nu(c_2^{-1}\mathfrak F c_2)-\nu(\mathfrak F)|&<2\epsilon
  \tag{6.2.\ensuremath{\mu}}\label{eq:6.2.mu}
\end{align}
obtain. Now (i$_1$) and \eqref{eq:6.2.zeta}, \eqref{eq:6.2.eta} and
\eqref{eq:6.2.kappa} give
\[
  \nu(\mathfrak F)+(\nu(\mathfrak F)+2\epsilon)>1,
\]
i.e.
\begin{equation}
  \nu(\mathfrak F)>\frac12-\epsilon.
  \tag{6.2.\ensuremath{\nu}}\label{eq:6.2.nu}
\end{equation}
On the other hand (i$_2$) and \eqref{eq:6.2.zeta},
\eqref{eq:6.2.lambda} and \eqref{eq:6.2.mu} give
\[
  \nu(\mathfrak F)+(\nu(\mathfrak F)-2\epsilon)
  +(\nu(\mathfrak F)-2\epsilon)<1,
\]
i.e.
\begin{equation}
  \nu(\mathfrak F)<\frac13+\frac43\epsilon.
  \tag{6.2.o}\label{eq:6.2.o}
\end{equation}
\eqref{eq:6.2.nu} and \eqref{eq:6.2.o} imply
\[
  \frac12-\epsilon<\frac13+\frac43\epsilon
  \quad\text{or}\quad \frac1{14}<\epsilon.
\]
Hence it suffices to choose $\epsilon=1/14$ in order to have a contradiction.
Thus we have shown that $\M$ cannot possess the property $\Gamma$.
\end{proof}

\begin{namedstatement}{Corollary}{1}\phantomsection\label{cor:6.2.1}
Under the above assumptions, $\M$ is not approximately finite.
\end{namedstatement}
\begin{proof}
Combine the above lemma with \lemref{lem:6.1.2}{6.1.2}.
\end{proof}

\begin{namedstatement}{Corollary}{2}\phantomsection\label{cor:6.2.2}
We may replace condition (i$_2$) on $\mathfrak G$ and $\mathfrak F$ by

\noindent(i$_2'$). There exist two elements $c_2$ and $c_3$ in $\mathfrak G$
such that $\mathfrak F,c_2\mathfrak F c_2^{-1},c_3\mathfrak F c_3^{-1}$ are
mutually exclusive.
\end{namedstatement}

\begin{proof}
In the first paragraph of the above proof we let $A_3=U_{c_3}$ and then take
$U=U(A_1,A_2,A_3,\epsilon)$ instead of $U=U(A_1,A_2,\epsilon)$. We have
$h=1,2,3$ in \eqref{eq:6.2.beta}, \eqref{eq:6.2.gamma},
\eqref{eq:6.2.epsilon}, \eqref{eq:6.2.theta} and \eqref{eq:6.2.iota}, instead
of $h=1,2$. Finally we obtain an equation
\begin{equation}
  |\nu(\mathfrak F)-\nu(c_3\mathfrak F c_3^{-1})|<2\epsilon
  \tag{6.2.\ensuremath{\mu'}}\label{eq:6.2.mu-prime}
\end{equation}
instead of \eqref{eq:6.2.mu}, by applying \eqref{eq:6.2.iota} to
$\mathfrak F,c_3$ instead of $c_2^{-1}\mathfrak F c_2,c_2$.
\eqref{eq:6.2.mu-prime} is used instead of \eqref{eq:6.2.mu} in order to
obtain \eqref{eq:6.2.o}, and the rest of the proof is the same.
\end{proof}

We conclude by giving a specific instance of an $\M$ which does not possess
the property $\Gamma$ and is not approximately finite. We do this by giving
an example of a group $\mathfrak G$ which fulfills the requirements of
\lemref{lem:6.2.1}{6.2.1}.

\begin{lemma}\phantomsection\label{lem:6.2.2}
Let $\mathfrak G$ be the free group of two generators $c_1,c_2$. Then
$\mathfrak G$ fulfills the requirements of \lemref{lem:6.2.1}{6.2.1}.
\end{lemma}

\begin{proof}
Ad (i) in \lemref{lem:5.3.4}{5.3.4}: Consider an $a\in\mathfrak G$,
$a\ne1$. Write $a$ as a power product of $c_1,c_2$ of minimum length. Then
it is easy to verify that $c_1^{-n}ac_1^n$, $n=1,2,\ldots$, are pairwise
different unless $a=c_1^m$. In the latter case we have that the
$c_2^{-n}ac_2^n=c_2^{-n}c_1^mc_2^n$ are distinct unless $m=0$, i.e. $a=1$.
Thus unless $a=1$, $\mathfrak L_a$ is infinite.

Ad (i) in \lemref{lem:6.2.1}{6.2.1}: Let $\mathfrak F$ be the set of
$a\in\mathfrak G$ which when written as a power product of $c_1,c_2$ of
minimum length end with a $c_1^n$, $n=\pm1,\pm2,\ldots$. Then (i$_1$) and
(i$_2$) in \lemref{lem:6.2.1}{6.2.1}, i.e. (i) eod., are immediately
verified.
\end{proof}

Combining \lemref{lem:6.2.2}{6.2.2} with Corollary 1 to
\lemref{lem:6.2.1}{6.2.1} gives
\begin{theorem}\phantomsection\label{thm:XVI}
There exist not approximately finite factors $\M$ in the case
$(\mathrm{II}_1)$.
\end{theorem}
Or, considering \thmref{thm:XV}{XV},
\begin{namedstatement}{Theorem}{XVI$'$}\phantomsection\label{thm:XVI-prime}
There exist in case $(\mathrm{II}_1)$ more than one algebraical type
$\overline{\M}$.
\end{namedstatement}

\PaperSection{sec:6.3}{Discussion of further examples}
We wish in this section to present a number of examples of groups $\mathfrak G$
for which the corresponding $\M$ of \secref{sec:5.3}{5.3} is not
approximately finite.

\begin{lemma}\phantomsection\label{lem:6.3.1}
Let $\mathfrak G$ be a group which contains two multiplicative subsets
$\mathfrak F_1$ and $\mathfrak F_2$. Suppose that $\mathfrak G$ is the free
product of $\mathfrak F_1$ and $\mathfrak F_2$ and that $\mathfrak F_1$
contains at least two elements and $\mathfrak F_2$ contains at least three.
Then $\mathfrak G$ satisfies condition (i) of \lemref{lem:5.3.4}{5.3.4},
condition (i$_1$) of \lemref{lem:6.2.1}{6.2.1} and condition (i$_2'$) of
Corollary 2 of that lemma (for the same $\mathfrak F$).
\end{lemma}

\begin{proof}
Let $c_1$ be an element of $\mathfrak F_1$ which is not $1$, and $c_2,c_3$
elements of $\mathfrak F_2$ which are each not $1$.

We prove first that (i) of \lemref{lem:5.3.4}{5.3.4} holds, i.e. that for
$a\in\mathfrak G$, $a\ne1$, the $\mathfrak L_a$ is infinite.

Suppose an $a\in\mathfrak G$, $a\ne1$, is given. Since $\mathfrak G$ is a
free product of $\mathfrak F_1$ and $\mathfrak F_2$, each $a$ has a unique
factorization $a_1\cdot a_2\cdots a_p$ into a product in which the factors are
alternately in $\mathfrak F_1$ and $\mathfrak F_2$ and no factor is $1$.

We note for later use that if $c\in\mathfrak F_i$, then
$c^{-1}\in\mathfrak F_i$. Suppose $i=1$. Then $c^{-1}$ will have a
factorization $b_1\cdot b_2\cdots b_r$. If $b_1\in\mathfrak F_2$ we have
$1=cc^{-1}=c\cdot b_1\cdot b_2\cdots b_r\ne1$. Thus
$b_1\in\mathfrak F_1$. But $1=cc^{-1}=(cb_1)\cdot b_2\cdots b_r$ implies
$r=1$. Consequently $c^{-1}=b_1\in\mathfrak F_1$.

Now suppose firstly that $a_1$ (the first factor) and $a_p$ (the last factor)
are both from $\mathfrak F_1$. Let $b_n=(c_2c_1)^n$. Then the set of
$b_n^{-1}ab_n$ are for $n=1,2,\ldots$ all distinct. (The factorization is
immediately apparent.) Hence $\mathfrak L_a$ is infinite in this case.

If $a_1$ and $a_p$ are both from $\mathfrak F_2$, we let
$b_n=(c_1c_2)^n$ and the result is again apparent.

If $a_1$ is from $\mathfrak F_1$ and $a_p$ is from $\mathfrak F_2$, we
consider $c_2$ and $c_3$ and denote by $c'$ that one of $c_2,c_3$ which is not
equal to $a_p^{-1}$. Let $b_n=(c'c_1)^n$. Then for $n=1,2,\ldots$ the
factorization of $b_n^{-1}ab_n$ is obtained by simply writing out the
expressions for these elements and coalescing $a_p$ and the first $c'$ in
$b_n$. $a_pc'$ is in $\mathfrak F_2$ but is not $1$ since $c'\ne a_p^{-1}$.
Hence the factorization has been obtained. These factorizations for
$n=1,2,\ldots$ show that the $b_n^{-1}ab_n$ are distinct, and hence
$\mathfrak L_a$ is infinite.

If $a_1$ is in $\mathfrak F_2$ and $a_p$ is in $\mathfrak F_1$ we define
$c'$ as that one of $c_2,c_3$ which is not equal to $a_1$. An argument
similar to that of the preceding paragraph will show that in this case also
$\mathfrak L_a$ is infinite.

Thus we have established (i) of \lemref{lem:5.3.4}{5.3.4}.

To establish (i$_1$) of \lemref{lem:6.2.1}{6.2.1} and (i$_2'$) of Corollary
2 of that lemma for the same $\mathfrak F$, we define $\mathfrak F$ as the
set of $a$'s, $a\ne1$, having a factorization $a_1\cdots a_p$ in which $a_p$
is in $\mathfrak F_1$.

We show firstly (i$_1$) in \lemref{lem:6.2.1}{6.2.1}, i.e.
\[
  \mathfrak F+c_1^{-1}\mathfrak F c_1=\mathfrak G-(1).
\]
If $a\in\mathfrak G$, $a\ne1$, then $a$ has a factorization
$a_1\cdots a_p$. Now if $a$ is not $\in\mathfrak F$, $a_p$ is not
$\in\mathfrak F_1$ but $a_p\in\mathfrak F_2$. In this case we consider
$c_1ac_1^{-1}=c_1a_1\cdots a_pc_1^{-1}$. If $p$ is even, i.e.
$p=2,4,\ldots$, then $a_1\in\mathfrak F_1$ (since
$a_p\in\mathfrak F_2$) and the factorization of $c_1ac_1^{-1}$ is either
$(c_1a_1)\cdot a_2\cdots a_p\cdot c_1^{-1}$ or (if $c_1a_1=1$)
$a_2\cdots a_p c_1^{-1}$. If $p$ is odd, i.e. $p=1,3,\ldots$,
$a_1\in\mathfrak F_2$ and $c_1ac_1^{-1}$ has the factorization
$c_1\cdot a_1\cdots a_p\cdot c_1^{-1}$. These factorizations all show that
$c_1ac_1^{-1}$ is $\in\mathfrak F$. Thus if $a\ne1$ and $a$ is not
$\in\mathfrak F$, $c_1ac_1^{-1}\in\mathfrak F$ or
$a\in c_1^{-1}\mathfrak F c_1$. This is equivalent to
$\mathfrak G-(1)\subseteq\mathfrak F+c_1^{-1}\mathfrak F c_1$. However this
inclusion is not proper. For $1$ is not in $\mathfrak F$. Consequently $1$
is not in $c_1^{-1}\mathfrak F c_1$ and
$\mathfrak F+c_1^{-1}\mathfrak F c_1$. We can now conclude
\[
  \mathfrak F+c_1^{-1}\mathfrak F c_1=\mathfrak G-(1).
\]
One can establish (i$_2'$) of Corollary 2 of \lemref{lem:6.2.1}{6.2.1} in
a similar (indeed easier) way.

Thus our Lemma has been proven.
\end{proof}

This lemma permits us to offer many other examples similar to the $\mathfrak G$
of \lemref{lem:6.2.2}{6.2.2}. For we may obtain $\mathfrak F_1$ and
$\mathfrak F_2$ in a number of ways, and then
$\mathfrak G=\mathfrak G(\mathfrak F_1,\mathfrak F_2)$ is constructed by
adding to $1$ the ``free'' products $a_1\cdot a_2\cdots a_p$ where the $a_i$'s
are never $1$ and are alternately from $\mathfrak F_1$ and $\mathfrak F_2$.
As an example we may take $\mathfrak F_1$ as the group generated by a single
generator $a$ of order $p=2,3,\ldots,\infty$ and $\mathfrak F_2$ as the group
generated by an element $b$ of order $q=3,4,\ldots,\infty$.
$\mathfrak G(a,b)=\mathfrak G(\mathfrak F_1,\mathfrak F_2)$ is then an
example. Of course we may take for $\mathfrak F_1$ any group of order $\geq2$
and $\mathfrak F_2$ of order $\geq3$.

The restriction that $\mathfrak F_2$ contains at least 3 members rules out the
case in which $\mathfrak F_1$ and $\mathfrak F_2$ are both groups of order 2,
i.e. $\mathfrak F_1$ consists of $(1,a)$ and $\mathfrak F_2$ consists of
$(1,b)$. Here of course we do not have the (i$_2$) of
\lemref{lem:6.2.1}{6.2.1} or the (i$_2'$) of its second corollary. However,
condition (i) of \lemref{lem:5.3.4}{5.3.4} also fails. For the element $ab$
has only $ba$ and itself as conjugates, i.e. $\mathfrak L_{ab}=(ab,ba)$. Thus
$\mathfrak L_{ab}$ is finite. Hence by \lemref{lem:5.3.4}{5.3.4}, the $\M$
for this $\mathfrak G$ is not even a factor.

Another lemma, which we shall use in an appendix, is readily shown now.
\begin{lemma}\phantomsection\label{lem:6.3.2}
If $\mathfrak G$ is such that to every $a\in\mathfrak G$, $a\ne1$, there is a
$b\in\mathfrak G$, $b$ of order $\geq3$, such that $\mathfrak G$ contains the
free group $\mathfrak G(a,b)$, then $\mathfrak G$ fulfills (i) of
\lemref{lem:5.3.4}{5.3.4}. (If $a$ is not of order 2, $b$ need only differ
from the identity.)
\end{lemma}

\begin{proof}
Let $\overline{\mathfrak L}_a$ denote the class of $a$ relative to the subgroup
$\mathfrak G(a,b)$. Our discussion following \lemref{lem:6.3.1}{6.3.1}
shows that $\mathfrak G(a,b)$ satisfies the hypotheses of
\lemref{lem:6.3.1}{6.3.1}. Consequently that Lemma and
\lemref{lem:5.3.4}{5.3.4} imply that $\overline{\mathfrak L}_a$ is infinite.
Since $\overline{\mathfrak L}_a\subseteq\mathfrak L_a$ we have
$\mathfrak L_a$ infinite for $\mathfrak G$ itself, and thus (i) of
\lemref{lem:5.3.4}{5.3.4} holds.
\end{proof}

\phantomsection\label{appendix}
\pdfbookmark[1]{Appendix}{bookmark.appendix}
\begin{center}\bfseries\MakeUppercase{Appendix}\footnote{This example was
originally part of a later investigation by the second-named author, but it
can be most easily presented here, using our present methods.}\end{center}

In this appendix we give examples of two Class $(\mathrm{II}_1)$ factors which
are not isomorphic and yet each of which is isomorphic to a subset of the
other. The factors will be the factors associated with certain groups
$\mathfrak G_1$ and $\mathfrak H_1$ as in
\secref{sec:5.3}{5.3}. These groups are best presented as subgroups of a group
$\mathfrak G$ which we shall now define.

$\mathfrak G$ is determined by generators $x_r$, one for each rational number
$r=p/q$. The commutation rules for these generators are determined as follows.
Let $A$ denote the set of dyadic rationals $q/2^n$ in the most reduced form
where $n=0,1,2,\ldots$ and $q$ is odd except possibly when $n=0$. Let $B$
denote the set of all other rationals.

The commutation rules for $x_r$ are the following:
\begin{enumerate}[label=(\roman*),leftmargin=3em]
\item If $r\in A$, then $x_r$ combines freely with any product
$x_{s_1}\cdot x_{s_2}\cdots x_{s_p}$ if $s_1,s_2,\ldots,s_p$ are all $<r$.
\item If $r\in B$, $x_r$ commutes with all products
$x_{s_1}\cdot x_{s_2}\cdots x_{s_p}$ if $s_1,s_2,\ldots,s_p$ are all $<r$.
\item All $x_r$ are of order $\infty$.
\end{enumerate}
One can show that the set of products
$x_{r_1}^{n_1}\cdots x_{r_p}^{n_p}$,
$n_i=\pm1,\pm2,\ldots$ (with a finite number of factors) constitutes a
denumerably infinite group $\mathfrak G$.\footnote{This can be most readily
shown by obtaining a ``normal form'' for each such product
$x_{r_1}^{n_1}\cdots x_{r_p}^{n_p}$. This normal form may be obtained as
follows. Let $r_i$ denote the largest subscript which appears and is in $B$.
We move $x_{r_i}^{n_i}$ as far to the left as possible. Let $r_{i'}$ be the
next largest subscript, which is also in $B$. We then move
$x_{r_{i'}}^{n_{i'}}$ as far to the left as possible. We continue on until we
have considered all the $r_i\in B$. The resulting product is then in a form in
which the $x_{r_j}^{n_j}$, $r_j\in A$, have the same order as before, but each
such $x_{r_j}^{n_j}$, $r_j\in A$, is followed by a product
$x_{r_{i_1}}^{n_{i_1}}\cdots x_{r_{i_k}}^{n_{i_k}}$ with the $r_{i_t}$ each
in $B$, each $<r_j$, and with $r_{i_1}<r_{i_2}<\cdots<r_{i_k}$. We may use
this form to establish the associativity of the composition rule. The
existence of an inverse is clear.} This is also true for the set of products
in which the $x_r$'s are restricted, so that $r$ is in a fixed non-empty
subset $S$ of the rational numbers.

Let $\mathfrak H_r$ be the subgroup of $\mathfrak G$ determined by
the generators $x_{r'}$, with $r'\leq r$. Let $\mathfrak G_r$ be the subgroup
of $\mathfrak G$ determined by the generators $x_{r'}$, with $r'<r$. We can
show
\begin{namedstatement}{Lemma}{A.1}\phantomsection\label{lem:A.1}
\begin{enumerate}[label=(\roman*),leftmargin=3em]
\item $\mathfrak G_r\subseteq\mathfrak H_r$.
\item If $r'<r$, $\mathfrak H_{r'}\subseteq\mathfrak G_r$.
\item If $n=\pm1,\pm2,\ldots$, then
$\mathfrak H_{r+n}$ is isomorphic to $\mathfrak H_r$.
\item $\mathfrak G_r$ has property (i) of \lemref{lem:5.3.4}{5.3.4}.
\item $\mathfrak G_r$ has property (i) of \lemref{lem:6.1.1}{6.1.1}.
\item Assume $r\in A$. Then $\mathfrak H_r$ satisfies the
hypotheses of \lemref{lem:6.3.1}{6.3.1}, and consequently it satisfies
(i$_1$) and (i$_2'$) of \lemref{lem:6.2.1}{6.2.1} and its Corollary 2 and (i)
in \lemref{lem:5.3.4}{5.3.4}.
\end{enumerate}
\end{namedstatement}

\begin{proof}
Ad (i) and (ii) are clear.

Ad (iii). The correspondence $x_{r+n}\sim x_r$ between the generators of
$\mathfrak H_{n+r}$ and $\mathfrak H_r$ is one-to-one
and preserves the properties $r\in A$ and $r\in B$. Consequently it
determines an isomorphism between $\mathfrak H_{n+r}$ and $\mathfrak H_r$.

Ad (iv). Suppose $a\in\mathfrak G_r$. Then
$a=x_{r_1}^{n_1}\cdots x_{r_p}^{n_p}$ with each $r_i<r$. Let $r'$ be the
$\operatorname{Max}_{i=1,\ldots,p}r_i$. Then of course $r'<r$ and there is an
$r''\in A$ such that $r'<r''<r$. By the first commutation rule above $x_{r''}$
combines freely with $a$. By the definition of $\mathfrak G_r$,
$\mathfrak G(a,x_{r''})\subseteq\mathfrak G_r$. The third rule above tells us
that $x_{r''}$ is of infinite order. These last two results show that the
hypotheses of \lemref{lem:6.3.2}{6.3.2} are satisfied. It follows from that
lemma that the condition (i) of \lemref{lem:5.3.4}{5.3.4} is satisfied.
% Editorial note E61: the source prints x_r at this point; x_{r''} is required
% by the immediately preceding choice and by the application of Lemma 6.3.2.

Ad (v). $\mathfrak G_r$ has property (i) of
\lemref{lem:6.1.1}{6.1.1}: For let $a_1,\ldots,a_k$ be any $k$ elements of
$\mathfrak G_r$ and let $x_{r_1},x_{r_2},\ldots,x_{r_l}$ denote the generators
which appear explicitly in $a_1,\ldots,a_k$. Let
$r'=\operatorname{Max}_{i=1,\ldots,l}r_i$. Then $r'<r$ by the definition of
$\mathfrak G_r$. Hence there is an $r''\in B$ such that $r'<r''<r$. Let
$c=x_{r''}$. Then $c\ne1$, $c\in\mathfrak G_r$ and $c$ commutes with
$a_1,\ldots,a_k$ by the second commutation rule above. This result shows that
(i) in \lemref{lem:6.1.1}{6.1.1} is satisfied.

Ad (vi). For $\mathfrak H_r$ we let
$\mathfrak F_1=\mathfrak G_r$ and $\mathfrak F_2=(x_r^n)$. Now $x_r$ is of
infinite order and the first commutation rule shows that
$\mathfrak H_r$ is the free product of $\mathfrak G_r$ and
$\mathfrak F_2$. These show that the hypotheses of
\lemref{lem:6.3.1}{6.3.1} are fulfilled. That lemma itself now shows that
(i$_1$) and (i$_2'$) of \lemref{lem:6.2.1}{6.2.1} and its Corollary 2 are
satisfied, and (i) in \lemref{lem:5.3.4}{5.3.4}.
\end{proof}

We next state a quite general lemma.
\begin{namedstatement}{Lemma}{A.2}\phantomsection\label{lem:A.2}
Let $\M$ be a ring which contains $(\alpha1)$ and suppose $f\in\Hilbert$ is
such that $Af=0$ and $A\in\M$ imply $A=0$. Let
$\mathfrak M=\mathfrak M_f^{\M}$. Then if $A$ is closed, has domain dense,
and $A\mathbin{\eta}\M$, $A$ is bounded if and only if it is bounded on
$\mathfrak M$.
\end{namedstatement}

\begin{proof}
It is clear that we need only show that if $A$ is unbounded, it is unbounded
on $\mathfrak M$. Let $A=WB$ where $W$ is partially isometric and $B$ is
definite (cf. \cite{ref5}, \S4.4). Since $A\mathbin{\eta}\M$, we have
$B\mathbin{\eta}\M$ and also for all $g$, $|Ag|=|Bg|$. (This may be $\infty$
of course.) Hence we may substitute $B$ for $A$, or what is the same thing
assume that $A$ is definite. But $A$ definite implies that there is a
resolution $\{E(\lambda)\}\subseteq\M$ such that
\[
  A=\int_0^\infty\lambda\,dE(\lambda).
\]
Since $A$ is unbounded, $1-E(\lambda)\ne0$ for $\lambda<\infty$. Let
$g=(1-E(\lambda))f$ for any fixed $\lambda$. By our hypothesis on $f$,
$g\ne0$ since $1-E(\lambda)\ne0$. From the formula for $A$ we see
$|Ag|>\lambda|g|$. Furthermore $g\in\mathfrak M_f^{\M}=\mathfrak M$. Thus if
$A$ is unbounded, it is unbounded on $\mathfrak M$.
\end{proof}

\begin{corollary}\phantomsection\label{cor:A.2}
If $A\in\M$ the bound of $A$ on $\mathfrak M$ is the same as the bound of
$A$ $(\OpNorm{A})$.
\end{corollary}
\begin{proof}
It is clear that the bound of $A$ on $\mathfrak M$ is $\leq\OpNorm{A}$. To
show that it can't be less one uses an argument similar to the above.
\end{proof}

We wish to apply \lemref{lem:A.2}{A.2} to the case in which $\M$ is a subring
of the set of group numbers in the sense of \secref{sec:5.3}{5.3}. We may do
this by taking $f=\varphi_1$ in the sense of (5.3.$\alpha$). For if $A\in\M$,
we have that $A\sim\{\eta_{a^{-1}b}\}$ by \lemref{lem:5.3.2}{5.3.2} and
\[
  A\varphi_1=\sum_{c\in\mathfrak G}\eta_c\varphi_c
\]
by \eqref{eq:49-gamma}. Thus $A\varphi_1=0$ implies $\eta_c=0$ for all
$c\in\mathfrak G$ and hence $A=0$.

It is also useful to notice in applying this to group numbers that if we have
any sequence $\{\eta_c\}$ with $\sum_{c\in\mathfrak G}|\eta_c|^2<\infty$ then
the matrix $\{\eta_{a^{-1}b}\}$ determines an $A$ with a closure and domain
dense. For it is clear from \eqref{eq:49-gamma} that $A$ is defined for all
$\varphi_a$, $a\in\mathfrak G$ of (5.3.$\alpha$) and hence for the linear
combinations of these. Since the latter are dense, $A$ has domain dense.
Furthermore the matrix $\{\overline{\eta}_{b^{-1}a}\}$ is adjoint to it. Since
this latter also has domain dense, $A$ must have a closure. We shall consider
$A$ to be the least closed transformation, with matrix
$\{\eta_{a^{-1}b}\}$. On its domain, $A$ is the limit of finite linear
combinations of the $U_c$ and hence is $\mathbin{\eta}\M$. Cf. Lemmas 5.3.1,
5.3.2 and 5.3.6.

\begin{namedstatement}{Lemma}{A.3}\phantomsection\label{lem:A.3}
Suppose $\mathfrak G^0$ and $\mathfrak G^{00}$ are two groups and $\M^0$ and
$\M^{00}$ their respective group numbers in the sense of
\secref{sec:5.3}{5.3}. Then if $\mathfrak G^0\subseteq\mathfrak G^{00}$,
$\M^0$ is algebraically isomorphic with a subring of $\M^{00}$.
\end{namedstatement}

\begin{proof}
For an $A\in\M^{00}$ we have a sequence
$\{\eta_c;c\in\mathfrak G^{00}\}$ as in \lemref{lem:5.3.2}{5.3.2}. Let
$\M^{0+}$ denote the set of $A$'s $\in\M^{00}$ for which $\eta_c=0$ if $c$
not $\in\mathfrak G^0$. Since $\mathfrak G^0$ is a subgroup, this set of
$A$'s is closed under the operations $\alpha A,A^*,A+B,AB$. Furthermore this
set of $A$'s is metrically closed, and as we mentioned at the end of
\secref{sec:5.5}{5.5}, for a subalgebra of group numbers, this implies ring
closure. Thus $\M^{0+}$ is a ring, and indeed is the ring determined by the
$U_c$, $c\in\mathfrak G^0$.

Let $\mathfrak M^0$ be the set $\subseteq\Hilbert^{00}$ determined by the
$\varphi_a$ with $a\in\mathfrak G^0$ (cf. (5.3.$\alpha$)). It is clear that we
can identify this with $\Hilbert^0$, the space on which we have $\M^0$
defined. Now for every $A\in\M^0$ we have a sequence
$\{\zeta_c;c\in\mathfrak G^0\}$ and corresponding matrix
$\{\zeta_{a^{-1}b}\}$. Correspondingly, we have a sequence
$\{\eta_c;c\in\mathfrak G^{00}\}$ with $\eta_c=\zeta_c$ for
$c\in\mathfrak G^0$ and $\eta_c=0$ otherwise. As we pointed out above, the
matrix $\{\eta_{a^{-1}b}\}$ determines a closed operator $\widetilde A$ with
domain dense. This is clearly $\mathbin{\eta}\M^{0+}$ and hence by the
application of \lemref{lem:A.2}{A.2} described above, it is bounded if and
only if it is bounded on $\mathfrak M_f^{\M^{0+}}$. Now one can show by our
above definition of $\M^{0+}$ that
$\mathfrak M_f^{\M^{0+}}=\mathfrak M^0$. Now $\widetilde A=A$ on
$\mathfrak M^0$ and thus is bounded there. Hence, by \lemref{lem:A.2}{A.2},
$\widetilde A$ is bounded and $\in\M^{0+}$. We now see that to every
$A\in\M^0$ there corresponds an $\widetilde A\in\M^{0+}$ such that
$\widetilde A$ agrees with $A$ on $\Hilbert^0=\mathfrak M^0$. On the other
hand, it is clear that to every $\widetilde A\in\M^{0+}$ we have an
$A\in\M^0$ with this property.

One can readily verify that this is an algebraic isomorphism.
\end{proof}

We are now ready for the final step. Let $\M_1$ be the set of group numbers
for $\mathfrak G_1$, $\M_2$ the set of group numbers for
$\mathfrak H_1$ (cf. \secref{sec:5.3}{5.3}).

\begin{namedstatement}{Lemma}{A.4}\phantomsection\label{lem:A.4}
\begin{enumerate}[label=(\roman*),leftmargin=3em]
\item $\M_1$ and $\M_2$ are factors in Case $(\mathrm{II}_1)$.
\item $\M_1$ is algebraically isomorphic with a subring of $\M_2$.
\item $\M_2$ is algebraically isomorphic with a subring of $\M_1$.
\item $\M_1$ has property $\Gamma$ of \defref{def:6.1.1}{6.1.1}.
\item $\M_2$ does not have property $\Gamma$.
\end{enumerate}
\end{namedstatement}

\begin{proof}
Ad (i). $\M_1$ is a factor by \lemref{lem:A.1}{A.1} (iv) and
\lemref{lem:5.3.4}{5.3.4}. $\M_2$ is a factor by \lemref{lem:A.1}{A.1} (vi)
and \lemref{lem:5.3.4}{5.3.4}. $\M_1$ and $\M_2$ are both in Case
$(\mathrm{II}_1)$ by \lemref{lem:5.3.5}{5.3.5} (i).

Ad (ii). This follows from \lemref{lem:A.1}{A.1} (i) and
\lemref{lem:A.3}{A.3}.

Ad (iii). Let $\M_0$ denote the set of group numbers for
$\mathfrak H_0$. \lemref{lem:A.1}{A.1} (iii) for $n=-1$ shows
that $\mathfrak H_0$ and $\mathfrak H_1$ are
isomorphic, and this of course implies that $\M_2$ is algebraically isomorphic
to $\M_0$. \lemref{lem:A.1}{A.1} (ii) and \lemref{lem:A.3}{A.3} imply that
$\M_0$ is isomorphic to a subring of $\M_1$. Hence $\M_2$ is algebraically
isomorphic to a subring of $\M_1$.

Ad (iv). This is implied by \lemref{lem:A.1}{A.1} (v) and
\lemref{lem:6.1.1}{6.1.1}.

Ad (v). This is implied by \lemref{lem:A.1}{A.1} (vi) and Corollary 2 of
\lemref{lem:6.2.1}{6.2.1}.

Since property $\Gamma$ is an isomorphism invariant, (iv) and (v) imply that
$\M_1$ and $\M_2$ are not algebraically isomorphic. Combining this with (ii),
(i) and (iii), we have
\end{proof}

\begin{namedstatement}{Theorem}{A}\phantomsection\label{thm:A}
There exist two rings $\M_1,\M_2$, each a factor in Case $(\mathrm{II}_1)$,
which are not algebraically isomorphic to each other but are such that each is
algebraically isomorphic to a subring of the other.
\end{namedstatement}

\bigskip
\noindent\textsc{The Institute for Advanced Study and}\\
\textsc{Columbia University}

\end{document}
